% ============================================================================
% Fichier     : APDC_main.tex
% Titre       : Audit de protection des données et conformité
% Auteur      : MINKA MI NGUIDJOÏ Thierry Emmanuel
% Institution : Chaire UNESCO d'Afrique Centrale Accès TIC, ENSP / LIMSI
% Version     : 1.0 -- novembre 2026 (traité complet : P00 à C16, annexes A à F)
% ============================================================================

\documentclass[12pt, twoside, openright, a4paper]{book}

\usepackage{Style_APDC}
\raggedbottom

\addbibresource{bibliographie/APDC_references.bib}

\hypersetup{
    pdftitle={Audit de protection des données et conformité},
    pdfauthor={MINKA MI NGUIDJOÏ Thierry Emmanuel},
    pdfsubject={MSN-2S4-SEC02 · Master Professionnel en Sécurité des Systèmes Numériques · Chaire UNESCO d'Afrique Centrale Accès TIC},
    pdfkeywords={protection des données personnelles, loi n° 2024/017, audit de conformité, registre des traitements, évaluation d'impact, ISO/IEC 27701, transferts, violation de données, NEK}
}

\begin{document}

% --------------------------------------------------------------------------
% ÉLÉMENTS PRÉLIMINAIRES (pages romaines)
% --------------------------------------------------------------------------
\frontmatter

    % ==============================================================================
% Fichier : front/00_page_titre.tex
% Titre   : Page de titre - Audit de Protection des Données et Conformité
% Auteur  : MINKA MI NGUIDJOÏ Thierry Emmanuel
% Version : 1.0 -- novembre 2026
% ==============================================================================

\begin{titlepage}
\thispagestyle{empty}

% --- Fond décoratif : bandeaux bleu marine et filets or ---
\begin{tikzpicture}[remember picture, overlay]
    \fill[CouleurPrimaire]
        (current page.north west) rectangle
        ([yshift=-3.6cm] current page.north east);
    \fill[CouleurAccent]
        ([yshift=-3.6cm] current page.north west) rectangle
        ([yshift=-4.0cm] current page.north east);
    \fill[CouleurPrimaire]
        (current page.south west) rectangle
        ([yshift=2.5cm] current page.south east);
    \fill[CouleurAccent]
        ([yshift=2.5cm] current page.south west) rectangle
        ([yshift=2.9cm] current page.south east);
    \node[align=center, text=white] at ([yshift=-1.8cm] current page.north) {%
        {\fontsize{11}{13}\selectfont
        \textsc{République du Cameroun}}\\[0.3ex]
        {\color{CouleurAccent!60!white}\rule{0.4\paperwidth}{0.4pt}}\\[0.6ex]
        {\fontsize{10}{12}\selectfont\textbf{\institutionprimaire}}\\[0.2ex]
        {\fontsize{8}{10}\selectfont \tutelle}\\[0.2ex]
        {\fontsize{9}{11}\selectfont \unitepedago}};
    \node[fill=white, rounded corners=4pt, inner sep=2pt]
        at ([xshift=1.75cm, yshift=-1.8cm] current page.north west)
        {\includegraphics[height=2.6cm]{figures/logos/enspy_logo.jpg}};
    \node[fill=white, rounded corners=4pt, inner sep=2pt]
        at ([xshift=-1.75cm, yshift=-1.8cm] current page.north east)
        {\includegraphics[height=2.6cm]{figures/logos/unesco_logo.jpeg}};
    \node[align=center] at ([yshift=1.25cm] current page.south) {%
        {\color{white}\bfseries\large Année académique \anneescolaire}\\[0.3ex]
        {\color{CouleurAccent!70!white}\small
        Version 1.0 | Document pédagogique de l'UE \codeUE}};
\end{tikzpicture}

\vspace*{0.6cm}

\begin{center}
    {\color{CouleurPrimaire}\rule{0.85\textwidth}{1.2pt}}\\[1.2ex]
    {\fontsize{26}{32}\selectfont\bfseries\color{CouleurPrimaire}
        Audit de Protection des Données\\[0.4ex]
        et Conformité}\\[0.8ex]
    {\color{CouleurAccent}\rule{0.85\textwidth}{2pt}}\\[0.5ex]
    {\color{CouleurPrimaire}\rule{0.85\textwidth}{0.5pt}}
\end{center}

\vspace{0.8cm}

\begin{center}
    {\fontsize{13}{16}\selectfont\itshape\color{CouleurSecondaire}
        Prouver la conformité, protéger les personnes~:\\
        auditer l'application de la loi n°~2024/017}
\end{center}

\vspace{0.5cm}

\begin{center}
    \begin{tikzpicture}
        \node[draw=CouleurPrimaire, fill=CouleurPrimaire, rounded corners=6pt,
              inner xsep=14pt, inner ysep=6pt]{%
            \color{white}\bfseries\large \sigleCours{} - UE~\codeUE
        };
    \end{tikzpicture}
\end{center}

\vspace{1.0cm}

% --- Schéma de la discipline : le cycle d'un audit de protection des données ---
\begin{center}
\begin{tikzpicture}[
    etape/.style = {draw=CouleurBordDef, fill=CouleurDef, rounded corners=3pt,
        minimum width=2.5cm, minimum height=0.8cm, align=center,
        font=\small\bfseries\color{CouleurPrimaire}, inner sep=2pt},
    arrow/.style = {-{Stealth[length=4pt]}, thick, color=CouleurSecondaire}
]
\node[etape] (n1) at (0, 0)      {Recenser};
\node[etape] (n2) at (3.2, 0)    {Apprécier\\le risque};
\node[etape] (n3) at (6.4, 0)    {Vérifier};
\node[etape] (n4) at (6.4, -1.4) {Conclure};
\node[etape] (n5) at (3.2, -1.4) {Corriger};
\node[etape] (n6) at (0, -1.4)   {Suivre};
\draw[arrow] (n1) -- (n2); \draw[arrow] (n2) -- (n3); \draw[arrow] (n3) -- (n4);
\draw[arrow] (n4) -- (n5); \draw[arrow] (n5) -- (n6); \draw[arrow] (n6) -- (n1);
\node[draw=CouleurPrimaire, fill=CouleurPrimaire, rounded corners=4pt,
      align=center, font=\footnotesize\bfseries\color{white}, inner xsep=8pt]
      at (3.2, -2.8)
      {Chaque exigence~: \OUI{} / \NON{} / \NEK{}, la constatation séparée de l'analyse};
\end{tikzpicture}
\end{center}

\vspace{0.6cm}

\begin{center}
    {\large\bfseries\color{CouleurPrimaire}
        \auteurcours}\\[0.5ex]
    {\normalsize\color{CouleurSecondaire}
       Enseignant Chercheur | \laboratoire{} | ENSPY}\\[0.2ex]
    {\small\color{CouleurBordInfo}
        \institutionprimaire}\\[0.2ex]
    {\small\color{CouleurBordInfo}
        \href{mailto:maletyon@proton.me}{maletyon@proton.me}}
\end{center}

\end{titlepage}

    \cleardoublepage

    % ============================================================================
% Fichier  : front/01_avant_propos.tex (Avant-propos)
% Version  : 1.0 -- novembre 2026
% ============================================================================

\chapter*{Avant-propos}
\addcontentsline{toc}{chapter}{Avant-propos}
\markboth{Avant-propos}{Avant-propos}

Le 23 juin 2026, le délai de dix-huit mois que la loi n°~2024/017 accordait
aux organisations camerounaises pour se mettre en conformité a expiré. Le
lendemain, une microfinance de Bafoussam reçoit la lettre d'un client dont
le dossier a fuité quinze mois plus tôt. Il demande ce que l'établissement
sait de lui, à qui il l'a transmis, et pourquoi personne ne l'a prévenu. Le
directeur général pose alors à ses équipes la question qui ouvre ce cours~:
sommes-nous en règle~?

Répondre exige un audit. Un audit de protection des données ne se contente
pas de lire des procédures~: il recense les traitements réels, les compare
à ce que la loi permet, vérifie que les personnes peuvent exercer leurs
droits, que les données sont protégées, que les sous-traitants et les
transferts sont encadrés, puis il conclut. Ce cours enseigne cette démarche,
dans le cadre de la loi camerounaise, avec les normes internationales qui
l'outillent.

Il s'adresse aux étudiants du Master Professionnel en Sécurité des Systèmes
Numériques de la Chaire UNESCO d'Afrique centrale, et aux praticiens qui
exercent ou exerceront comme auditeurs de la protection des données,
délégués à la protection des données, juristes conformité, RSSI ou
responsables de traitement. Il suppose acquis le droit du numérique (UE
MSN-2S3-JUR01) et la méthode générale de l'audit de cybersécurité (UE
MSN-2S3-SEC08). Il laisse au cours de forensique (UE MSN-2S3-SEC07)
l'enquête sur l'incident lui-même.

Trois partis pris le traversent. Le premier est de lire la loi dans son
texte officiel, article par article, sans lui prêter ce qu'elle ne dit pas~:
elle ne crée pas de délégué à la protection des données, elle ne fixe pas de
délai chiffré de notification, et de nombreuses dispositions attendent leurs
textes d'application. Le deuxième est de séparer, à chaque constat, le fait
établi de son interprétation, avec un niveau de confiance. Le troisième est
de coter ce que l'on ignore~: l'instrument \OUI{}~/ \NON{}~/ \NEK{} du
programme de recherche de l'auditologie sert ici à dire, pour chaque
exigence, si la conformité est établie, si l'écart est établi, ou si les
preuves manquent pour trancher. Une loi dont certains décrets ne sont pas
publiés produit des \NEK{} que l'auditeur doit assumer.

Les trois organisations qui accompagnent le lecteur, NYANGA CRÉDIT, le
CENTRE MÉDICAL MBOLO et NOUN PAY, sont fictives, comme leurs dirigeants,
leurs prestataires, les dates, les adresses et les montants. Les mécanismes,
les normes et les textes cités sont exacts. Les articles de loi sont cités
dans leur texte officiel~; les normes sous licence sont paraphrasées.

Le livre est accompagné d'un livret des corrigés et de NDAP SUKULU, où les
fiches de questions et les dix laboratoires du cours permettent de pratiquer
chaque chapitre.

\bigskip
\hfill\begin{minipage}{0.5\textwidth}
\raggedleft
\auteurcours\\
\laboratoire, Yaoundé, novembre 2026
\end{minipage}

    \cleardoublepage

    % ============================================================================
% Fichier  : front/02_mode_emploi.tex (Mode d'emploi)
% Version  : 1.0 -- novembre 2026
% ============================================================================
% 30 h de cours magistral = 15 séances de 2 h ; TD (fiches, exercices) et TP
% (douze labs) faits par l'étudiant chez lui, hors des 30 h.
% ============================================================================

\chapter*{Mode d'emploi du cours}
\addcontentsline{toc}{chapter}{Mode d'emploi du cours}
\markboth{Mode d'emploi}{Mode d'emploi}

\section*{Structure générale}

Le cours compte un prologue et sept parties, soit dix-neuf chapitres, pour
30~heures de cours magistral réparties en quinze séances de deux heures.
Les travaux dirigés et les travaux pratiques se font à la maison, hors de
ces 30~heures~: pour les travaux dirigés, la fiche de questions de chaque
chapitre et ses exercices de fin de chapitre (environ dix heures au
total)~; pour les travaux pratiques, les douze laboratoires de NDAP SUKULU,
d'une heure chacun. Les corrigés des exercices sont réunis dans un livret
séparé.

\begin{tableaucours}{cL{6.4cm}cL{3.2cm}}{%
    \tableauHeader{Partie} &
    \tableauHeader{Thème} &
    \tableauHeader{Chap.} &
    \tableauHeader{Niveau}
}
0 & Prologue~: le 24 juin 2026 & P & introductif \\
I & Fondements~: la donnée, la personne, le traitement & 1--2 & fondamental \\
II & Le contexte international & 3--5 & fondamental \\
III & Le contexte national & 6--7 & fondamental \\
IV & Méthode~: auditer la protection des données & 8--12 & fondamental à
intermédiaire \\
V & Pratique~: de l'écart à la conformité & 13--16 & intermédiaire \\
VI & Défis avancés & 17--18 & intermédiaire à expert \\
VII & Recherche~: mesurer la conformité, articuler les normes & 19 & recherche \\
\end{tableaucours}

\section*{Les quinze séances}

\begin{tableaucours}{cL{6.8cm}L{4.2cm}}{%
    \tableauHeader{Séance} &
    \tableauHeader{Contenu} &
    \tableauHeader{Travail à la maison}
}
1 & Prologue et Chap.~1~: données, traitements, acteurs & Fiches P et 1 \\
2 & Chap.~2~: les principes du traitement & Fiche 2, lab 01 \\
3 & Chap.~3 et~4~: sources internationales et africaines, RGPD & Fiches 3
et 4 \\
4 & Chap.~5~: NIST et normes ISO & Fiche 5, lab 12 \\
5 & Chap.~6 et~7~: loi n°~2024/017, loi n°~2010/012, choix des critères &
Fiches 6 et 7, lab 11 \\
6 & Chap.~8~: conduire un audit de protection des données & Fiche 8 \\
7 & Chap.~9~: cartographier les traitements & Fiche 9, lab 02 \\
8 & Chap.~10~: évaluer le risque pour les personnes & Fiche 10, lab 03 \\
9 & Chap.~11~: auditer les droits des personnes & Fiche 11, lab 04 \\
10 & Chap.~12~: auditer la sécurité des données & Fiche 12, lab 05 \\
11 & Chap.~13 et~14~: sous-traitants, transferts internationaux & Fiches 13
et 14, labs 06 et 07 \\
12 & Chap.~15~: violations de données & Fiche 15, lab 08 \\
13 & Chap.~16~: programme de conformité et rapport & Fiche 16, lab 09 \\
14 & Chap.~17~: profilage, décisions automatisées, IA & Fiche 17, lab 10 \\
15 & Chap.~18 et~19~: secteurs sensibles, mesurer et articuler & Fiches 18
et 19 \\
\end{tableaucours}

Cette répartition est indicative~; l'enseignant l'adapte au calendrier et au
rythme du groupe.

\section*{Le gabarit des chapitres}

Chaque chapitre suit le même plan~: une épigraphe, le niveau et les
prérequis, les objectifs, une situation d'ouverture tirée de l'un des trois
cas, des sections numérotées, ce qu'il faut retenir, des points de
contrôle, trois à quatre exercices gradués, des lectures, et une
transition.

\section*{Les encadrés}

\begin{tableaucours}{L{3.6cm}L{8.6cm}}{%
    \tableauHeader{Encadré} &
    \tableauHeader{Rôle}
}
Objectifs & Ce que l'on doit savoir faire à la fin du chapitre \\
Définition & Les notions, en gras dans le texte et indexées \\
Violation typique & Une erreur fréquente, et la règle qui l'évite \\
Information & Un complément, un renvoi, une nuance \\
Mission (NYANGA, MBOLO, NOUN PAY) & Le cas~: les constatations, puis
l'analyse, avec un niveau de confiance \\
Traitement & La fiche d'un traitement audité~: finalité, base légale,
données, destinataires, durée, ce que l'audit en a établi \\
Démarche d'audit & Les étapes d'une procédure d'audit, dans l'ordre \\
Lecture & Un décompte de cotes \OUI{}~/ \NON{}~/ \NEK{} et sa lecture \\
Métier & Ce que la notion change pour chacun des métiers du cours \\
Texte juridique & Un article cité dans son texte officiel, et sa portée \\
Ce qu'il faut retenir & Les faits et les règles du chapitre \\
Exercice & Trois niveaux de difficulté~: \niveaubase, \niveaumoyen,
\niveauexpert \\
\end{tableaucours}

\section*{Les conventions}

\begin{itemize}
    \item Les cotes~: \OUI{} (conformité établie par une preuve), \NON{}
          (écart établi par une preuve), \NEK{} (connaissance insuffisante
          pour trancher). Chaque exigence auditée reçoit l'une de ces trois
          cotes.
    \item Constatation et analyse sont toujours séparées~; l'analyse porte
          un niveau de confiance (faible, moyenne, élevée).
    \item « La loi » désigne la loi n°~2024/017 du 23~décembre~2024
          relative à la protection des données à caractère personnel au
          Cameroun~; ses articles sont cités sans autre précision.
    \item Les textes~: seuls les articles lus dans leur version officielle
          sont cités entre guillemets~; les normes sous licence sont
          paraphrasées~; le degré de vérification figure à
          l'annexe~\ref{ann:textes}.
    \item Les cas sont fictifs~: noms en majuscules, montants en francs
          CFA, villes réelles. Les données personnelles des exemples sont
          inventées.
\end{itemize}

    \cleardoublepage

    % ==============================================================================
% Fichier : front/03_ndap_sukulu.tex
% Titre   : NDAP SUKULU, l'environnement compagnon du cours
% Auteur  : MINKA MI NGUIDJOÏ Thierry Emmanuel
% Version : 1.0 -- novembre 2026
% ==============================================================================
% Texte commun aux cours de l'auteur, mis à jour en novembre 2026 (version en
% ligne, banques de 100 questions, compte de démonstration) ; le tableau des
% labs est propre au cours (front/03b_labs.tex).
% ==============================================================================

\chapter*{NDAP SUKULU, la salle de classe compagnon}
\addcontentsline{toc}{chapter}{NDAP SUKULU, la salle de classe compagnon}
\markboth{NDAP SUKULU}{NDAP SUKULU}

\begin{center}
    \includegraphics[width=0.62\textwidth]{figures/logos/logo_ndap_sukulu.png}
\end{center}

\section*{Un nom venu du pays bassa}

\begin{quote}\itshape
Le soir tombe sur Éséka. Un enfant revient des champs avec son grand-père et
lui demande pourquoi il doit marcher si loin, chaque matin, pour aller à
l'école. Le vieil homme s'arrête au bord du chemin et répond, en bassa~:
« \textup{Ndap sukulu}, la salle de classe, ce ne sont pas les murs. Ce sont
les murs qui viennent après. La salle de classe, c'est l'endroit où quelqu'un
t'attend pour t'apprendre ce qu'il sait. » L'enfant insiste~: « Et qui
m'attend~? » Le grand-père sourit~: « \textup{Malet yon}. Ton enseignant. »
\end{quote}

En bassa, la langue paternelle de l'auteur, \textit{ndap sukulu} désigne la
salle de classe et \textit{malet yon} signifie « ton enseignant ».
L'environnement compagnon de ce cours porte ce nom parce qu'il prolonge la
salle de classe au-delà des heures de cours~: il attend l'étudiant le soir,
le week-end, à distance, avec des questions et des pièces à examiner. Il
est édité par MaletYon et mis à la disposition des étudiants et des
praticiens qui souhaitent se perfectionner
(\href{mailto:maletyon@proton.me}{maletyon@proton.me}).

\section*{Ce que l'on y trouve}

\begin{itemize}
    \item Une fiche de questions par chapitre. Chaque tentative tire vingt
          questions d'une banque de cent (QCM, vrai ou faux, textes à
          trous) et dure trente minutes.
    \item Dix laboratoires. Chaque lab place l'étudiant devant les pièces
          d'une mission d'audit du cours (registres, exports, journaux
          d'accès, contrats, tableaux de flux) et se termine par des points
          de contrôle.
    \item Un espace enseignant, qui corrige les remises, établit la
          synthèse de la classe et publie les résultats.
\end{itemize}

Les fiches et les labs sont des travaux à faire hors des heures de cours.
Ils ne remplacent ni le livre ni les exercices de fin de chapitre~: ils en
vérifient la maîtrise.

\section*{Comment y accéder}

\begin{enumerate}
    \item En ligne, à l'adresse
          \url{https://maletyon.github.io/cours-apdc/ndap-sukulu/}, dans un
          navigateur récent~; sur téléphone, ajouter \texttt{telephone/} à
          la fin de l'adresse.
    \item Hors ligne, en téléchargeant l'archive de l'édition ordinateur
          depuis le dépôt \texttt{MaletYon/cours-apdc}, dossier
          \texttt{ndap-sukulu}.
    \item Se connecter avec le matricule et le mot de passe remis par
          l'enseignant. Pour essayer avant d'être inscrit, un compte de
          démonstration est ouvert à tous~: matricule \texttt{ETU001}, mot
          de passe \texttt{eTU001}~; ses essais ne sont jamais transmis à
          l'enseignant.
    \item Pour un lab, vérifier d'abord l'empreinte SHA-256 des pièces
          (\texttt{sha256sum -c MANIFESTE\_SHA256.txt}), puis travailler sur
          une copie. Les données personnelles des pièces sont inventées.
\end{enumerate}

\begin{boxinfo}
Les fiches de questions sont aussi fournies au format texte dans le dossier
\texttt{fiches} du dépôt du cours, prêtes à être distribuées. Leurs
corrections figurent dans le livret de correction, à la suite des corrigés
des exercices de chaque chapitre.
\end{boxinfo}

\section*{Les laboratoires du cours}

% Tableau des labs du cours APDC dans NDAP SUKULU (propre au cours).
Le cours compte douze laboratoires, disponibles sous le code \codeUE. Chacun
place l'étudiant dans l'une des trois missions d'audit du cours, devant les
pièces que l'auditeur réunit~: exports, registres, journaux, contrats.

\begin{tableaucours}{lL{6.2cm}L{1.6cm}L{2.2cm}}{%
    \tableauHeader{Lab} &
    \tableauHeader{Sujet} &
    \tableauHeader{Chapitre} &
    \tableauHeader{Mission}
}
APDC-LAB-01 & Qualifier les données d'un export & 1, 2 & NYANGA \\
APDC-LAB-02 & Registre déclaré, fichiers découverts & 9 & NYANGA \\
APDC-LAB-03 & Évaluation d'impact de la téléconsultation & 10 & MBOLO \\
APDC-LAB-04 & Demandes des personnes et délais & 11 & NYANGA \\
APDC-LAB-05 & Journaux d'accès au dossier patient & 12 & MBOLO \\
APDC-LAB-06 & Contrats de sous-traitance & 13 & NYANGA, MBOLO \\
APDC-LAB-07 & Flux de données vers l'étranger & 14 & NOUN PAY \\
APDC-LAB-08 & Violation de données et notification & 15 & NYANGA \\
APDC-LAB-09 & Durées de conservation & 16 & NYANGA \\
APDC-LAB-10 & Audit intégrateur du profilage & 17, 19 & NOUN PAY \\
APDC-LAB-11 & La loi de 2010 à NYANGA & 7 & NYANGA \\
APDC-LAB-12 & Double cote et profil NIST & 4, 5 & NOUN PAY \\
\end{tableaucours}
\vspace{-0.8em}
\begin{center}{\small\itshape Les laboratoires du cours dans NDAP SUKULU}\end{center}


    \cleardoublepage

    \tableofcontents
    \cleardoublepage

    \listoffigures
    \listoftables
    \cleardoublepage

% CORPS (pages arabes). Numérotation continue : le prologue est le chapitre 0,
% puis C01 à C16, un numéro par fichier.
% --------------------------------------------------------------------------
\mainmatter

    % PARTIE 0 · PROLOGUE OPÉRATIONNEL
    \part*{Prologue Opérationnel}
    \addcontentsline{toc}{part}{Prologue Opérationnel}

        \setcounter{chapter}{-1}
        % ============================================================================
% Fichier  : partie0/P00_vingt_quatre_juin.tex
% Titre    : Le 24 Juin 2026 (prologue opérationnel)
% Auteur   : MINKA MI NGUIDJOÏ Thierry Emmanuel
% Version  : 1.0 -- novembre 2026
% Corrigés : corriges/Corr_P00.tex
% ============================================================================
% Sources : loi n° 2024/017, art. 2, 5, 19, 22, 29, 37-39, 62, 73 (texte
% officiel, scan certifié conforme) ; FAV pour la fuite du 14/03/2025.
% Faits posés (bible) : lettre de M. TAGNE du 24/06/2026 ; demande de l'OPJ du
% 20/03/2025 de différer l'information des clients ; mandat d'audit du
% 30/06/2026 à KORO AUDIT & CONFORMITÉ.
% ============================================================================

\chapter{Le 24 Juin 2026}
\label{chap:prologue}
\chaptermark{Le 24 juin 2026}

\epigraph{%
    « Être en règle, c'est pouvoir le montrer.\\
    Le dire ne suffit pas. »%
}{Principe du cours}

\niveauChapitre{Introductif}{%
    Aucun, sinon la lecture de l'avant-propos
}

\begin{boxobjectifs}
\begin{itemize}
    \item Reconnaître, dans une situation réelle, les questions qui relèvent
          de la protection des données à caractère personnel.
    \item Situer une organisation par rapport à l'échéance du 23~juin~2026
          fixée par l'article~73 de la loi n°~2024/017.
    \item Distinguer une déclaration de conformité d'une conformité
          démontrée.
    \item Coter une exigence \OUI{}, \NON{} ou \NEK{} et justifier la cote.
    \item Situer les métiers qui interviennent, de la réclamation au
          rapport d'audit.
\end{itemize}
\end{boxobjectifs}

\bigskip

% ============================================================================
\section*{Avant de commencer : une lettre, un mercredi matin}
% ============================================================================

Bafoussam, mercredi 24~juin~2026, 8~h~15. Le courrier de NYANGA
CRÉDIT\footnote{NYANGA CRÉDIT, ses dirigeants, ses clients, ses prestataires
et le cabinet KORO AUDIT \& CONFORMITÉ sont fictifs, comme les dates, les
adresses et les montants des cas de ce cours. Les mécanismes et les textes
cités sont exacts.}, établissement de microfinance de vingt-deux agences et
de 210\,000 clients, contient une lettre recommandée de M.~TAGNE Paul,
commerçant au marché A, client depuis 2016.

M.~TAGNE a assisté, en février, à l'audience du tribunal de grande instance
de la Mifi qui jugeait l'affaire de la nuit du 14~mars~2025. Il y a appris
que son dossier figurait parmi les 18\,412 fichiers du dossier
\texttt{Clients} copiés cette nuit-là vers l'étranger. Personne ne l'en
avait informé. Sa lettre demande trois choses~: savoir quelles données
l'établissement détient sur lui et à qui il les a communiquées, obtenir
l'effacement de celles qui ne sont plus nécessaires, et comprendre pourquoi
il n'a pas été prévenu. Elle annonce qu'à défaut de réponse il saisira le
juge des référés.

À 10~heures, le directeur général, M.~TCHOUA, réunit la RSSI, Mme~NDOUMBE,
le DSI, M.~KAMGA, et la juriste, Me~ETOUNDI. Il pose une seule question~:
« Sommes-nous en règle~? » Me~ETOUNDI rappelle que la veille, le 23~juin,
a expiré le délai de dix-huit mois que la loi accordait aux organisations
pour se mettre en conformité. M.~KAMGA assure que « tout est sécurisé
depuis l'incident ». Mme~NDOUMBE répond qu'elle ne sait pas, et qu'elle ne
peut pas le savoir sans un examen.

\begin{boxalert}
\textbf{Violation typique.} Répondre « nous sommes conformes » sur la foi
d'une politique écrite, d'un pare-feu renforcé ou d'une absence de
plainte.

\smallskip
\textbf{Règle~:} la conformité se démontre traitement par traitement, par
des preuves~: un registre, des autorisations, des contrats, des journaux,
des réponses aux personnes. Ce qui n'a pas été examiné reste inconnu, et se
dit inconnu.
\end{boxalert}

% ============================================================================
\section{Ce que l'on sait le 24 juin}
\label{sec:p00-situation}
% ============================================================================

La réunion permet de réunir quelques faits. Ce sont des
\termecle{constatations}~: chacune se rattache à une pièce que quelqu'un
d'autre pourrait vérifier.

\begin{itemize}
    \item Le 14~mars~2025, 3,2~Go du dossier \texttt{Clients} ont été copiés
          vers une adresse étrangère~; l'enquête technique l'a établi, et le
          jugement de février~2026 le reprend.
    \item Le 20~mars~2025, l'officier de police judiciaire chargé de
          l'enquête a demandé par écrit à l'établissement de ne pas informer
          les clients, pour ne pas compromettre l'enquête. Aucun client n'a
          été informé depuis.
    \item L'établissement n'a jamais adressé de demande d'autorisation de
          traitement à l'Autorité prévue par la loi.
    \item Me~ETOUNDI n'a trouvé aucune liste des traitements de
          l'établissement. Le DSI tient un inventaire des applications, qui
          n'est pas la même chose.
\end{itemize}

Les autres phrases entendues pendant la réunion sont des opinions~: « tout
est sécurisé », « les clients ont signé les conditions générales », « le
prestataire de recouvrement est sérieux ». Elles ne sont ni vraies ni
fausses à ce stade~; elles attendent leur preuve.

\begin{boxdef}
\termecle{Donnée à caractère personnel}~: selon l'article~5 de la loi,
« information se rapportant à une personne permettant de l'identifier
directement ou indirectement, notamment par référence à toute forme
d'identifiant ou à un ou plusieurs éléments, propres à son identité
physique, psychologique, génétique, psychique, culturelle,
socio-professionnelle ou économique ». La loi cite parmi les exemples un
nom, une photo, une adresse, un numéro de téléphone, un matricule interne,
une adresse IP et un enregistrement vocal.

\smallskip
\termecle{Traitement}~: toute opération appliquée à de telles données,
« notamment la collecte, l'enregistrement, l'organisation, la conservation
[\ldots], la consultation, l'utilisation, la communication [\ldots], ainsi
que le verrouillage, l'effacement ou la destruction » (article~5).
\end{boxdef}

% ============================================================================
\section{L'échéance du 23 juin}
\label{sec:p00-echeance}
% ============================================================================

La loi n°~2024/017 a été promulguée le 23~décembre~2024~\cite{loi2024017}.
Elle a laissé aux organisations un délai pour adapter leurs pratiques.

\begin{boxjuris}
\textbf{\loicam{n°~2024/017}, article~73}~\cite{loi2024017}~:
\textit{« Les personnes physiques ou morales en charge du traitement des
données à caractère personnel disposent d'un délai de dix-huit (18) mois, à
compter de la date de promulgation de la présente loi, pour se conformer aux
dispositions de celle-ci. »}

\smallskip
\textbf{Portée.} Le délai court du 23~décembre~2024 au 23~juin~2026. À
partir du 24~juin, une organisation qui traite des données ne peut plus
invoquer la période d'adaptation. L'article ne dit pas si les traitements
commencés avant la loi doivent être régularisés rétroactivement~; la
question est discutée au Chapitre~\ref{chap:recherche}.
\end{boxjuris}

La loi s'applique largement. Son article~2 vise tout traitement effectué
par l'État, les collectivités territoriales décentralisées ou toute
personne physique ou morale, tout traitement de données d'une personne
établie, résidant ou en transit au Cameroun, et tout traitement effectué
par un responsable ou un sous-traitant établi au Cameroun. NYANGA CRÉDIT
entre dans chacun de ces cas.

\begin{boxinfo}
\textbf{Ce que la loi ne dit pas.} Elle ne crée pas de délégué à la
protection des données~; elle ne fixe pas de délai chiffré pour informer
d'une violation (l'article~22 dit « sans délai »)~; et de nombreuses
modalités sont renvoyées à des textes d'application. Elle confie enfin
la plus grande part du contrôle à une Autorité de protection des données à
caractère personnel dont la création et le fonctionnement relèvent d'un
décret (article~53, alinéa~2). Le cours lit la loi telle qu'elle est écrite,
et cote \NEK{} ce que ses textes d'application ne permettent pas encore de
trancher.
\end{boxinfo}

% ============================================================================
\section{Trois questions pour un audit}
\label{sec:p00-questions}
% ============================================================================

La question du directeur général, « sommes-nous en règle », est trop large
pour recevoir une réponse. Mme~NDOUMBE la décompose en trois questions,
chacune vérifiable.

\begin{enumerate}
    \item \textbf{Que traitons-nous~?} Quelles données, sur quelles
          personnes, pour quelles finalités, avec quels destinataires et
          pendant combien de temps. C'est l'objet du registre des activités
          de traitement que l'article~29 impose de tenir
          (Chapitre~\ref{chap:cartographie}).
    \item \textbf{Ce que nous faisons est-il permis~?} Le consentement
          ou l'une des trois dérogations de l'article~9, l'autorisation de
          l'article~19, l'information de l'article~21, la durée de
          conservation de l'article~13 (Chapitres~\ref{chap:donnees}
          et~\ref{chap:principes}).
    \item \textbf{Pouvons-nous le prouver~?} Pour chaque exigence, une
          pièce qui l'établit, ou l'aveu qu'elle manque
          (Chapitre~\ref{chap:audit}).
\end{enumerate}

Le 30~juin, le conseil d'administration confie la mission au cabinet KORO
AUDIT \& CONFORMITÉ, de Yaoundé. Son auditrice principale, Mme~BELLO
Aïssatou, demande deux mois de terrain. Elle refuse de commencer par la
sécurité informatique, que M.~KAMGA voulait lui montrer en premier~: « Nous
commencerons par savoir ce que vous traitez. Protéger des données que l'on
n'a pas recensées, c'est protéger au hasard. »

% ============================================================================
\section{Répondre à M. TAGNE}
\label{sec:p00-reponse}
% ============================================================================

La lettre n'attend pas l'audit. Me~ETOUNDI prépare la réponse avec trois
textes en main.

\begin{boxjuris}
\textbf{\loicam{n°~2024/017}, article~39}~\cite{loi2024017}~:
\textit{« La personne concernée peut demander au responsable du
traitement~: les informations permettant de connaitre et de contester le
traitement~; la confirmation que les données à caractère personnel la
concernant font ou ne font pas l'objet de traitement~; la communication des
données à caractère personnel qui la concerne ainsi que toute information
disponible quant à l'origine de celle-ci~; [\ldots] une copie de ses données
à caractère personnel sous une forme intelligible contre paiement d'une
somme ne pouvant excéder le coût de la reproduction. »}

\smallskip
\textbf{Portée.} M.~TAGNE a droit à la liste de ses données, à leur origine
et à leurs destinataires. Le délai de réponse doit être fixé par voie
réglementaire (article~47)~; en attendant, l'auditrice recommande de
répondre dans un délai raisonnable et de le dater.
\end{boxjuris}

La demande d'effacement est plus délicate. L'article~38 oblige à effacer
« sans délai » les données qui ne sont plus nécessaires, mais « sous réserve
des exigences de conservation » prévues par d'autres textes. Or les données
de M.~TAGNE sont aussi des pièces d'une procédure judiciaire et de la
comptabilité de l'établissement. Me~ETOUNDI ne peut ni tout effacer, ni tout
garder~: elle doit séparer, donnée par donnée, ce qui reste nécessaire et ce
qui ne l'est plus. Cette séparation suppose de savoir où se trouvent les
données de M.~TAGNE, ce que personne ne sait encore.

Reste la troisième question~: pourquoi n'a-t-il pas été prévenu~? L'article~22
impose d'informer « sans délai » l'Autorité et la personne concernée dès
connaissance d'une violation. En mars~2025, l'officier de police judiciaire
a demandé le silence pour protéger l'enquête. Deux obligations se sont
rencontrées~; laquelle l'emportait, et à partir de quand la loi
s'appliquait-elle à cette violation antérieure à l'échéance~? Le
Chapitre~\ref{chap:violations} examine la question pour la pratique, et le
Chapitre~\ref{chap:recherche} la présente comme un objet de recherche en
cours.

% ============================================================================
\section{Premières cotes}
\label{sec:p00-cotes}
% ============================================================================

Mme~BELLO ouvre son premier papier de travail le 6~juillet, premier jour de
terrain. Elle cote huit propositions avec ce qu'elle a vu.

\begin{boxdef}
\termecle{Cote \OUI{}, \NON{}, \NEK{}}~: réponse à une question d'audit
fermée sur une exigence. \OUI{}~: la conformité est établie par une preuve.
\NON{}~: l'écart est établi par une preuve. \NEK{} (\textit{not enough
knowledge})~: les preuves disponibles ne permettent pas de trancher.
L'instrument vient du programme de recherche de
l'auditologie~\cite{minka2026nek}.
\end{boxdef}

\begin{boxlecture}{Les huit premières propositions, 6 juillet 2026}
\begin{tableaucours}{L{7.4cm}cL{3.6cm}}{%
    \tableauHeader{Proposition} &
    \tableauHeader{Cote} &
    \tableauHeader{Pièce}
}
NYANGA CRÉDIT traite des données à caractère personnel & \OUI{} & Fichiers
clients, paie \\
NYANGA CRÉDIT est soumise à la loi n°~2024/017 & \OUI{} & Art.~2 \\
Un registre des activités de traitement existe & \NON{} & Aucun produit,
attestation de la juriste \\
Une autorisation de l'Autorité a été demandée & \NON{} & Attestation de la
direction \\
Les clients ont consenti aux traitements & \NEK{} & Formulaires non encore
examinés \\
Les données des clients sont protégées & \NEK{} & Aucune mesure encore
testée \\
Les prestataires sont liés par un contrat conforme & \NEK{} & Contrats non
encore reçus \\
Des données sont transférées hors du Cameroun & \NEK{} & Hébergement du
logiciel bancaire non encore examiné \\
\end{tableaucours}

\medskip
\textbf{Décompte.} Deux \OUI{}, deux \NON{}, quatre \NEK{} sur huit.
\hfill\lecture{0{,}25}{0{,}50}{1{,}50}

\medskip
\textbf{Lecture.} La moitié des propositions est inconnue. Ce n'est pas un
mauvais résultat~: c'est l'état normal d'un audit qui commence. Les deux
\NON{} sont, eux, établis dès le premier jour, et ils suffisent à répondre au
directeur général que l'établissement n'est pas en règle sur ces deux
points.
\end{boxlecture}

Les trois mesures se lisent ensemble. $\mu$ est la part des \OUI{}, $U$
celle des \NEK{}, $H$ l'entropie de la répartition, en bits. Le
Chapitre~\ref{chap:audit} en donne la règle de lecture, et le
Chapitre~\ref{chap:recherche} en discute les limites.

% ============================================================================
\section{Deux autres organisations}
\label{sec:p00-autres}
% ============================================================================

Deux autres organisations accompagnent NYANGA CRÉDIT dans ce cours.

Le \textbf{CENTRE MÉDICAL MBOLO}, à Bertoua, établissement privé de
120~lits, traite des données de santé, qui comptent parmi les données
sensibles de l'article~5. Son dossier patient informatisé est édité et
hébergé par un prestataire de Douala~; il prépare un service de
téléconsultation pour les villages de la Kadey, et il a découvert en mai
qu'un dossier de patiente avait été consulté par des agents sans lien avec
ses soins. Sa directrice, le Dr~ABENA, a commandé un audit de préparation
avant l'échéance.

\textbf{NOUN PAY}, à Bafoussam, fintech de paiement mobile et de
microcrédit, héberge tout son système d'information chez un fournisseur de
nuage à l'étranger et décide de l'octroi de ses crédits par une notation
automatique. Un client conteste un refus. Les transferts de données hors du
Cameroun et les décisions automatisées sont au centre de son audit.

% ============================================================================
\section{Les métiers de l'affaire}
\label{sec:p00-metiers}
% ============================================================================

\begin{boxmetier}{Auditeur de la protection des données}
Mme~BELLO examine, prouve et conclut. Elle ne corrige pas elle-même~: elle
dit ce qui est établi, ce qui manque, et ce qu'il faut faire. Son
indépendance à l'égard de la direction qu'elle audite conditionne la valeur
de ses conclusions.
\end{boxmetier}

\begin{boxmetier}{Délégué à la protection des données}
La loi n'impose pas cette fonction. Beaucoup d'organisations la créent
pourtant, sur le modèle européen, pour qu'une personne tienne le registre,
conseille les équipes et réponde aux personnes. NYANGA CRÉDIT n'en a pas~:
la lettre de M.~TAGNE est arrivée sur le bureau du directeur général.
\end{boxmetier}

\begin{boxmetier}{Juriste conformité}
Me~ETOUNDI lit les textes, répond aux personnes et mesure le risque de
contentieux. Elle doit savoir dire où la loi est claire, où elle renvoie à
des textes non publiés, et où deux obligations se contredisent.
\end{boxmetier}

\begin{boxmetier}{RSSI}
Mme~NDOUMBE tient la sécurité des données, exigée par les articles~22, 24
et~27. Elle sait aussi que la sécurité n'est qu'une partie de la
conformité~: des données bien protégées peuvent être collectées sans droit.
\end{boxmetier}

\begin{boxmetier}{Responsable de traitement}
M.~TCHOUA, pour l'établissement, détermine les finalités et les moyens des
traitements. C'est lui que la loi oblige, et lui qui répond des sanctions
des articles~54 à~71.
\end{boxmetier}

% ============================================================================
\section{Cinq réflexes}
\label{sec:p00-reflexes}
% ============================================================================

\begin{enumerate}
    \item \textbf{Lire le texte officiel.} On cite la loi telle qu'elle est
          écrite, pas telle qu'un résumé ou le modèle européen la décrit.
    \item \textbf{Recenser avant de juger.} On ne peut pas dire si un
          traitement est permis tant qu'on ne sait pas qu'il existe.
    \item \textbf{Demander la preuve.} Une déclaration n'établit rien~; une
          pièce établit quelque chose, et l'on dit quoi.
    \item \textbf{Coter l'inconnu.} Ce qui n'a pas été examiné est \NEK{},
          jamais conforme par défaut.
    \item \textbf{Protéger les données de l'audit.} L'auditeur manipule
          lui-même des données personnelles~: il en prend le moins
          possible et les protège comme il l'exige des autres.
\end{enumerate}

\begin{boxresume}
\begin{itemize}
    \item Depuis le 24~juin~2026, le délai de mise en conformité de
          l'article~73 est expiré~: la loi n°~2024/017 s'applique
          pleinement.
    \item La conformité se démontre par des preuves, traitement par
          traitement~; une déclaration ne la démontre pas.
    \item La question « sommes-nous en règle » se décompose en trois~: que
          traitons-nous, est-ce permis, pouvons-nous le prouver.
    \item À NYANGA CRÉDIT, l'absence de registre et de demande
          d'autorisation est établie dès le premier jour~; la moitié des
          propositions reste \NEK{}.
    \item Une violation antérieure à l'échéance, couverte par une demande de
          silence de l'enquête pénale, pose une question d'articulation des
          normes que le cours traite à part.
\end{itemize}
\end{boxresume}

\begin{boxinfo}
\textbf{Points de contrôle.} Vous avez compris ce prologue si vous savez~:
\begin{itemize}
    \item dater l'échéance de l'article~73 et dire ce qu'elle change~;
    \item séparer, dans un compte rendu de réunion, les constatations des
          opinions~;
    \item coter une exigence \OUI{}, \NON{} ou \NEK{} et citer la pièce qui
          fonde la cote.
\end{itemize}
\end{boxinfo}

\begin{boxexo}[title={\ding{45}\quad Exercice~0.1. Donnée personnelle ou non \niveaubase}]
Dites si chaque élément est une donnée à caractère personnel au sens de
l'article~5, et pourquoi~: (a) le numéro de compte de M.~TAGNE~; (b) le
nombre total de clients de l'agence de Dschang~; (c) l'adresse IP d'où un
client se connecte à l'application mobile~; (d) l'enregistrement d'un appel
au service client~; (e) le taux d'intérêt moyen des crédits de 2025.
\end{boxexo}

\begin{boxexo}[title={\ding{45}\quad Exercice~0.2. Constatation ou opinion \niveaumoyen}]
Classez chaque phrase de la réunion du 24~juin comme constatation ou
opinion, et proposez une cote \OUI{}, \NON{} ou \NEK{} avec la pièce qui la
fonderait~: (a) « tout est sécurisé depuis l'incident »~; (b) « les clients
ont signé les conditions générales »~; (c) « aucune demande d'autorisation
n'a été envoyée »~; (d) « le prestataire de recouvrement est sérieux ».
\end{boxexo}

\begin{boxexo}[title={\ding{45}\quad Exercice~0.3. Répondre à M. TAGNE \niveaumoyen}]
\begin{enumerate}[(a)]
    \item Rédigez la réponse à la demande d'accès de M.~TAGNE, en vous
          appuyant sur l'article~39.
    \item Quelles données pouvez-vous effacer, et lesquelles devez-vous
          garder~? Justifiez avec l'article~38.
    \item Que répondez-vous à sa troisième question sans engager
          imprudemment l'établissement~?
\end{enumerate}
\end{boxexo}

\begin{boxexo}[title={\ding{45}\quad Exercice~0.4. Le silence demandé \niveauexpert}]
Le 20~mars~2025, l'officier de police judiciaire demande à NYANGA CRÉDIT de
ne pas informer ses clients.
\begin{enumerate}[(a)]
    \item Quelles obligations entrent en concurrence~?
    \item La loi n°~2024/017 s'appliquait-elle déjà à cette violation~?
          Distinguez la promulgation et l'échéance de l'article~73.
    \item Quelle trace écrite l'établissement aurait-il dû garder pour se
          défendre aujourd'hui~?
\end{enumerate}
\end{boxexo}

\bigskip

\begin{boxinfo}
\textbf{Pour aller plus loin.}
\begin{itemize}
    \item La loi n°~2024/017~\cite{loi2024017}, dans son texte officiel.
    \item D.~J. Solove, \textit{Understanding Privacy}~\cite{solove2008},
          pour la diversité de ce que l'on appelle la vie privée.
\end{itemize}
\end{boxinfo}

\begin{boxinfo}
\textbf{Transition.} Avant de recenser les traitements de NYANGA CRÉDIT, il
faut savoir ce qu'est exactement une donnée, un traitement, un responsable
et un sous-traitant, et quand une donnée devient sensible. C'est l'objet du
Chapitre~\ref{chap:donnees}.
\end{boxinfo}


    % PARTIE 1 · FONDEMENTS (loi n° 2024/017, art. 1 à 18 ; ISO/IEC 29100)
    \part{Fondements : la Donnée, la Personne, le Traitement}
    \label{part:fondements}

        % ============================================================================
% Fichier  : partie1/P01_C01_donnees_traitements.tex
% Titre    : Données, Traitements, Acteurs
% Auteur   : MINKA MI NGUIDJOÏ Thierry Emmanuel
% Version  : 1.0 -- novembre 2026
% Corrigés : corriges/Corr_C01.tex
% ============================================================================
% Sources : loi n° 2024/017, art. 1 à 5, 20, 31, 48, 72 (texte officiel) ;
% ISO/IEC 29100 (cadre de protection de la vie privée, paraphrasé) ;
% Solove (2008) ; Nissenbaum (2010).
% Faits nouveaux (bible) : inventaire des familles de données de NYANGA,
% 7 juillet 2026.
% ============================================================================

\chapter{Données, Traitements, Acteurs}
\label{chap:donnees}
\chaptermark{Données, traitements, acteurs}

\epigraph{%
    « On ne protège bien que ce que l'on sait nommer. »%
}{Principe du cours}

\niveauChapitre{Fondamental}{%
    Lecture du Chapitre~\ref{chap:prologue}
}

\begin{boxobjectifs}
\begin{itemize}
    \item Reconnaître une donnée à caractère personnel, directe ou
          indirecte, selon l'article~5 de la loi n°~2024/017.
    \item Distinguer anonymisation et pseudonymisation, et dire ce que
          chacune change pour l'application de la loi.
    \item Identifier les données sensibles, dont la liste camerounaise est
          plus large que la liste européenne.
    \item Qualifier les acteurs d'un traitement~: personne concernée,
          responsable, sous-traitant, responsables conjoints, destinataire.
    \item Délimiter le champ d'application de la loi et ses exclusions.
\end{itemize}
\end{boxobjectifs}

\bigskip

% ============================================================================
\section*{Avant de commencer : « nos données personnelles, ce sont les clients »}
% ============================================================================

Bafoussam, mardi 7~juillet~2026. Pour sa deuxième journée de terrain,
Mme~BELLO demande à chaque directeur de NYANGA CRÉDIT de lister les données
personnelles que traite son service. Le directeur commercial répond
aussitôt~: « Nos données personnelles, ce sont les clients. » La directrice
des ressources humaines n'a rien noté~; elle pensait que la question ne la
concernait pas.

En fin de journée, l'auditrice a recensé onze familles de données, dont
huit ne figuraient sur aucune liste~: les cautions des emprunteurs, les
prospects démarchés par les agents, les candidats à l'embauche, les
salariés, les images des caméras des agences, les enregistrements du
service client, les journaux de connexion de l'application mobile et les
listes de débiteurs remises au prestataire de recouvrement.

\begin{boxalert}
\textbf{Violation typique.} Réduire les données personnelles aux fichiers
clients, et oublier les salariés, les tiers, les images, les voix et les
journaux techniques.

\smallskip
\textbf{Règle~:} une donnée est personnelle dès qu'elle permet d'identifier
une personne, directement ou indirectement, quel que soit son support et
quelle que soit la personne. Le recensement part des personnes concernées,
pas des applications.
\end{boxalert}

% ============================================================================
\section{La Donnée à Caractère Personnel}
\label{sec:c01-donnee}
% ============================================================================

L'article~5 de la loi définit la donnée à caractère personnel par sa
relation à une personne identifiable, et non par sa nature.

\begin{boxjuris}
\textbf{\loicam{n°~2024/017}, article~5}~\cite{loi2024017}~:
\textit{« Donnée à caractère personnel~: information se rapportant à une
personne permettant de l'identifier directement ou indirectement, notamment
par référence à toute forme d'identifiant ou à un ou plusieurs éléments,
propres à son identité physique, psychologique, génétique, psychique,
culturelle, socio-professionnelle ou économique, notamment un nom, une
photo, une empreinte, une adresse postale, une adresse mail, un numéro de
téléphone, un numéro de sécurité sociale, un matricule interne, un
identifiant numérique, une adresse IP, un identifiant de connexion
informatique, un enregistrement vocal. »}

\smallskip
\textbf{Portée.} La liste des exemples est longue et précise~: une adresse
IP, un matricule interne et un enregistrement vocal sont des données
personnelles par la lettre même de la loi. L'auditeur n'a pas à le
démontrer.
\end{boxjuris}

\subsection{Identifier directement ou indirectement}

Une donnée identifie \termecle{directement} quand elle désigne la personne à
elle seule~: un nom complet, une photographie, une empreinte. Elle identifie
\termecle{indirectement} quand il faut la rapprocher d'autres informations~:
un numéro de compte, rapproché du fichier des comptes~; une adresse IP,
rapprochée du journal de l'opérateur~; la combinaison « femme, 34~ans,
pharmacienne, quartier Tamdja » dans une petite ville. La question n'est pas
de savoir si l'organisation fait ce rapprochement, mais s'il est
raisonnablement possible.

\subsection{Anonymiser n'est pas pseudonymiser}

\begin{boxdef}
\termecle{Anonymisation}~: selon l'article~5, « processus de modification ou
de suppression d'informations à caractère personnel identifiables dans un
ensemble de données, afin de rendre ces données non attribuables à des
individus spécifiques ».

\smallskip
\termecle{Pseudonymisation}~: « technique utilisée pour remplacer des
données d'identification directe par des identifiants artificiels dans le
but de protéger la vie privée des individus » (article~5).
\end{boxdef}

La différence est décisive pour l'auditeur. Des données réellement
anonymisées ne se rapportent plus à une personne identifiable~: elles
sortent du champ de la loi. Des données pseudonymisées restent
personnelles, puisque celui qui détient la table de correspondance peut
rétablir le lien. Un fichier où le nom du client est remplacé par son
numéro de compte n'est pas anonyme~: il est pseudonymisé, et le
numéro de compte est lui-même un identifiant.

\begin{boxinfo}
\textbf{L'anonymisation se prouve.} Une organisation qui déclare « nos
statistiques sont anonymes » doit montrer qu'aucun rapprochement raisonnable
ne permet de retrouver une personne. Supprimer les noms ne suffit pas quand
il reste la date de naissance, le quartier et la profession. Tant que la
preuve manque, l'auditeur cote la proposition « ces données sont anonymes »
\NEK{}, et traite les données comme personnelles.
\end{boxinfo}

% ============================================================================
\section{Les Données Sensibles}
\label{sec:c01-sensibles}
% ============================================================================

Certaines données exposent davantage la personne~: leur divulgation peut
conduire à une discrimination, une persécution ou une atteinte grave à la
réputation. La loi les nomme et les encadre plus strictement.

\begin{boxjuris}
\textbf{\loicam{n°~2024/017}, article~5}~\cite{loi2024017}~:
\textit{« Donnée sensible~: information relative notamment aux opinions et
activités religieuses, philosophiques, politiques, syndicales, aux
transactions bancaires, à l'origine raciale ou ethnique, linguistique ou
régionale, à la vie sexuelle, à la génétique, à la biométrie, à la santé,
aux poursuites judiciaires et aux sanctions pénales. »}

\smallskip
\textbf{Portée.} Le mot « notamment » rend la liste ouverte. Elle comprend
trois catégories que le droit européen ne range pas parmi les données
sensibles~: les transactions bancaires, l'origine linguistique ou régionale,
et les poursuites judiciaires.
\end{boxjuris}

La présence des \textbf{transactions bancaires} dans la liste change tout
pour un établissement financier~: la plupart des données de NYANGA CRÉDIT
et de NOUN PAY est sensible. L'article~48, alinéa~2, en tire la
conséquence~: le traitement des données dans le cadre des transactions
bancaires est interdit « sans autorisation préalable des administrations et
structures compétentes ». L'origine \textbf{linguistique ou régionale}
répond à une préoccupation propre au Cameroun~; un formulaire qui demande la
région d'origine ou la langue maternelle collecte une donnée sensible.

L'article~48, alinéa~1, interdit le traitement des données relatives aux
opinions, à l'origine, à la vie sexuelle, « à la génétique et à la
biométrie en matière de santé ». La lecture de cette fin de phrase n'est
pas tranchée~: la biométrie est-elle interdite seulement en matière de
santé, ou dans tous les cas~? L'auditeur ne tranche pas à la place du juge~;
il signale la difficulté et cote \NEK{} la licéité d'un pointage
biométrique tant qu'aucune autorisation ni aucun texte d'application ne
l'éclaire (Chapitre~\ref{chap:profilage}).

% ============================================================================
\section{Le Traitement et le Fichier}
\label{sec:c01-traitement}
% ============================================================================

\begin{boxdef}
\termecle{Traitement}~: selon l'article~5, « opération ou ensemble
d'opérations effectuées ou non à l'aide de procédés entièrement ou
partiellement automatisés ou non automatisés et appliqués à des données à
caractère personnel, notamment la collecte, l'enregistrement,
l'organisation, la conservation, l'adaptation ou la modification,
l'extraction, la consultation, l'utilisation, la communication par
transmission, diffusion ou toute autre forme de mise à disposition, le
rapprochement, ainsi que le verrouillage, l'effacement ou la destruction ».

\smallskip
\termecle{Fichier}~: « ensemble de données structurées et à caractère
personnel accessibles selon des critères déterminés, que cet ensemble soit
centralisé, décentralisé ou réparti de manière fonctionnelle ou
géographique » (article~5).
\end{boxdef}

Deux conséquences pratiques en découlent. D'abord, le papier compte~: la
loi vise aussi les procédés « non automatisés », et les dossiers de crédit
rangés dans les armoires des agences sont des fichiers. Ensuite, une simple
consultation est un traitement~: l'agent qui ouvre le dossier d'un client
sans raison de service traite des données, et c'est précisément ce que le
CENTRE MÉDICAL MBOLO découvrira dans ses journaux (Chapitre~\ref{chap:securite}).

Un traitement se définit par sa \termecle{finalité}~: l'ensemble des
opérations faites dans un même but. « Gérer les comptes des clients » est
un traitement~; « tenir la base de données Oracle » n'en est pas un, c'est
un moyen technique qui sert plusieurs traitements. Cette distinction
organise tout le recensement du Chapitre~\ref{chap:cartographie}.

\begin{boxobjet}{Gestion des comptes des clients, NYANGA CRÉDIT}
\textbf{Finalité~:} ouvrir, tenir et clôturer les comptes d'épargne et de
crédit. \textbf{Personnes concernées~:} clients, mandataires, cautions.
\textbf{Données~:} identité, pièce d'identité, adresse, téléphone,
profession, revenus déclarés, mouvements des comptes (données sensibles,
article~5). \textbf{Destinataires~:} agences, service crédit, éditeur du
logiciel bancaire, prestataire de recouvrement pour les impayés.
\textbf{Ce que l'audit établit au 7~juillet~:} le traitement existe
(\OUI{})~; son autorisation n'a pas été demandée (\NON{})~; sa durée de
conservation n'est fixée nulle part (\NON{})~; le fondement du traitement
des cautions reste à examiner (\NEK{}).
\end{boxobjet}

% ============================================================================
\section{Les Acteurs du Traitement}
\label{sec:c01-acteurs}
% ============================================================================

\begin{boxdef}
\termecle{Personne concernée}~: « toute personne dont les données à
caractère personnel font l'objet d'un traitement » (article~5).

\smallskip
\termecle{Responsable du traitement}~: « personne physique ou morale qui,
seule ou conjointement, collecte et traite des données à caractère
personnel et en détermine les moyens et les finalités » (article~5).

\smallskip
\termecle{Sous-traitant}~: « toute personne physique ou morale qui traite des
données à caractère personnel pour le compte du responsable du traitement
et sous ses instructions » (article~5).

\smallskip
\termecle{Destinataire}~: « toute personne physique ou morale, habilitée à
recevoir une communication des données à caractère personnel »
(article~5).
\end{boxdef}

La qualification se fait \textbf{opération par opération}, pas acteur par
acteur. L'éditeur du logiciel bancaire de NYANGA CRÉDIT est sous-traitant
quand il héberge les comptes selon les instructions de l'établissement~; il
deviendrait responsable s'il utilisait ces données pour améliorer ses
propres produits sans instruction. Un même acteur peut donc porter deux
qualités pour deux opérations différentes~; le Chapitre~\ref{chap:recherche}
montre que cette règle pose des questions difficiles dans la chaîne
pénale.

\begin{boxjuris}
\textbf{\loicam{n°~2024/017}, article~20}~\cite{loi2024017}~:
\textit{« Le responsable du traitement et le sous-traitant sont assujettis
aux mêmes obligations au titre de l'activité de traitement des données à
caractère personnel. »}

\smallskip
\textbf{Portée.} Le texte camerounais va plus loin que le droit européen,
qui distingue nettement les obligations des deux acteurs. Pour l'auditeur,
le prestataire de recouvrement et l'éditeur du logiciel bancaire répondent
des mêmes obligations que NYANGA CRÉDIT pour les traitements qu'ils
réalisent.
\end{boxjuris}

L'article~31 ajoute les \termecle{responsables conjoints}~: lorsque deux
organisations déterminent ensemble les finalités et les moyens, elles sont
« solidairement responsables » et doivent fixer leurs obligations dans une
convention « rendue publique » (Chapitre~\ref{chap:soustraitants}).

\begin{tableaucours}{L{3.6cm}L{4.4cm}L{4.0cm}}{%
    \tableauHeader{Loi n°~2024/017} &
    \tableauHeader{ISO/IEC 29100} &
    \tableauHeader{RGPD (UE)}
}
Personne concernée & Personne à qui se rapportent les données (\textit{PII
principal}) & Personne concernée \\
Responsable du traitement & \textit{PII controller} & Responsable du
traitement \\
Sous-traitant & \textit{PII processor} & Sous-traitant \\
Destinataire, tiers & Tiers (\textit{third party}) & Destinataire, tiers \\
Autorité de protection & Autorité de contrôle (hors norme) & Autorité de
contrôle \\
\end{tableaucours}
\vspace{-0.8em}
\begin{center}{\small\itshape Tableau~1.1. Les acteurs dans trois vocabulaires}\end{center}

La norme ISO/IEC~29100~\cite{iso29100} fournit un vocabulaire et des
principes communs à toutes les normes de la famille, en particulier
l'ISO/IEC~27701 (Chapitre~\ref{chap:normes}). Le RGPD~\cite{ue2016rgpd}
sert de point de comparaison~: beaucoup d'organisations camerounaises, comme
NOUN PAY, y sont exposées parce que leurs prestataires sont européens.

% ============================================================================
\section{Le Champ d'Application}
\label{sec:c01-champ}
% ============================================================================

L'article~2 soumet à la loi tout traitement effectué par l'État, les
collectivités territoriales décentralisées ou toute personne physique ou
morale~; tout traitement des données de personnes établies, résidant ou en
transit au Cameroun~; tout traitement réalisé par un responsable ou un
sous-traitant établi au Cameroun~; et tout traitement réalisé dans un
territoire où le droit camerounais s'applique. Le critère des personnes
« en transit » est large~: un voyageur dont les données sont traitées
pendant son passage au Cameroun est protégé.

L'article~3 exclut trois catégories~: les traitements d'une personne
physique pour ses activités exclusivement personnelles ou domestiques, à
condition que les données ne soient pas destinées à une communication
systématique, à un tiers ou à la diffusion~; les copies temporaires faites
pour la transmission sur un réseau~; et les traitements à des fins
littéraires, artistiques, archivistiques d'intérêt public, de recherche,
de statistique ou de journalisme, dans le respect des règles
déontologiques de ces professions.

Deux articles renvoient enfin à des textes particuliers~: l'article~4 pour
les traitements des autorités de sécurité et de défense, l'article~72 pour
les traitements relatifs « notamment à la sécurité, à la défense, à la
santé, à la justice et à l'état civil ». La portée de cette réserve est
discutée aux Chapitres~\ref{chap:secteurs} et~\ref{chap:recherche}~: tant
que les textes spécifiques ne sont pas identifiés, l'auditeur applique la
loi générale et signale la réserve.

% ============================================================================
\section{L'Inventaire de NYANGA}
\label{sec:c01-nyanga}
% ============================================================================

\begin{boxscenario}{Mission NYANGA. Onze familles de données, 7 juillet 2026}
\begin{tableaucours}{L{3.6cm}L{4.4cm}cc}{%
    \tableauHeader{Famille} &
    \tableauHeader{Exemples de données} &
    \tableauHeader{Sensible} &
    \tableauHeader{Connue}
}
Clients & Identité, revenus, mouvements & Oui & Oui \\
Cautions & Identité, revenus & Oui & Non \\
Prospects & Nom, téléphone, quartier & Non & Non \\
Débiteurs remis au recouvrement & Identité, montant dû & Oui & Non \\
Salariés & Identité, paie, santé (arrêts) & Oui & Oui \\
Candidats & CV, région d'origine & Oui & Non \\
Images des agences & Visages, heures de passage & Non & Non \\
Appels du service client & Voix, motif d'appel & Non & Non \\
Journaux de l'application & Adresse IP, identifiant & Non & Non \\
Utilisateurs du SI & Comptes, journaux d'accès & Non & Non \\
Fournisseurs (personnes physiques) & Identité, coordonnées bancaires & Oui &
Oui \\
\end{tableaucours}

\medskip
\textbf{Constatations.} Onze familles, dont huit absentes des listes de la
direction. Sept contiennent des données sensibles au sens de l'article~5,
dont quatre à cause des transactions bancaires et une à cause de la
région d'origine demandée aux candidats.

\textbf{Analyse} (confiance élevée). La vision de l'établissement était
organisée par applications et par clients. Les familles oubliées sont
celles dont la finalité n'est pas commerciale~: elles n'ont pas de
« propriétaire » évident dans l'organisation.
\end{boxscenario}

% ============================================================================
\section{Ce que la Notion demande à chaque Métier}
\label{sec:c01-metiers}
% ============================================================================

\begin{boxmetier}{Auditeur de la protection des données}
Recenser à partir des personnes concernées, pas des applications. Pour
chaque famille, demander qui la collecte, pour quoi faire, et où elle va.
Ne jamais accepter « anonyme » sans preuve.
\end{boxmetier}

\begin{boxmetier}{Juriste conformité}
Qualifier chaque acteur pour chaque opération, et relire les contrats à la
lumière de l'article~20~: le prestataire porte les mêmes obligations que
l'établissement.
\end{boxmetier}

\begin{boxmetier}{Responsable de traitement}
Désigner, pour chaque famille de données, un propriétaire dans
l'organisation. Une famille sans propriétaire est une famille que personne
ne protège.
\end{boxmetier}

\begin{boxresume}
\begin{itemize}
    \item Une donnée est personnelle dès qu'elle identifie une personne,
          directement ou indirectement~; la loi cite l'adresse IP, le
          matricule interne et la voix.
    \item Les données pseudonymisées restent personnelles~; seules des
          données dont l'anonymisation est prouvée sortent du champ de la
          loi.
    \item La liste camerounaise des données sensibles inclut les
          transactions bancaires, l'origine linguistique ou régionale et les
          poursuites judiciaires.
    \item Le responsable et le sous-traitant sont assujettis aux mêmes
          obligations (article~20)~; la qualification se fait opération par
          opération.
    \item À NYANGA CRÉDIT, onze familles de données ont été recensées, dont
          huit inconnues de la direction et sept sensibles.
\end{itemize}
\end{boxresume}

\begin{boxinfo}
\textbf{Points de contrôle.} Vous avez compris ce chapitre si vous savez~:
\begin{itemize}
    \item dire pourquoi un fichier pseudonymisé reste soumis à la loi~;
    \item citer trois données sensibles propres à la liste camerounaise~;
    \item qualifier un prestataire de responsable ou de sous-traitant pour
          une opération donnée.
\end{itemize}
\end{boxinfo}

\begin{boxexo}[title={\ding{45}\quad Exercice~1.1. Personnelle, sensible, ni l'une ni l'autre \niveaubase}]
Classez~: (a) le relevé de compte d'un client~; (b) la langue maternelle
d'un candidat~; (c) le chiffre d'affaires global d'une agence~; (d) la
photographie d'un client sur sa carte d'identité~; (e) l'adresse IP d'un
utilisateur de l'application~; (f) le casier judiciaire d'un candidat.
\end{boxexo}

\begin{boxexo}[title={\ding{45}\quad Exercice~1.2. Anonyme ou pseudonyme \niveaumoyen}]
Le service marketing transmet à un cabinet d'études un fichier où chaque
client est désigné par un code, avec son âge, son sexe, son agence, sa
profession et son encours. Le fichier est-il anonyme~? Que faudrait-il
faire pour qu'il le soit, et comment le prouver~?
\end{boxexo}

\begin{boxexo}[title={\ding{45}\quad Exercice~1.3. Qualifier les acteurs \niveaumoyen}]
Pour chaque opération, dites qui est responsable, qui est sous-traitant~:
(a) l'éditeur héberge les comptes et applique les sauvegardes demandées~;
(b) le prestataire de recouvrement appelle les débiteurs selon un script
fourni par NYANGA~; (c) le même prestataire utilise les numéros pour
proposer ses propres services~; (d) un opérateur de téléphonie envoie les
SMS de relance.
\end{boxexo}

\begin{boxexo}[title={\ding{45}\quad Exercice~1.4. Le champ de la loi \niveauexpert}]
\begin{enumerate}[(a)]
    \item Un agent de NYANGA conserve sur son téléphone personnel les
          numéros de ses clients pour les relancer. L'exclusion des
          activités domestiques (article~3) s'applique-t-elle~?
    \item Un journaliste publie des extraits des fichiers exfiltrés en
          mars~2025. Quel régime s'applique à son traitement~?
    \item L'expert judiciaire conserve une copie du dossier \texttt{Clients}
          pour son rapport. Quels articles faut-il lire, et quelle
          difficulté signalez-vous~?
\end{enumerate}
\end{boxexo}

\bigskip

\begin{boxinfo}
\textbf{Pour aller plus loin.}
\begin{itemize}
    \item L'ISO/IEC~29100~\cite{iso29100}, pour le vocabulaire et les onze
          principes de protection de la vie privée.
    \item H.~Nissenbaum, \textit{Privacy in Context}~\cite{nissenbaum2010},
          pour l'idée qu'une donnée est protégée par le contexte dans lequel
          elle circule.
\end{itemize}
\end{boxinfo}

\begin{boxinfo}
\textbf{Transition.} Savoir qu'un traitement existe ne dit pas s'il est
permis. Le Chapitre~\ref{chap:principes} expose les principes que tout
traitement doit respecter, et la règle du consentement, plus stricte au
Cameroun qu'en Europe.
\end{boxinfo}

        % ============================================================================
% Fichier  : partie1/P01_C02_principes.tex
% Titre    : Les Principes du Traitement
% Auteur   : MINKA MI NGUIDJOÏ Thierry Emmanuel
% Version  : 1.0 -- novembre 2026
% Corrigés : corriges/Corr_C02.tex
% ============================================================================
% Sources : loi n° 2024/017, art. 5 (consentement), 6 à 18, 21, 36, 50
% (texte officiel) ; ISO/IEC 29100 (principes, paraphrasés) ; RGPD art. 6
% (comparaison).
% Faits nouveaux (bible) : formulaire d'ouverture de compte de 2016 ; comptes
% clos depuis 2012 ; prospects collectés par les agents ; 14 juillet 2026.
% ============================================================================

\chapter{Les Principes du Traitement}
\label{chap:principes}
\chaptermark{Les principes du traitement}

\epigraph{%
    « Une case cochée n'est pas un consentement.\\
    C'est, au mieux, un indice. »%
}{Principe du cours}

\niveauChapitre{Fondamental}{%
    Chapitre~\ref{chap:donnees}
}

\begin{boxobjectifs}
\begin{itemize}
    \item Énoncer les principes des articles~6 à~18 de la loi n°~2024/017
          et les traduire en questions d'audit.
    \item Apprécier la validité d'un consentement au regard des
          articles~5 et~9.
    \item Identifier les trois cas où l'on peut se passer du consentement,
          et les comparer aux bases légales du droit européen.
    \item Vérifier la finalité, l'exactitude et la durée de conservation
          d'un traitement.
    \item Contrôler l'information fournie aux personnes selon
          l'article~21.
\end{itemize}
\end{boxobjectifs}

\bigskip

% ============================================================================
\section*{Avant de commencer : « ils ont tous signé »}
% ============================================================================

Bafoussam, mardi 14~juillet~2026. Le directeur commercial de NYANGA CRÉDIT
pose sur la table de Mme~BELLO une liasse de formulaires d'ouverture de
compte. « Ils ont tous signé. » Le formulaire, inchangé depuis 2016, tient
sur une page. En bas, au-dessus de la signature, une case pré-imprimée~:
« J'accepte les conditions générales de NYANGA CRÉDIT. » Les conditions
générales, huit pages, ne sont pas remises au client~; elles sont
« disponibles en agence ».

L'auditrice note trois questions. Les clients ont-ils consenti à quelque
chose de précis~? Les cautions, qui ne signent pas ce formulaire, ont-elles
consenti à quoi que ce soit~? Et les prospects, dont les agents notent le
nom et le téléphone sur un cahier lors des tournées de marché, savent-ils
seulement qu'ils figurent dans un fichier~?

\begin{boxalert}
\textbf{Violation typique.} Tenir une signature au bas de conditions
générales pour un consentement au traitement des données.

\smallskip
\textbf{Règle~:} le consentement exigé par la loi est préalable, libre,
éclairé, spécifique et univoque. Il porte sur un traitement identifié, après
une information claire~; il se prouve, et il peut être retiré.
\end{boxalert}

% ============================================================================
\section{Six Principes, Six Questions d'Audit}
\label{sec:c02-principes}
% ============================================================================

Le chapitre~I du titre~II de la loi pose les principes du traitement. Chacun
se traduit en une question fermée que l'auditeur peut coter.

\begin{tableaucours}{L{3.0cm}cL{8.0cm}}{%
    \tableauHeader{Principe} &
    \tableauHeader{Article} &
    \tableauHeader{Question d'audit}
}
Respect de la vie privée & 6 & Le traitement porte-t-il atteinte à la vie
privée au-delà de ce que sa finalité exige~? \\
Confidentialité & 7, 15 & Les données sont-elles protégées contre la
divulgation, y compris en transmission~? \\
Consentement & 9 & Un consentement valable, ou l'une des trois dérogations,
fonde-t-il le traitement~? \\
Licéité et loyauté & 10 & Le traitement est-il permis, et les personnes
ont-elles été traitées loyalement~? \\
Finalité & 11 & La finalité est-elle déterminée, explicite et légitime, et
l'usage ultérieur compatible~? \\
Exactitude & 12 & Les données sont-elles complètes, fiables, tenues à jour,
et les erreurs corrigées~? \\
Conservation limitée & 13 & La durée de conservation est-elle fixée et
respectée~? \\
Information & 14, 21 & Les personnes reçoivent-elles, à la collecte, les
informations de l'article~21~? \\
Non-discrimination & 17 & Le traitement est-il indifférent à la classe
sociale, à l'origine, aux opinions~? \\
Mineurs & 18 & Le traitement des données de mineurs est-il adéquat,
pertinent et limité~? \\
\end{tableaucours}
\vspace{-0.8em}
\begin{center}{\small\itshape Tableau~2.1. Les principes de la loi et leurs questions d'audit}\end{center}

\begin{boxinfo}
\textbf{La minimisation n'est pas un principe général de la loi.} Le mot
n'apparaît qu'à l'article~36, pour l'interconnexion des fichiers, et l'idée
de données « adéquates, pertinentes et limitées » n'est formulée qu'à
l'article~18, pour les mineurs. L'auditeur ne peut donc pas coter un écart à
un principe général de minimisation au titre de la loi. Il peut en revanche
s'appuyer sur la finalité (article~11)~: une donnée collectée sans lien avec
la finalité annoncée n'a pas de fondement. Les normes, elles, posent la
minimisation comme principe~\cite{iso29100}~; un audit contre l'ISO/IEC~27701
peut donc la retenir comme critère, à condition de le dire.
\end{boxinfo}

% ============================================================================
\section{Le Consentement}
\label{sec:c02-consentement}
% ============================================================================

\begin{boxdef}
\termecle{Consentement}~: selon l'article~5, « toute manifestation de volonté
expresse, non équivoque, non contrainte, libre et éclairée par une
information claire, précise et complète, par laquelle la personne
concernée, ou son représentant légal, accepte que ses données à caractère
personnel fassent l'objet d'un traitement ».
\end{boxdef}

\begin{boxjuris}
\textbf{\loicam{n°~2024/017}, article~9}~\cite{loi2024017}~:
\textit{« (1) Le traitement des données à caractère personnel est subordonné
au consentement préalable, libre, éclairé, spécifique et univoque de la
personne concernée. (2) Toutefois, il peut être dérogé à cette exigence du
consentement préalable lorsque le traitement est nécessaire~: au respect
d'une obligation légale à laquelle le responsable du traitement est soumis~;
à l'exécution d'une mission d'intérêt public ou relevant de l'exercice de
l'Autorité de protection des données à caractère personnel, dont est
investi le responsable du traitement ou le tiers auquel les données sont
communiquées~; à la préservation de la santé de la personne concernée ou
d'une autre personne physique. (3) Le consentement d'une personne mineure de
dix-huit (18) ans n'est valable que s'il est appuyé par celui de ses père
et mère, ou de son représentant légal. »}

\smallskip
\textbf{Portée.} Le consentement est la règle~; les dérogations sont trois,
et elles sont étroites. L'article~50 le confirme de façon négative~:
« le traitement des données à caractère personnel sans le consentement
préalable de la personne concernée est interdit ».
\end{boxjuris}

\subsection{Les cinq qualités d'un consentement}

Pour coter \OUI{} la proposition « le traitement repose sur un consentement
valable », l'auditeur vérifie cinq qualités, chacune par une pièce.

\begin{enumerate}
    \item \textbf{Préalable}~: recueilli avant le traitement. Pièce~: la date
          du recueil, comparée à la date de la première donnée.
    \item \textbf{Libre}~: sans contrainte. Un consentement exigé comme
          condition d'un service qui n'en a pas besoin n'est pas libre.
    \item \textbf{Éclairé}~: après une information claire, précise et
          complète. Pièce~: le texte remis au moment du recueil.
    \item \textbf{Spécifique}~: pour une finalité identifiée. Un
          consentement global à « toutes les opérations » ne l'est pas.
    \item \textbf{Univoque}~: un acte positif, sans ambiguïté. Une case
          pré-cochée ou un silence ne le sont pas.
\end{enumerate}

À ces cinq qualités s'ajoute la \textbf{preuve}~: c'est au responsable de
montrer qu'il a recueilli le consentement. Un consentement dont on ne
retrouve pas la trace est, pour l'auditeur, un consentement \NEK{}, et
souvent un \NON{} si la procédure de recueil n'existait pas.

\subsection{Les trois dérogations, comparées au droit européen}

Le droit européen connaît six bases légales~: le consentement, le contrat,
l'obligation légale, les intérêts vitaux, la mission d'intérêt public et
l'intérêt légitime~\cite{ue2016rgpd}. La loi camerounaise n'en retient que
quatre~: le consentement et trois dérogations. Elle ne connaît ni
l'\textbf{exécution d'un contrat}, ni l'\textbf{intérêt légitime}. Le
Chapitre~\ref{chap:rgpd} reprend les six bases une à une et montre ce que
cet écart change pour une organisation soumise aux deux textes.

\begin{tableaucours}{L{4.4cm}L{3.6cm}L{4.0cm}}{%
    \tableauHeader{Situation à NYANGA} &
    \tableauHeader{Fondement possible} &
    \tableauHeader{Commentaire}
}
Vérifier l'identité d'un client à l'ouverture & Obligation légale & Si un
texte de lutte contre le blanchiment l'impose \\
Tenir le compte et exécuter les opérations & Consentement & Pas de base
« contrat » dans la loi \\
Vidéosurveillance des agences & Consentement~? & Pas d'intérêt légitime~:
fondement fragile \\
Relancer un débiteur par un prestataire & Consentement & À recueillir
expressément \\
Démarcher un prospect au marché & Consentement & Préalable à l'inscription
au cahier \\
Transmettre les fichiers à l'expert judiciaire & Obligation légale &
Réquisition judiciaire \\
\end{tableaucours}
\vspace{-0.8em}
\begin{center}{\small\itshape Tableau~2.2. Fondements des traitements de NYANGA CRÉDIT}\end{center}

L'absence de l'intérêt légitime a des conséquences lourdes. La
vidéosurveillance, la sécurité informatique ou la prévention de la fraude,
qui reposent en Europe sur l'intérêt légitime, doivent trouver au Cameroun
un consentement, une obligation légale ou une mission d'intérêt public.
L'auditeur ne comble pas ce vide en important le droit européen~: il
signale le traitement dont le fondement est incertain, le cote \NEK{}, et
recommande d'obtenir l'autorisation de l'Autorité, qui en fixera les
conditions (article~19).

% ============================================================================
\section{Finalité, Exactitude, Conservation}
\label{sec:c02-finalite}
% ============================================================================

\begin{boxjuris}
\textbf{\loicam{n°~2024/017}, article~11}~\cite{loi2024017}~:
\textit{« Les données à caractère personnel doivent être collectées pour des
finalités déterminées, explicites et légitimes. Elles ne peuvent pas être
ultérieurement traitées de manière incompatible avec ces finalités. »}
\end{boxjuris}

Le détournement de finalité est l'écart le plus fréquent, et l'un des plus
graves~: l'article~67 en fait un délit. Le fichier des débiteurs, constitué
pour le recouvrement, devient incompatible s'il sert à démarcher les mêmes
personnes pour un nouveau crédit~; les images des caméras, prises pour la
sécurité, le deviennent si elles servent à contrôler les horaires des
agents.

L'article~12 impose des données « complètes, fiables et tenues à jour » et
oblige à corriger ou effacer les données inexactes. L'auditeur le teste par
sondage~: il tire des dossiers au hasard et compare les données aux pièces
justificatives.

\begin{boxjuris}
\textbf{\loicam{n°~2024/017}, article~13}~\cite{loi2024017}~:
\textit{« (1) Les données à caractère personnel doivent être conservées
pendant une durée qui n'excède pas la période nécessaire à leur traitement.
(2) La période de conservation des données à caractère personnel visée à
l'alinéa 1 ci-dessus est précisée dans l'autorisation de traitement des
données et suivant les modalités fixées par voie réglementaire. (3) Au-delà
de la période visée à l'alinéa 1 ci-dessus, les données à caractère
personnel ne peuvent être conservées qu'en vue de répondre spécifiquement à
un traitement à des fins historiques, statistiques ou de recherche
conformément aux dispositions légales en vigueur. »}

\smallskip
\textbf{Portée.} La durée se fixe dans l'autorisation~; tant que celle-ci
n'existe pas, l'organisation doit fixer elle-même une durée justifiée, et
l'auditeur vérifie qu'elle l'applique. L'article~28 renvoie aussi au
référentiel de l'Autorité.
\end{boxjuris}

% ============================================================================
\section{L'Information des Personnes}
\label{sec:c02-information}
% ============================================================================

L'article~14 exige une information « claire et transparente », et
l'article~21 en fixe le contenu~: elle doit être fournie « au plus tard, lors
de la collecte et quels que soient les moyens employés ».

\begin{tableaucours}{L{6.0cm}L{6.2cm}}{%
    \tableauHeader{Élément exigé (article~21)} &
    \tableauHeader{Ce que l'auditeur cherche}
}
Identité du responsable & Nom et coordonnées sur le formulaire \\
Finalités déterminées & Une phrase par finalité, sans formule globale \\
Catégories de données & La liste des données collectées \\
Destinataires & Les catégories de destinataires, prestataires compris \\
Possibilité de refuser de figurer au fichier & Une mention explicite \\
Droits (accès, rectification, effacement, opposition, limitation,
portabilité, décision automatisée, directives post-mortem, réclamation) &
La liste des droits et la manière de les exercer \\
Durée de conservation & Une durée ou un critère \\
Transfert éventuel vers un État étranger & La mention du pays ou de la
région \\
\end{tableaucours}
\vspace{-0.8em}
\begin{center}{\small\itshape Tableau~2.3. Le contenu de l'information (article~21)}\end{center}

La liste est longue, et elle se contrôle mécaniquement~: l'auditeur pose le
formulaire à côté du tableau et coche. L'information peut être donnée en
deux niveaux, une mention courte sur le formulaire et une notice complète
remise ou accessible, pourvu que la personne la reçoive vraiment au moment
de la collecte. « Disponible en agence » ne suffit pas.

% ============================================================================
\section{Les Principes à NYANGA}
\label{sec:c02-nyanga}
% ============================================================================

\begin{boxscenario}{Mission NYANGA. Le formulaire de 2016 et les autres, 14 juillet 2026}
\textbf{Constatations.}
\begin{itemize}
    \item Le formulaire d'ouverture de compte porte une seule case,
          pré-imprimée, d'acceptation des conditions générales~; aucune
          finalité n'y figure. Les conditions générales ne sont pas remises.
    \item Les cautions signent un acte de cautionnement qui ne mentionne
          pas le traitement de leurs données.
    \item Les cahiers de prospection de trois agences tirées au hasard
          contiennent 1\,214 noms et numéros de téléphone, sans aucune
          mention d'information ni de consentement.
    \item La base des comptes contient 61\,380 comptes clos, dont 22\,950
          clos avant 2016, conservés avec toutes leurs données.
    \item Un sondage de 40 dossiers actifs trouve 9 numéros de téléphone
          erronés et 6 adresses périmées.
\end{itemize}

\textbf{Analyse} (confiance élevée). Le consentement des clients n'est ni
spécifique, ni éclairé~: il ne fonde pas les traitements au sens de
l'article~9. Une partie des données de compte peut relever d'une obligation
légale d'identification, que l'audit n'a pas encore établie. Les cautions et les prospects n'ont
consenti à rien. Les comptes clos depuis dix ans et plus n'ont plus de
finalité identifiée.
\end{boxscenario}

\begin{boxlecture}{Les principes à NYANGA CRÉDIT}
\begin{tableaucours}{L{7.6cm}cL{3.4cm}}{%
    \tableauHeader{Proposition} &
    \tableauHeader{Cote} &
    \tableauHeader{Pièce}
}
Les clients ont donné un consentement valable & \NON{} & Formulaire 2016 \\
Les cautions ont consenti au traitement & \NON{} & Acte de cautionnement \\
Les prospects ont consenti au traitement & \NON{} & Cahiers de prospection \\
Une obligation légale fonde l'identification des clients & \NEK{} & Texte non produit \\
Une durée de conservation est fixée & \NON{} & Aucune règle écrite \\
Les données des comptes actifs sont exactes & \NON{} & Sondage, 15 erreurs
sur 40 \\
L'information de l'article~21 est fournie & \NON{} & Formulaire 2016 \\
La vidéosurveillance repose sur un fondement valable & \NEK{} & Pas de
consentement, pas d'autorisation \\
\end{tableaucours}

\medskip
\textbf{Décompte.} Zéro \OUI{}, six \NON{}, deux \NEK{} sur huit.
\hfill\lecture{0}{0{,}25}{0{,}81}

\medskip
\textbf{Lecture.} Le domaine des principes est clair~: $U$ est faible et
les écarts sont établis. L'entropie est basse parce que les réponses se
concentrent sur \NON{}. Une conclusion de non-conformité peut être donnée
dès maintenant, avec deux points à instruire.
\end{boxlecture}

% ============================================================================
\section{Ce que les Principes demandent à chaque Métier}
\label{sec:c02-metiers}
% ============================================================================

\begin{boxmetier}{Auditeur de la protection des données}
Pour chaque traitement, demander le fondement, puis la pièce qui le prouve.
Tester les cinq qualités du consentement une par une, et ne jamais
importer une base légale européenne que la loi ne connaît pas.
\end{boxmetier}

\begin{boxmetier}{Juriste conformité}
Réécrire les formulaires~: une finalité par phrase, une case non cochée par
consentement distinct, la notice de l'article~21 remise au moment de la
collecte. Identifier les textes qui imposent une obligation légale.
\end{boxmetier}

\begin{boxmetier}{Délégué à la protection des données}
Tenir, pour chaque traitement fondé sur le consentement, la preuve de son
recueil et la procédure de retrait.
\end{boxmetier}

\begin{boxresume}
\begin{itemize}
    \item Les principes des articles~6 à~18 se traduisent chacun en une
          question d'audit fermée.
    \item Le consentement est la règle~; il doit être préalable, libre,
          éclairé, spécifique et univoque, et sa preuve incombe au
          responsable.
    \item La loi ne connaît que trois dérogations au consentement~; ni le
          contrat ni l'intérêt légitime n'en font partie.
    \item La durée de conservation se fixe dans l'autorisation~; à défaut,
          l'organisation doit la fixer et l'appliquer.
    \item À NYANGA CRÉDIT, six écarts sont établis sur huit propositions~:
          consentements, durée, exactitude, information.
\end{itemize}
\end{boxresume}

\begin{boxinfo}
\textbf{Points de contrôle.} Vous avez compris ce chapitre si vous savez~:
\begin{itemize}
    \item tester les cinq qualités d'un consentement sur un formulaire~;
    \item expliquer pourquoi la vidéosurveillance pose un problème de
          fondement au Cameroun~;
    \item contrôler une notice d'information contre l'article~21.
\end{itemize}
\end{boxinfo}

\begin{boxexo}[title={\ding{45}\quad Exercice~2.1. Valable ou non \niveaubase}]
Le consentement est-il valable~? (a) une case pré-cochée « J'accepte de
recevoir des offres »~; (b) un client qui dicte son accord par téléphone,
enregistré~; (c) un consentement exigé pour ouvrir un compte d'épargne à
une utilisation de ses données pour la prospection~; (d) un élève de
16~ans qui s'inscrit seul à un service de microcrédit scolaire.
\end{boxexo}

\begin{boxexo}[title={\ding{45}\quad Exercice~2.2. Réécrire le formulaire \niveaumoyen}]
Réécrivez le formulaire d'ouverture de compte de NYANGA CRÉDIT pour qu'il
respecte les articles~9 et~21. Distinguez ce qui relève d'une obligation
légale et ce qui exige un consentement.
\end{boxexo}

\begin{boxexo}[title={\ding{45}\quad Exercice~2.3. Les comptes clos \niveaumoyen}]
\begin{enumerate}[(a)]
    \item Proposez une règle de conservation pour les comptes clos, en
          justifiant chaque durée.
    \item Que faire des 22\,950 comptes clos avant 2016~?
    \item Que faire des données de M.~TAGNE si son compte était clos~?
\end{enumerate}
\end{boxexo}

\begin{boxexo}[title={\ding{45}\quad Exercice~2.4. Sans intérêt légitime \niveauexpert}]
NOUN PAY veut analyser les transactions de ses clients pour détecter la
fraude. En Europe, elle invoquerait l'intérêt légitime. Quels fondements
peut-elle invoquer au Cameroun~? Quelles pièces l'auditeur demandera-t-il,
et comment cotez-vous la proposition « la détection de fraude repose sur un
fondement valable » si aucune autorisation n'a été délivrée~?
\end{boxexo}

\bigskip

\begin{boxinfo}
\textbf{Pour aller plus loin.}
\begin{itemize}
    \item L'ISO/IEC~29100~\cite{iso29100}, pour les principes de
          consentement, de limitation de la collecte et de minimisation.
    \item Le RGPD~\cite{ue2016rgpd}, article~6, pour comparer les bases
          légales.
\end{itemize}
\end{boxinfo}

\begin{boxinfo}
\textbf{Transition.} Les écarts de NYANGA commencent à se dessiner, mais
une liste de défauts n'est pas un audit. Le Chapitre~\ref{chap:audit}
présente la démarche qui transforme des constats en une conclusion que l'on peut défendre.
\end{boxinfo}


    % PARTIES 2 ET 3 · CONTEXTE INTERNATIONAL PUIS NATIONAL (version 3)
    \part{Le Contexte International}
    \label{part:international}

        % ============================================================================
% Fichier  : partie2/P02_C03_sources_internationales.tex
% Titre    : Les Sources Internationales et Africaines
% Auteur   : MINKA MI NGUIDJOÏ Thierry Emmanuel
% Version  : 3.0 -- novembre 2026
% Corrigés : corriges/Corr_C03.tex
% ============================================================================
% Sources : lignes directrices de l'OCDE (1980, révisées en 2013) ; principes
% directeurs de l'ONU (résolution 45/95) ; Convention 108 et protocole STCE
% n° 223 ; Convention de Malabo ; cadre stratégique de l'UA en matière de
% données (2022) ; acte additionnel CEDEAO A/SA.1/01/10 ; Convention de
% Budapest. Tous paraphrasés. Statuts vérifiés en ligne le 12/11/2026 pour
% Malabo (en vigueur le 08/06/2023) et la Convention 108+ (34 ratifications
% sur 38), la directive CEMAC n° 07/08-UEAC-133-CM-18 et la loi ivoirienne
% n° 2013-450. Loi marocaine n° 09-08 et liste des Parties africaines à la
% Convention 108 : citées de mémoire, à confirmer.
% ============================================================================

\chapter{Les Sources Internationales et Africaines}
\label{chap:international}
\chaptermark{Les sources internationales}

\epigraph{%
    « Une loi nationale de protection des données\\
    est rarement une invention. C'est une traduction. »%
}{Principe du cours}

\niveauChapitre{Fondamental}{%
    Chapitres~\ref{chap:donnees} et~\ref{chap:principes}~; droit du numérique
    (UE MSN-2S3-JUR01)
}

\begin{boxobjectifs}
\begin{itemize}
    \item Retracer la filiation des principes de protection des données, des
          lignes directrices de l'OCDE à la loi n°~2024/017.
    \item Situer le droit à la vie privée dans les textes de droits
          fondamentaux qui lient le Cameroun.
    \item Donner le statut exact, pour le Cameroun, de chaque instrument
          africain et international.
    \item Dresser la fiche du droit d'un pays où vont des données.
    \item Classer un instrument comme opposable, contractuel ou comparatif
          avant de s'en servir dans un audit.
\end{itemize}
\end{boxobjectifs}

\bigskip

% ============================================================================
\section*{Avant de commencer : quatre pays, quatre droits}
% ============================================================================

Yaoundé, vendredi 3~juillet~2026. Dans la salle de réunion de KORO AUDIT
\& CONFORMITÉ, M.~NGANDO punaise une carte du monde au mur. Il y plante
quatre épingles pour les trois missions du cabinet. Abidjan, où l'éditeur
du logiciel bancaire de NYANGA CRÉDIT héberge les comptes. Casablanca, où
un centre d'appels répond la nuit aux clients de NOUN PAY. Francfort, où se
trouvent les serveurs de la même NOUN PAY. La Virginie, d'où lui parviennent
sa messagerie et sa mesure d'audience.

Un stagiaire demande quel droit s'applique à tout cela. « Le nôtre d'abord,
répond M.~NGANDO, puisque les personnes sont ici. Mais chaque épingle a aussi
son droit et certains de ces droits suivent les données jusqu'à nous. » Il
ajoute une cinquième épingle sur Strasbourg et une sixième sur Addis-Abeba.
« Celles-là n'hébergent rien. Elles expliquent pourquoi nos lois se
ressemblent. »

\begin{boxalert}
\textbf{Violation typique.} Citer dans un rapport d'audit un instrument
international comme s'il obligeait l'organisation, sans avoir vérifié que le
Cameroun l'a ratifié ni qu'il vise ce type d'acteur.

\textbf{Règle~:} avant de s'appuyer sur un texte, l'auditeur établit son
statut. Il est opposable, parce que le Cameroun y est Partie ou qu'une loi
le reprend. Il est contractuel, parce qu'un partenaire l'impose par contrat.
Ou il est comparatif et ne sert alors qu'à éclairer et à proposer.
\end{boxalert}

% ============================================================================
\section{La Filiation des Principes}
\label{sec:c03-filiation}
% ============================================================================

Les principes que le Chapitre~\ref{chap:principes} a lus dans les
articles~6 à~18 de la loi ont une histoire. Ils ont été formulés pour la
première fois à l'échelle internationale en 1980 et n'ont guère changé de
substance depuis. Les lois nationales les internalisent, chacune avec ses
ajouts.

\subsection{Les lignes directrices de l'OCDE}

En 1980, l'Organisation de coopération et de développement économiques
adopte des lignes directrices sur la protection de la vie privée et les flux
transfrontières de données~\cite{ocde1980}. Elles ne sont pas un traité. Ce
sont des recommandations, révisées en 2013, que les États membres
s'engagent à suivre. Elles énoncent huit principes qui forment encore le
socle de presque toutes les lois du monde.

\begin{tableaucours}{L{4.2cm}L{4.6cm}L{3.2cm}}{%
    \tableauHeader{Principe de l'OCDE (paraphrasé)} &
    \tableauHeader{Idée} &
    \tableauHeader{Loi n°~2024/017}
}
Limitation de la collecte & Collecter loyalement, avec le consentement ou
l'information de la personne & Art.~9, 10 \\
Qualité des données & Des données exactes, complètes et à jour & Art.~12 \\
Spécification des finalités & Dire pourquoi l'on collecte, au plus tard à la
collecte & Art.~11, 21 \\
Limitation de l'utilisation & Ne pas détourner les données de leur finalité
& Art.~11, 67 \\
Garanties de sécurité & Protéger contre la perte, l'accès et la divulgation
& Art.~7, 15, 27 \\
Transparence & Faire connaître l'existence et la nature des traitements &
Art.~14, 29 \\
Participation individuelle & Permettre à la personne d'accéder à ses données
et de les contester & Art.~37 à~46 \\
Responsabilité & Rendre le responsable comptable du respect des principes &
Art.~20, 54 à~71 \\
\end{tableaucours}
\vspace{-0.8em}
\begin{center}{\small\itshape Tableau~3.1. Les huit principes de l'OCDE et
leur écho dans la loi}\end{center}

Le Cameroun n'est pas membre de l'OCDE. Le tableau ne décrit donc pas une
obligation, il décrit une parenté. Un auditeur qui la connaît lit plus vite
une loi étrangère et repère plus vite ce que la sienne a de particulier.

\subsection{Les principes directeurs des Nations Unies}

L'Assemblée générale des Nations Unies a adopté en 1990 des principes
directeurs pour la réglementation des fichiers informatisés contenant des
données à caractère personnel~\cite{onu1990}. Le texte reprend la licéité, la
loyauté, l'exactitude, la finalité, l'accès et la sécurité. Il ajoute deux
idées que l'on retrouve dans la loi camerounaise~: la non-discrimination, à
l'article~17 et l'existence d'une autorité de contrôle indépendante, à
l'article~53. C'est une résolution, qui recommande sans lier.

\subsection{La Convention 108 et sa version modernisée}

La Convention du Conseil de l'Europe du 28~janvier~1981, dite
Convention~108~\cite{coe1981conv108}, est le premier traité contraignant
consacré à la protection des données. Elle est ouverte aux États qui ne sont
pas membres du Conseil de l'Europe et plusieurs États africains y ont
adhéré, parmi lesquels le Maroc, le Sénégal, la Tunisie, Maurice et le
Cap-Vert. Le Cameroun n'y est pas Partie.

Un protocole de 2018 la modernise~: c'est la Convention~108+~\cite{coe2018conv108plus}.
Il ajoute la notification des violations, la protection des données dès la
conception et des droits face aux décisions automatisées. Il n'entrera en
vigueur qu'à la trente-huitième ratification. En 2026, le Conseil de
l'Europe en comptait trente-quatre.

\begin{boxinfo}
\textbf{Pourquoi un traité européen intéresse l'auditeur camerounais.} Le
centre d'appels de NOUN PAY est au Maroc, qui est Partie à la
Convention~108. Cette adhésion est un fait vérifiable sur le site du
dépositaire. Elle ne vaut pas reconnaissance par l'Autorité camerounaise,
mais elle entre dans la description du pays de destination que l'auditeur
joint au dossier de transfert (section~\ref{sec:c03-pays}).
\end{boxinfo}

% ============================================================================
\section{Le Droit à la Vie Privée}
\label{sec:c03-vieprivee}
% ============================================================================

La protection des données procède d'un droit plus ancien, le droit au
respect de la vie privée. Quatre textes le portent pour le Cameroun.

\begin{tableaucours}{L{4.4cm}L{4.4cm}L{3.2cm}}{%
    \tableauHeader{Texte} &
    \tableauHeader{Ce qu'il protège} &
    \tableauHeader{Statut}
}
Déclaration universelle des droits de l'homme (1948), article~12 & Vie
privée, famille, domicile, correspondance & Référence, visée par le préambule
de la Constitution \\
Pacte international relatif aux droits civils et politiques (1966),
article~17 & Les mêmes, contre les immixtions arbitraires ou illégales &
Traité, opposable \\
Constitution, préambule & Inviolabilité du domicile, secret de toute
correspondance & Norme suprême \\
Charte africaine des droits et du bien-être de l'enfant (1990), article~10 &
Vie privée de l'enfant & Traité, opposable \\
\end{tableaucours}
\vspace{-0.8em}
\begin{center}{\small\itshape Tableau~3.2. Le droit à la vie privée dans les
textes de droits fondamentaux}\end{center}

La Charte africaine des droits de l'homme et des peuples de 1981 ne contient
pas d'article consacré à la vie privée. Le fait est souvent ignoré. En
Afrique, la protection des données s'est donc construite par des textes
spécialisés plutôt que par la Charte.

Ces textes ne fournissent pas de critères d'audit~: ils sont trop généraux
pour être cotés. Ils servent au juge, qui peut y rattacher une atteinte que
la loi spéciale ne prévoit pas et ils rappellent à l'auditeur que la
personne concernée est titulaire d'un droit fondamental avant d'être une
ligne dans un registre.

% ============================================================================
\section{Le Cadre Africain}
\label{sec:c03-afrique}
% ============================================================================

\subsection{La Convention de Malabo}

La Convention de l'Union africaine sur la cybersécurité et la protection des
données à caractère personnel a été adoptée à Malabo le
27~juin~2014~\cite{ua2014malabo}. Elle traite trois sujets~: les transactions
électroniques, la protection des données et la cybersécurité. Son chapitre~II
demande à chaque État de se doter d'un cadre légal, d'une autorité nationale
indépendante, de formalités préalables aux traitements, de principes et de
droits proches de ceux du Tableau~3.1.

Elle est entrée en vigueur le 8~juin~2023, trente jours après le dépôt du
quinzième instrument de ratification. Le Cameroun l'a signée et ne l'avait
pas ratifiée à la date de ce cours. Elle ne l'oblige donc pas encore.

Sa lecture éclaire pourtant deux choix de la loi n°~2024/017. Le régime
d'autorisation préalable de l'article~19 prolonge les formalités préalables
de la Convention. La liste ouverte des données sensibles reprend l'esprit de
la sienne. Le législateur camerounais a internalisé un texte qu'il n'avait
pas ratifié.

\subsection{Le cadre stratégique de l'Union africaine}

En 2022, l'Union africaine a adopté un cadre stratégique en matière de
données~\cite{ua2022donnees}. Ce n'est pas un texte normatif. Il fixe une
orientation~: des données qui circulent entre États africains, des autorités
qui coopèrent, des règles nationales qui convergent. L'auditeur le cite
quand une organisation invoque la souveraineté des données pour refuser tout
hébergement hors du pays, ou au contraire le marché africain pour justifier
un transfert. Dans les deux cas, le cadre donne le vocabulaire et non la
règle.

\subsection{Les textes régionaux}

L'Afrique de l'Ouest a pris de l'avance. L'acte additionnel de la CEDEAO du
16~février~2010~\cite{cedeao2010acte} oblige ses États membres à légiférer
et la plupart l'ont fait. Le Cameroun n'appartient pas à la CEDEAO~; l'acte
ne lui est pas applicable.

En Afrique centrale, il n'existe pas de texte équivalent consacré aux
données personnelles. La CEMAC a adopté le 19~décembre~2008 cinq directives sur les
communications électroniques. L'une d'elles, la directive
n°~07/08-UEAC-133-CM-18, protège les droits des utilisateurs et demande aux
États de garantir la confidentialité des communications et des données de
trafic~\cite{cemac2008}. La loi camerounaise n°~2010/013 l'a transposée. Le secteur bancaire et celui des paiements
relèvent en outre de règlements communautaires et de la Commission bancaire
de l'Afrique centrale, qui encadrent le secret et la conservation des
données des clients (Chapitre~\ref{chap:loi2010}). Pour une microfinance ou
une fintech, ces textes sont opposables et s'appliquent directement dans
leur domaine.

\begin{boxinfo}
\textbf{Les autorités se parlent.} Les autorités africaines de protection
des données se sont réunies en 2016 dans un réseau, le RAPDP. Quand
l'Autorité camerounaise sera organisée, c'est là qu'elle trouvera ses
homologues. La loi lui donne mission de coopérer avec elles (article~53).
\end{boxinfo}

% ============================================================================
\section{Le Droit des Pays où Vont les Données}
\label{sec:c03-pays}
% ============================================================================

L'article~32 de la loi subordonne tout transfert à la vérification que le
pays de destination « offre un niveau de protection suffisant ». C'est à
l'Autorité de l'apprécier et le Chapitre~\ref{chap:transferts} dira comment
coter tant qu'elle ne l'a pas fait. L'organisation, elle, doit préparer la
description du pays. L'auditeur vérifie que cette description existe et
qu'elle repose sur des sources primaires.

\begin{boxprocedure}
\textbf{Dresser la fiche d'un pays de destination.}
\begin{enumerate}
    \item La loi générale de protection des données~: référence, date, champ.
    \item L'autorité de contrôle~: nom, indépendance, pouvoirs de sanction.
    \item Les engagements internationaux~: Convention~108, Convention de
          Malabo, autres.
    \item Les reconnaissances reçues d'autres autorités, par exemple une
          décision d'adéquation de l'Union européenne.
    \item L'accès des autorités publiques aux données et les recours ouverts
          aux étrangers.
    \item Les sources consultées et la date de consultation.
\end{enumerate}
\end{boxprocedure}

\begin{tableaucours}{L{2.2cm}L{3.6cm}L{2.8cm}L{3.4cm}}{%
    \tableauHeader{Pays} &
    \tableauHeader{Loi générale} &
    \tableauHeader{Autorité} &
    \tableauHeader{Engagements}
}
Côte d'Ivoire & Loi n°~2013-450 du 19~juin~2013~\cite{ci2013loi} & Autorité de régulation des télécommunications
(ARTCI) & Malabo ratifiée~; CEDEAO \\
Maroc & Loi n°~09-08 de 2009~\cite{ma2009loi} & Commission nationale
(CNDP) & Convention~108 \\
Union européenne & RGPD (Chapitre~\ref{chap:rgpd}) & Une autorité par État,
un comité européen & Convention~108 \\
États-Unis & Pas de loi fédérale générale~; lois sectorielles et lois
d'États & Pas d'autorité unique & Aucun traité de protection des données \\
\end{tableaucours}
\vspace{-0.8em}
\begin{center}{\small\itshape Tableau~3.3. Les quatre destinations des
données des trois missions (à vérifier sur les textes avant tout usage)}\end{center}

La dernière ligne mérite un mot. Les États-Unis n'ont pas de loi générale
comparable aux trois autres. La protection y dépend du secteur (santé,
finance, enfants) et de l'État (la Californie a la loi la plus connue). Les
entreprises américaines s'appuient souvent sur des cadres volontaires, dont
celui du NIST, que le Chapitre~\ref{chap:normes} présente. Une fiche pays
honnête l'écrit sans juger~: l'absence de loi générale est un fait, pas une
cote.

Aucun pays d'Afrique subsaharienne ne bénéficiait en 2026 d'une décision
d'adéquation de l'Union européenne. Les données qui partent d'Europe vers le
Cameroun voyagent donc sous clauses contractuelles, ce qui aura des
conséquences pour les sous-traitants camerounais de clients européens
(Chapitre~\ref{chap:rgpd}).

% ============================================================================
\section{La Coopération Pénale}
\label{sec:c03-budapest}
% ============================================================================

Les données franchissent aussi les frontières à la demande des enquêteurs.
La Convention de Budapest sur la cybercriminalité de
2001~\cite{budapest2001} organise cette coopération~: conservation rapide de
données stockées à l'étranger, entraide, points de contact disponibles en
permanence. Le Cameroun y est Partie depuis le 15~décembre~2023 et elle est
entrée en vigueur à son égard le 1\textsuperscript{er}~avril~2024, d'après
les travaux de l'auteur et de ses coauteurs~\cite{minka2026chainepenale}.

Pour l'auditeur de la protection des données, la Convention a deux effets.
Elle fonde des demandes que l'organisation auditée peut recevoir et qu'elle
doit savoir traiter sans violer l'article~32. Elle rencontre aussi le régime
d'autorisation préalable des transferts, rencontre que le
Chapitre~\ref{chap:recherche} étudie comme une tension transfrontalière.

% ============================================================================
\section{Le Statut de Chaque Instrument}
\label{sec:c03-statut}
% ============================================================================

\begin{tableaucours}{L{3.6cm}L{4.2cm}L{4.2cm}}{%
    \tableauHeader{Instrument} &
    \tableauHeader{Statut pour le Cameroun} &
    \tableauHeader{Usage dans l'audit}
}
Pacte de 1966, article~17 & Opposable & Fondement général, non coté \\
Convention de Budapest (2001) & Opposable depuis 2024 & Demandes des
autorités, coopération pénale \\
Textes CEMAC et COBAC & Opposables dans leur domaine & Critères sectoriels
(banque, paiement, télécommunications) \\
Convention de Malabo (2014) & Signée, non ratifiée & Comparatif et
prospectif \\
Convention~108 et 108+ & Non Partie & Comparatif~; fiche des pays Parties \\
Acte additionnel CEDEAO (2010) & Non applicable & Comparatif~; fiche des
pays de la CEDEAO \\
Lignes directrices de l'OCDE & Non membre & Comparatif \\
Principes directeurs de l'ONU (1990) & Recommandation & Comparatif \\
Cadre de l'UA sur les données (2022) & Orientation politique & Vocabulaire \\
RGPD (UE, 2016) & Voir Chapitre~\ref{chap:rgpd} & Opposable à certains
acteurs, contractuel pour d'autres \\
\end{tableaucours}
\vspace{-0.8em}
\begin{center}{\small\itshape Tableau~3.4. Les instruments internationaux et
leur statut}\end{center}

Trois statuts suffisent à l'auditeur.

\begin{itemize}
    \item \textbf{Opposable.} Le texte lie l'organisation. Un écart se cote
          \NON{} et se rattache à l'article.
    \item \textbf{Contractuel.} Le texte ne lie pas l'organisation par
          lui-même, mais un contrat le lui impose. L'écart est un manquement
          au contrat, que l'auditeur cote si le contrat figure parmi ses
          critères.
    \item \textbf{Comparatif.} Le texte n'oblige pas. Il sert à interpréter,
          à combler un silence par un critère de substitution déclaré, ou à
          proposer une bonne pratique. Il ne fonde jamais un \NON{} à lui
          seul.
\end{itemize}

Un statut change. La ratification de Malabo ferait passer la Convention de
comparative à opposable. L'auditeur date donc le tableau et le revoit à
chaque mission.

% ============================================================================
\section{La Carte des Trois Missions}
\label{sec:c03-cas}
% ============================================================================

\begin{boxscenario}{Missions NYANGA, MBOLO et NOUN PAY. La carte de M. NGANDO, 3 juillet 2026}
\textbf{Constatations.}
\begin{itemize}
    \item NYANGA CRÉDIT~: un flux vers la Côte d'Ivoire (hébergement du
          logiciel bancaire). Aucune description du droit ivoirien au
          dossier. Établissement soumis aux textes de la COBAC.
    \item CENTRE MÉDICAL MBOLO~: aucun flux hors du Cameroun identifié à ce
          stade. Prestataire à Douala.
    \item NOUN PAY~: flux vers l'Union européenne, les États-Unis et le
          Maroc, un quatrième vers un pays inconnu (les selfies). Aucune
          fiche pays. Établissement de paiement soumis aux textes
          communautaires.
\end{itemize}

\textbf{Analyse} (confiance élevée sur les flux déclarés, faible sur leur
exhaustivité). Aucune des trois organisations n'a décrit le droit des pays
où vont ses données. Deux d'entre elles relèvent de textes communautaires
qu'elles n'ont pas rapprochés de la loi n°~2024/017.
\end{boxscenario}

\begin{boxlecture}{Le dossier international de NOUN PAY}
\begin{tableaucours}{L{7.6cm}cL{3.4cm}}{%
    \tableauHeader{Proposition} &
    \tableauHeader{Cote} &
    \tableauHeader{Pièce}
}
Les pays de destination sont identifiés & \NON{} & Selfies, pays inconnu \\
Une fiche existe pour chaque pays de destination & \NON{} & Aucune fiche \\
Les fiches reposent sur des sources primaires datées & \NON{} & Sans objet,
aucune fiche \\
Les textes communautaires applicables sont recensés & \NON{} & Aucune
liste \\
Le Maroc est Partie à la Convention~108 & \OUI{} & Liste du dépositaire \\
Le niveau de protection de chaque pays est reconnu suffisant & \NEK{} &
Liste de l'Autorité non publiée \\
\end{tableaucours}

\medskip
\textbf{Décompte.} Un \OUI{}, quatre \NON{}, un \NEK{} sur six.
\hfill\lecture{0{,}17}{0{,}17}{1{,}25}

\medskip
\textbf{Lecture.} Le seul \OUI{} ne doit rien à l'entreprise~: c'est un fait
de droit international que l'auditeur a vérifié lui-même. Le \NEK{} est de
critère. Les quatre \NON{} relèvent tous de la préparation du dossier, qui
dépend de NOUN PAY seule.
\end{boxlecture}

% ============================================================================
\section{Ce que les Sources demandent à chaque Métier}
\label{sec:c03-metiers}
% ============================================================================

\begin{boxmetier}{Juriste conformité}
Tenir le tableau des instruments avec leur statut et leur date de
vérification. Aller au dépositaire~: l'Union africaine pour Malabo, le
Conseil de l'Europe pour la Convention~108 et pour Budapest.
\end{boxmetier}

\begin{boxmetier}{Auditeur de la protection des données}
Écrire dans le plan d'audit le statut de chaque texte retenu comme critère.
Ne jamais coter \NON{} sur la seule foi d'un texte comparatif.
\end{boxmetier}

\begin{boxmetier}{Délégué à la protection des données}
Tenir une fiche par pays de destination et la joindre au registre. La
mettre à jour quand un prestataire change de région d'hébergement.
\end{boxmetier}

\begin{boxresume}
\begin{itemize}
    \item Les principes de la loi n°~2024/017 descendent des huit principes
          de l'OCDE de 1980, repris par l'ONU, le Conseil de l'Europe et
          l'Union africaine.
    \item Le droit à la vie privée est garanti par le Pacte de 1966 et par
          le préambule de la Constitution~; la Charte africaine de 1981 ne
          le mentionne pas.
    \item La Convention de Malabo est en vigueur depuis le 8~juin~2023. Le
          Cameroun l'a signée sans la ratifier~: elle reste comparative.
    \item Un instrument est opposable, contractuel ou comparatif. Seuls les
          deux premiers fondent une cote \NON{}.
    \item L'organisation qui transfère des données tient une fiche du droit
          de chaque pays de destination, fondée sur des sources primaires.
\end{itemize}
\end{boxresume}

\begin{boxinfo}
\textbf{Points de contrôle.} Vous avez compris ce chapitre si vous savez~:
\begin{itemize}
    \item rattacher un article de la loi à un principe de l'OCDE~;
    \item dire pourquoi la Convention de Malabo ne fonde pas un constat
          d'audit au Cameroun~;
    \item dresser la fiche d'un pays de destination.
\end{itemize}
\end{boxinfo}

\begin{boxexo}[title={\ding{45}\quad Exercice~3.1. Quel statut \niveaubase}]
Donnez le statut (opposable, contractuel, comparatif) de chaque texte pour
NYANGA CRÉDIT~: (a) la Convention de Malabo~; (b) le Pacte de 1966~;
(c) l'acte additionnel de la CEDEAO~; (d) un règlement de la COBAC~;
(e) les lignes directrices de l'OCDE.
\end{boxexo}

\begin{boxexo}[title={\ding{45}\quad Exercice~3.2. La filiation \niveaumoyen}]
Pour chacun des articles suivants de la loi n°~2024/017, nommez le principe
de l'OCDE qu'il internalise, puis dites ce que la loi camerounaise y ajoute~:
(a) article~9~; (b) article~13~; (c) article~19~; (d) article~27.
\end{boxexo}

\begin{boxexo}[title={\ding{45}\quad Exercice~3.3. La fiche pays \niveaumoyen}]
Dressez, en suivant les six rubriques de la procédure, la fiche de la Côte
d'Ivoire pour le dossier de transfert de NYANGA CRÉDIT. Indiquez pour chaque
rubrique la source primaire que vous consulteriez et ce que vous écrivez si
vous ne la trouvez pas.
\end{boxexo}

\begin{boxexo}[title={\ding{45}\quad Exercice~3.4. Si le Cameroun ratifiait \niveauexpert}]
Supposez que le Cameroun dépose son instrument de ratification de la
Convention de Malabo.
\begin{enumerate}[(a)]
    \item À quelle date la Convention entrerait-elle en vigueur à son égard~?
          Quelle source consultez-vous pour le savoir~?
    \item Quelles lignes du Tableau~3.4 changent~?
    \item Un audit en cours doit-il modifier ses critères~? Rédigez la note
          au commanditaire.
\end{enumerate}
\end{boxexo}

\bigskip

\begin{boxinfo}
\textbf{Pour aller plus loin.}
\begin{itemize}
    \item Les lignes directrices de l'OCDE~\cite{ocde1980}, pour la
          formulation d'origine des huit principes.
    \item La Convention de Malabo~\cite{ua2014malabo}, chapitre~II et sa
          liste d'état tenue par l'Union africaine.
    \item La Convention~108+~\cite{coe2018conv108plus} et son rapport
          explicatif.
\end{itemize}
\end{boxinfo}

\begin{boxinfo}
\textbf{Transition.} Parmi tous ces textes, un seul suit les données
au-delà de ses frontières avec assez de force pour atteindre une entreprise
de Bafoussam~: le règlement européen. C'est l'objet du
Chapitre~\ref{chap:rgpd}.
\end{boxinfo}

        % ============================================================================
% Fichier  : partie2/P02_C04_rgpd.tex
% Titre    : Le RGPD et son Effet hors d'Europe
% Auteur   : MINKA MI NGUIDJOÏ Thierry Emmanuel
% Version  : 3.0 -- novembre 2026
% Corrigés : corriges/Corr_C04.tex
% ============================================================================
% Sources : règlement (UE) 2016/679 (Journal officiel de l'UE), paraphrasé ;
% lignes directrices 3/2018 du Comité européen de la protection des données
% sur le champ territorial ; décision d'exécution (UE) 2021/914 (clauses
% contractuelles types). Loi n° 2024/017, texte officiel. Le service NOUN
% DIASPORA et le contrat bruxellois de MEDISOFT sont des ajouts fictifs de
% la version 3.
% ============================================================================

\chapter{Le RGPD et son Effet hors d'Europe}
\label{chap:rgpd}
\chaptermark{Le RGPD hors d'Europe}

\epigraph{%
    « Le règlement européen ne s'arrête pas aux frontières de l'Europe.\\
    Il ne remplace pas pour autant la loi du pays où l'on est. »%
}{Principe du cours}

\niveauChapitre{Fondamental}{%
    Chapitres~\ref{chap:principes} et~\ref{chap:international}
}

\begin{boxobjectifs}
\begin{itemize}
    \item Présenter l'économie du RGPD en une quinzaine d'articles utiles à
          l'auditeur.
    \item Dire quand le RGPD oblige une organisation établie au Cameroun, en
          appliquant son article~3.
    \item Distinguer les trois voies d'exposition~: directe, contractuelle,
          en cascade.
    \item Comparer la loi n°~2024/017 et le RGPD point par point, sans
          prêter à l'une le contenu de l'autre.
    \item Coter une même exigence contre les deux textes quand les deux
          s'appliquent.
\end{itemize}
\end{boxobjectifs}

\bigskip

% ============================================================================
\section*{Avant de commencer : « c'est le RGPD camerounais »}
% ============================================================================

Bafoussam, lundi 3~août~2026. Le directeur général de NYANGA CRÉDIT a reçu
d'un cabinet de conseil une offre de « mise en conformité RGPD camerounais »
en douze semaines. Le document promet la désignation d'un délégué à la
protection des données « obligatoire », une procédure de notification des
violations « sous 72~heures » et la constitution d'un dossier d'« intérêt
légitime » pour la vidéosurveillance.

M.~NGANDO pose la proposition à côté du texte de la loi. Aucune de ces trois
affirmations ne s'y trouve. La loi ne connaît pas de délégué obligatoire,
elle dit « sans délai » et non 72~heures et elle ignore l'intérêt légitime.
L'offre ne dit pas un mot, en revanche, de l'autorisation préalable de
l'article~19.

Le même jour, à quelques rues de là, M.~NJOYA pose à Mme~BELLO la question
inverse. NOUN PAY\footnote{NOUN PAY, son service NOUN DIASPORA, MEDISOFT
SARL et la clinique bruxelloise citée plus bas sont fictifs.} a ouvert en
mars un service de transfert d'argent pour les Camerounais de France et de
Belgique. « Eux sont en Europe. Est-ce que leur règlement nous concerne~? »
La réponse est oui.

\begin{boxalert}
\textbf{Violation typique.} Deux erreurs symétriques. Lire la loi
camerounaise à travers le RGPD et lui attribuer des obligations européennes
qu'elle ne contient pas. Ou croire qu'une entreprise camerounaise échappe au
RGPD parce qu'elle n'a aucun bureau en Europe.

\textbf{Règle~:} l'auditeur détermine d'abord, traitement par traitement,
quel texte s'applique. La loi n°~2024/017 s'applique à tout traitement visé
par son article~2. Le RGPD s'y ajoute quand son article~3 est rempli ou
qu'un contrat l'impose. Aucun des deux ne dispense de l'autre.
\end{boxalert}

% ============================================================================
\section{Le RGPD en Quinze Articles}
\label{sec:c04-economie}
% ============================================================================

Le règlement (UE) 2016/679~\cite{ue2016rgpd} est applicable depuis le
25~mai~2018 dans tous les États de l'Union européenne. Il compte
quatre-vingt-dix-neuf articles. L'auditeur en retient une quinzaine.

\begin{tableaucours}{cL{4.0cm}L{6.6cm}}{%
    \tableauHeader{Art.} &
    \tableauHeader{Objet} &
    \tableauHeader{Ce qu'il faut en retenir (paraphrasé)}
}
3 & Champ territorial & Établissement dans l'Union, ou ciblage de personnes
qui s'y trouvent \\
5 & Principes & Sept principes, dont la minimisation et la responsabilité \\
6 & Bases légales & Six bases, dont le contrat et l'intérêt légitime \\
9 & Catégories particulières & Liste fermée de données sensibles \\
12 à 22 & Droits & Information, accès, rectification, effacement,
limitation, portabilité, opposition~; réponse en un mois \\
24, 25 & Responsabilité & Démontrer la conformité~; protection dès la
conception \\
27 & Représentant & Obligatoire pour l'acteur hors Union visé par
l'article~3 \\
28 & Sous-traitant & Contrat écrit au contenu imposé \\
30 & Registre & Tenu par le responsable et par le sous-traitant \\
32 & Sécurité & Mesures adaptées au risque \\
33, 34 & Violations & Autorité en 72~heures~; personnes si le risque est
élevé \\
35, 36 & Analyse d'impact & Obligatoire si le risque est élevé~;
consultation de l'autorité \\
37 à 39 & Délégué & Obligatoire dans trois cas \\
44 à 49 & Transferts & Adéquation, garanties appropriées, dérogations \\
83 & Amendes & Jusqu'à 20~millions d'euros ou 4~\% du chiffre d'affaires
mondial \\
\end{tableaucours}
\vspace{-0.8em}
\begin{center}{\small\itshape Tableau~4.1. Quinze articles du RGPD pour
l'auditeur}\end{center}

Le règlement repose sur une idée que la loi camerounaise n'a pas retenue~:
la \termecle{responsabilité}. Il a supprimé les formalités préalables auprès
de l'autorité. En échange, le responsable doit pouvoir démontrer à tout
moment qu'il respecte le texte. La loi n°~2024/017 fait le choix opposé,
celui du contrôle préalable par l'autorisation (Chapitre~\ref{chap:loi}).

% ============================================================================
\section{Quand le RGPD Oblige une Organisation Camerounaise}
\label{sec:c04-article3}
% ============================================================================

L'article~3 du règlement fixe son champ territorial. Il retient deux
critères, dont le second vise précisément les acteurs établis hors de
l'Union.

\begin{enumerate}
    \item \textbf{L'établissement.} Le règlement s'applique aux traitements
          faits dans le cadre des activités d'un établissement situé dans
          l'Union, que le traitement ait lieu ou non dans l'Union. Une
          filiale, une succursale ou un simple bureau stable suffisent.
    \item \textbf{Le ciblage.} Il s'applique aussi à un responsable ou à un
          sous-traitant non établi dans l'Union, lorsque le traitement est
          lié à l'offre de biens ou de services à des personnes qui se
          trouvent dans l'Union, ou au suivi de leur comportement dans
          l'Union.
\end{enumerate}

Le critère du ciblage tient à la localisation de la personne au moment du
traitement, pas à sa nationalité. Un Camerounais qui vit à Lyon est protégé
par le RGPD. Un Français en visite à Kribi ne l'est pas pour les traitements
d'une entreprise camerounaise qui ne vise pas l'Europe~; il l'est en
revanche par la loi camerounaise, dont l'article~2 protège les personnes
« en transit ».

Le Comité européen de la protection des données a précisé ce qu'est une
offre~\cite{edpb2018territorial}. Un site simplement accessible depuis
l'Europe ne suffit pas. Il faut des indices d'une intention de s'adresser à
des personnes dans l'Union.

\begin{tableaucours}{L{6.0cm}L{6.0cm}}{%
    \tableauHeader{Indice de ciblage} &
    \tableauHeader{À NOUN DIASPORA}
}
Prix affichés dans une monnaie de l'Union & Frais affichés en euros \\
Langue ou pays cités dans l'offre & « Envoyez depuis la France et la
Belgique » \\
Publicité achetée pour un public de l'Union & Campagnes sur les réseaux
sociaux, zone France et Belgique \\
Moyens de paiement de l'Union acceptés & Cartes bancaires européennes \\
Mention de clients de l'Union & Témoignages de clients de Lyon et de Liège \\
\end{tableaucours}
\vspace{-0.8em}
\begin{center}{\small\itshape Tableau~4.2. Les indices de ciblage, appliqués
au service de NOUN PAY}\end{center}

Les cinq indices sont réunis. NOUN PAY est un responsable de traitement non
établi dans l'Union qui offre un service à des personnes qui s'y trouvent.
Le règlement s'applique aux traitements de NOUN DIASPORA et à eux seuls~:
le microcrédit proposé aux clients de Bafoussam n'en relève pas.

L'article~27 en tire une première conséquence. L'acteur ainsi visé désigne
par écrit un représentant établi dans l'un des États où se trouvent les
personnes. Le représentant reçoit les demandes des personnes et de
l'autorité.

% ============================================================================
\section{Trois Voies d'Exposition}
\label{sec:c04-voies}
% ============================================================================

\begin{tableaucours}{L{2.4cm}L{4.6cm}L{5.0cm}}{%
    \tableauHeader{Voie} &
    \tableauHeader{Situation} &
    \tableauHeader{Conséquence}
}
Directe & L'organisation camerounaise cible des personnes dans l'Union
(NOUN DIASPORA) & Le RGPD lui est opposable pour ces traitements~;
représentant à désigner \\
Contractuelle & Elle est sous-traitante d'un responsable européen & Contrat
de l'article~28 et clauses contractuelles types pour le transfert vers le
Cameroun \\
En cascade & Elle confie ses données à un prestataire établi dans l'Union &
Le prestataire applique ses obligations de sous-traitant et les répercute
dans son contrat \\
\end{tableaucours}
\vspace{-0.8em}
\begin{center}{\small\itshape Tableau~4.3. Trois manières de rencontrer le
RGPD depuis le Cameroun}\end{center}

\subsection{La voie contractuelle}

MEDISOFT SARL, l'éditeur de Douala qui héberge le dossier patient de MBOLO,
a signé en 2026 un contrat de télémaintenance avec une clinique de
Bruxelles. Ses techniciens consultent depuis Douala des données de patients
belges. La clinique est responsable de traitement, MEDISOFT sous-traitant.

La clinique ne peut pas transférer ces données au Cameroun sans garantie~:
aucune décision d'adéquation ne couvre le pays
(Chapitre~\ref{chap:international}). Elle recourt donc aux clauses
contractuelles types adoptées par la Commission européenne en
2021~\cite{ue2021clauses}. MEDISOFT les signe telles quelles, car elles ne
se négocient pas. Elles lui imposent de n'agir que sur instruction,
d'assurer la sécurité, de signaler toute violation, d'accepter des audits et
de dire si le droit camerounais permet à une autorité d'accéder aux données.

Le RGPD est ici \textbf{contractuel} au sens du
Chapitre~\ref{chap:international}. L'auditeur de MEDISOFT cote ces
engagements si le contrat figure dans ses critères.

\subsection{La voie en cascade}

L'hébergeur de NOUN PAY est établi dans l'Union. Il est soumis au RGPD pour
ses propres obligations de sous-traitant, quel que soit le pays de son
client. Son contrat type reprend donc les clauses de l'article~28 et NOUN
PAY les a acceptées en ouvrant son compte. L'effet est favorable~: NOUN PAY
dispose d'un droit d'audit, d'un engagement de notification et d'une liste
des sous-traitants ultérieurs qu'elle n'aurait peut-être pas obtenus seule.

Un point reste à surveiller. Quand l'hébergeur renvoie les données vers le
Cameroun, par exemple pour qu'un agent de Bafoussam les consulte, il fait
lui-même un transfert hors de l'Union. Il peut exiger pour cela des
garanties de son client.

% ============================================================================
\section{La Loi et le RGPD, Point par Point}
\label{sec:c04-comparaison}
% ============================================================================

\begin{tableaucours}{L{2.8cm}L{4.6cm}L{4.6cm}}{%
    \tableauHeader{Point} &
    \tableauHeader{Loi n°~2024/017} &
    \tableauHeader{RGPD}
}
Champ & Personnes établies, résidant ou en transit au Cameroun~; acteurs
établis au Cameroun (art.~2) & Établissement dans l'Union~; ciblage de
personnes dans l'Union (art.~3) \\
Formalité préalable & Autorisation pour tout traitement (art.~19) & Aucune,
sauf consultation préalable (art.~36) \\
Bases du traitement & Consentement et trois dérogations (art.~9) & Six
bases, dont contrat et intérêt légitime (art.~6) \\
Minimisation & Pas de principe général & Principe (art.~5) \\
Délégué & Non prévu & Obligatoire dans trois cas (art.~37) \\
Registre & Responsable et sous-traitant (art.~29) & Les deux, avec une
exemption limitée pour les petites structures (art.~30) \\
Évaluation d'impact & Si risque élevé~; modalités réglementaires (art.~33) &
Si risque élevé~; critères publiés (art.~35) \\
Violation & « Sans délai », à l'Autorité et à la personne, pour toute
violation (art.~22) & 72~heures à l'autorité~; personnes si risque élevé
(art.~33, 34) \\
Sous-traitant & Mêmes obligations que le responsable (art.~20) & Obligations
propres, plus limitées (art.~28) \\
Responsables conjoints & Convention rendue publique (art.~31) & Grandes
lignes mises à disposition (art.~26) \\
Transferts & Autorisation, quatre vérifications (art.~32) & Adéquation,
garanties, dérogations (art.~44 à~49) \\
Décision automatisée & Droit d'opposition~; profilage puni (art.~44, 65) &
Encadrée, non interdite (art.~22) \\
Données sensibles & Liste ouverte, dont transactions bancaires et origine
régionale (art.~5) & Liste fermée (art.~9) \\
Délai de réponse aux demandes & Renvoyé à un texte réglementaire (art.~47) &
Un mois, prolongeable (art.~12) \\
Sanction maximale & Un milliard de francs CFA, personne morale (art.~71) &
4~\% du chiffre d'affaires mondial (art.~83) \\
\end{tableaucours}
\vspace{-0.8em}
\begin{center}{\small\itshape Tableau~4.4. Quinze points de comparaison}\end{center}

Trois lectures de ce tableau sont utiles.

La loi camerounaise est \textbf{plus stricte} sur six points~: l'autorisation
préalable, le petit nombre de fondements, la notification de toute
violation, les obligations du sous-traitant, la publicité de la convention
entre responsables conjoints et la liste des données sensibles. Une
organisation conforme au RGPD n'est donc pas conforme à la loi.

Elle est \textbf{moins précise} sur quatre~: la minimisation, les délais,
les critères du risque élevé et le contenu de la notification. Là, le RGPD
fournit des repères que l'auditeur peut emprunter.

Elle est \textbf{simplement différente} sur le reste. Le délégué n'existe
pas dans la loi et sa désignation ne peut être qu'une bonne pratique.

% ============================================================================
\section{Emprunter sans Confondre}
\label{sec:c04-emprunter}
% ============================================================================

Le Chapitre~\ref{chap:audit} montrera qu'un critère renvoyant à un texte non
publié se remplace par un \termecle{critère de substitution} déclaré. Le
droit européen est la première source de ces critères, parce qu'il est
précis, publié et commenté. Trois emprunts reviennent dans ce cours.

\begin{tableaucours}{L{3.6cm}L{4.2cm}L{4.2cm}}{%
    \tableauHeader{Silence de la loi} &
    \tableauHeader{Repère européen} &
    \tableauHeader{Formulation dans le rapport}
}
Délai de réponse aux demandes (art.~47) & Un mois (RGPD, art.~12) & « 30
jours (critère de l'audit) » \\
Critères du risque élevé (art.~33) & Liste du Comité européen, méthode de la
CNIL~\cite{cnil2018pia} & « Critères déclarés au plan d'audit » \\
Contenu de la notification (art.~22) & Rubriques des articles~33 et~34 &
« Contenu minimal retenu par l'audit » \\
\end{tableaucours}
\vspace{-0.8em}
\begin{center}{\small\itshape Tableau~4.5. Trois emprunts et leur
formulation}\end{center}

La règle tient en une phrase. Un emprunt se déclare dans le plan d'audit, se
nomme dans le rapport et ne se présente jamais comme une obligation légale
camerounaise. L'emprunt ne joue que dans un sens~: il précise un silence, il
n'allège jamais une exigence de la loi. On n'importe pas l'intérêt légitime.

% ============================================================================
\section{La Double Conformité}
\label{sec:c04-double}
% ============================================================================

Pour NOUN DIASPORA, les deux textes s'appliquent au même traitement. Le
client est à Lyon, donc protégé par le RGPD. NOUN PAY est établie au
Cameroun, donc soumise à la loi par son article~2. L'auditeur cote chaque
exigence deux fois.

\begin{tableaucours}{L{3.4cm}L{3.4cm}L{3.4cm}L{1.6cm}}{%
    \tableauHeader{Exigence} &
    \tableauHeader{Critère camerounais} &
    \tableauHeader{Critère européen} &
    \tableauHeader{Règle}
}
Fondement du traitement & Consentement ou dérogation (art.~9) & L'une des
six bases (art.~6) & Le plus strict \\
Notification d'une violation & Sans délai, Autorité et personne (art.~22) &
72~heures, autorité~; personne si risque élevé & Les deux \\
Formalité préalable & Autorisation (art.~19) & Aucune & Camerounais \\
Représentant & Aucun & Désignation écrite (art.~27) & Européen \\
\end{tableaucours}
\vspace{-0.8em}
\begin{center}{\small\itshape Tableau~4.6. Une exigence, deux critères}\end{center}

Quand les deux textes demandent la même chose à des degrés différents,
l'organisation satisfait le plus strict et satisfait ainsi les deux. Quand
ils demandent des choses différentes, elle fait les deux. Les deux textes ne
se contredisent sur aucun de ces points~: aucun n'interdit ce que l'autre
impose.

\begin{boxscenario}{Mission NOUN PAY. NOUN DIASPORA face au RGPD, septembre 2026}
\textbf{Constatations.}
\begin{itemize}
    \item 6\,840 clients inscrits depuis la France et la Belgique entre mars
          et août~2026~; frais en euros, publicité ciblée sur ces deux pays.
    \item Aucun représentant désigné dans l'Union. Aucune mention du RGPD
          dans la notice de l'application.
    \item La vérification d'identité de ces clients passe par le même
          prestataire de selfies, dont le pays reste inconnu.
    \item Le contrat de l'hébergeur contient les clauses de l'article~28.
    \item Aucune analyse d'impact~; aucune procédure de notification.
\end{itemize}

\textbf{Analyse} (confiance élevée). Le service remplit le critère du
ciblage. NOUN PAY l'a lancé sans se demander quel droit protégeait ses
clients. Le seul point conforme lui vient de son hébergeur.
\end{boxscenario}

\begin{boxlecture}{NOUN DIASPORA, exigences du RGPD}
\begin{tableaucours}{L{7.6cm}cL{3.4cm}}{%
    \tableauHeader{Exigence} &
    \tableauHeader{Cote} &
    \tableauHeader{Preuve}
}
Un représentant est désigné dans l'Union (art.~27) & \NON{} & Aucune
désignation \\
Chaque traitement a une base légale identifiée (art.~6) & \NON{} & Case
unique \\
L'information de l'article~13 est fournie & \NON{} & Notice de
l'application \\
Un registre est tenu (art.~30) & \NON{} & Aucun registre \\
Une analyse d'impact a été conduite (art.~35) & \NON{} & Aucune \\
Une procédure permet de notifier en 72~heures (art.~33) & \NON{} & Aucune \\
Le contrat de l'hébergeur contient les clauses de l'article~28 & \OUI{} &
Contrat, annexe données \\
Un délégué doit être désigné (art.~37) & \NEK{} & Suivi « à grande échelle »
à apprécier \\
\end{tableaucours}

\medskip
\textbf{Décompte.} Un \OUI{}, six \NON{}, un \NEK{} sur huit.
\hfill\lecture{0{,}125}{0{,}125}{1{,}06}

\medskip
\textbf{Lecture.} Le domaine est clair et non conforme. Ces huit cotes
s'ajoutent à celles de la loi camerounaise, elles ne les remplacent pas~: le
même service reçoit par ailleurs un \NON{} au titre de l'article~19.
\end{boxlecture}

% ============================================================================
\section{Ce que le RGPD demande à chaque Métier}
\label{sec:c04-metiers}
% ============================================================================

\begin{boxmetier}{Auditeur de la protection des données}
Poser la question de l'article~3 pour chaque traitement et consigner la
réponse au plan d'audit. Tenir deux colonnes de cotes quand les deux textes
s'appliquent.
\end{boxmetier}

\begin{boxmetier}{Juriste conformité}
Lire avant signature les clauses qu'un client ou un prestataire européen
impose. Elles créent des obligations d'audit, de notification et de
coopération que l'organisation devra tenir.
\end{boxmetier}

\begin{boxmetier}{Responsable de traitement}
Décider en connaissance de cause d'ouvrir un service à des personnes en
Europe. Le marché de la diaspora emporte le RGPD avec lui.
\end{boxmetier}

\begin{boxresume}
\begin{itemize}
    \item Le RGPD repose sur la responsabilité du responsable, la loi
          n°~2024/017 sur l'autorisation préalable.
    \item Son article~3 l'applique à une organisation camerounaise qui cible
          des personnes se trouvant dans l'Union~; l'article~27 lui impose
          alors un représentant.
    \item Trois voies d'exposition~: directe, contractuelle, en cascade.
    \item La loi camerounaise est plus stricte que le RGPD sur six points et
          moins précise sur quatre~; être conforme à l'un ne dit rien de
          l'autre.
    \item Un repère européen s'emprunte comme critère de substitution
          déclaré, pour préciser un silence, jamais pour alléger la loi.
\end{itemize}
\end{boxresume}

\begin{boxinfo}
\textbf{Points de contrôle.} Vous avez compris ce chapitre si vous savez~:
\begin{itemize}
    \item appliquer les indices de ciblage à un service en ligne~;
    \item citer trois points où la loi camerounaise est plus stricte que le
          RGPD~;
    \item formuler dans un rapport un critère emprunté au droit européen.
\end{itemize}
\end{boxinfo}

\begin{boxexo}[title={\ding{45}\quad Exercice~4.1. Le RGPD s'applique-t-il \niveaubase}]
Dites si le RGPD s'applique et par quelle voie~: (a) NYANGA CRÉDIT gère les
comptes de ses clients de l'Ouest~; (b) NOUN PAY propose NOUN DIASPORA~;
(c) MEDISOFT assure la télémaintenance d'une clinique de Bruxelles~;
(d) un touriste allemand est soigné à MBOLO~; (e) NOUN PAY héberge ses
données à Francfort.
\end{boxexo}

\begin{boxexo}[title={\ding{45}\quad Exercice~4.2. L'offre du cabinet \niveaumoyen}]
Reprenez l'offre de « mise en conformité RGPD camerounais » reçue par
NYANGA. Corrigez ses trois affirmations et ajoutez les trois obligations
camerounaises qu'elle omet, articles à l'appui.
\end{boxexo}

\begin{boxexo}[title={\ding{45}\quad Exercice~4.3. Deux colonnes \niveaumoyen}]
NOUN PAY découvre le 2~octobre qu'un fichier de 1\,200 clients de NOUN
DIASPORA a été envoyé par erreur à un partenaire. Dressez le tableau des
obligations de notification au titre de l'article~22 de la loi et des
articles~33 et~34 du RGPD~: destinataires, délais, contenu. Que fait NOUN
PAY pour satisfaire les deux textes~?
\end{boxexo}

\begin{boxexo}[title={\ding{45}\quad Exercice~4.4. Les clauses de MEDISOFT \niveauexpert}]
Les clauses signées par MEDISOFT l'obligent à dire si le droit camerounais
permet à une autorité publique d'accéder aux données des patients belges.
\begin{enumerate}[(a)]
    \item Quels textes camerounais MEDISOFT doit-il examiner pour répondre~?
          Appuyez-vous sur le Chapitre~\ref{chap:loi2010}.
    \item Le transfert de la Belgique vers le Cameroun est-il aussi un
          traitement soumis à l'article~19 de la loi n°~2024/017~?
    \item Comment l'auditeur de MEDISOFT cote-t-il la proposition « MEDISOFT
          respecte ses engagements envers la clinique »~?
\end{enumerate}
\end{boxexo}

\bigskip

\begin{boxinfo}
\textbf{Pour aller plus loin.}
\begin{itemize}
    \item Le RGPD~\cite{ue2016rgpd}, articles~3, 27, 28 et~44 à~49.
    \item Les lignes directrices 3/2018 du Comité européen sur le champ
          territorial~\cite{edpb2018territorial}.
    \item Les recommandations 01/2020 sur les transferts~\cite{edpb2020transferts}.
\end{itemize}
\end{boxinfo}

\begin{boxinfo}
\textbf{Transition.} Le RGPD est une loi. D'autres référentiels n'obligent
personne et organisent pourtant le travail de milliers d'entreprises~: le
cadre du NIST et les normes ISO. C'est l'objet du
Chapitre~\ref{chap:normes}.
\end{boxinfo}

        % ============================================================================
% Fichier  : partie2/P02_C05_nist_iso.tex
% Titre    : Le NIST Privacy Framework et les Normes ISO
% Auteur   : MINKA MI NGUIDJOÏ Thierry Emmanuel
% Version  : 3.0 -- novembre 2026
% Corrigés : corriges/Corr_C05.tex
% ============================================================================
% Sources : NIST Privacy Framework 1.0 (16 janvier 2020) et projet public de
% la version 1.1 (14 avril 2025) ; NISTIR 8062 ; SP 800-53 rév. 5 ;
% SP 800-122 ; AI RMF 1.0. Normes ISO/IEC 29100, 27701 (éditions 2019 et
% 2025), 29134, 27018, sous licence, paraphrasées. Les intitulés anglais des
% fonctions du NIST sont conservés. État de la version 1.1 du cadre du NIST
% à revérifier à chaque édition du cours. Les sections sur les normes ISO et
% les deux certifications viennent de l'ancien chapitre 9 (version 1).
% ============================================================================

\chapter{Le NIST Privacy Framework et les Normes ISO}
\label{chap:normes}
\chaptermark{NIST et normes ISO}

\epigraph{%
    « Un certificat dit qu'un système de management existe.\\
    Il ne dit pas que chaque traitement est licite. »%
}{Principe du cours}

\niveauChapitre{Fondamental}{%
    Chapitres~\ref{chap:international} et~\ref{chap:rgpd}~; ISO/IEC~27001
    (UE MSN-2S3-SEC08)
}

\begin{boxobjectifs}
\begin{itemize}
    \item Distinguer trois natures de référentiel~: la loi, la norme
          certifiable, le cadre volontaire.
    \item Décrire le \textit{Privacy Framework} du NIST~: noyau, profils,
          niveaux de mise en œuvre.
    \item Situer les normes ISO/IEC~29100, 27701, 29134 et 27018 et dire ce
          que l'édition 2025 de l'ISO/IEC~27701 a changé.
    \item Distinguer la certification prévue par la loi (article~34) de la
          certification d'un système de management.
    \item Établir une correspondance entre la loi, le RGPD, le cadre du NIST
          et la norme et en dire les limites.
\end{itemize}
\end{boxobjectifs}

\bigskip

% ============================================================================
\section*{Avant de commencer : « avec le certificat, nous serons conformes »}
% ============================================================================

Bafoussam, mardi 1\textsuperscript{er}~septembre~2026. NOUN PAY\footnote{NOUN
PAY, ses dirigeants et ses clients sont fictifs. L'entreprise figure aussi
dans le cours de gouvernance des SI de l'auteur.} prépare une levée de fonds.
Un investisseur américain demande une preuve de conformité en matière de
données et joint à son questionnaire une grille tirée du \textit{Privacy
Framework} du NIST. Le directeur général, M.~NJOYA, a trouvé la solution~: « Nous
allons nous faire certifier ISO/IEC~27701. Avec le certificat, nous serons conformes. »

Mme~BELLO, sollicitée pour un diagnostic, lui pose deux questions. Le
périmètre du certificat couvrira-t-il la notation automatique des
crédits, ou seulement l'application de paiement~? Et l'organisme
certificateur vérifiera-t-il que chaque traitement a reçu l'autorisation de
l'article~19~? À la seconde, la réponse est non~: un auditeur de
certification vérifie le système de management contre la norme, il ne
rend pas un avis juridique. Quant à la grille du NIST, personne à NOUN PAY
ne sait ce qu'elle mesure.

\begin{boxalert}
\textbf{Violation typique.} Présenter un certificat de système de
management comme une preuve de conformité à la loi.

\smallskip
\textbf{Règle~:} un certificat ISO/IEC~27701 atteste qu'un système de
management de la protection des données existe et fonctionne sur un
périmètre donné. La conformité à la loi se démontre par l'audit de
conformité légale, exigence par exigence. Les deux se complètent~; l'un ne
remplace pas l'autre.
\end{boxalert}

% ============================================================================
\section{Trois Natures de Référentiel}
\label{sec:c05-natures}
% ============================================================================

Les textes rencontrés depuis le début de cette partie n'ont pas la même
nature. L'auditeur les range en trois familles avant de s'en servir.

\begin{tableaucours}{L{2.6cm}L{3.2cm}L{3.0cm}L{3.2cm}}{%
    \tableauHeader{Nature} &
    \tableauHeader{Exemples} &
    \tableauHeader{Qui l'impose} &
    \tableauHeader{Ce que prouve la conformité}
}
Loi & Loi n°~2024/017, loi n°~2010/012, RGPD & L'État, sous peine de
sanction & Le respect d'obligations \\
Norme certifiable & ISO/IEC~27701, ISO/IEC~27001 & Un client, un marché,
l'organisation elle-même & L'existence d'un système de management, attestée
par un tiers \\
Cadre volontaire & \textit{Privacy Framework} du NIST & Personne~; un
partenaire peut le demander & Une maturité déclarée, sans certificat \\
\end{tableaucours}
\vspace{-0.8em}
\begin{center}{\small\itshape Tableau~5.1. Trois natures de référentiel}\end{center}

Un cadre volontaire devient contraignant de la même manière qu'un texte
étranger~: par contrat. L'investisseur de NOUN PAY peut faire du profil NIST
une condition de son entrée au capital. Le cadre reste alors contractuel au
sens du Chapitre~\ref{chap:international}.

% ============================================================================
\section{Le Privacy Framework du NIST}
\label{sec:c05-nist}
% ============================================================================

Le \textit{National Institute of Standards and Technology} est l'institut de
normalisation du département américain du commerce. Il a publié en
janvier~2020 la version~1.0 de son \textit{Privacy Framework}~\cite{nist2020pf},
un outil de gestion du risque d'atteinte à la vie privée. Le cadre est
volontaire, gratuit et sans certification. Il ne dépend d'aucune loi~: il a
été conçu dans un pays qui n'a pas de loi générale de protection des données
(Chapitre~\ref{chap:international}), pour des organisations qui doivent en
respecter plusieurs à la fois.

Il se compose de trois éléments, repris du cadre de cybersécurité du même
institut.

\subsection{Le noyau}

Le \termecle{noyau} (\textit{Core}) est un catalogue de résultats attendus,
rangés en cinq fonctions. Chaque fonction se divise en catégories, puis en
sous-catégories rédigées comme des résultats et non comme des moyens.

\begin{tableaucours}{L{2.8cm}L{5.0cm}L{4.2cm}}{%
    \tableauHeader{Fonction} &
    \tableauHeader{Question posée (paraphrasée)} &
    \tableauHeader{Écho dans la loi n°~2024/017}
}
\textit{Identify-P} & Quelles données traite-t-on, pour qui, avec quel
risque pour les personnes~? & Registre (art.~29), évaluation du risque
(art.~25, 33) \\
\textit{Govern-P} & Qui décide, selon quelles règles et quelles
obligations~? & Responsabilité des acteurs (art.~20), autorisation
(art.~19) \\
\textit{Control-P} & Sait-on gérer les données finement~: accès, correction,
suppression, conservation~? & Droits (art.~37 à~46), durée (art.~13) \\
\textit{Communicate-P} & Les personnes comprennent-elles ce que l'on fait de
leurs données~? & Information (art.~14, 21) \\
\textit{Protect-P} & Les données sont-elles protégées contre les
incidents~? & Sécurité (art.~27), violations (art.~22) \\
\end{tableaucours}
\vspace{-0.8em}
\begin{center}{\small\itshape Tableau~5.5. Les cinq fonctions du noyau}\end{center}

La fonction \textit{Protect-P} recouvre le cadre de cybersécurité du NIST,
dont les fonctions de détection, de réponse et de reprise complètent le
dispositif en cas de violation. Le cadre distingue ainsi deux risques que ce
cours a séparés dès le Chapitre~\ref{chap:principes}~: celui qui vient d'un
incident de sécurité et celui qui vient du traitement lui-même, même
parfaitement sécurisé. Le NIST nomme le second « action problématique sur
les données ».

\subsection{Les profils}

Un \termecle{profil} est la sélection des résultats du noyau qui comptent
pour une organisation. On en dresse deux. Le profil actuel décrit ce qui est
atteint aujourd'hui. Le profil cible décrit ce que l'organisation veut
atteindre, compte tenu de ses obligations et de ses risques. L'écart entre
les deux donne le plan d'action.

Le profil actuel se prête à la cotation du cours. Chaque résultat retenu
devient une proposition fermée~: atteint (\OUI{}), non atteint (\NON{}), ou
pas assez d'éléments pour le dire (\NEK{}). Un profil que l'organisation
remplit seule, sans preuve, vaut une déclaration au sens du
Chapitre~\ref{chap:audit}.

\subsection{Les niveaux de mise en œuvre}

Les \termecle{niveaux} (\textit{Tiers}) décrivent la manière dont
l'organisation gère le risque, de~1 à~4~: partielle, informée par le risque,
reproductible, adaptative. Ce ne sont pas des niveaux de maturité à
atteindre tous. Une petite structure peut viser légitimement le niveau~2.

\begin{boxinfo}
\textbf{La version 1.1.} Le NIST a publié en avril~2025 le projet d'une
version~1.1, qui réaligne le cadre sur la version~2.0 de son cadre de
cybersécurité et ajoute des éléments sur l'intelligence artificielle. À la
date de ce cours, l'institut annonçait la version finale pour 2026. La
structure en noyau, profils et niveaux ne change pas. L'auditeur à qui l'on
présente un profil vérifie la version qu'il utilise.
\end{boxinfo}

\subsection{Les autres publications utiles}

\begin{tableaucours}{L{3.0cm}L{4.8cm}L{4.2cm}}{%
    \tableauHeader{Publication} &
    \tableauHeader{Objet (paraphrasé)} &
    \tableauHeader{Usage dans ce cours}
}
NISTIR 8062 & Ingénierie de la vie privée~: prévisibilité, maîtrise,
dissociabilité & Évaluer le risque (Chapitre~\ref{chap:risque}) \\
SP 800-53, révision 5 & Catalogue de mesures de sécurité et de vie privée &
Détailler une mesure (Chapitre~\ref{chap:securite}) \\
SP 800-122 & Protection de la confidentialité des données
d'identification & Classer les données par niveau d'impact \\
SP 800-61 & Gestion des incidents & Violations
(Chapitre~\ref{chap:violations}) \\
AI RMF 1.0 & Gestion des risques de l'intelligence artificielle & Profilage
(Chapitre~\ref{chap:profilage}) \\
\end{tableaucours}
\vspace{-0.8em}
\begin{center}{\small\itshape Tableau~5.6. Cinq publications du NIST pour
l'auditeur}\end{center}

Les trois objectifs du NISTIR~8062 complètent utilement la triade de la
sécurité. La \textit{prévisibilité}~: la personne et l'organisation peuvent
anticiper ce qui sera fait des données. La \textit{maîtrise}~: on peut
modifier, supprimer ou communiquer une donnée à un niveau fin. La
\textit{dissociabilité}~: on peut traiter sans rattacher les données à une
personne quand ce lien n'est pas nécessaire.

\subsection{Ce que le cadre n'apporte pas}

Le cadre du NIST ne dit pas ce qui est permis. Il ne connaît ni
l'autorisation préalable, ni les trois dérogations au consentement, ni la
liste camerounaise des données sensibles. Une organisation peut présenter un
profil cible entièrement atteint et traiter des données sans fondement.

% ============================================================================
\section{La Famille des Normes}
\label{sec:c05-famille}
% ============================================================================

\begin{tableaucours}{L{2.6cm}L{5.0cm}L{4.4cm}}{%
    \tableauHeader{Norme} &
    \tableauHeader{Objet (paraphrasé)} &
    \tableauHeader{Usage dans l'audit}
}
ISO/IEC 29100 & Cadre de la protection de la vie privée~: vocabulaire,
acteurs, onze principes & Référentiel de principes, comparaison avec la
loi \\
ISO/IEC 27701 & Système de management de la protection de la vie privée &
Audit contre une norme, certification \\
ISO/IEC 29134 & Lignes directrices pour l'évaluation d'impact & Méthode de
l'évaluation d'impact (Chapitre~\ref{chap:risque}) \\
ISO/IEC 27018 & Protection des données personnelles dans le nuage public,
pour les sous-traitants & Évaluer un hébergeur
(Chapitre~\ref{chap:transferts}) \\
ISO/IEC 27001 et 27002 & Système de management et mesures de sécurité de
l'information & Domaine sécurité (Chapitre~\ref{chap:securite}) \\
ISO 19011, ISO/IEC 27007 & Conduite des audits & Méthode
(Chapitre~\ref{chap:audit}) \\
\end{tableaucours}
\vspace{-0.8em}
\begin{center}{\small\itshape Tableau~5.2. Les normes utiles à l'audit de protection des données}\end{center}

Toutes ces normes sont publiées sous licence~: ce cours les paraphrase et
n'en reproduit aucun passage. L'auditeur qui fonde un constat sur une
exigence normative doit disposer de l'édition en vigueur et la citer avec
son année.

\subsection{Les onze principes de l'ISO/IEC 29100}

L'ISO/IEC~29100~\cite{iso29100} énonce onze principes~: le consentement et le
choix~; la légitimité et la spécification des finalités~; la limitation de
la collecte~; la minimisation des données~; la limitation de l'utilisation,
de la conservation et de la divulgation~; l'exactitude et la qualité~;
l'ouverture, la transparence et l'information~; la participation et l'accès
des personnes~; la responsabilité~; la sécurité~; et le respect des
obligations de protection de la vie privée. La comparaison avec la loi
camerounaise fait apparaître la minimisation, absente comme principe
général de la loi (Chapitre~\ref{chap:principes}).

\subsection{L'ISO/IEC 27701}

L'ISO/IEC~27701~\cite{iso27701} étend la logique de l'ISO/IEC~27001 aux
données personnelles~: contexte, leadership, planification, support,
fonctionnement, évaluation et amélioration, plus des mesures propres aux
responsables de traitement et aux sous-traitants. Publiée en 2019 comme
extension de l'ISO/IEC~27001, elle a été révisée en octobre~2025 pour
devenir une norme de système de management autonome~: la certification
ISO/IEC~27001 n'est plus un préalable. Les deux éditions coexistent pendant
la période de transition~; l'auditeur vérifie l'édition sur laquelle porte
le certificat qu'on lui présente.

Elle apporte à l'audit trois choses que la loi ne donne pas~: une méthode
de management (un cycle d'amélioration, des revues de direction, des
audits internes), des mesures détaillées propres aux données personnelles
et un vocabulaire partagé avec les partenaires étrangers. Elle n'apporte
pas les obligations propres au Cameroun~: l'autorisation préalable, les
trois dérogations au consentement, le rapport annuel de l'article~27, la
liste camerounaise des données sensibles.

% ============================================================================
\section{Deux Certifications}
\label{sec:c05-certification}
% ============================================================================

\begin{boxjuris}
\textbf{\loicam{n°~2024/017}, article~34}~\cite{loi2024017}~:
\textit{« L'Autorité de protection des données à caractère personnel met en
place un mécanisme de certification du traitement des données à caractère
personnel conforme aux principes et exigences énoncés par la présente
loi. »}

\smallskip
\textbf{Portée.} L'article~5 définit la certification comme un outil
permettant de communiquer sur le niveau de protection « au regard des
critères d'un référentiel approuvé au préalable par l'Autorité ». Il s'agit
d'une certification de conformité à la loi, distincte des certifications
ISO. L'article~59 punit le responsable qui « opère sans certification
prévue dans la présente loi »~: tant que le mécanisme n'existe pas, la
portée de cette sanction est un \NEK{} de critère.
\end{boxjuris}

\begin{tableaucours}{L{3.2cm}L{4.4cm}L{4.4cm}}{%
    \tableauHeader{} &
    \tableauHeader{Certification de l'article~34} &
    \tableauHeader{Certification ISO/IEC 27701}
}
Délivrée par & L'Autorité, ou selon son mécanisme & Un organisme de
certification accrédité \\
Référentiel & Approuvé par l'Autorité & Norme internationale \\
Objet & Traitement, processus, produit & Système de management sur un
périmètre \\
Portée juridique & Prévue par la loi (articles~34, 59) & Aucune au regard de
la loi \\
État en novembre~2026 & Mécanisme non publié & Disponible \\
\end{tableaucours}
\vspace{-0.8em}
\begin{center}{\small\itshape Tableau~5.3. Deux certifications à ne pas confondre}\end{center}

% ============================================================================
\section{Correspondances entre les Référentiels}
\label{sec:c05-correspondances}
% ============================================================================

Une organisation exposée à plusieurs référentiels établit une table de
correspondance. Elle évite de faire plusieurs fois le même travail~: le
registre de l'article~29 sert aussi le RGPD, le cadre du NIST et la norme.

\begin{tableaucours}{L{2.6cm}L{2.2cm}L{2.2cm}L{2.4cm}L{2.6cm}}{%
    \tableauHeader{Thème} &
    \tableauHeader{Loi n°~2024/017} &
    \tableauHeader{RGPD} &
    \tableauHeader{NIST} &
    \tableauHeader{ISO/IEC 27701}
}
Recensement & Art.~29 & Art.~30 & \textit{Identify-P} & Registre des
traitements \\
Fondement & Art.~9 & Art.~6 & \textit{Govern-P} & Base légale identifiée \\
Autorisation & Art.~19 & Aucun & Aucun & Aucun \\
Information & Art.~21 & Art.~13, 14 & \textit{Communicate-P} & Information
des personnes \\
Droits & Art.~37 à~47 & Art.~15 à~22 & \textit{Control-P} & Obligations
envers les personnes \\
Risque et impact & Art.~25, 33 & Art.~35 & \textit{Identify-P} & Évaluation
d'impact \\
Sécurité & Art.~27 & Art.~32 & \textit{Protect-P} & Mesures de sécurité \\
Sous-traitants & Art.~30 & Art.~28 & \textit{Govern-P} & Mesures pour les
sous-traitants \\
Transferts & Art.~32 & Art.~44 à~49 & \textit{Govern-P} & Transferts,
pays \\
Violations & Art.~22 & Art.~33, 34 & \textit{Protect-P} et cadre de
cybersécurité & Gestion des incidents \\
Rapport annuel & Art.~27(2) & Aucun & Aucun & Aucun \\
\end{tableaucours}
\vspace{-0.8em}
\begin{center}{\small\itshape Tableau~5.4. Correspondances indicatives entre
quatre référentiels (thèmes paraphrasés)}\end{center}

Deux lignes n'ont de contenu que dans la première colonne~: l'autorisation
préalable et le rapport annuel. Aucun référentiel étranger ne les porte. Une
organisation qui bâtit sa conformité sur le RGPD, le NIST ou l'ISO seuls les
oubliera.

Une correspondance n'est jamais une équivalence. Une ligne dit que le
dispositif exigé par un référentiel peut porter l'exigence légale~; elle ne
dit pas qu'il le fait. L'audit de conformité vérifie toujours l'exigence
légale elle-même. Le Chapitre~\ref{chap:loi2010} le complétera par les
règles propres à la loi n°~2010/012 et donnera la règle pour choisir ses
critères.

% ============================================================================
\section{Le Diagnostic de NOUN PAY}
\label{sec:c05-cas}
% ============================================================================

\begin{boxscenario}{Mission NOUN PAY. Le diagnostic du 8 septembre 2026}
\textbf{Constatations.} Le périmètre envisagé pour la certification couvre
l'application de paiement et exclut la notation automatique « NOUN SCORE »,
« en cours de refonte ». Aucune demande d'autorisation n'a été préparée.
L'hébergeur, établi dans l'Union européenne, présente un certificat
ISO/IEC~27018.

\textbf{Analyse} (confiance élevée). Le certificat visé laisserait hors
périmètre le traitement le plus risqué de l'entreprise, celui qui
décide des crédits. Il serait exact et trompeur s'il était présenté comme
une preuve de conformité. Le certificat de l'hébergeur porte sur
l'hébergeur~: il ne dit rien des traitements de NOUN PAY ni du transfert
lui-même (Chapitre~\ref{chap:transferts}).
\end{boxscenario}

Pour répondre à l'investisseur, Mme~BELLO dresse avec l'équipe un profil
actuel de huit résultats, deux au moins par fonction qu'elle a pu examiner.

\begin{boxlecture}{NOUN PAY, profil actuel coté, 8 septembre 2026}
\begin{tableaucours}{L{2.4cm}L{5.4cm}cL{2.8cm}}{%
    \tableauHeader{Fonction} &
    \tableauHeader{Résultat attendu (paraphrasé)} &
    \tableauHeader{Cote} &
    \tableauHeader{Preuve}
}
\textit{Identify-P} & Les traitements et les flux sont inventoriés & \NON{} &
Aucun registre \\
\textit{Identify-P} & Le risque pour les personnes est évalué & \NON{} &
Aucune évaluation \\
\textit{Govern-P} & Les obligations légales sont recensées & \NON{} & Aucune
liste \\
\textit{Govern-P} & Les rôles sont attribués & \NON{} & Aucune
désignation \\
\textit{Control-P} & Les données sont supprimées à l'échéance & \NEK{} & Non
examiné \\
\textit{Communicate-P} & Les personnes sont informées des traitements &
\NON{} & Notice de l'application \\
\textit{Protect-P} & Les données sont chiffrées au repos et en transit &
\OUI{} & Configuration de l'hébergeur \\
\textit{Protect-P} & Les accès sont nominatifs et tracés & \OUI{} & Journaux
de la console \\
\end{tableaucours}

\medskip
\textbf{Décompte.} Deux \OUI{}, cinq \NON{}, un \NEK{} sur huit.
\hfill\lecture{0{,}25}{0{,}125}{1{,}30}

\medskip
\textbf{Lecture.} Les deux résultats atteints sont dans la seule fonction de
protection. NOUN PAY a traité la vie privée comme un problème de sécurité.
Le profil montre à l'investisseur ce qu'un certificat de sécurité n'aurait
pas montré.
\end{boxlecture}

% ============================================================================
\section{Ce que les Référentiels demandent à chaque Métier}
\label{sec:c05-metiers}
% ============================================================================

\begin{boxmetier}{Délégué à la protection des données}
Bâtir le système de management sur les obligations de la loi, puis le
rapprocher des référentiels demandés par les partenaires. Tenir la table de
correspondance et signaler les exigences camerounaises qu'aucun référentiel
étranger ne porte.
\end{boxmetier}

\begin{boxmetier}{Auditeur de la protection des données}
Lire le périmètre et l'édition de tout certificat, la version de tout
profil. Ne jamais coter \OUI{} une exigence légale sur la seule foi d'un
certificat ou d'un profil déclaré.
\end{boxmetier}

\begin{boxmetier}{Responsable de traitement}
Choisir le périmètre d'une certification en fonction des risques, pas de la
facilité. Un certificat qui exclut le traitement le plus risqué expose à un
reproche de présentation trompeuse.
\end{boxmetier}

\begin{boxresume}
\begin{itemize}
    \item Une loi oblige, une norme se certifie, un cadre volontaire se
          déclare. Les deux derniers ne deviennent contraignants que par
          contrat.
    \item Le \textit{Privacy Framework} du NIST comprend un noyau de cinq
          fonctions, des profils actuel et cible et quatre niveaux de mise
          en œuvre. Il ne dit pas ce qui est permis.
    \item L'ISO/IEC~27701 est depuis 2025 une norme autonome~; l'auditeur
          vérifie l'édition du certificat.
    \item La certification de l'article~34 est une certification de
          conformité à la loi, distincte d'une certification ISO.
    \item L'autorisation préalable et le rapport annuel n'ont d'équivalent
          dans aucun référentiel étranger.
\end{itemize}
\end{boxresume}

\begin{boxinfo}
\textbf{Points de contrôle.} Vous avez compris ce chapitre si vous savez~:
\begin{itemize}
    \item nommer les cinq fonctions du noyau du NIST et la question que pose
          chacune~;
    \item dire ce qu'un certificat ISO/IEC~27701 prouve et ce qu'il ne
          prouve pas~;
    \item citer deux obligations camerounaises absentes des référentiels
          étrangers.
\end{itemize}
\end{boxinfo}

\begin{boxexo}[title={\ding{45}\quad Exercice~5.1. Quelle norme \niveaubase}]
Quelle norme consultez-vous pour~: (a) conduire une évaluation d'impact~;
(b) évaluer un hébergeur en nuage~; (c) mettre en place un système de
management de la protection des données~; (d) définir le vocabulaire d'un
contrat~?
\end{boxexo}

\begin{boxexo}[title={\ding{45}\quad Exercice~5.2. Le périmètre de NOUN PAY \niveaumoyen}]
Rédigez l'avis de Mme~BELLO à M.~NJOYA sur le périmètre de la
certification, en séparant constatations et analyse. Proposez un périmètre
et un calendrier qui tiennent compte de la notation automatique.
\end{boxexo}

\begin{boxexo}[title={\ding{45}\quad Exercice~5.3. Compléter la correspondance \niveaumoyen}]
Ajoutez au Tableau~5.4 les lignes correspondant aux articles~22
(violations), 30 (sous-traitance) et 32 (transferts), avec le thème de la
norme et la couverture.
\end{boxexo}

\begin{boxexo}[title={\ding{45}\quad Exercice~5.4. Le profil cible \niveauexpert}]
L'investisseur demande à NOUN PAY un profil cible à douze mois.
\begin{enumerate}[(a)]
    \item Pour chacun des huit résultats du profil actuel, proposez la cible
          et la preuve qui permettra de la coter \OUI{}.
    \item Ajoutez deux résultats que la loi n°~2024/017 impose et que le
          noyau du NIST ne contient pas. Comment les présentez-vous à un
          investisseur qui ne connaît que le cadre du NIST~?
    \item Quel niveau de mise en œuvre NOUN PAY peut-elle viser
          raisonnablement~? Justifiez.
\end{enumerate}
\end{boxexo}

\bigskip

\begin{boxinfo}
\textbf{Pour aller plus loin.}
\begin{itemize}
    \item Le \textit{Privacy Framework} du NIST~\cite{nist2020pf}, librement
          accessible et le NISTIR~8062~\cite{nistir8062}.
    \item L'ISO/IEC~27701~\cite{iso27701} et l'ISO/IEC~29100~\cite{iso29100}.
    \item L'ISO/IEC~27018~\cite{iso27018}, pour les sous-traitants en nuage.
\end{itemize}
\end{boxinfo}

\begin{boxinfo}
\textbf{Transition.} Les sources internationales, le règlement européen et
les référentiels volontaires forment le paysage. Il reste à lire de près le
texte qui oblige toute organisation camerounaise et qui a internalisé une
grande partie de ce paysage~: la loi n°~2024/017, avec son Autorité et ses
sanctions. C'est l'objet du Chapitre~\ref{chap:loi}.
\end{boxinfo}


    \part{Le Contexte National}
    \label{part:national}

        % ============================================================================
% Fichier  : partie3/P03_C06_loi_2024.tex
% Titre    : La Loi n° 2024/017 : Autorisation, Autorité, Sanctions
% Auteur   : MINKA MI NGUIDJOÏ Thierry Emmanuel
% Version  : 3.0 -- novembre 2026
% Corrigés : corriges/Corr_C06.tex
% ============================================================================
% Sources : loi n° 2024/017 (texte officiel, titres I à VII) ; loi
% n° 2010/012, art. 41 à 47 (texte officiel) ; Convention de Malabo ;
% Convention de Budapest ; acte additionnel CEDEAO A/SA.1/01/10 ; RGPD.
% Statut des conventions d'après Minka Mi Nguidjoï, Massoussi Songo et
% Bialo (2026, manuscrit) : Budapest, Cameroun Partie le 15/12/2023, en
% vigueur le 01/04/2024 ; Malabo, signée le 12/08/2021, non ratifiée.
% ============================================================================

\chapter{La Loi n° 2024/017 : Autorisation, Autorité, Sanctions}
\label{chap:loi}
\chaptermark{La loi n° 2024/017}

\epigraph{%
    « On cite une loi pour ce qu'elle dit,\\
    pas pour ce que l'on aimerait qu'elle dise. »%
}{Principe du cours}

\niveauChapitre{Fondamental}{%
    Chapitres~\ref{chap:donnees}, \ref{chap:principes} et
    \ref{chap:international} à~\ref{chap:normes}~; droit du numérique
    (UE MSN-2S3-JUR01)
}

\begin{boxobjectifs}
\begin{itemize}
    \item Présenter l'architecture de la loi n°~2024/017 en sept titres.
    \item Expliquer le régime d'autorisation préalable et ses conséquences
          tant que l'Autorité n'est pas organisée.
    \item Décrire les missions de l'Autorité et l'échelle des sanctions
          administratives, civiles et pénales.
    \item Dire ce que la loi a internalisé des sources internationales et
          ce qu'elle a ajouté en propre.
    \item Recenser les dispositions qui attendent un texte d'application et
          en tirer les conséquences pour la cotation.
    \item Lire les dispositions transitoires et la clause d'abrogation.
\end{itemize}
\end{boxobjectifs}

\bigskip

% ============================================================================
\section*{Avant de commencer : « où dépose-t-on la demande ? »}
% ============================================================================

Bafoussam, jeudi 25~juin~2026. Le lendemain de la lettre de M.~TAGNE, Me~ETOUNDI
lit la loi au directeur général, article par article. À l'article~19,
M.~TCHOUA l'arrête~: « Une autorisation pour chaque traitement. Très bien.
Où dépose-t-on la demande~? »

Me~ETOUNDI tourne les pages jusqu'à l'article~53. L'Autorité de protection
existe dans la loi. Sa création, son organisation et son fonctionnement
relèvent d'un décret qu'elle n'a pas trouvé. Les modalités de l'autorisation
relèvent d'un autre texte. « Je ne sais pas, répond-elle. Mais je sais ce
qu'il y aura dans le dossier le jour où le guichet ouvrira. »

\begin{boxalert}
\textbf{Violation typique.} Conclure que la loi ne s'applique pas encore
parce que l'Autorité n'est pas organisée, et attendre.

\smallskip
\textbf{Règle~:} la loi est en vigueur et son délai de mise en conformité est
expiré. Les obligations qui ne dépendent d'aucun texte d'application
s'imposent dès maintenant. Pour les autres, l'organisation prépare ce qui
dépend d'elle et l'auditeur cote séparément l'exigence et la diligence.
\end{boxalert}

% ============================================================================
\section{Ce que la Loi a Internalisé}
\label{sec:c06-internalise}
% ============================================================================

Les chapitres précédents ont décrit le paysage international. La loi
n°~2024/017 en est l'héritière sans en être la copie. Le tableau
ci-dessous sépare ce qu'elle reprend de ce qu'elle ajoute.

\begin{tableaucours}{L{2.8cm}L{4.4cm}L{4.8cm}}{%
    \tableauHeader{Titre de la loi} &
    \tableauHeader{Ce qu'elle reprend} &
    \tableauHeader{Ce qu'elle ajoute ou change}
}
I. Définitions, champ & Vocabulaire commun à la Convention~108, à Malabo et
au RGPD & Protection des personnes « en transit »~; liste ouverte de données
sensibles, dont les transactions bancaires et l'origine régionale \\
II. Principes & Les huit principes de l'OCDE & Consentement comme règle,
trois dérogations seulement~; pas de principe général de minimisation \\
II. Formalités & Formalités préalables de Malabo & Autorisation pour tout
traitement, sans régime de déclaration \\
II. Obligations & Registre, évaluation d'impact, sous-traitance, transferts
du RGPD & Mêmes obligations pour le sous-traitant~; convention publique~;
quatre vérifications cumulées pour un transfert~; rapport annuel de
sécurité \\
III. Droits & Droits du RGPD, portabilité comprise & Directives après le
décès~; délais renvoyés au règlement \\
V. Autorité & Autorité indépendante des principes de l'ONU et de Malabo &
Création et organisation par décret \\
VI. Sanctions & Amendes administratives & Délits nombreux, dont le
profilage~; responsabilité pénale des personnes morales \\
\end{tableaucours}
\vspace{-0.8em}
\begin{center}{\small\itshape Tableau~6.1. Ce que la loi reprend et ce
qu'elle ajoute}\end{center}

La colonne de droite est celle que l'auditeur doit connaître par cœur. Elle
contient tout ce qu'un consultant formé au seul RGPD oubliera
(Chapitre~\ref{chap:rgpd}).

% ============================================================================
\section{L'Architecture de la Loi}
\label{sec:c06-architecture}
% ============================================================================

\begin{tableaucours}{cL{5.6cm}cL{4.0cm}}{%
    \tableauHeader{Titre} &
    \tableauHeader{Objet} &
    \tableauHeader{Art.} &
    \tableauHeader{Chapitres du cours}
}
I & Dispositions générales~: objet, champ, définitions & 1--5 &
\ref{chap:donnees} \\
II & Régime juridique~: principes, formalités préalables, obligations,
interconnexion & 6--36 & \ref{chap:principes}, \ref{chap:cartographie},
\ref{chap:risque}, \ref{chap:securite}, \ref{chap:soustraitants},
\ref{chap:transferts}, \ref{chap:violations} \\
III & Droits de la personne concernée & 37--47 & \ref{chap:droits},
\ref{chap:profilage} \\
IV & Interdictions & 48--52 & \ref{chap:donnees}, \ref{chap:secteurs} \\
V & Autorité de protection & 53 & ce chapitre \\
VI & Sanctions administratives, civiles et pénales & 54--71 & ce chapitre \\
VII & Dispositions diverses, transitoires et finales & 72--75 &
\ref{chap:prologue}, \ref{chap:secteurs} \\
\end{tableaucours}
\vspace{-0.8em}
\begin{center}{\small\itshape Tableau~6.2. Les sept titres de la loi n°~2024/017}\end{center}

La loi compte soixante-quinze articles. Elle « abroge toutes les
dispositions antérieures contraires » (article~75) et renvoie à des textes
particuliers les modalités d'application (article~74). Une quinzaine de ses
dispositions dépendent d'un texte réglementaire ou d'un référentiel de
l'Autorité~: critères de risque, évaluation d'impact, durées de
conservation, délais des droits, modalités d'autorisation et de contrôle.

% ============================================================================
\section{Ce qui Attend un Texte d'Application}
\label{sec:c06-application}
% ============================================================================

\begin{tableaucours}{L{4.6cm}cL{5.6cm}}{%
    \tableauHeader{Objet} &
    \tableauHeader{Art.} &
    \tableauHeader{Conséquence pour l'audit}
}
Création et organisation de l'Autorité & 53 & Pas de guichet, pas de
contrôle administratif \\
Modalités de l'autorisation & 19, 32 & Double cote~: exigence et diligence \\
Critères et modalités de l'évaluation du risque & 25 & Méthode déclarée par
l'organisation \\
Évaluation d'impact & 33 & Critères de risque élevé déclarés \\
Durées de conservation & 13, 28 & Durées fixées et justifiées par
l'organisation \\
Délais d'examen des demandes & 47 & Critère de substitution déclaré \\
Référentiel des mesures & 24, 27, 57 & Contenu du rapport annuel inconnu \\
Mécanisme de certification & 34, 59 & Sanction de l'article~59 sans objet
pour l'instant \\
Modalités du contrôle & 35 & Sans objet pour l'instant \\
Liste des pays, clauses types & 32, 53 & Deux vérifications de l'article~32
non cotables \\
Textes des secteurs & 4, 72 & Loi générale appliquée, réserve signalée \\
\end{tableaucours}
\vspace{-0.8em}
\begin{center}{\small\itshape Tableau~6.3. Les dispositions suspendues à un
texte (état en novembre~2026, tableau à refaire à chaque mission)}\end{center}

Chaque ligne de ce tableau produira dans les chapitres suivants un \NEK{} de
critère, c'est-à-dire une incertitude qui tient au droit et non à
l'organisation (Chapitre~\ref{chap:audit}). L'auditeur refait le tableau au
début de chaque mission. Un texte publié entre-temps transforme un \NEK{} en
exigence cotable.

% ============================================================================
\section{Le Régime d'Autorisation Préalable}
\label{sec:c06-autorisation}
% ============================================================================

\begin{boxjuris}
\textbf{\loicam{n°~2024/017}, article~19}~\cite{loi2024017}~:
\textit{« (1) Le traitement des données à caractère personnel est soumis à
l'obtention préalable d'une autorisation délivrée par l'Autorité de
protection des données à caractère personnel. (2) Tout processus
d'interconnexion et d'interopérabilité des fichiers de données sensibles,
relatives aux mineurs, est soumis à l'obtention préalable d'une
autorisation délivrée par l'Autorité de protection des données à caractère
personnel. (3) Les modalités de délivrance de l'autorisation visée aux
alinéas 1 et 2 ci-dessus sont fixées par voie réglementaire. »}

\smallskip
\textbf{Portée.} La loi ne prévoit ni déclaration simplifiée ni dispense~:
tout traitement est soumis à autorisation. L'article~49 interdit le
traitement sans autorisation préalable, et l'article~55 le punit d'une
pénalité de 5 à 50~millions de francs CFA.
\end{boxjuris}

Le régime est l'inverse du modèle européen, qui a supprimé les formalités
préalables au profit de la responsabilisation du responsable. Il rapproche
la loi camerounaise des premières lois africaines et françaises, fondées
sur un contrôle \textit{a priori}.

Ce régime pose une difficulté pratique~: l'autorisation est délivrée par une
Autorité dont « la création, ainsi que l'organisation et le
fonctionnement » sont fixés par décret (article~53, alinéa~2), et selon des
modalités elles-mêmes réglementaires (article~19, alinéa~3). Tant que ces
textes ne sont pas publiés et que l'Autorité ne reçoit pas de demandes,
l'organisation ne peut pas obtenir d'autorisation. Que doit coter
l'auditeur~?

\begin{itemize}
    \item « Le traitement est autorisé »~: \NON{}, puisqu'aucune autorisation
          n'existe. La cote est exacte~; elle ne dit pas que l'écart est
          imputable à l'organisation.
    \item « L'organisation a fait ce qui dépendait d'elle »~: \OUI{} si elle
          a préparé le dossier (registre, évaluations, mesures) et peut
          prouver qu'elle le déposera dès que la procédure sera ouverte~;
          \NON{} sinon.
\end{itemize}

Le rapport présente les deux cotes ensemble. À NYANGA, la seconde est
\NON{}~: aucun dossier n'a été préparé.

% ============================================================================
\section{L'Autorité de Protection}
\label{sec:c06-autorite}
% ============================================================================

\begin{boxjuris}
\textbf{\loicam{n°~2024/017}, article~53, alinéa~1 (extraits)}~\cite{loi2024017}~:
\textit{« L'Autorité de protection des données à caractère personnel est un
organisme public indépendant chargé, notamment~: de veiller à l'application
de la présente loi, des textes pris pour son application ainsi que des
conventions internationales en la matière~; de délivrer les autorisations
prévues par la présente loi, assorties des cahiers de charges
conséquents~; d'élaborer, publier et mettre à jour le référentiel des
mesures techniques et organisationnelles des données à caractère
personnel~; [\ldots] de traiter les réclamations, pétitions et plaintes
[\ldots]~; d'établir et publier la liste des pays reconnus comme offrant un
niveau de protection des données à caractère personnel équivalent aux
exigences prévues par le droit camerounais~; [\ldots] »}
\end{boxjuris}

Trois instruments de l'Autorité conditionnent l'audit~: le
\textbf{référentiel des mesures} (articles~24, 27, 28, 57), les
\textbf{cahiers des charges} attachés aux autorisations (article~61), et la
\textbf{liste des pays} offrant une protection équivalente
(Chapitre~\ref{chap:transferts}). Tant qu'ils n'existent pas, les exigences
qui en dépendent sont des \NEK{} de critère.

\begin{boxinfo}
\textbf{État de l'Autorité au moment de la rédaction.} À la date de ce cours
(novembre~2026), l'auteur n'a pas trouvé de texte publié organisant
l'Autorité, et ses coauteurs et lui relèvent dans leurs travaux en cours que
ce texte était attendu après l'échéance du 23~juin~2026. Le lecteur
vérifiera la situation au moment où il lit~: si l'Autorité est organisée,
une partie des \NEK{} de critère de ce livre se lève.
\end{boxinfo}

% ============================================================================
\section{Contrôle et Certification}
\label{sec:c06-certification}
% ============================================================================

L'article~35 confie à l'Autorité le contrôle des traitements, selon des
modalités réglementaires. L'article~34 lui demande en outre de mettre en
place une certification.

\begin{boxjuris}
\textbf{\loicam{n°~2024/017}, article~34}~\cite{loi2024017}~:
\textit{« L'Autorité de protection des données à caractère personnel met en
place un mécanisme de certification du traitement des données à caractère
personnel conforme aux principes et exigences énoncés par la présente
loi. »}

\smallskip
\textbf{Portée.} L'article~5 définit la certification comme un outil
permettant de communiquer sur le niveau de protection « au regard des
critères d'un référentiel approuvé au préalable par l'Autorité ». Il s'agit
d'une certification de conformité à la loi, distincte des certifications
ISO. L'article~59 punit le responsable qui « opère sans certification
prévue dans la présente loi »~: tant que le mécanisme n'existe pas, la
portée de cette sanction est un \NEK{} de critère.
\end{boxjuris}

\begin{tableaucours}{L{3.2cm}L{4.4cm}L{4.4cm}}{%
    \tableauHeader{} &
    \tableauHeader{Certification de l'article~34} &
    \tableauHeader{Certification ISO/IEC 27701}
}
Délivrée par & L'Autorité, ou selon son mécanisme & Un organisme de
certification accrédité \\
Référentiel & Approuvé par l'Autorité & Norme internationale \\
Objet & Traitement, processus, produit & Système de management sur un
périmètre \\
Portée juridique & Prévue par la loi (articles~34, 59) & Aucune au regard de
la loi \\
État en novembre~2026 & Mécanisme non publié & Disponible \\
\end{tableaucours}
\vspace{-0.8em}
\begin{center}{\small\itshape Tableau~6.4. Deux certifications à ne pas confondre (voir aussi le Chapitre~\ref{chap:normes})}\end{center}

% ============================================================================
\section{Les Sanctions}
\label{sec:c06-sanctions}
% ============================================================================

\subsection{Sanctions administratives}

\begin{boxjuris}
\textbf{\loicam{n°~2024/017}, article~54}~\cite{loi2024017}~:
\textit{« (1) Lorsque l'Autorité de protection des données à caractère
personnel constate le non-respect des obligations mises à la charge du
responsable du traitement ou de son sous-traitant, celui-ci est mis en
demeure de s'y conformer dans un délai maximum de dix (10) jours. (2) Passé
le délai prévu à l'alinéa 1 ci-dessus, l'Autorité de protection des données
à caractère personnel adresse une injonction de mise en conformité du
traitement, sous astreinte dont le montant ne peut excéder cent mille
(100~000) francs CFA par jour de retard. [\ldots] »}

\smallskip
\textbf{Portée.} À défaut de mise en conformité, l'Autorité peut suspendre
l'activité, retirer l'autorisation ou interdire tout traitement
(alinéa~3).
\end{boxjuris}

\begin{tableaucours}{cL{6.4cm}L{4.0cm}}{%
    \tableauHeader{Art.} &
    \tableauHeader{Manquement} &
    \tableauHeader{Montant (F CFA)}
}
55 & Traitement sans autorisation préalable & 5 à 50~millions \\
56 & Refus de communiquer à la personne les informations la concernant & 1
à 10~millions \\
57 & Non-respect du référentiel & 5 à 20~millions \\
58 & Non-respect des règles d'interconnexion & 1 à 5~millions \\
59 & Traitement sans la certification prévue & 5 à 15~millions \\
60 & Transfert vers un pays tiers sans autorisation & 10 à 50~millions \\
61 & Non-respect du cahier des charges & 10 à 100~millions \\
\end{tableaucours}
\vspace{-0.8em}
\begin{center}{\small\itshape Tableau~6.5. Pénalités et amendes administratives (articles~55 à~61)}\end{center}

\subsection{Sanctions civiles et pénales}

L'article~62 ouvre à la personne concernée le référé « en cas d'atteinte
grave aux droits », éventuellement sous astreinte, et l'action en
réparation. Les articles~63 à~70 créent des délits~; l'article~71 rend les
personnes morales pénalement responsables, avec une amende de 50~millions à
un milliard de francs CFA.

\begin{tableaucours}{cL{6.6cm}L{3.8cm}}{%
    \tableauHeader{Art.} &
    \tableauHeader{Délit} &
    \tableauHeader{Peine d'emprisonnement}
}
63 & Collecte ou accès frauduleux, déloyal ou illicite (peines doublées
avec verrouillage ou cryptage) & 2 à 5~ans \\
64 & Traitement maintenu malgré l'opposition, refus de rectifier ou
d'effacer & 1 à 3~ans \\
65 & Traitement à des fins de profilage & 3 à 10~ans \\
66 & Données sensibles sans consentement exprès~; récupération
frauduleuse de données supprimées & 6~mois à 2~ans \\
67 & Détournement de finalité, traitement ultérieur incompatible & 6~mois
à 2~ans \\
68 & Divulgation portant atteinte à la dignité ou à la vie privée (sur
plainte) & 6~mois à 2~ans \\
69 & Transfert vers l'étranger en méconnaissance de la loi & 3 à 10~ans \\
70 & Entrave à l'action de l'Autorité & 6~mois à 3~ans \\
\end{tableaucours}
\vspace{-0.8em}
\begin{center}{\small\itshape Tableau~6.6. Les délits de la loi (amendes assorties, non reproduites)}\end{center}

L'article~65 mérite attention~: il punit « le responsable de traitement ou
le sous-traitant qui procède ou fait procéder à un traitement de données à
caractère personnel à des fins de profilage », sans condition apparente.
Sa portée exacte, pour une notation de crédit par exemple, est discutée au
Chapitre~\ref{chap:profilage}.

% ============================================================================
\section{Dispositions Transitoires et Abrogation}
\label{sec:c06-transitoires}
% ============================================================================

Trois articles ferment la loi. L'article~73 a donné dix-huit mois aux
organisations pour se conformer~; ce délai a expiré le 23~juin~2026
(Chapitre~\ref{chap:prologue}). L'article~74 renvoie les modalités
d'application à des textes particuliers. L'article~75 abroge « toutes les
dispositions antérieures contraires ».

La formule de l'article~75 ne nomme aucun texte. Elle laisse donc en vigueur
tout ce qui, dans les lois antérieures, n'est pas contraire à la loi
nouvelle. La loi n°~2010/012 relative à la cybersécurité et à la
cybercriminalité contient une quinzaine d'articles sur la vie privée, la
conservation des données et les infractions qui s'y rapportent. Lesquels
subsistent~? Aucune liste officielle ne le dit. C'est à l'auditeur et au
juriste de le déterminer article par article, et c'est l'objet du
Chapitre~\ref{chap:loi2010}.

% ============================================================================
\section{Ce que la Loi demande à chaque Métier}
\label{sec:c06-metiers}
% ============================================================================

\begin{boxmetier}{Juriste conformité}
Lire le texte officiel, distinguer les obligations applicables, les
obligations suspendues à un texte d'application et les bonnes pratiques
importées, et le dire dans chaque avis.
\end{boxmetier}

\begin{boxmetier}{Auditeur de la protection des données}
Coter séparément « l'exigence est satisfaite » et « l'organisation a fait ce
qui dépendait d'elle » quand l'écart tient à l'absence de l'Autorité.
\end{boxmetier}

\begin{boxmetier}{Responsable de traitement}
Préparer dès maintenant les dossiers d'autorisation~: registre, fondements,
évaluations, mesures. Le jour où la procédure ouvrira, ce sont eux qui
seront demandés.
\end{boxmetier}

\begin{boxresume}
\begin{itemize}
    \item La loi compte sept titres et soixante-quinze articles~; une
          quinzaine de dispositions attendent un texte d'application.
    \item Tout traitement est soumis à autorisation préalable de
          l'Autorité (article~19)~; tant qu'elle n'est pas organisée,
          l'auditeur cote aussi ce qui dépendait de l'organisation.
    \item Les sanctions vont de la mise en demeure à dix jours jusqu'à un
          milliard de francs CFA pour une personne morale.
    \item La loi internalise les principes internationaux et y ajoute
          l'autorisation préalable, un consentement de principe, une liste
          large de données sensibles et de nombreux délits.
    \item L'article~75 n'abroge que les dispositions contraires~: la loi
          n°~2010/012 subsiste pour le reste.
\end{itemize}
\end{boxresume}

\begin{boxinfo}
\textbf{Points de contrôle.} Vous avez compris ce chapitre si vous savez~:
\begin{itemize}
    \item citer trois ajouts de la loi par rapport à ses sources~;
    \item expliquer comment coter l'absence d'autorisation quand l'Autorité
          n'est pas organisée~;
    \item rattacher un écart à l'article qui le sanctionne.
\end{itemize}
\end{boxinfo}

\begin{boxexo}[title={\ding{45}\quad Exercice~6.1. Quel article \niveaubase}]
Indiquez l'article qui sanctionne~: (a) un transfert vers un hébergeur
étranger sans autorisation~; (b) un agent qui récupère des fichiers
effacés de la corbeille d'un collègue~; (c) un établissement qui refuse de
communiquer à un client ses données~; (d) une entreprise qui refuse de
recevoir les agents de l'Autorité.
\end{boxexo}

\begin{boxexo}[title={\ding{45}\quad Exercice~6.2. Sans Autorité \niveauexpert}]
Un client de MBOLO dépose plainte auprès de l'Autorité en septembre~2026.
L'Autorité n'est pas encore organisée. (a) Quelles autres voies s'offrent
à lui~? (b) Quel est l'intérêt, pour MBOLO, d'avoir préparé ses dossiers
d'autorisation~? (c) Comment l'auditeur rend-il compte de cette
situation dans son rapport~?
\end{boxexo}

\begin{boxexo}[title={\ding{45}\quad Exercice~6.3. Reprise ou ajout \niveaumoyen}]
Pour chaque disposition, dites si la loi reprend une source internationale
(laquelle) ou innove~: (a) la portabilité de l'article~43~; (b) les
transactions bancaires parmi les données sensibles~; (c) l'autorisation
préalable de l'article~19~; (d) la convention publique de l'article~31~;
(e) le droit d'accès de l'article~39.
\end{boxexo}

\begin{boxexo}[title={\ding{45}\quad Exercice~6.4. Le tableau des textes attendus \niveauexpert}]
Un décret organisant l'Autorité paraît pendant votre mission.
\begin{enumerate}[(a)]
    \item Quelles lignes du Tableau~6.3 changent, et lesquelles ne changent
          pas~?
    \item Quelles cotes \NEK{} de vos papiers de travail devez-vous
          reprendre~?
    \item Rédigez le paragraphe du rapport qui explique au commanditaire
          pourquoi des cotes ont changé sans que l'organisation ait rien
          fait.
\end{enumerate}
\end{boxexo}

\bigskip

\begin{boxinfo}
\textbf{Pour aller plus loin.}
\begin{itemize}
    \item La loi n°~2024/017~\cite{loi2024017}, dans son texte officiel,
          titres~V à~VII.
    \item La Convention de Malabo~\cite{ua2014malabo}, chapitre~II, pour
          mesurer ce que la loi en a repris.
\end{itemize}
\end{boxinfo}

\begin{boxinfo}
\textbf{Transition.} La loi de 2024 n'est pas seule. Une loi de 2010, écrite
pour la sécurité des réseaux, impose encore de conserver, de protéger et de
ne pas intercepter. C'est l'objet du Chapitre~\ref{chap:loi2010}.
\end{boxinfo}

        % ============================================================================
% Fichier  : partie3/P03_C07_loi_2010.tex
% Titre    : La Loi n° 2010/012 et les Textes Connexes
% Auteur   : MINKA MI NGUIDJOÏ Thierry Emmanuel
% Version  : 3.0 -- novembre 2026
% Corrigés : corriges/Corr_C07.tex
% ============================================================================
% Sources : loi n° 2010/012 du 21 décembre 2010 (texte de la bibliothèque du
% cours), articles 13, 14, 24 à 32, 35, 41 à 51, 74, 91 : cités mot pour mot
% quand ils sont entre guillemets. Loi n° 2024/017, texte officiel. Les
% textes connexes du Tableau 7.4 portent chacun leur degré de vérification ;
% ceux qui sont marqués « à vérifier » sont cités de mémoire et ne doivent
% pas servir de critère d'audit avant lecture du texte.
% ============================================================================

\chapter{La Loi n° 2010/012 et les Textes Connexes}
\label{chap:loi2010}
\chaptermark{La loi n° 2010/012}

\epigraph{%
    « Une loi nouvelle n'efface pas les anciennes.\\
    Elle oblige à les relire. »%
}{Principe du cours}

\niveauChapitre{Fondamental}{%
    Chapitre~\ref{chap:loi}~; droit du numérique (UE MSN-2S3-JUR01)
}

\begin{boxobjectifs}
\begin{itemize}
    \item Présenter ce que la loi n°~2010/012 impose en matière de sécurité,
          de conservation et de vie privée.
    \item Confronter la conservation décennale de la loi de 2010 à la durée
          limitée de la loi de 2024.
    \item Rapprocher les délits de l'article~74 de ceux de la loi de 2024.
    \item Recenser les textes connexes applicables à une organisation, avec
          leur degré de vérification.
    \item Choisir les critères d'un audit et dresser la carte
          d'applicabilité d'une organisation.
\end{itemize}
\end{boxobjectifs}

\bigskip

% ============================================================================
\section*{Avant de commencer : « dix ans, c'est la loi »}
% ============================================================================

Bafoussam, mercredi 5~août~2026. Mme~BELLO demande à M.~KAMGA combien de
temps NYANGA CRÉDIT garde les journaux de connexion de son application
mobile. « Dix ans. C'est la loi de 2010. » Il sort le texte et montre
l'article~29.

Me~ETOUNDI, qui assiste à l'entretien, ouvre l'autre loi. L'article~13 de la
loi de 2024 limite la conservation à « la période nécessaire » au
traitement. Les journaux contiennent l'adresse IP et l'identifiant de chaque
client à chaque connexion. Dix ans de ces données sont-ils nécessaires~?

Les deux ont un texte en main. L'auditrice note que la question n'est pas de
savoir qui a raison, mais de savoir si l'article~29 vise bien un
établissement de microfinance et ce que l'établissement a écrit pour
justifier son choix. Il n'a rien écrit.

\begin{boxalert}
\textbf{Violation typique.} Auditer la protection des données contre la
seule loi n°~2024/017, comme si elle avait remplacé tout ce qui la
précédait.

\smallskip
\textbf{Règle~:} l'article~75 de la loi de 2024 n'abroge que les dispositions
contraires. L'auditeur recense les textes antérieurs encore applicables, les
inscrit parmi ses critères et signale les points où deux textes se
rencontrent sans qu'une règle dise lequel l'emporte.
\end{boxalert}

% ============================================================================
\section{Une Loi de Sécurité qui Parle de Données}
\label{sec:c07-architecture}
% ============================================================================

La loi n°~2010/012 du 21~décembre~2010~\cite{loi2010012} régit la sécurité des
réseaux de communications électroniques et des systèmes d'information, et
réprime la cybercriminalité. Elle n'est pas une loi de protection des
données. Elle en a pourtant tenu lieu pendant quatorze ans, par quelques
articles dispersés.

\begin{tableaucours}{cL{5.4cm}cL{4.0cm}}{%
    \tableauHeader{Titre} &
    \tableauHeader{Objet} &
    \tableauHeader{Art.} &
    \tableauHeader{Intérêt pour ce cours}
}
I & Dispositions générales, définitions & 1--5 & Vocabulaire~: données de
connexion, données de trafic \\
II & Cybersécurité~: régulation, audit, certification électronique,
protection des réseaux, des systèmes et de la vie privée & 6--51 & Sécurité
obligatoire, conservation, vie privée, interception \\
III & Cybercriminalité~: procédure, infractions & 52--89 & Enquête, délits
sur les données (art.~74) \\
IV & Coopération et entraide judiciaire internationales & 90--94 & Demandes
venues de l'étranger \\
V & Dispositions transitoires et finales & 95 et s. & Textes
d'application \\
\end{tableaucours}
\vspace{-0.8em}
\begin{center}{\small\itshape Tableau~7.1. La loi n°~2010/012 en cinq
titres}\end{center}

La loi confie la régulation de la sécurité électronique à l'Agence nationale
des technologies de l'information et de la communication (article~7). Tant
que l'Autorité de protection des données n'est pas organisée, l'Agence est
le seul régulateur en place dont les compétences touchent aux données
personnelles, par le biais de la sécurité.

% ============================================================================
\section{La Sécurité Obligatoire}
\label{sec:c07-securite}
% ============================================================================

Trois obligations de la loi de 2010 s'ajoutent à l'article~27 de la loi de
2024.

L'article~26 demande aux exploitants des systèmes d'information de se doter
« de systèmes normalisés leur permettant d'identifier, d'évaluer, de traiter
et de gérer continûment les risques liés à la sécurité ». C'est une
obligation de gestion des risques, que l'ISO/IEC~27001 permet de satisfaire.

L'article~30 leur demande d'évaluer et de réviser leurs mesures « en fonction
de l'évolution des technologies ».

Les articles~13 et~32 instituent un audit de sécurité obligatoire. Il vise
les réseaux et les systèmes d'information des opérateurs, des autorités de
certification et des fournisseurs de services de communications
électroniques. Il est périodique, conduit par l'Agence ou sous son contrôle
et ses rapports sont confidentiels (article~32).

\begin{boxinfo}
\textbf{Ce que l'auditeur de protection des données en fait.} Il ne refait
pas l'audit de sécurité. Il demande si l'organisation y est assujettie, la
date du dernier audit et la suite donnée à ses recommandations. Le rapport
est confidentiel~: l'auditeur peut ne recevoir qu'une attestation. Il cote
alors l'existence de l'audit, pas son contenu. Le cours d'audit de
cybersécurité (UE MSN-2S3-SEC08) traite l'audit lui-même.
\end{boxinfo}

% ============================================================================
\section{La Conservation Décennale}
\label{sec:c07-conservation}
% ============================================================================

Quatre articles de la loi de 2010 imposent de conserver des données pendant
dix ans.

\begin{tableaucours}{cL{4.2cm}L{6.0cm}}{%
    \tableauHeader{Art.} &
    \tableauHeader{Qui} &
    \tableauHeader{Quoi}
}
25 & Opérateurs de réseaux, fournisseurs de services de communications
électroniques & Données de connexion et de trafic \\
29 & Exploitants des systèmes d'information & Données de connexion et de
trafic de leurs systèmes \\
35 & Fournisseurs d'accès et hébergeurs & Données d'identification des
auteurs de contenus \\
46 & Fournisseurs de contenus & Contenus et données stockées \\
\end{tableaucours}
\vspace{-0.8em}
\begin{center}{\small\itshape Tableau~7.2. Les quatre obligations de
conservation pendant dix ans}\end{center}

\begin{boxjuris}
\textbf{\loicam{n°~2010/012}, article~29, alinéa~1}~\cite{loi2010012}~:
\textit{« Les exploitants des systèmes d'information ont l'obligation de
conserver les données de connexion et de trafic de leurs systèmes
d'information pendant une période de dix (10) ans. »}

\smallskip
\textbf{Portée.} L'alinéa~2 ajoute que ces données « peuvent être
accessibles lors des investigations judiciaires ». La finalité est donc la
preuve pénale. La loi ne définit pas l'exploitant de système d'information~;
on ne sait pas avec certitude si une microfinance ou une clinique en est un
au sens de ce texte.
\end{boxjuris}

Une adresse IP et un identifiant de connexion sont des données personnelles
par la lettre de l'article~5 de la loi de 2024 (Chapitre~\ref{chap:donnees}).
Les journaux visés par l'article~29 sont donc des données personnelles, que
l'article~13 de la loi de 2024 demande de ne pas garder au-delà du
nécessaire.

Les deux textes ne sont pas incompatibles. L'article~9 de la loi de 2024
admet le traitement « nécessaire au respect d'une obligation légale », et
l'article~38 réserve « les exigences de conservation » prévues par d'autres
textes. Une obligation de conserver dix ans, si elle s'applique, fonde la
conservation. La difficulté est en amont~: savoir à qui elle s'applique.

L'auditeur procède en trois temps.

\begin{enumerate}
    \item Il demande à l'organisation si elle se considère comme opérateur,
          fournisseur, hébergeur ou exploitant au sens de la loi de 2010
          et sur quel fondement.
    \item Si elle répond oui, il vérifie que seules les données visées
          (connexion et trafic) sont gardées dix ans, que leur accès est
          restreint et que leur finalité reste la preuve.
    \item Si elle répond non ou ne sait pas, il cote \NEK{} de critère
          l'application de l'article~29 et vérifie que la durée retenue est
          écrite et justifiée au regard de l'article~13.
\end{enumerate}

Dans les trois cas, des journaux gardés dix ans et consultés pour surveiller
les clients ou le personnel sont détournés de leur finalité. Cet écart-là ne
dépend d'aucune interprétation.

% ============================================================================
\section{La Vie Privée dans les Communications}
\label{sec:c07-vieprivee}
% ============================================================================

La section~IV du chapitre~IX (articles~41 à~48) protège la vie privée. Elle
reste applicable en ce qu'elle n'est pas contraire à la loi de 2024 et elle
contient des règles que celle-ci n'a pas reprises.

\begin{tableaucours}{cL{5.6cm}L{4.8cm}}{%
    \tableauHeader{Art.} &
    \tableauHeader{Règle (paraphrasée)} &
    \tableauHeader{Traitement concerné}
}
41 & Droit au respect de la vie privée~; mesures conservatoires du juge &
Tous \\
42 & Confidentialité des communications et des données de trafic &
Messagerie, application, téléphonie \\
44 & Interdiction d'écouter, d'intercepter ou de stocker les communications
sans consentement, sauf autorisation légale & Enregistrement des appels,
supervision de la messagerie \\
45 & Enregistrement autorisé dans le cadre professionnel pour fournir la
preuve d'une communication & Service client, ordres par téléphone \\
46 (2) & Filtres contre les atteintes aux données personnelles & Sites et
applications ouverts au public \\
47 & Consentement préalable pour lire ou écrire dans l'équipement terminal &
Traceurs, kits de mesure d'audience \\
48 & Prospection électronique~: émetteur identifié, adresse valide pour
demander l'arrêt & SMS et courriels commerciaux \\
\end{tableaucours}
\vspace{-0.8em}
\begin{center}{\small\itshape Tableau~7.3. La vie privée dans la loi
n°~2010/012}\end{center}

\begin{boxjuris}
\textbf{\loicam{n°~2010/012}, article~47}~\cite{loi2010012}~:
\textit{« L'utilisation des réseaux de communications électroniques et des
systèmes d'information aux fins de stocker les informations ou d'accéder à
des informations stockées dans un équipement terminal d'une personne
physique ou morale, ne peut se faire qu'avec son consentement préalable. »}

\smallskip
\textbf{Portée.} Un kit de mesure d'audience intégré à une application
mobile, qui lit l'identifiant du téléphone, accède à des informations
stockées dans l'équipement terminal. Il exige un consentement préalable,
distinct de celui que la loi de 2024 demande pour le traitement lui-même.
C'est le cas du traitement que l'audit découvrira dans l'application de
NYANGA (Chapitre~\ref{chap:cartographie}).
\end{boxjuris}

L'article~45 autorise l'enregistrement des communications « effectué dans le
cadre professionnel en vue de fournir la preuve numérique d'une
communication électronique ». L'enregistrement des appels du service client
de NYANGA peut s'appuyer sur ce texte s'il sert à prouver les opérations
convenues par téléphone. Il ne le peut pas s'il sert à évaluer les agents ou
à entraîner un logiciel. La finalité décide, comme à l'article~11 de la loi
de 2024.

% ============================================================================
\section{Interception, Réquisitions et Entraide}
\label{sec:c07-enquete}
% ============================================================================

L'organisation auditée détient des données que les enquêteurs peuvent
demander. La loi de 2010 fixe le cadre de ces demandes. Les officiers de
police judiciaire et les agents assermentés de l'Agence conduisent les
enquêtes selon le Code de procédure pénale (article~52). En cas de crime ou
de délit prévu par la loi, l'officier de police judiciaire peut intercepter
et enregistrer des communications (article~49) et les personnels des
opérateurs sont tenus au secret sur les réquisitions reçues (article~51).
Les demandes venues de l'étranger passent par le ministère chargé des
relations extérieures, sauf convention contraire (article~91)~; la
Convention de Budapest en est une (Chapitre~\ref{chap:international}).

L'auditeur ne juge pas la régularité d'une enquête. Il vérifie que
l'organisation sait recevoir une demande~: qui la reçoit, comment son auteur
et son fondement sont vérifiés, quelles données sont remises, quelle trace
en est gardée. Une remise de données sur simple appel téléphonique d'un
prétendu enquêteur est une violation de données
(Chapitre~\ref{chap:violations}).

% ============================================================================
\section{Deux Lois, des Délits Voisins}
\label{sec:c07-delits}
% ============================================================================

L'article~74 de la loi de 2010 punit plusieurs atteintes aux données
personnelles. Il a été écrit à une époque où aucune formalité préalable
n'existait encore au Cameroun et il en parle pourtant.

\begin{boxjuris}
\textbf{\loicam{n°~2010/012}, article~74, alinéas~3 et~7}~\cite{loi2010012}~:
\textit{« (3) Est puni d'un emprisonnement d'un (01) à trois (03) ans et
d'une amende de 1.000.000 (un million) à 5.000.000 (cinq millions) F CFA ou
de l'une de ces deux peines seulement, quiconque procède ou fait procéder,
même par négligence au traitement des données à caractère personnel en
violation des formalités préalables à leur mise en œuvre. [\ldots]
(7) Est puni d'un emprisonnement de six (06) mois à deux (02) ans et d'une
amende de 5.000.000 (cinq millions) à 50.000.000 (cinquante millions) F
CFA, ou de l'une de ces deux peines seulement, celui qui conserve des
informations sous une forme nominative ou chiffrée au-delà de la durée
légale indiquée dans la demande d'avis ou la déclaration préalable à la mise
en œuvre du traitement automatisé. »}

\smallskip
\textbf{Portée.} L'alinéa~3 punit le traitement fait sans formalité
préalable, « même par négligence ». Depuis 2024, cette formalité existe~:
c'est l'autorisation de l'article~19. L'alinéa~7 punit la conservation
au-delà de la durée déclarée, délit que la loi de 2024 ne reprend pas en ces
termes.
\end{boxjuris}

\begin{tableaucours}{L{4.6cm}cL{4.8cm}c}{%
    \tableauHeader{Comportement} &
    \tableauHeader{2010} &
    \tableauHeader{Loi n°~2024/017} &
    \tableauHeader{2024}
}
Traitement sans formalité préalable & 74 (3) & Traitement sans autorisation
(pénalité administrative) & 55 \\
Collecte par des moyens illicites & 74 (4) & Collecte frauduleuse, déloyale
ou illicite & 63 \\
Données révélant l'origine tribale, les opinions, les mœurs, sans accord
exprès & 74 (5) & Données sensibles sans consentement exprès & 66 \\
Détournement des informations & 74 (6) & Détournement de finalité & 67 \\
Conservation au-delà de la durée déclarée & 74 (7) & Pas de délit
identique~; durée limitée & 13 \\
Divulgation portant atteinte à la considération & 74 (8) & Divulgation
portant atteinte à la dignité ou à la vie privée & 68 \\
Interception de données personnelles en transmission & 74 (2) & Accès
frauduleux & 63 \\
\end{tableaucours}
\vspace{-0.8em}
\begin{center}{\small\itshape Tableau~7.4. Concordance indicative des délits
des deux lois}\end{center}

Les peines diffèrent d'une colonne à l'autre. Quand deux textes punissent le
même fait de peines différentes, lequel s'applique~? L'article~75 de la loi
de 2024 abroge les dispositions « contraires », sans dire si une peine
différente pour le même fait est une contrariété. La réponse appartient au
juge pénal.

L'auditeur n'a pas à la donner. Il rattache l'écart aux deux articles,
signale la concurrence et classe l'incertitude sur la peine parmi les
\NEK{} de critère. L'écart lui-même reste un \NON{}~: les deux lois le
réprouvent.

% ============================================================================
\section{Les Textes Connexes}
\label{sec:c07-connexes}
% ============================================================================

D'autres textes touchent aux données personnelles d'une organisation
camerounaise. Le tableau en donne la liste de travail du cours, avec pour
chacun ce que l'auteur a pu vérifier.

\begin{tableaucours}{L{4.6cm}L{4.4cm}L{3.0cm}}{%
    \tableauHeader{Texte} &
    \tableauHeader{Ce qu'il apporte} &
    \tableauHeader{Vérification}
}
Loi n°~2010/013 sur les communications électroniques, modifiée en 2015 &
Secret des correspondances, obligations des opérateurs envers les abonnés &
Texte en bibliothèque, à lire \\
Décret n°~2015/3759 sur l'identification des abonnés & Collecte de
l'identité de tout abonné & Texte en bibliothèque, à lire \\
Loi-cadre n°~2011/012 sur la protection du consommateur & Information et
loyauté envers le consommateur & Texte en bibliothèque, à lire \\
Code de procédure pénale, articles~102 et~154 & Secret de l'enquête et de
l'information judiciaire & D'après~\cite{minka2026chainepenale} \\
Code pénal & Secret professionnel, violation de correspondance & À
vérifier \\
Loi sur le secret bancaire & Confidentialité des informations des clients &
À vérifier \\
Règlements de la CEMAC et de la COBAC & Conservation et communication des
données des clients, lutte contre le blanchiment, services de paiement & À
identifier \\
Texte sur la protection des enfants en ligne & Obligations envers les
mineurs & À vérifier \\
Code de déontologie médicale & Secret médical & À vérifier \\
\end{tableaucours}
\vspace{-0.8em}
\begin{center}{\small\itshape Tableau~7.5. Les textes connexes et leur degré
de vérification (novembre~2026)}\end{center}

La dernière colonne n'est pas une précaution de style. Un texte « à
vérifier » ne peut pas figurer parmi les critères d'un audit~: l'auditeur le
lit d'abord dans sa version officielle. S'il ne le trouve pas, il l'écrit et
cote \NEK{} de critère les exigences qui en dépendraient. Le
Chapitre~\ref{chap:secteurs} reprend ces textes secteur par secteur.

% ============================================================================
\section{Choisir ses Critères}
\label{sec:c07-criteres}
% ============================================================================

Cinq chapitres ont présenté des textes de nature et de portée différentes.
Avant d'ouvrir une mission, l'auditeur décide lesquels serviront de
critères. La règle tient en quatre lignes.

\begin{boxprocedure}
\textbf{Choisir les critères d'un audit de protection des données.}
\begin{enumerate}
    \item \textbf{L'obligatoire d'abord.} La loi n°~2024/017, la loi
          n°~2010/012 pour ce qui subsiste, les textes sectoriels et
          communautaires applicables à l'organisation. Ils entrent toujours
          dans les critères.
    \item \textbf{L'étranger opposable ensuite.} Le RGPD pour les seuls
          traitements qui remplissent son article~3.
    \item \textbf{Le contractuel.} Les clauses imposées par un client, un
          prestataire ou un investisseur, si le commanditaire veut qu'elles
          soient auditées.
    \item \textbf{Le volontaire déclaré.} Une norme, un cadre ou un repère
          étranger, retenu comme référentiel ou comme critère de
          substitution et nommé comme tel.
\end{enumerate}
\end{boxprocedure}

Le résultat se présente dans une \termecle{carte d'applicabilité}~: un
tableau qui dit, pour chaque texte, son statut à l'égard de l'organisation.

\begin{tableaucours}{L{4.0cm}L{2.6cm}L{2.6cm}L{2.8cm}}{%
    \tableauHeader{Texte} &
    \tableauHeader{NYANGA} &
    \tableauHeader{MBOLO} &
    \tableauHeader{NOUN PAY}
}
Loi n°~2024/017 & Obligatoire & Obligatoire & Obligatoire \\
Loi n°~2010/012, art.~41 à~48 et~74 & Obligatoire & Obligatoire &
Obligatoire \\
Loi n°~2010/012, art.~29 & À qualifier & À qualifier & À qualifier \\
Textes CEMAC et COBAC & Obligatoires & Sans objet & Obligatoires \\
Secret médical & Sans objet & Obligatoire & Sans objet \\
RGPD & Sans objet & Sans objet & Obligatoire pour NOUN DIASPORA~;
contractuel avec l'hébergeur \\
Cadre du NIST & Sans objet & Sans objet & Contractuel (investisseur) \\
ISO/IEC~27701 & Sans objet & Sans objet & Volontaire (projet de
certification) \\
Convention de Malabo & Comparatif & Comparatif & Comparatif \\
\end{tableaucours}
\vspace{-0.8em}
\begin{center}{\small\itshape Tableau~7.6. La carte d'applicabilité des
trois missions}\end{center}

La carte se joint au plan d'audit (Chapitre~\ref{chap:audit}). Elle protège
l'auditeur contre deux reproches~: avoir oublié un texte applicable, ou
avoir reproché à l'organisation de ne pas respecter un texte qui ne la lie
pas.

% ============================================================================
\section{La Loi de 2010 à NYANGA}
\label{sec:c07-cas}
% ============================================================================

\begin{boxscenario}{Mission NYANGA. Les exigences de la loi n° 2010/012, août 2026}
\textbf{Constatations.}
\begin{itemize}
    \item L'application mobile intègre un kit de mesure d'audience qui lit
          l'identifiant du téléphone dès le premier lancement, avant tout
          écran d'information.
    \item Tous les appels au service client sont enregistrés et gardés deux
          ans. Le chef du service en écoute un échantillon chaque mois pour
          noter les agents.
    \item Les SMS de prospection n'indiquent aucun moyen de demander leur
          arrêt.
    \item Des relevés de compte sont envoyés aux clients depuis la
          messagerie personnelle des agents.
    \item Une analyse des risques de sécurité a été refaite après l'incident
          de 2025 et mise à jour en mars~2026.
    \item Les journaux de connexion de l'application sont gardés dix ans.
          Aucune note ne dit pourquoi. L'établissement n'a jamais fait
          l'objet d'un audit de l'Agence et ne sait pas s'il y est
          assujetti.
\end{itemize}

\textbf{Analyse} (confiance élevée sur les faits, moyenne sur
l'assujettissement). Quatre écarts ne dépendent d'aucune interprétation. Les
trois points incertains tiennent tous à la même question~: NYANGA CRÉDIT
est-elle un exploitant de système d'information au sens de la loi de 2010~?
\end{boxscenario}

\begin{boxlecture}{NYANGA CRÉDIT, exigences de la loi n° 2010/012}
\begin{tableaucours}{L{7.6cm}cL{3.4cm}}{%
    \tableauHeader{Exigence} &
    \tableauHeader{Cote} &
    \tableauHeader{Preuve}
}
Le kit de mesure d'audience recueille un consentement préalable (art.~47) &
\NON{} & Analyse de l'application \\
L'enregistrement des appels sert la seule preuve (art.~45) & \NON{} &
Entretien, grille de notation \\
Les messages de prospection permettent de demander l'arrêt (art.~48) &
\NON{} & Échantillon de 20~SMS \\
Les communications avec les clients sont confidentielles (art.~42) & \NON{} &
Messageries personnelles \\
Les risques de sécurité sont gérés de façon continue (art.~26) & \OUI{} &
Analyse de mars~2026 \\
L'audit de sécurité obligatoire a été réalisé (art.~32) & \NEK{} &
Assujettissement non établi \\
L'obligation de conserver dix ans s'applique (art.~29) & \NEK{} & Qualité
d'exploitant non établie \\
La durée de dix ans est justifiée au regard de l'article~13 de la loi de
2024 & \NEK{} & Dépend de la ligne précédente \\
\end{tableaucours}

\medskip
\textbf{Décompte.} Un \OUI{}, quatre \NON{}, trois \NEK{} sur huit.
\hfill\lecture{0{,}125}{0{,}375}{1{,}41}

\medskip
\textbf{Lecture.} Domaine opaque. Les trois \NEK{} sont de critère et se
lèveraient ensemble si la notion d'exploitant était précisée. Les quatre
\NON{} appellent des mesures immédiates, dont le retrait du kit de mesure
d'audience.
\end{boxlecture}

% ============================================================================
\section{Ce que les Textes demandent à chaque Métier}
\label{sec:c07-metiers}
% ============================================================================

\begin{boxmetier}{Juriste conformité}
Tenir la liste des textes applicables avec leur degré de vérification.
Écrire, pour chaque point où deux textes se rencontrent, la lecture retenue
par l'organisation et sa justification.
\end{boxmetier}

\begin{boxmetier}{RSSI}
Séparer les journaux gardés pour la preuve de ceux qui servent
l'exploitation. Restreindre l'accès aux premiers et documenter leur durée.
\end{boxmetier}

\begin{boxmetier}{Auditeur de la protection des données}
Joindre la carte d'applicabilité au plan d'audit. Rattacher chaque écart à
tous les articles qui le visent, dans les deux lois.
\end{boxmetier}

\begin{boxresume}
\begin{itemize}
    \item La loi n°~2010/012 reste applicable pour tout ce qui n'est pas
          contraire à la loi de 2024~: sécurité, conservation, vie privée
          dans les communications, délits.
    \item Elle impose une gestion continue des risques de sécurité
          (article~26) et un audit de sécurité obligatoire pour certains
          acteurs (articles~13 et~32).
    \item Quatre articles imposent de conserver des données dix ans. Ils
          fondent la conservation s'ils s'appliquent~; il faut d'abord
          établir à qui.
    \item L'article~47 exige un consentement préalable pour tout accès à
          l'équipement terminal~; l'article~45 limite l'enregistrement des
          communications à la preuve.
    \item L'auditeur choisit ses critères dans l'ordre~: obligatoire,
          étranger opposable, contractuel, volontaire déclaré. Il les
          présente dans une carte d'applicabilité.
\end{itemize}
\end{boxresume}

\begin{boxinfo}
\textbf{Points de contrôle.} Vous avez compris ce chapitre si vous savez~:
\begin{itemize}
    \item citer trois règles de la loi de 2010 absentes de la loi de 2024~;
    \item expliquer comment coter la conservation de journaux pendant dix
          ans~;
    \item dresser la carte d'applicabilité d'une organisation.
\end{itemize}
\end{boxinfo}

\begin{boxexo}[title={\ding{45}\quad Exercice~7.1. Quel article \niveaubase}]
Indiquez l'article de la loi n°~2010/012 qui s'applique~: (a) une
application installe un traceur sans rien demander~; (b) un établissement
envoie des SMS publicitaires sans numéro pour se désinscrire~; (c) un
responsable écoute les appels de ses agents à leur insu~; (d) une entreprise
garde les enregistrements d'appels comme preuve des ordres reçus.
\end{boxexo}

\begin{boxexo}[title={\ding{45}\quad Exercice~7.2. Les appels enregistrés \niveaumoyen}]
Le service client de NYANGA enregistre tous les appels, les conserve deux
ans et les écoute pour évaluer les agents. Analysez ce traitement au regard
des articles~44 et~45 de la loi n°~2010/012 et des articles~9 et~11 de la
loi n°~2024/017. Proposez une règle conforme.
\end{boxexo}

\begin{boxexo}[title={\ding{45}\quad Exercice~7.3. La carte de MBOLO \niveaumoyen}]
Le CENTRE MÉDICAL MBOLO envisage d'accueillir des patients envoyés par une
organisation humanitaire établie à Genève, qui lui impose par contrat ses
propres règles de protection des données. Refaites la colonne MBOLO du
Tableau~7.6. Quels textes changent de statut et lesquels restent à
vérifier avant d'entrer dans les critères~?
\end{boxexo}

\begin{boxexo}[title={\ding{45}\quad Exercice~7.4. Dix ans ou le nécessaire \niveauexpert}]
M.~KAMGA vous demande une réponse écrite~: NYANGA doit-elle garder ses
journaux de connexion dix ans~?
\begin{enumerate}[(a)]
    \item Rédigez la note, en séparant ce qui est établi par les textes de
          ce qui relève d'une interprétation.
    \item Proposez une politique de conservation des journaux qui soit
          défendable dans les deux lectures de l'article~29.
    \item Dites quelles cotes de l'encadré de lecture changeraient si
          l'Agence confirmait par écrit que l'établissement est un
          exploitant au sens de la loi.
\end{enumerate}
\end{boxexo}

\bigskip

\begin{boxinfo}
\textbf{Pour aller plus loin.}
\begin{itemize}
    \item La loi n°~2010/012~\cite{loi2010012}, articles~24 à~51 et~74, dans
          son texte officiel.
    \item T.~E.~Minka Mi Nguidjoï, R.~T.~Massoussi Songo et
          D.~Bialo~\cite{minka2026chainepenale}, pour l'articulation des
          deux lois dans la chaîne pénale.
\end{itemize}
\end{boxinfo}

\begin{boxinfo}
\textbf{Transition.} Le contexte est posé~: ses sources, le règlement
européen, les référentiels volontaires, les deux lois camerounaises et la
manière de choisir entre eux. Il reste à s'en servir. La Partie~IV ouvre la
mission et commence par la manière de conduire un audit. C'est l'objet du
Chapitre~\ref{chap:audit}.
\end{boxinfo}


    \part{Méthode : Auditer la Protection des Données}
    \label{part:methode}

        % ============================================================================
% Fichier  : partie4/P04_C08_conduire_audit.tex
% Titre    : Conduire un Audit de Protection des Données
% Auteur   : MINKA MI NGUIDJOÏ Thierry Emmanuel
% Version  : 3.0 -- novembre 2026
% Corrigés : corriges/Corr_C08.tex
% ============================================================================
% Sources : ISO 19011:2018 (principes et cycle, paraphrasés) ; ISO/IEC 27007
% (paraphrasée) ; auditologie (hal-05701007, hal-05701270) ; loi n° 2024/017,
% art. 35, 70 ; NRAC C06-C07 (quadrant, règle de décision).
% Faits nouveaux (bible) : plan d'audit de KORO AUDIT, six domaines,
% 2 juillet 2026.
% ============================================================================

\chapter{Conduire un Audit de Protection des Données}
\label{chap:audit}
\chaptermark{Conduire un audit}

\epigraph{%
    « Un audit ne dit pas si l'organisation est bonne.\\
    Il dit ce qui est établi, et sur quelle preuve. »%
}{Principe du cours}

\niveauChapitre{Fondamental}{%
    Chapitres~\ref{chap:donnees} et~\ref{chap:principes}~; méthode générale
    de l'audit (UE MSN-2S3-SEC08)
}

\begin{boxobjectifs}
\begin{itemize}
    \item Distinguer audit de conformité légale, audit contre une norme et
          audit de préparation.
    \item Rédiger le mandat, les critères et le périmètre d'une mission.
    \item Choisir des preuves selon leur force et documenter chaque constat
          dans un papier de travail.
    \item Décomposer la conformité en domaines et en exigences cotées
          \OUI{}, \NON{} ou \NEK{}.
    \item Lire les mesures $\mu$, $U$ et $H$ et appliquer une règle de
          décision écrite à l'avance.
    \item Protéger les données personnelles que l'audit lui-même manipule.
\end{itemize}
\end{boxobjectifs}

\bigskip

% ============================================================================
\section*{Avant de commencer : « montrez-nous la sécurité »}
% ============================================================================

Yaoundé, jeudi 2~juillet~2026. Dans les bureaux de KORO AUDIT \&
CONFORMITÉ, Mme~BELLO prépare la mission de NYANGA CRÉDIT avec son associé
juriste, M.~NGANDO. Le DSI de l'établissement a envoyé la veille un dossier
de quarante pages~: schémas du réseau, règles du pare-feu, rapport du test
d'intrusion de 2025. « Voilà notre conformité », écrit-il.

M.~NGANDO le parcourt et le referme. Le dossier ne dit rien des
consentements, des durées, des droits des personnes, des contrats des
prestataires ni des transferts. Il répond à une question que personne n'a
posée. Mme~BELLO commence par écrire ce que la mission doit établir, avec
quels critères, sur quel périmètre, et ce qu'elle ne couvrira pas.

\begin{boxalert}
\textbf{Violation typique.} Confondre l'audit de protection des données avec
un audit de sécurité informatique, et conclure à la conformité parce que les
systèmes sont bien protégés.

\smallskip
\textbf{Règle~:} les critères de l'audit sont les obligations de la loi et,
le cas échéant, d'une norme désignée. La sécurité n'en est qu'un domaine
parmi six~; des données bien gardées peuvent avoir été collectées sans
droit.
\end{boxalert}

% ============================================================================
\section{Trois Sortes d'Audit}
\label{sec:c03-types}
% ============================================================================

\begin{tableaucours}{L{3.2cm}L{3.8cm}L{5.0cm}}{%
    \tableauHeader{Audit} &
    \tableauHeader{Critères} &
    \tableauHeader{Usage}
}
De conformité légale & Loi n°~2024/017, textes d'application, autres lois
applicables & Dire si l'organisation respecte ses obligations~; préparer un
contrôle ou un contentieux \\
Contre une norme & ISO/IEC~27701, ISO/IEC~27001 & Préparer ou maintenir une
certification (Chapitre~\ref{chap:normes}) \\
De préparation & Loi et bonnes pratiques & Mesurer l'écart avant une
échéance~; prioriser les chantiers \\
\end{tableaucours}
\vspace{-0.8em}
\begin{center}{\small\itshape Tableau~8.1. Trois audits, trois jeux de critères}\end{center}

Ces audits peuvent être conduits par l'organisation elle-même (audit
interne), par un cabinet qu'elle mandate (NYANGA, MBOLO), ou par un tiers qui
certifie. Le contrôle de l'Autorité de protection, prévu par la loi
(article~53) et dont les modalités relèvent de textes réglementaires
(article~35), n'est pas un audit au sens de ce chapitre~: c'est un contrôle
administratif, qui peut conduire à des sanctions (Chapitre~\ref{chap:loi}).

% ============================================================================
\section{Les Principes de l'Auditeur}
\label{sec:c03-principes}
% ============================================================================

La norme ISO~19011~\cite{iso19011_2018} pose sept principes pour tout audit
de système de management. Ils s'appliquent sans changement à l'audit de
protection des données~; on les résume ici en paraphrase.

\begin{enumerate}
    \item \textbf{Intégrité}~: honnêteté, diligence, respect des lois.
    \item \textbf{Présentation impartiale}~: rendre compte avec exactitude,
          y compris des difficultés rencontrées.
    \item \textbf{Conscience professionnelle}~: le soin et le jugement
          proportionnés à l'importance de la mission.
    \item \textbf{Confidentialité}~: la protection des informations
          obtenues, et ici des données personnelles consultées.
    \item \textbf{Indépendance}~: l'absence de conflit d'intérêts, fondement
          de l'objectivité des conclusions.
    \item \textbf{Approche fondée sur la preuve}~: des conclusions
          vérifiables, tirées d'un échantillon défini.
    \item \textbf{Approche par les risques}~: concentrer l'effort là où
          l'écart serait le plus grave pour les personnes.
\end{enumerate}

Le quatrième principe prend ici un relief particulier. Pour vérifier un
consentement, l'auditeur ouvre un dossier client~; pour tester les droits,
il lit des réponses adressées à des personnes~; pour contrôler les
journaux, il voit qui a consulté quoi. L'audit est lui-même un traitement~:
il a une finalité (la mission), une durée (celle de la conservation des
papiers de travail) et des destinataires (le commanditaire). L'auditeur
consulte sur place, copie le moins possible, masque les identités dans ses
papiers et détruit les copies à l'échéance prévue au mandat
(section~\ref{sec:c08-deontologie}).

% ============================================================================
\section{La Déontologie de l'Auditeur}
\label{sec:c08-deontologie}
% ============================================================================

L'auditeur de la protection des données est soumis aux principes communs de
l'audit de la section précédente. Son objet ajoute trois exigences.

\subsection{Minimiser les données de l'audit}

L'auditeur ne collecte que les données nécessaires à la preuve. Il consulte
sur place plutôt que de copier, relève des indicateurs plutôt que des
contenus, masque les identités dans ses papiers de travail (« client~A,
dossier 2018-0412 »), conserve les copies indispensables dans un espace
chiffré et les détruit à la date prévue au mandat. Un auditeur qui
emporterait un extrait du fichier clients sur une clé non chiffrée
commettrait l'écart même qu'il est chargé de relever.

\subsection{L'indépendance}

L'indépendance interdit d'auditer son propre travail. KORO AUDIT a proposé
au CENTRE MÉDICAL MBOLO, en juin, de rédiger ses procédures de protection
des données~; s'il accepte, il ne pourra pas, en octobre, auditer ces
mêmes procédures sans menacer son objectivité. Mme~BELLO a refusé la
mission de rédaction et conservé la mission d'audit.

\subsection{Ce que l'auditeur découvre}

L'auditeur peut découvrir des faits graves~: un agent qui vend des fichiers,
un traitement manifestement délictueux. Il en rend compte au commanditaire,
par écrit et sans délai, comme son mandat doit le prévoir. Il ne mène pas
d'enquête qui excéderait sa mission, et il ne saisit pas lui-même la
justice, sauf obligation légale qui lui serait propre. Le
Chapitre~\ref{chap:violations} montre comment un tel fait devient une
violation à traiter, et le cours de forensique comment il devient une
enquête.

% ============================================================================
\section{Mandat, Critères, Périmètre}
\label{sec:c03-mandat}
% ============================================================================

\begin{boxprocedure}
\textbf{Préparer une mission d'audit de protection des données.}
\begin{enumerate}
    \item Obtenir un mandat écrit~: commanditaire, objectif, accès accordés,
          destinataires du rapport, date de remise.
    \item Fixer les critères~: articles de la loi retenus, textes
          d'application publiés, norme éventuelle, politiques internes.
    \item Délimiter le périmètre~: entités, sites, traitements, période~;
          écrire ce qui est exclu, et pourquoi.
    \item Décomposer en domaines et en exigences, chacune formulée en
          question fermée.
    \item Choisir, pour chaque exigence, le test et la preuve attendue.
    \item Prévoir le traitement des données de l'audit~: consultation sur
          place, masquage, durée de conservation des papiers.
    \item Fixer à l'avance la règle de décision qui permettra de conclure à
          partir des cotes.
\end{enumerate}
\end{boxprocedure}

Les critères doivent être \textbf{publiés et applicables}. Un critère qui
renvoie à un texte réglementaire non publié, comme la durée de conservation
fixée « suivant les modalités fixées par voie réglementaire »
(article~13), ne peut pas être appliqué tel quel~: l'auditeur le remplace
par un critère de substitution déclaré (« une durée fixée, justifiée et
appliquée par l'organisation »), ou il cote \NEK{} en précisant que
l'incertitude tient au critère et non à la preuve. Les deux sources
d'incertitude se distinguent dans le rapport~: un \NEK{} \textit{de
preuve} se lève en cherchant une pièce, un \NEK{} \textit{de critère} ne se
lève qu'avec un texte.

% ============================================================================
\section{Les Preuves et leur Force}
\label{sec:c03-preuves}
% ============================================================================

\begin{tableaucours}{cL{3.4cm}L{4.4cm}L{3.4cm}}{%
    \tableauHeader{Rang} &
    \tableauHeader{Type de preuve} &
    \tableauHeader{Exemple} &
    \tableauHeader{Limite}
}
1 & Réexécution, test & Demande d'accès test suivie jusqu'à la réponse &
Coûteuse, à encadrer \\
2 & Trace produite par le système & Journal d'accès, horodatage de
consentement & Dépend de la configuration \\
3 & Document signé et daté & Contrat de sous-traitance, autorisation &
Prouve l'engagement, pas son exécution \\
4 & Observation & Affichage en agence, armoire ouverte & Instantanée \\
5 & Déclaration & Entretien, questionnaire & À corroborer \\
\end{tableaucours}
\vspace{-0.8em}
\begin{center}{\small\itshape Tableau~8.2. Hiérarchie indicative des preuves}\end{center}

Une déclaration seule ne permet jamais de coter \OUI{}~; elle oriente la
recherche d'une preuve de rang supérieur. Un document de rang~3 établit
l'existence d'une règle, pas son application~: un contrat qui impose des
mesures de sécurité au prestataire ne prouve pas que le prestataire les
applique.

Chaque constat se documente dans un \termecle{papier de travail}~:
l'exigence, le critère, le test, l'échantillon, la preuve (référence de la
pièce), la cote, la date, l'auteur et le réviseur. Le papier de travail
permet à un autre auditeur de refaire le raisonnement, et à l'audité de
contester un constat précis plutôt qu'une impression.

\begin{boxinfo}
\textbf{La demande test.} Pour vérifier le traitement des demandes de
droits, l'auditeur peut adresser, avec l'accord écrit du commanditaire, une
demande d'accès au nom d'un membre de l'équipe qui est réellement client.
Une demande au nom d'une personne fictive, ou d'un client à son insu,
serait une collecte déloyale. Le mandat doit prévoir ce test.
\end{boxinfo}

% ============================================================================
\section{Décomposer et Coter}
\label{sec:c03-decomposer}
% ============================================================================

La conformité à une loi de soixante-quinze articles ne se cote pas d'un
bloc. On la décompose en \termecle{domaines}, puis en \termecle{exigences},
chacune posée comme une question fermée à laquelle une preuve peut répondre
\OUI{}, \NON{} ou \NEK{}.

\begin{tableaucours}{cL{3.6cm}L{3.4cm}L{4.4cm}}{%
    \tableauHeader{Code} &
    \tableauHeader{Domaine} &
    \tableauHeader{Articles} &
    \tableauHeader{Exemple d'exigence}
}
D1 & Recensement et fondements & 9--14, 19, 29 & Le registre existe-t-il~? \\
D2 & Droits des personnes & 21, 37--47 & Les demandes reçoivent-elles une
réponse~? \\
D3 & Sécurité et confidentialité & 15, 24, 25, 27 & Les accès sont-ils
tracés~? \\
D4 & Sous-traitants et transferts & 16, 26, 30--32 & Chaque prestataire
a-t-il un contrat conforme~? \\
D5 & Violations & 22 & Les violations sont-elles notifiées~? \\
D6 & Gouvernance & 27(2), 33, 35 & L'évaluation d'impact est-elle faite
quand le risque est élevé~? \\
\end{tableaucours}
\vspace{-0.8em}
\begin{center}{\small\itshape Tableau~8.3. Les six domaines d'un audit de protection des données}\end{center}

Une exigence formulée avec « et » ou « ou » se scinde~: « le registre existe
et il est à jour » contient deux questions, qui peuvent recevoir deux cotes
différentes. La règle d'atomicité évite qu'un \NON{} ne désigne deux défauts
à la fois.

% ============================================================================
\section{Lire les Mesures, Conclure}
\label{sec:c03-conclure}
% ============================================================================

\begin{boxdef}
Pour un ensemble de $n$ exigences cotées, avec $n_O$ réponses \OUI{}, $n_N$
\NON{} et $n_K$ \NEK{}~:

\smallskip
\termecle{Taux de conformité}~: $\mu = n_O / n$.
\termecle{Taux d'incertitude}~: $U = n_K / n$.
\termecle{Entropie}~: $H = -\sum p \log_2 p$, sur les trois cotes, en bits,
entre~0 et $\log_2 3 \approx 1{,}585$.
\end{boxdef}

\begin{boxprocedure}
\textbf{La règle de décision du cours.}
\begin{enumerate}
    \item \textbf{$U$ d'abord.} Si $U \geq 0{,}20$, le domaine est opaque~:
          aucune conclusion de conformité. On recherche les preuves, ou l'on
          conclut sur l'opacité elle-même.
    \item \textbf{$\mu$ ensuite.} Si $U < 0{,}20$, le domaine est clair~:
          conforme si $\mu \geq 0{,}5$ et sans \NON{} sur une exigence
          pénalement sanctionnée~; non conforme sinon.
    \item \textbf{$H$ pour contrôler.} Une entropie élevée avec $U$ faible
          signale un désaccord entre \OUI{} et \NON{}~: on décompose avant
          de conclure.
    \item \textbf{Les cotes prévalent.} Un seul \NON{} grave (collecte
          déloyale, transfert illicite, données sensibles sans
          consentement) suffit à motiver une recommandation prioritaire,
          quelles que soient les mesures.
\end{enumerate}
\end{boxprocedure}

Les seuils sont ceux du cours~; chaque mission peut en fixer d'autres, à
condition de les écrire dans le plan d'audit avant de coter. La quatrième
règle est propre à la protection des données~: la loi sanctionne
pénalement certains écarts (articles~63 à~71), et un taux de conformité
élevé ne compense pas une infraction.

% ============================================================================
\section{Le Plan d'Audit de NYANGA}
\label{sec:c03-nyanga}
% ============================================================================

\begin{boxscenario}{Mission NYANGA. Le plan d'audit de KORO AUDIT, 2 juillet 2026}
\textbf{Mandat.} Conseil d'administration de NYANGA CRÉDIT, 30~juin~2026~:
« dire si l'établissement respecte la loi n°~2024/017 et recommander les
mesures nécessaires ». Rapport au conseil le 25~septembre.

\textbf{Critères.} Les articles~2 à~52 de la loi n°~2024/017~; les
articles~26, 41 à~48 et~74 de la loi n°~2010/012~; les textes d'application
publiés à la date de la mission (aucun identifié)~; la politique de sécurité
interne de 2025. La carte d'applicabilité du Chapitre~\ref{chap:loi2010} est
jointe au plan.

\textbf{Périmètre.} Siège et trois agences tirées au sort (Dschang,
Mbouda, Foumban)~; les onze familles de données du
Chapitre~\ref{chap:donnees}~; période du 1\textsuperscript{er}~janvier~2025
au 31~août~2026. \textbf{Exclusions déclarées~:} l'enquête pénale sur
l'incident de mars~2025 et le contrôle de l'éditeur du logiciel bancaire
dans ses propres locaux, à Abidjan.

\textbf{Données de l'audit.} Consultation dans les locaux de NYANGA~;
papiers de travail avec identités masquées~; conservation trois ans, puis
destruction attestée.

\textbf{Règle de décision.} Celle du cours, seuils inchangés.
\end{boxscenario}

\begin{boxlecture}{Le domaine D1 après la première semaine}
\begin{tableaucours}{L{7.8cm}cL{3.2cm}}{%
    \tableauHeader{Exigence} &
    \tableauHeader{Cote} &
    \tableauHeader{Preuve}
}
Les familles de données sont recensées & \OUI{} & Inventaire du 7~juillet \\
Un registre des activités de traitement est tenu & \NON{} & Aucun registre \\
Chaque traitement a un fondement valable & \NON{} & Formulaires, cahiers \\
L'autorisation de l'Autorité est demandée & \NON{} & Attestation \\
Une durée de conservation est fixée & \NON{} & Aucune règle \\
L'information de l'article~21 est fournie & \NON{} & Formulaire 2016 \\
L'identification des clients repose sur une obligation légale & \NEK{} &
Texte non produit \\
La vidéosurveillance repose sur un fondement valable & \NEK{} & Fondement
incertain \\
\end{tableaucours}

\medskip
\textbf{Décompte.} Un \OUI{}, cinq \NON{}, deux \NEK{} sur huit.
\hfill\lecture{0{,}125}{0{,}25}{1{,}30}

\medskip
\textbf{Lecture.} $U = 0{,}25 \geq 0{,}20$~: selon la première règle, le
domaine est opaque et ne permet aucune conclusion de conformité. Mais la
quatrième règle s'applique déjà~: cinq écarts sont établis, dont
l'absence d'autorisation, sanctionnée par l'article~55. La conclusion
provisoire est une non-conformité établie sur cinq exigences, avec deux
exigences à instruire.
\end{boxlecture}

% ============================================================================
\section{Ce que la Méthode demande à chaque Métier}
\label{sec:c03-metiers}
% ============================================================================

\begin{boxmetier}{Auditeur de la protection des données}
Écrire le mandat, les critères, le périmètre et la règle de décision avant
le premier entretien. Distinguer dans le rapport les \NEK{} de preuve et les
\NEK{} de critère.
\end{boxmetier}

\begin{boxmetier}{Responsable de traitement}
Donner à l'auditeur un accès réel aux pièces, et nommer un interlocuteur
par domaine. L'entrave à un contrôle de l'Autorité est un délit
(article~70)~; l'entrave à un audit interne prive surtout l'organisation
de sa propre information.
\end{boxmetier}

\begin{boxmetier}{RSSI}
Préparer les preuves de sécurité sous une forme exploitable~: journaux
exportés et datés, listes d'accès, rapports de sauvegarde. Un schéma de
réseau ne prouve aucune mesure de protection des données.
\end{boxmetier}

\begin{boxresume}
\begin{itemize}
    \item L'audit de protection des données a pour critères la loi et, le
          cas échéant, une norme~; la sécurité n'en est qu'un domaine.
    \item Les sept principes de l'ISO~19011 s'appliquent~; la
          confidentialité impose de protéger les données que l'audit
          manipule.
    \item Un critère renvoyant à un texte non publié produit un \NEK{} de
          critère, distinct d'un \NEK{} de preuve.
    \item La conformité se décompose en six domaines et en exigences
          atomiques cotées \OUI{}, \NON{} ou \NEK{}.
    \item La règle de décision lit $U$ d'abord, $\mu$ ensuite, et un seul
          écart grave suffit à motiver une recommandation prioritaire.
\end{itemize}
\end{boxresume}

\begin{boxinfo}
\textbf{Points de contrôle.} Vous avez compris ce chapitre si vous savez~:
\begin{itemize}
    \item rédiger le mandat et le périmètre d'une mission~;
    \item classer des preuves selon leur force~;
    \item appliquer la règle de décision à un domaine coté.
\end{itemize}
\end{boxinfo}

\begin{boxexo}[title={\ding{45}\quad Exercice~8.1. Quelle preuve \niveaubase}]
Pour chaque exigence, proposez la preuve la plus forte possible~: (a) les
salariés ont été informés du traitement de leurs données~; (b) le
prestataire de recouvrement efface les dossiers soldés~; (c) les accès aux
dossiers clients sont limités aux agents de l'agence~; (d) les demandes
d'accès reçoivent une réponse.
\end{boxexo}

\begin{boxexo}[title={\ding{45}\quad Exercice~8.2. Rendre atomique \niveaumoyen}]
Réécrivez en exigences atomiques~: (a) « le registre existe et il est tenu à
jour par un responsable désigné »~; (b) « les données sont chiffrées ou
pseudonymisées »~; (c) « les contrats des prestataires prévoient la
sécurité, la confidentialité et la restitution des données ».
\end{boxexo}

\begin{boxexo}[title={\ding{45}\quad Exercice~8.3. Conclure sur D2 \niveaumoyen}]
Le domaine D2 de NYANGA compte dix exigences~: 2~\OUI{}, 5~\NON{}, 3~\NEK{}.
Calculez $\mu$, $U$ et $H$, appliquez la règle de décision et rédigez la
phrase de conclusion du domaine.
\end{boxexo}

\begin{boxexo}[title={\ding{45}\quad Exercice~8.4. L'audit comme traitement \niveauexpert}]
Rédigez la clause du mandat qui encadre les données personnelles consultées
par l'auditeur~: finalité, accès, copies, masquage, conservation,
destruction. Dites comment l'auditeur prouvera lui-même qu'il l'a respectée.
\end{boxexo}

\begin{boxexo}[title={\ding{45}\quad Exercice~8.5. Une découverte \niveauexpert}]
Pendant l'audit de MBOLO, un auditeur junior découvre qu'un agent exporte
chaque semaine la liste des patientes de la maternité vers une adresse de
messagerie personnelle. Rédigez la note au commanditaire, et dites ce que
l'auditeur doit faire et ne pas faire.
\end{boxexo}

\bigskip

\begin{boxinfo}
\textbf{Pour aller plus loin.}
\begin{itemize}
    \item L'ISO~19011~\cite{iso19011_2018} et l'ISO/IEC~27007~\cite{iso27007},
          pour la conduite des audits.
    \item T.~E.~Minka Mi Nguidjoï, « Vers l'auditologie »~\cite{minka2026auditologie},
          pour la décomposition et les mesures d'incertitude.
\end{itemize}
\end{boxinfo}

\begin{boxinfo}
\textbf{Transition.} La première exigence de D1 est le registre. Comment
recenser tous les traitements d'une organisation, y compris ceux que la
direction ignore~? C'est l'objet du Chapitre~\ref{chap:cartographie}.
\end{boxinfo}

        % ============================================================================
% Fichier  : partie4/P04_C09_cartographier.tex
% Titre    : Cartographier les Traitements
% Auteur   : MINKA MI NGUIDJOÏ Thierry Emmanuel
% Version  : 1.0 -- novembre 2026
% Corrigés : corriges/Corr_C09.tex
% ============================================================================
% Sources : loi n° 2024/017, art. 5 (fichier, interconnexion), 29, 36
% (texte officiel) ; ISO/IEC 27701 (inventaire et registre, paraphrasés).
% Faits nouveaux (bible) : 47 traitements recensés, dont 9 inconnus de la
% direction ; découverte sur les partages et les postes, 20-24 juillet 2026.
% ============================================================================

\chapter{Cartographier les Traitements}
\label{chap:cartographie}
\chaptermark{Cartographier les traitements}

\epigraph{%
    « Le registre dit ce que l'organisation croit faire.\\
    La découverte dit ce qu'elle fait. »%
}{Principe du cours}

\niveauChapitre{Fondamental}{%
    Chapitres~\ref{chap:donnees} et~\ref{chap:audit}
}

\begin{boxobjectifs}
\begin{itemize}
    \item Énoncer le contenu obligatoire du registre des activités de
          traitement (article~29).
    \item Construire un registre à partir des processus, et non des
          applications.
    \item Confronter le registre déclaré à une découverte des données
          réellement stockées.
    \item Cartographier les flux de données, y compris vers les prestataires
          et l'étranger.
    \item Mesurer la complétude d'un registre et coter ce qui reste
          inconnu.
\end{itemize}
\end{boxobjectifs}

\bigskip

% ============================================================================
\section*{Avant de commencer : la liste noire de Foumban}
% ============================================================================

Foumban, mardi 21~juillet~2026. Dans l'agence, l'équipe de KORO AUDIT
examine avec l'accord du chef d'agence le contenu du partage réseau
commun. Elle y trouve un classeur intitulé \texttt{mauvais\_payeurs.xlsx}~:
2\,140~noms, numéros de téléphone, quartiers, montants impayés, et une
colonne « remarques » où figurent des mentions comme « famille du chef de
quartier, ne pas relancer » ou « commerçant bamoun, fiable malgré
retard ». Le fichier est partagé entre six agences par messagerie. Il ne
figure dans aucune liste de l'établissement.

Ce classeur est un traitement à part entière~: il a une finalité (décider à
qui prêter), des personnes concernées, des destinataires. Il contient des
données sensibles (une origine régionale et des informations bancaires), et
il sert à décider sans que les personnes le sachent.

\begin{boxalert}
\textbf{Violation typique.} Établir le registre par questionnaire auprès des
directions, puis le tenir pour complet.

\smallskip
\textbf{Règle~:} un registre déclaratif se confronte à une découverte des
données réellement stockées et des flux réellement observés. Ce que la
découverte trouve et que le registre ignore mesure l'incomplétude du
registre.
\end{boxalert}

% ============================================================================
\section{Le Registre de l'Article 29}
\label{sec:c04-registre}
% ============================================================================

\begin{boxjuris}
\textbf{\loicam{n°~2024/017}, article~29}~\cite{loi2024017}~:
\textit{« (1) Le responsable du traitement ou le sous-traitant tient un
registre physique ou numérique des activités de traitement effectuées sous
leur responsabilité. (2) Le registre visé à l'alinéa 1 ci-dessus comporte
les informations ci-après~: le nom et les coordonnées du responsable du
traitement et, le cas échéant, le nom du sous-traitant~; les finalités du
traitement~; la description des catégories de personnes concernées et des
catégories des données à caractère personnel~; les catégories de
destinataires auxquels les données à caractère personnel ont été ou seront
communiquées~; les documents attestant de l'existence de garanties
appropriées ou le numéro de l'autorisation de l'Autorité de protection des
données à caractère personnel. »}

\smallskip
\textbf{Portée.} L'obligation pèse sur le responsable \textit{et} sur le
sous-traitant. Le registre peut être papier ou numérique. Son dernier champ
relie le registre aux autorisations~: un traitement sans numéro
d'autorisation ni garanties documentées est, par construction, un écart.
\end{boxjuris}

Le contenu légal est un minimum. Un registre utile à l'audit et à la gestion
ajoute la base du traitement (consentement ou dérogation de l'article~9), la
durée de conservation, les transferts hors du Cameroun, les mesures de
sécurité, la date de la dernière revue et le propriétaire du traitement dans
l'organisation.

\begin{tableaucours}{L{4.0cm}cL{6.6cm}}{%
    \tableauHeader{Champ} &
    \tableauHeader{Obligatoire} &
    \tableauHeader{Ce que l'auditeur vérifie}
}
Responsable, sous-traitant & Oui & Coordonnées exactes~; sous-traitants
cohérents avec les contrats \\
Finalités & Oui & Une finalité par traitement, déterminée et explicite \\
Personnes et données & Oui & Catégories complètes, données sensibles
signalées \\
Destinataires & Oui & Comparés aux flux observés \\
Garanties ou autorisation & Oui & Numéro d'autorisation ou pièces
justificatives \\
Fondement & Recommandé & Cohérent avec l'article~9 \\
Durée de conservation & Recommandé & Cohérente avec l'article~13 \\
Transferts & Recommandé & Pays, garanties (Chapitre~\ref{chap:transferts}) \\
Mesures de sécurité & Recommandé & Renvoi aux mesures de l'article~27 \\
Propriétaire, date de revue & Recommandé & Une personne, une date récente \\
\end{tableaucours}
\vspace{-0.8em}
\begin{center}{\small\itshape Tableau~9.1. Les champs du registre}\end{center}

% ============================================================================
\section{Recenser par les Processus}
\label{sec:c04-processus}
% ============================================================================

L'erreur la plus courante consiste à partir de l'inventaire des
applications. Une application sert souvent plusieurs traitements (le
logiciel bancaire sert la tenue des comptes, le crédit, le recouvrement et
la lutte contre le blanchiment)~; à l'inverse, un traitement peut ne
reposer sur aucune application (un cahier, un classeur partagé, un groupe de
messagerie).

On part donc des \termecle{processus} de l'organisation~: pour chaque
direction, quelles activités, quelles personnes concernées, quelles
décisions prises sur ces personnes. Le Chapitre~\ref{chap:donnees} a fourni
les familles de données~; on les croise ici avec les processus pour obtenir
les traitements, puis on rattache à chacun les applications et les supports
qui le portent.

\begin{boxprocedure}
\textbf{Établir le registre d'une organisation.}
\begin{enumerate}
    \item Lister les processus de chaque direction, avec un interlocuteur.
    \item Pour chaque processus, identifier les personnes concernées et les
          décisions prises sur elles.
    \item En déduire les traitements, un par finalité.
    \item Rattacher à chaque traitement ses supports~: applications,
          fichiers bureautiques, papier, messagerie.
    \item Renseigner les champs du registre, avec la pièce qui fonde chaque
          champ.
    \item Confronter le registre à la découverte et aux flux
          (sections~\ref{sec:c04-decouverte} et~\ref{sec:c04-flux}).
    \item Faire valider chaque fiche par son propriétaire, et dater la
          revue.
\end{enumerate}
\end{boxprocedure}

% ============================================================================
\section{La Découverte des Données}
\label{sec:c04-decouverte}
% ============================================================================

La \termecle{découverte des données} consiste à chercher, dans les supports
réels, des données personnelles que le registre ne mentionne pas. Elle se
pratique sur quatre types de supports.

\begin{itemize}
    \item \textbf{Les partages de fichiers et les postes}~: recherche de
          motifs (numéros de téléphone, numéros de compte, numéros de carte
          nationale d'identité) dans les documents bureautiques.
    \item \textbf{Les bases de données}~: lecture des schémas, à la
          recherche de colonnes nommées \texttt{nom}, \texttt{tel},
          \texttt{cni}, \texttt{date\_naissance}.
    \item \textbf{La messagerie et les outils de collaboration}~: groupes,
          pièces jointes récurrentes, formulaires en ligne.
    \item \textbf{Le papier}~: armoires, cahiers, archives d'agence.
\end{itemize}

La découverte est elle-même un traitement de données personnelles. Elle se
fait avec l'accord écrit du commanditaire, sur les seuls supports de
l'organisation, en ne relevant que des indicateurs (nom du fichier,
emplacement, nombre d'occurrences, types de données) et jamais le contenu
lui-même, sauf échantillon nécessaire à la qualification.

\begin{boxinfo}
\textbf{Les faux positifs.} Une recherche de motifs trouve des numéros qui
n'en sont pas~: un numéro de facture peut ressembler à un numéro de
téléphone. On qualifie donc chaque fichier signalé par un examen rapide d'un
échantillon, et l'on cote la proposition « ce fichier contient des données
personnelles » \OUI{}, \NON{} ou \NEK{}. Le laboratoire APDC-LAB-02 exerce
cette qualification.
\end{boxinfo}

% ============================================================================
\section{Cartographier les Flux}
\label{sec:c04-flux}
% ============================================================================

Un registre décrit des traitements~; une \termecle{cartographie des flux}
décrit le chemin des données entre les traitements, les systèmes, les
prestataires et les pays. Elle répond à trois questions que le registre
laisse ouvertes~: par où les données entrent-elles, où sont-elles stockées,
et vers qui sortent-elles.

On l'établit à partir de trois sources indépendantes~: les entretiens
(ce que les équipes décrivent), les configurations (connecteurs,
exportations planifiées, interfaces de programmation) et les traces
réseau ou les factures (ce qui sort réellement, et vers qui l'organisation
paie un service). Une facture mensuelle d'un service de formulaires en
ligne révèle souvent un flux que personne n'avait déclaré.

L'article~36 encadre l'\termecle{interconnexion des fichiers}, définie à
l'article~5 comme la mise en relation de données traitées pour une
finalité avec d'autres données. Elle n'est licite que si elle est
« nécessaire au respect d'une obligation légale », et ne doit entraîner ni
discrimination ni réduction des droits. Le classeur des mauvais payeurs,
alimenté par six agences et rapproché du fichier des clients, est une
interconnexion~; l'auditeur en cherche l'obligation légale, et n'en trouve
pas.

% ============================================================================
\section{Le Registre de NYANGA}
\label{sec:c04-nyanga}
% ============================================================================

\begin{boxscenario}{Mission NYANGA. Quarante-sept traitements, 20 au 24 juillet 2026}
\textbf{Constatations.}
\begin{itemize}
    \item Les entretiens avec les huit directions produisent un registre
          déclaratif de 38~traitements.
    \item La découverte sur les partages du siège et des trois agences signale
          1\,312 fichiers contenant des motifs de données personnelles~;
          l'examen d'un échantillon de 120 fichiers en confirme 97.
    \item Neuf traitements apparaissent que personne n'a déclarés~: le
          classeur des mauvais payeurs~; deux groupes de messagerie où les
          agents échangent des photos de pièces d'identité~; un questionnaire
          de satisfaction hébergé par un service de formulaires étranger~; un
          kit de mesure d'audience intégré à l'application mobile, qui envoie
          les identifiants des téléphones à un tiers~; l'enregistrement des
          appels du service client~; le fichier des clients « privilégiés »
          du service marketing~; les badges d'accès du siège~; les CV reçus
          par messagerie.
    \item Les factures de 2025 révèlent trois prestataires absents de la
          liste des sous-traitants.
\end{itemize}

\textbf{Analyse} (confiance élevée). Le registre complet compte
47~traitements. Le registre déclaratif en couvrait 38, soit
81~\%~; les neuf manquants sont ceux qui transfèrent des données à des
tiers ou qui prennent des décisions sans procédure, c'est-à-dire les plus
risqués.
\end{boxscenario}

\begin{boxobjet}{Liste des mauvais payeurs, agences de l'Ouest}
\textbf{Finalité réelle~:} décider de l'octroi de crédit et de la relance.
\textbf{Personnes~:} 2\,140 débiteurs, dont des clients d'autres
établissements cités par des tiers. \textbf{Données~:} identité, téléphone,
quartier, montant impayé, remarques incluant l'origine régionale (donnée
sensible). \textbf{Destinataires~:} six agences, par messagerie.
\textbf{Ce que l'audit établit~:} aucun fondement (\NON{}), interconnexion
sans obligation légale (\NON{}), aucune information des personnes
(\NON{})~; usage pour des décisions de crédit (\OUI{}, entretiens et
échantillon de 12 refus).
\end{boxobjet}

\begin{boxlecture}{Le registre de NYANGA, 24 juillet 2026}
\begin{tableaucours}{L{7.8cm}cL{3.2cm}}{%
    \tableauHeader{Exigence} &
    \tableauHeader{Cote} &
    \tableauHeader{Preuve}
}
Un registre existe & \OUI{} & Établi par l'audit \\
Le registre est complet & \NON{} & 9 traitements découverts \\
Chaque fiche comporte les champs de l'article~29 & \NON{} & Garanties
absentes \\
Chaque traitement a un propriétaire & \NON{} & 14 sans propriétaire \\
Les interconnexions sont fondées sur une obligation légale & \NON{} &
Classeur des mauvais payeurs \\
Les flux vers les tiers sont tous connus & \NEK{} & Trois prestataires
découverts, recherche inachevée \\
Les fichiers papier des agences non visitées sont recensés & \NEK{} & 19
agences hors échantillon \\
\end{tableaucours}

\medskip
\textbf{Décompte.} Un \OUI{}, quatre \NON{}, deux \NEK{} sur sept.
\hfill\lecture{0{,}14}{0{,}29}{1{,}38}

\medskip
\textbf{Lecture.} Le registre existe désormais, mais c'est l'audit qui l'a
établi~: le \OUI{} ne crédite pas l'établissement. Avec $U = 0{,}29$, la
complétude reste opaque~; l'incertitude se localise dans les flux et dans
les dix-neuf agences non visitées.
\end{boxlecture}

% ============================================================================
\section{Ce que la Cartographie demande à chaque Métier}
\label{sec:c04-metiers}
% ============================================================================

\begin{boxmetier}{Auditeur de la protection des données}
Mesurer la complétude du registre par la découverte et les flux, et
rapporter le taux de couverture. Un registre à 81~\% n'est pas un registre
« presque complet »~: les traitements manquants sont rarement les moins
risqués.
\end{boxmetier}

\begin{boxmetier}{Délégué à la protection des données}
Tenir le registre comme un outil vivant~: une fiche par traitement, un
propriétaire, une date de revue, et une procédure qui oblige tout nouveau
projet à créer sa fiche avant de collecter.
\end{boxmetier}

\begin{boxmetier}{RSSI}
Fournir les sources de la découverte~: exports des schémas de bases,
journaux de flux sortants, inventaire des services en ligne payés. Bloquer
les outils de partage qui contournent les traitements déclarés.
\end{boxmetier}

\begin{boxresume}
\begin{itemize}
    \item L'article~29 impose un registre au responsable et au
          sous-traitant, avec cinq catégories d'informations au minimum.
    \item On recense par les processus et les personnes concernées, puis on
          rattache les applications et les supports.
    \item Le registre déclaratif se confronte à la découverte des données et
          aux flux réels~; l'écart mesure son incomplétude.
    \item L'interconnexion de fichiers n'est licite que si une obligation
          légale l'impose (article~36).
    \item À NYANGA, 47~traitements, dont 9 inconnus de la direction et parmi
          eux un fichier de mauvais payeurs mentionnant l'origine régionale.
\end{itemize}
\end{boxresume}

\begin{boxinfo}
\textbf{Points de contrôle.} Vous avez compris ce chapitre si vous savez~:
\begin{itemize}
    \item citer les champs obligatoires du registre~;
    \item expliquer pourquoi un inventaire d'applications n'est pas un
          registre~;
    \item mesurer la complétude d'un registre à partir d'une découverte.
\end{itemize}
\end{boxinfo}

\begin{boxexo}[title={\ding{45}\quad Exercice~9.1. Traitement ou application \niveaubase}]
Dans la liste suivante, distinguez les traitements des moyens
techniques~: (a) la paie~; (b) le serveur de messagerie~; (c) le
recrutement~; (d) la base Oracle~; (e) la vidéosurveillance des agences~;
(f) le logiciel bancaire.
\end{boxexo}

\begin{boxexo}[title={\ding{45}\quad Exercice~9.2. Une fiche de registre \niveaumoyen}]
Rédigez la fiche de registre du traitement « enregistrement des appels du
service client », avec les champs obligatoires et recommandés. Signalez les
champs que vous ne pouvez pas remplir sans pièce supplémentaire.
\end{boxexo}

\begin{boxexo}[title={\ding{45}\quad Exercice~9.3. Mesurer la complétude \niveaumoyen}]
La découverte d'une autre organisation signale 800 fichiers, dont un
échantillon de 50 donne 35 confirmations. Le registre déclaratif compte
20~traitements~; la découverte en ajoute 6. (a) Quel est le taux de
couverture du registre déclaratif~? (b) Combien de fichiers contiennent
probablement des données personnelles, et avec quelle réserve~?
\end{boxexo}

\begin{boxexo}[title={\ding{45}\quad Exercice~9.4. La liste noire \niveauexpert}]
Le directeur des crédits défend le classeur des mauvais payeurs~: « Sans
lui, nous prêterions à des fraudeurs. » Analysez sa licéité au regard des
articles~5, 9, 11, 17, 36 et~48. Proposez une solution qui réponde au
besoin de l'établissement dans le respect de la loi.
\end{boxexo}

\bigskip

\begin{boxinfo}
\textbf{Pour aller plus loin.}
\begin{itemize}
    \item L'ISO/IEC~27701~\cite{iso27701}, pour la tenue des inventaires et
          des registres dans un système de management de la vie privée.
\end{itemize}
\end{boxinfo}

\begin{boxinfo}
\textbf{Transition.} Quarante-sept traitements ne présentent pas le même
danger pour les personnes. Lesquels examiner d'abord, et lesquels exigent
une évaluation d'impact~? C'est l'objet du Chapitre~\ref{chap:risque}.
\end{boxinfo}

        % ============================================================================
% Fichier  : partie4/P04_C10_risque_personnes.tex
% Titre    : Évaluer le Risque pour les Personnes
% Auteur   : MINKA MI NGUIDJOÏ Thierry Emmanuel
% Version  : 1.0 -- novembre 2026
% Corrigés : corriges/Corr_C10.tex
% ============================================================================
% Sources : loi n° 2024/017, art. 5 (évaluation d'impact), 25, 33 (texte
% officiel) ; ISO/IEC 29134 (étapes, paraphrasées) ; méthode PIA de la CNIL
% (échelles à quatre niveaux, paraphrasées) ; ISO 31000:2018.
% Faits nouveaux (bible) : évaluation d'impact du projet de téléconsultation
% du CENTRE MÉDICAL MBOLO, avril 2026.
% ============================================================================

\chapter{Évaluer le Risque pour les Personnes}
\label{chap:risque}
\chaptermark{Le risque pour les personnes}

\epigraph{%
    « Le risque que mesure la protection des données\\
    est celui que court la personne, pas l'entreprise. »%
}{Principe du cours}

\niveauChapitre{Intermédiaire}{%
    Chapitres~\ref{chap:principes} et~\ref{chap:cartographie}~; notions de
    gestion des risques (UE MSN-2S3-SEC08)
}

\begin{boxobjectifs}
\begin{itemize}
    \item Distinguer le risque pour les personnes du risque pour
          l'organisation.
    \item Appliquer les articles~25 et~33~: évaluation du risque pour tout
          traitement, évaluation d'impact pour les traitements à risque
          élevé.
    \item Repérer les traitements à risque élevé à l'aide de critères
          déclarés.
    \item Conduire une évaluation d'impact selon les étapes de
          l'ISO/IEC~29134~: description, nécessité, risques, mesures,
          décision.
    \item Auditer une évaluation d'impact existante.
\end{itemize}
\end{boxobjectifs}

\bigskip

% ============================================================================
\section*{Avant de commencer : « le risque, c'est la panne »}
% ============================================================================

Bertoua, lundi 6~avril~2026. Le CENTRE MÉDICAL MBOLO\footnote{Le CENTRE
MÉDICAL MBOLO, sa directrice, son personnel, ses patients et son prestataire
MEDISOFT SARL sont fictifs.} prépare un service de téléconsultation pour
les villages de la Kadey~: un infirmier équipé d'une tablette examine le
patient au centre de santé du village, et un médecin de Bertoua pose le
diagnostic par visioconférence. La directrice, le Dr~ABENA, a demandé au
responsable informatique, M.~YOUSSOUFA, une analyse des risques. Il revient
avec un tableau de quinze lignes~: coupure du réseau, panne de la tablette,
vol du matériel, perte de recettes, retard du prestataire.

Aucune ligne ne parle des patients~: de la consultation filmée qu'un tiers
pourrait voir, du diagnostic d'une maladie stigmatisante transmis par un
réseau partagé, de l'enfant examiné hors de la présence de ses parents. Le
tableau mesure ce que l'établissement peut perdre. La loi demande ce que
les personnes peuvent subir.

\begin{boxalert}
\textbf{Violation typique.} Présenter une analyse des risques pour
l'organisation (pannes, pertes financières) comme l'évaluation exigée par la
loi.

\smallskip
\textbf{Règle~:} l'évaluation d'impact apprécie « la probabilité et la
gravité des risques encourus par les droits et libertés des individus »
(article~5). Ses événements redoutés concernent les personnes~: la
divulgation, la modification, la disparition de leurs données, et leurs
conséquences sur leur vie.
\end{boxalert}

% ============================================================================
\section{Deux Obligations}
\label{sec:c05-obligations}
% ============================================================================

La loi impose deux évaluations distinctes.

\begin{boxjuris}
\textbf{\loicam{n°~2024/017}, article~25, alinéa~1}~\cite{loi2024017}~:
\textit{« Le responsable du traitement et le sous-traitant effectuent
l'évaluation du risque présenté par le traitement des données à caractère
personnel, suivant les critères d'évaluation, ainsi que les modalités
d'examen et de validation fixées par voie règlementaire. »}

\medskip
\textbf{\loicam{n°~2024/017}, article~33}~\cite{loi2024017}~:
\textit{« (1) Tout type de traitement susceptible d'engendrer un risque
élevé pour les droits et libertés des personnes physiques est subordonné à
la réalisation préalable d'une évaluation d'impact des opérations de
traitement envisagées sur la protection des données à caractère
personnel. (2) Les conditions et les modalités de la réalisation de
l'évaluation d'impact visées à l'alinéa 1 ci-dessus sont fixées par voie
règlementaire. »}

\smallskip
\textbf{Portée.} Tout traitement doit être évalué (article~25)~; seuls ceux
qui présentent un risque élevé exigent une évaluation d'impact préalable
(article~33). Dans les deux cas, critères et modalités relèvent de textes
d'application. En leur absence, l'auditeur vérifie que l'organisation a
déclaré et appliqué une méthode reconnue~; il cote \NEK{} de critère la
conformité aux modalités réglementaires elles-mêmes.
\end{boxjuris}

\begin{boxdef}
\termecle{Évaluation d'impact sur la protection des données}~: selon
l'article~5, « procédure visant à analyser la probabilité et la gravité des
risques encourus par les droits et libertés des individus résultant du
traitement de données à caractère personnel ».
\end{boxdef}

% ============================================================================
\section{Repérer le Risque Élevé}
\label{sec:c05-eleve}
% ============================================================================

La loi ne définit pas le « risque élevé ». En attendant le texte
d'application, le cours retient une liste de critères déclarés, inspirée des
pratiques internationales. Un traitement qui en réunit deux ou plus est
présumé à risque élevé et soumis à évaluation d'impact.

\begin{enumerate}
    \item Des données sensibles au sens de l'article~5 (santé, biométrie,
          transactions bancaires, origine régionale\ldots).
    \item Des personnes vulnérables~: mineurs, patients, salariés face à
          leur employeur, personnes en difficulté financière.
    \item Une décision automatisée ou un profilage produisant des effets sur
          la personne (Chapitre~\ref{chap:profilage}).
    \item Une surveillance systématique~: vidéosurveillance, géolocalisation,
          contrôle des communications.
    \item Un traitement à grande échelle, au regard du nombre de personnes
          ou de la durée.
    \item Le croisement ou l'interconnexion de fichiers (article~36).
    \item Une technologie nouvelle ou un usage nouveau d'une technologie.
    \item Un transfert hors du Cameroun (Chapitre~\ref{chap:transferts}).
\end{enumerate}

Ces critères sont une convention du cours, déclarée comme telle dans le
plan d'audit~; ils ne remplacent pas le texte d'application quand il sera
publié. La téléconsultation de MBOLO en réunit cinq~: données de santé,
patients vulnérables dont des mineurs, technologie nouvelle pour
l'établissement, transmission par un réseau d'opérateur, hébergement chez
un prestataire.

\begin{boxinfo}
\textbf{Le risque de vie privée selon le NIST et l'ISO/IEC~27557.} Deux
référentiels volontaires aident à décrire ce risque sans le confondre avec
le risque de sécurité.

\smallskip
Le NIST part des \textit{actions problématiques sur les données}~: une
opération peut nuire à une personne alors même que le système fonctionne
comme prévu et qu'aucune attaque n'a eu lieu~\cite{nistir8062}. Le dossier
d'une patiente lu par une secrétaire qui en avait le droit technique en est
un exemple. Le NIST en tire trois objectifs d'ingénierie qui complètent
confidentialité, intégrité et disponibilité~: la prévisibilité (la personne
peut s'attendre à ce qui est fait de ses données), la maîtrise (les données
peuvent être corrigées, effacées, communiquées de façon sélective) et la
dissociabilité (le traitement n'associe les données à une personne que si
c'est nécessaire).

\smallskip
L'ISO/IEC~27557:2022~\cite{iso27557} applique à la vie privée la démarche
de gestion du risque de l'ISO~31000. Elle demande de tenir ensemble deux
lectures~: la conséquence pour la personne et la conséquence qui en
découle pour l'organisation.

\smallskip
\textbf{Usage dans l'audit.} Ni l'un ni l'autre n'est un critère de
conformité au Cameroun. L'auditeur peut s'en servir pour vérifier qu'une
évaluation n'a oublié aucune famille de risques, à condition de le
déclarer dans son plan (Chapitre~\ref{chap:normes}).
\end{boxinfo}

% ============================================================================
\section{Les Étapes d'une Évaluation d'Impact}
\label{sec:c05-etapes}
% ============================================================================

L'ISO/IEC~29134~\cite{iso29134} décrit le déroulement d'une évaluation
d'impact~; la méthode publiée par l'autorité française~\cite{cnil2018pia}
en donne une version outillée. On les résume ici en cinq étapes.

\begin{boxprocedure}
\textbf{Conduire une évaluation d'impact.}
\begin{enumerate}
    \item \textbf{Décrire} le traitement~: finalités, données, personnes,
          flux, supports, durées, acteurs.
    \item \textbf{Apprécier la nécessité et la proportionnalité}~: fondement,
          finalités, données strictement utiles, information, droits.
    \item \textbf{Apprécier les risques} pour les personnes~: sources de
          risque, événements redoutés, gravité, vraisemblance.
    \item \textbf{Choisir les mesures} qui ramènent chaque risque à un niveau
          acceptable, et évaluer le risque résiduel.
    \item \textbf{Décider et documenter}~: validation par le responsable,
          avis du délégué s'il existe, consultation des personnes ou de
          leurs représentants si possible, date de revue.
\end{enumerate}
\end{boxprocedure}

\subsection{Trois événements redoutés}

Les événements redoutés se rangent en trois familles, qui correspondent aux
trois propriétés de la sécurité vues du côté de la personne~:
\begin{itemize}
    \item l'\textbf{accès illégitime} à ses données (confidentialité)~: un
          voisin apprend un diagnostic~;
    \item la \textbf{modification non désirée} (intégrité)~: un résultat
          d'analyse attribué au mauvais patient~;
    \item la \textbf{disparition} (disponibilité)~: le dossier introuvable le
          jour d'une urgence.
\end{itemize}

\subsection{Gravité et vraisemblance}

\begin{tableaucours}{cL{3.0cm}L{8.0cm}}{%
    \tableauHeader{Niveau} &
    \tableauHeader{Gravité} &
    \tableauHeader{Effet sur la personne}
}
1 & Négligeable & Désagrément surmonté sans difficulté \\
2 & Limitée & Désagréments significatifs, surmontables avec quelques
efforts \\
3 & Importante & Conséquences importantes, surmontables avec de sérieuses
difficultés (exclusion, perte d'emploi, préjudice financier) \\
4 & Maximale & Conséquences irrémédiables (atteinte à la santé, à la
sécurité physique, stigmatisation durable) \\
\end{tableaucours}
\vspace{-0.8em}
\begin{center}{\small\itshape Tableau~10.1. Échelle de gravité (d'après la méthode de la CNIL, paraphrasée)}\end{center}

La vraisemblance se cote de même sur quatre niveaux (négligeable, limitée,
importante, maximale), d'après les sources de risque et les vulnérabilités
des supports. On reporte chaque risque sur une grille gravité ×
vraisemblance~; un risque dont la gravité ou la vraisemblance reste à~3
ou~4 après mesures n'est pas acceptable sans décision motivée.

\begin{boxinfo}
\textbf{Coter l'inconnu dans une évaluation d'impact.} Une vraisemblance
qu'on ne sait pas apprécier (par exemple, la sécurité réelle d'un
prestataire jamais audité) ne se cote pas « limitée » par défaut. On la
note \NEK{}, et l'on traite le risque comme s'il était au niveau le plus
élevé plausible tant que la preuve manque.
\end{boxinfo}

% ============================================================================
\section{La Téléconsultation de MBOLO}
\label{sec:c05-mbolo}
% ============================================================================

\begin{boxscenario}{Mission MBOLO. L'évaluation d'impact, avril 2026}
\textbf{Description.} Patients des villages de la Kadey, dont des enfants~;
données d'identité, symptômes, images de l'examen, diagnostic,
prescription~; flux tablette, réseau mobile, serveur de MEDISOFT à Douala,
poste du médecin à Bertoua~; enregistrement des séances envisagé « pour la
formation ».

\textbf{Nécessité.} Le fondement est la préservation de la santé
(article~9(2)) pour la consultation elle-même. L'enregistrement des séances
pour la formation n'a pas ce fondement~: il exige un consentement
spécifique, et des parents pour les mineurs (article~9(3)).

\textbf{Risques.} Trois risques prioritaires~:
\begin{itemize}
    \item R1~: un tiers présent dans la salle du centre de santé voit ou
          entend la consultation (gravité~3, vraisemblance~3)~;
    \item R2~: interception ou accès indu aux images stockées chez le
          prestataire (gravité~4, vraisemblance \NEK{}, prestataire jamais
          audité)~;
    \item R3~: erreur d'attribution d'une image au mauvais patient
          (gravité~4, vraisemblance~2).
\end{itemize}

\textbf{Mesures.} Salle fermée et procédure d'accueil~; chiffrement de bout
en bout et contrat avec MEDISOFT révisé (Chapitre~\ref{chap:soustraitants})~;
double identification du patient~; abandon de l'enregistrement des
séances.

\textbf{Décision} (8~mai~2026). R1 et R3 ramenés à un niveau acceptable.
R2 reste \NEK{} tant que MEDISOFT n'a pas été audité~: le Dr~ABENA lance le
service sur deux villages seulement, en attendant l'audit du prestataire.
\end{boxscenario}

\begin{boxlecture}{L'évaluation d'impact de MBOLO, auditée en octobre 2026}
\begin{tableaucours}{L{7.8cm}cL{3.2cm}}{%
    \tableauHeader{Exigence} &
    \tableauHeader{Cote} &
    \tableauHeader{Preuve}
}
L'évaluation a précédé la mise en service & \OUI{} & Datée du 8~mai \\
Le traitement est décrit avec ses flux & \OUI{} & Schéma, liste des
données \\
Le fondement de chaque finalité est apprécié & \OUI{} & Section
nécessité \\
Les risques sont appréciés du point de vue des patients & \OUI{} & Trois
risques, échelles \\
Les mesures décidées sont en place & \NON{} & Chiffrement de bout en bout
non activé \\
Le risque lié au prestataire est apprécié & \NEK{} & Audit de MEDISOFT
non fait \\
La conformité aux modalités réglementaires est établie & \NEK{} & Texte
d'application non publié \\
\end{tableaucours}

\medskip
\textbf{Décompte.} Quatre \OUI{}, un \NON{}, deux \NEK{} sur sept.
\hfill\lecture{0{,}57}{0{,}29}{1{,}38}

\medskip
\textbf{Lecture.} L'évaluation est bien faite sur le papier, mais
$U = 0{,}29$~: le domaine reste opaque. Les deux \NEK{} sont de nature
différente~: l'un est de preuve (l'audit du prestataire le lèvera), l'autre
de critère (seul le texte d'application le lèvera). Le \NON{} est un écart
entre la décision et l'exécution, le plus fréquent dans les évaluations
d'impact.
\end{boxlecture}

% ============================================================================
\section{Auditer une Évaluation d'Impact}
\label{sec:c05-auditer}
% ============================================================================

L'auditeur ne refait pas l'évaluation~; il vérifie qu'elle existe quand elle
est exigée, qu'elle a été faite \textit{avant} le traitement, qu'elle
porte sur les personnes, que ses cotes sont justifiées, et surtout que les
mesures décidées sont en place. Ce dernier point est le plus souvent en
défaut~: une évaluation d'impact décrit un traitement idéal, et le
traitement réel s'en écarte.

Pour les traitements qui n'en ont pas, l'auditeur applique les critères de
la section~\ref{sec:c05-eleve} aux fiches du registre. À NYANGA CRÉDIT,
douze des quarante-sept traitements réunissent au moins deux critères~;
aucun n'a fait l'objet d'une évaluation d'impact. Le classeur des mauvais
payeurs en réunit quatre.

% ============================================================================
\section{Ce que le Risque demande à chaque Métier}
\label{sec:c05-metiers}
% ============================================================================

\begin{boxmetier}{Responsable de traitement}
Décider, au vu de l'évaluation, de lancer, d'adapter ou de renoncer, et
signer cette décision. Le Dr~ABENA a limité le service à deux villages~:
c'est une décision proportionnée à un risque encore inconnu.
\end{boxmetier}

\begin{boxmetier}{Délégué à la protection des données}
Conduire ou animer l'évaluation d'impact, rendre un avis écrit, et suivre la
mise en place des mesures jusqu'à la preuve.
\end{boxmetier}

\begin{boxmetier}{RSSI}
Apporter l'appréciation de la vraisemblance~: vulnérabilités des supports,
menaces réelles, état des mesures. Signaler quand une mesure décidée n'est
pas activée.
\end{boxmetier}

\begin{boxresume}
\begin{itemize}
    \item La loi exige une évaluation du risque pour tout traitement
          (article~25) et une évaluation d'impact préalable pour les
          traitements à risque élevé (article~33).
    \item Le risque évalué est celui des personnes~: accès illégitime,
          modification non désirée, disparition de leurs données.
    \item En l'absence de texte d'application, le risque élevé se repère à
          l'aide de critères déclarés, au moins deux réunis.
    \item Une évaluation d'impact se juge surtout à la mise en place
          effective des mesures qu'elle décide.
    \item À MBOLO, l'évaluation est sérieuse mais une mesure n'est pas
          activée et le risque lié au prestataire reste \NEK{}.
\end{itemize}
\end{boxresume}

\begin{boxinfo}
\textbf{Points de contrôle.} Vous avez compris ce chapitre si vous savez~:
\begin{itemize}
    \item distinguer les articles~25 et~33~;
    \item dire si un traitement exige une évaluation d'impact, critères à
          l'appui~;
    \item coter la gravité et la vraisemblance d'un risque pour une
          personne.
\end{itemize}
\end{boxinfo}

\begin{boxexo}[title={\ding{45}\quad Exercice~10.1. Évaluation d'impact ou non \niveaubase}]
Dites si ces traitements exigent une évaluation d'impact, critères à
l'appui~: (a) la paie de 210 salariés~; (b) le pointage biométrique du
personnel~; (c) la liste des fournisseurs~; (d) la notation automatique du
crédit chez NOUN PAY~; (e) la vidéosurveillance d'une salle d'attente de
pédiatrie.
\end{boxexo}

\begin{boxexo}[title={\ding{45}\quad Exercice~10.2. Réécrire le tableau \niveaumoyen}]
Reprenez le tableau de M.~YOUSSOUFA. Pour trois de ses lignes, écrivez
l'événement redouté du point de vue du patient, sa gravité et sa
vraisemblance.
\end{boxexo}

\begin{boxexo}[title={\ding{45}\quad Exercice~10.3. Le classeur des mauvais payeurs \niveaumoyen}]
Conduisez les trois premières étapes de l'évaluation d'impact du classeur
des mauvais payeurs de NYANGA. Quelle conclusion tirez-vous de l'étape de
nécessité~?
\end{boxexo}

\begin{boxexo}[title={\ding{45}\quad Exercice~10.4. Un risque inconnu \niveauexpert}]
Le Dr~ABENA vous demande s'il peut étendre la téléconsultation à vingt
villages avant l'audit de MEDISOFT. Rédigez votre avis, en séparant les
constatations de l'analyse, et en indiquant ce qui ferait passer R2 de
\NEK{} à une cote établie.
\end{boxexo}

\bigskip

\begin{boxinfo}
\textbf{Pour aller plus loin.}
\begin{itemize}
    \item L'ISO/IEC~29134~\cite{iso29134}, pour les étapes de l'évaluation
          d'impact.
    \item La méthode PIA de la CNIL~\cite{cnil2018pia}, librement
          accessible, pour ses échelles et ses modèles.
\end{itemize}
\end{boxinfo}

\begin{boxinfo}
\textbf{Transition.} Parmi les risques pour les personnes, l'un tient à leur
capacité d'agir sur leurs propres données. Peuvent-elles savoir ce que l'on
sait d'elles, le corriger, s'y opposer~? C'est l'objet du
Chapitre~\ref{chap:droits}.
\end{boxinfo}

        % ============================================================================
% Fichier  : partie4/P04_C11_droits_personnes.tex
% Titre    : Auditer les Droits des Personnes
% Auteur   : MINKA MI NGUIDJOÏ Thierry Emmanuel
% Version  : 1.0 -- novembre 2026
% Corrigés : corriges/Corr_C11.tex
% ============================================================================
% Sources : loi n° 2024/017, art. 21, 37 à 47, 56, 62, 64 (texte officiel) ;
% ISO/IEC 27701 (traitement des demandes, paraphrasé).
% Faits nouveaux (bible) : 31 demandes du 23/06 au 31/08/2026, 12 traitées,
% délai médian 26 jours ; demande test de KORO AUDIT, 27 juillet 2026.
% ============================================================================

\chapter{Auditer les Droits des Personnes}
\label{chap:droits}
\chaptermark{Les droits des personnes}

\epigraph{%
    « Un droit qu'on ne peut pas exercer\\
    n'est qu'une phrase dans une loi. »%
}{Principe du cours}

\niveauChapitre{Intermédiaire}{%
    Chapitres~\ref{chap:principes} et~\ref{chap:audit}
}

\begin{boxobjectifs}
\begin{itemize}
    \item Énoncer les droits des articles~37 à~46 et leurs limites.
    \item Vérifier qu'un canal de demande existe, qu'il est connu et que
          l'identité du demandeur est contrôlée sans excès.
    \item Mesurer les délais de réponse en l'absence de délai réglementaire,
          avec un critère de substitution déclaré.
    \item Tester le traitement d'une demande de bout en bout.
    \item Relier les défaillances aux sanctions administratives et pénales
          encourues.
\end{itemize}
\end{boxobjectifs}

\bigskip

% ============================================================================
\section*{Avant de commencer : trente et une lettres}
% ============================================================================

Bafoussam, lundi 27~juillet~2026. La lettre de M.~TAGNE n'est pas restée
seule. L'audience de février a été relayée par la radio locale, et d'autres
clients ont écrit. Me~ETOUNDI a ouvert un classeur~: entre le 23~juin et le
27~juillet, vingt-deux lettres, courriels et messages vocaux laissés au
standard. Certaines demandent l'accès, d'autres l'effacement, deux
s'opposent aux appels de prospection, une demande la rectification d'un
nom mal orthographié qui bloque un virement.

Il n'existe ni procédure, ni délai, ni modèle de réponse. Les demandes
arrivent par six canaux différents~; trois ont été perdues entre l'agence
et le siège. La seule qui ait reçu une réponse complète est celle d'un
client qui connaissait personnellement le directeur général.

\begin{boxalert}
\textbf{Violation typique.} Considérer que l'organisation respecte les droits
des personnes parce qu'elle n'a « jamais refusé » une demande.

\smallskip
\textbf{Règle~:} le respect des droits se prouve par une procédure écrite,
un registre des demandes, des réponses datées et un délai mesuré. Une
demande perdue est un refus de fait~; une demande traitée par faveur
n'est pas un droit exercé.
\end{boxalert}

% ============================================================================
\section{Les Droits de la Personne Concernée}
\label{sec:c06-droits}
% ============================================================================

Le titre~III de la loi (articles~37 à~47) organise les droits. L'article~21
impose d'en informer la personne dès la collecte~; l'auditeur vérifie donc
à la fois l'information et l'exercice.

\begin{tableaucours}{L{3.0cm}cL{8.0cm}}{%
    \tableauHeader{Droit} &
    \tableauHeader{Article} &
    \tableauHeader{Contenu et limites}
}
Accès & 39 & Confirmation du traitement, communication des données et de
leur origine, finalités, destinataires~; copie contre paiement limité au
coût de reproduction \\
Rectification & 42 & Rectifier, compléter, mettre à jour, verrouiller ou
supprimer les données inexactes ou incomplètes, sur justification de
l'identité \\
Effacement & 37, 38 & Cessation de la diffusion et effacement si les
données ne sont plus nécessaires, si le consentement est vicié, retiré ou
expiré, ou sans base légale~; sous réserve des obligations de conservation \\
Opposition & 40 & À tout moment, pour les catégories particulières de
données, pour la prospection, et pour tout motif prévu par la loi~; droit
propre des parents d'un mineur \\
Prospection & 41 & Communication à des tiers pour la prospection soumise à
consentement préalable \\
Portabilité & 43 & Recevoir gratuitement ses données dans un format
structuré, et les faire transmettre à un autre responsable~; exclue pour
les missions d'intérêt public \\
Décision automatisée & 44 & S'opposer à une décision fondée exclusivement
sur un traitement automatisé (Chapitre~\ref{chap:profilage}) \\
Décès & 45 & Fin du traitement au décès, sauf exceptions~; directives
post-mortem~; mise à jour par les ayants droit \\
Limitation & 46 & Geler le traitement quand l'exactitude ou la finalité est
contestée \\
\end{tableaucours}
\vspace{-0.8em}
\begin{center}{\small\itshape Tableau~11.1. Les droits de la personne concernée}\end{center}

\begin{boxjuris}
\textbf{\loicam{n°~2024/017}, article~43}~\cite{loi2024017}~:
\textit{« (1) La personne concernée a le droit de recevoir gratuitement les
données à caractère personnel la concernant qu'elle a fournies à un
responsable du traitement, dans un format structuré, couramment utilisé et
lisible par machine. (2) La personne concernée a le droit à la portabilité
de ses données. (3) Lorsque la personne concernée exerce son droit à la
portabilité des données en application de l'alinéa 1 ci-dessus, elle a le
droit d'obtenir que les données à caractère personnel soient transmises
directement d'un responsable du traitement à un autre, si cela est
techniquement possible. [\ldots] »}

\smallskip
\textbf{Portée.} L'accès peut être facturé au coût de reproduction
(article~39)~; la portabilité est gratuite. Pour une microfinance, la
portabilité permet à un client de faire transmettre son historique à un
autre établissement~: l'organisation doit être capable de l'exporter.
\end{boxjuris}

L'article~45 contient une disposition originale~: le traitement des données
d'une personne « cesse dès constatation de son décès », sauf obligation
légale, défense en justice ou directives que la personne a laissées, et les
ayants droit peuvent demander la mise à jour des informations du défunt.
Pour un établissement financier, cela suppose une procédure de succession
qui distingue ce que la loi impose de garder et ce qui doit cesser.

% ============================================================================
\section{Le Délai, un Critère à Déclarer}
\label{sec:c06-delai}
% ============================================================================

L'article~47 renvoie à la voie réglementaire « les délais et les modalités
d'examen des demandes ». Tant que ce texte n'est pas publié, la loi ne fixe
aucun délai. L'auditeur a deux possibilités, à choisir dans le plan
d'audit~:

\begin{itemize}
    \item coter \NEK{} de critère toute exigence de délai réglementaire, ce
          qui est exact mais ne dit rien de la qualité du service~;
    \item déclarer un \textbf{critère de substitution}, par exemple
          « réponse complète dans les 30 jours calendaires », en précisant
          qu'il s'agit d'une référence de bonne pratique et non d'une
          obligation légale, et mesurer les délais contre ce critère.
\end{itemize}

KORO AUDIT a retenu la seconde, en la déclarant dans le plan d'audit. Le
rapport dira donc « 30~jours (critère de l'audit) », jamais « 30~jours
(délai légal) ».

% ============================================================================
\section{Auditer le Traitement des Demandes}
\label{sec:c06-auditer}
% ============================================================================

\begin{boxprocedure}
\textbf{Auditer l'exercice des droits.}
\begin{enumerate}
    \item Vérifier l'information~: les droits et la manière de les exercer
          figurent-ils sur les formulaires et la notice (article~21)~?
    \item Recenser les canaux réels~: courrier, courriel, agence, téléphone,
          application, réseaux sociaux.
    \item Obtenir le registre des demandes, ou le reconstituer à partir des
          canaux.
    \item Pour chaque demande, relever la date de réception, le type, la
          vérification d'identité, la date et le contenu de la réponse.
    \item Calculer les délais, la part des demandes traitées et la part
          traitées dans le délai retenu.
    \item Réaliser une demande test, prévue au mandat, et la suivre jusqu'à
          la réponse.
    \item Vérifier que les effacements et les oppositions ont été exécutés
          dans les systèmes, et pas seulement acceptés par lettre.
\end{enumerate}
\end{boxprocedure}

La vérification d'identité mérite attention. L'article~42 vise la personne
« justifiant de son identité »~: l'organisation doit s'assurer qu'elle ne
communique pas les données d'un client à un imposteur, sans exiger plus que
nécessaire. Demander une copie de la carte nationale d'identité pour une
simple opposition aux SMS est excessif~; la demander pour communiquer un
relevé de compte est justifié.

Le dernier point de la procédure est le plus souvent oublié. Une lettre
« nous avons bien pris en compte votre demande d'effacement » ne prouve pas
l'effacement~: l'auditeur cherche le client dans la base, dans les
sauvegardes de moins d'un mois, dans le fichier du prestataire.

% ============================================================================
\section{Les Demandes de NYANGA}
\label{sec:c06-nyanga}
% ============================================================================

\begin{boxscenario}{Mission NYANGA. Le registre reconstitué, 31 août 2026}
\textbf{Constatations.}
\begin{itemize}
    \item 31~demandes reçues entre le 23~juin et le 31~août~: 14 d'accès,
          8 d'effacement, 5 d'opposition à la prospection, 3 de
          rectification, 1 de portabilité.
    \item 12 ont reçu une réponse complète, 7 une réponse d'attente, 12
          aucune réponse. Délai médian des réponses complètes~: 26~jours~;
          4 au-delà de 30~jours.
    \item Les 5 oppositions ont été acceptées par lettre~; 3 des 5 numéros
          figurent encore dans la liste d'appels de prospection du 28~août.
    \item La demande de portabilité a été refusée au motif que « le
          logiciel ne permet pas l'export ».
    \item La demande test de KORO AUDIT, envoyée le 27~juillet par courriel à
          l'adresse du service client, a reçu un accusé de réception
          automatique et aucune réponse au 31~août.
\end{itemize}

\textbf{Analyse} (confiance élevée). L'organisation n'a pas de procédure~;
les réponses dépendent de la personne qui reçoit la demande. Les
oppositions non exécutées exposent l'établissement à l'article~64, qui
punit le traitement maintenu malgré l'opposition à la prospection.
\end{boxscenario}

\begin{boxlecture}{Le domaine D2, droits des personnes}
\begin{tableaucours}{L{7.8cm}cL{3.2cm}}{%
    \tableauHeader{Exigence} &
    \tableauHeader{Cote} &
    \tableauHeader{Preuve}
}
Les personnes sont informées de leurs droits & \NON{} & Formulaire 2016 \\
Une procédure de traitement des demandes existe & \NON{} & Aucune \\
Un registre des demandes est tenu & \NON{} & Reconstitué par l'audit \\
Les demandes d'accès reçoivent une réponse & \NON{} & 12 sur 31 \\
Les réponses respectent le critère de 30 jours & \NON{} & 8 sur 12 \\
Les oppositions sont exécutées dans les systèmes & \NON{} & 3 numéros
encore appelés \\
La portabilité est possible & \NON{} & Refus du 14~août \\
L'identité est vérifiée de façon proportionnée & \NEK{} & Pratiques
variables, non documentées \\
Les effacements acceptés sont exécutés & \NEK{} & Sauvegardes non
examinées \\
La procédure de succession respecte l'article~45 & \OUI{} & Note de 2025,
appliquée sur 3 dossiers \\
\end{tableaucours}

\medskip
\textbf{Décompte.} Un \OUI{}, sept \NON{}, deux \NEK{} sur dix.
\hfill\lecture{0{,}10}{0{,}20}{1{,}16}

\medskip
\textbf{Lecture.} $U = 0{,}20$ atteint juste le seuil~: le domaine est
opaque selon la règle, mais de peu. La non-conformité est établie sur sept
exigences, dont une, l'opposition non exécutée, relève d'une sanction
pénale. La recommandation prioritaire ne dépend pas des deux \NEK{}.
\end{boxlecture}

\begin{boxjuris}
\textbf{\loicam{n°~2024/017}, article~56}~\cite{loi2024017}~:
\textit{« Est passible d'une amende d'un million (1~000~000) à dix millions
(10~000~000) de francs CFA, tout responsable de traitement ou sous-traitant
qui refuse de mettre à la disposition de la personne concernée, les
informations sollicitées la concernant. »}

\smallskip
\textbf{Portée.} Une demande d'accès laissée sans réponse peut être
qualifiée de refus. L'article~62 ouvre en outre à la personne la voie du
référé « en cas d'atteinte grave aux droits », et la réparation devant la
juridiction compétente~: c'est la voie qu'annonçait M.~TAGNE.
\end{boxjuris}

% ============================================================================
\section{Ce que les Droits demandent à chaque Métier}
\label{sec:c06-metiers}
% ============================================================================

\begin{boxmetier}{Délégué à la protection des données}
Centraliser les demandes sur un canal connu, tenir le registre, fixer un
délai cible, vérifier l'exécution dans les systèmes, et rendre compte
chaque mois du nombre de demandes et des délais.
\end{boxmetier}

\begin{boxmetier}{Juriste conformité}
Préparer des modèles de réponse par type de droit, et une grille des
données à conserver malgré une demande d'effacement (comptabilité,
procédure judiciaire, lutte contre le blanchiment).
\end{boxmetier}

\begin{boxmetier}{Auditeur de la protection des données}
Toujours vérifier l'exécution, pas seulement l'acceptation~; et déclarer
dans le plan d'audit tout critère de délai qui n'est pas légal.
\end{boxmetier}

\begin{boxresume}
\begin{itemize}
    \item La loi reconnaît l'accès, la rectification, l'effacement,
          l'opposition, la portabilité gratuite, la limitation, des droits
          face aux décisions automatisées et des règles en cas de décès.
    \item Les délais de réponse relèvent d'un texte d'application~; en son
          absence, l'auditeur déclare un critère de substitution.
    \item Un droit exercé se prouve par un registre, une réponse datée et
          l'exécution dans les systèmes.
    \item Le refus de communiquer les informations (article~56) et le
          traitement maintenu malgré l'opposition (article~64) sont
          sanctionnés.
    \item À NYANGA, 12 demandes sur 31 ont reçu une réponse complète, et des
          oppositions acceptées n'ont pas été exécutées.
\end{itemize}
\end{boxresume}

\begin{boxinfo}
\textbf{Points de contrôle.} Vous avez compris ce chapitre si vous savez~:
\begin{itemize}
    \item citer les limites du droit à l'effacement~;
    \item expliquer pourquoi un délai de 30~jours ne peut pas être présenté
          comme légal~;
    \item vérifier l'exécution d'une opposition dans les systèmes.
\end{itemize}
\end{boxinfo}

\begin{boxexo}[title={\ding{45}\quad Exercice~11.1. Quel droit \niveaubase}]
Quel droit chaque demande exerce-t-elle~? (a) « Donnez-moi la liste des
personnes à qui vous avez transmis mon numéro. »~; (b) « Je ne veux plus de
vos SMS. »~; (c) « Transmettez mon historique de crédit à ma nouvelle
banque. »~; (d) « Mon père est décédé~; supprimez son compte de vos
campagnes. »~; (e) « Mon salaire est faux dans votre dossier. »
\end{boxexo}

\begin{boxexo}[title={\ding{45}\quad Exercice~11.2. Mesurer les délais \niveaumoyen}]
Voici les délais de réponse complète, en jours, de dix demandes~: 5, 12,
18, 21, 26, 26, 31, 38, 44, 60. Calculez le délai médian et la part des
réponses dans le critère de 30~jours. Rédigez la phrase du rapport.
\end{boxexo}

\begin{boxexo}[title={\ding{45}\quad Exercice~11.3. L'effacement de M. TAGNE \niveaumoyen}]
Dressez la liste des données de M.~TAGNE que NYANGA peut effacer et de
celles qu'elle doit conserver, avec pour chacune le texte qui le justifie
ou la mention \NEK{} si vous ne connaissez pas le texte.
\end{boxexo}

\begin{boxexo}[title={\ding{45}\quad Exercice~11.4. La demande test \niveauexpert}]
Rédigez la clause du mandat qui autorise la demande test, puis le papier de
travail qui en rend compte. Quelles précautions prenez-vous pour que le
test ne soit pas lui-même un traitement déloyal~?
\end{boxexo}

\bigskip

\begin{boxinfo}
\textbf{Pour aller plus loin.}
\begin{itemize}
    \item L'ISO/IEC~27701~\cite{iso27701}, pour l'organisation du traitement
          des demandes des personnes.
\end{itemize}
\end{boxinfo}

\begin{boxinfo}
\textbf{Transition.} Les droits protègent la personne contre l'organisation.
La sécurité la protège contre tous les autres~: intrus, agents curieux,
prestataires négligents. C'est l'objet du Chapitre~\ref{chap:securite}.
\end{boxinfo}

        % ============================================================================
% Fichier  : partie4/P04_C12_securite_donnees.tex
% Titre    : Auditer la Sécurité des Données
% Auteur   : MINKA MI NGUIDJOÏ Thierry Emmanuel
% Version  : 1.0 -- novembre 2026
% Corrigés : corriges/Corr_C12.tex
% ============================================================================
% Sources : loi n° 2024/017, art. 7, 15, 22(2), 24, 25(2-3), 27, 28, 57
% (texte officiel) ; ISO/IEC 27002:2022 et ISO/IEC 27701 (mesures,
% paraphrasées).
% Faits nouveaux (bible) : consultations du dossier d'une patiente les 11 et
% 12 mai 2026 au CENTRE MÉDICAL MBOLO ; 14 consultations, 6 comptes, 4 sans
% lien avec la prise en charge.
% ============================================================================

\chapter{Auditer la Sécurité des Données}
\label{chap:securite}
\chaptermark{La sécurité des données}

\epigraph{%
    « La menace la plus fréquente sur un dossier médical\\
    n'est pas le pirate. C'est le collègue curieux. »%
}{Principe du cours}

\niveauChapitre{Intermédiaire}{%
    Chapitres~\ref{chap:audit} et~\ref{chap:risque}~; bases de la sécurité
    des systèmes
}

\begin{boxobjectifs}
\begin{itemize}
    \item Énoncer les obligations de sécurité des articles~22, 24, 25
          et~27 de la loi.
    \item Traduire les dix exigences de l'article~27 en contrôles
          vérifiables.
    \item Exploiter des journaux d'accès pour établir des consultations
          sans lien avec la finalité.
    \item Distinguer la sécurité des systèmes de la protection des données
          contre les usages internes illégitimes.
    \item Auditer le rapport annuel de l'article~27, alinéa~2.
\end{itemize}
\end{boxobjectifs}

\bigskip

% ============================================================================
\section*{Avant de commencer : quatorze consultations en quarante-huit heures}
% ============================================================================

Bertoua, mercredi 13~mai~2026. Une chanteuse connue de la région, admise
deux jours plus tôt au CENTRE MÉDICAL MBOLO pour une intervention, apprend
par un message vocal qui circule en ville le nom de son affection. Elle
appelle la directrice. Le Dr~ABENA demande à M.~YOUSSOUFA d'extraire du
dossier patient informatisé la liste des consultations de ce dossier.

L'extraction tient sur une page~: quatorze consultations en quarante-huit
heures, par six comptes. Deux sont ceux du chirurgien et de l'infirmière du
service. Les quatre autres sont ceux d'une secrétaire de la maternité, d'un
agent de la facturation, d'un brancardier dont le compte est partagé par
l'équipe de nuit, et d'un stagiaire dont le stage s'est terminé en mars.

\begin{boxalert}
\textbf{Violation typique.} Auditer la sécurité des données comme la sécurité
d'un réseau~: pare-feu, antivirus, mots de passe, sans regarder qui, à
l'intérieur, consulte quoi.

\smallskip
\textbf{Règle~:} la loi exige que les personnes autorisées n'accèdent
« qu'aux données relevant de leur autorisation » et que l'on puisse
vérifier « a posteriori » qui a consulté quoi (article~27). L'audit teste
les droits d'accès et exploite les journaux.
\end{boxalert}

% ============================================================================
\section{Les Obligations de Sécurité}
\label{sec:c07-obligations}
% ============================================================================

Les obligations de sécurité sont dispersées dans la loi. Les articles~7
et~15 posent la confidentialité~; l'article~24 réserve le traitement aux
personnes « qui agissent sous l'autorité du responsable du traitement et
seulement sur ses instructions », et renvoie au référentiel de mesures de
l'Autorité~; l'article~25, alinéa~2, impose d'empêcher que les données
soient « déformées, endommagées, indisponibles » ou que des tiers non
autorisés y aient accès~; l'article~22, alinéa~2, cite le chiffrement et
« l'instauration des autorisations restrictives d'accès aux données »~; l'article~27 en
donne la liste la plus complète.

\begin{boxjuris}
\textbf{\loicam{n°~2024/017}, article~27, alinéa~1 (extraits)}~\cite{loi2024017}~:
\textit{« Le responsable du traitement et le sous-traitant des données à
caractère personnel sont tenus~: de garantir la disponibilité des données à
caractère personnel notamment grâce à la mise en œuvre des mesures visant à
garantir la résilience des systèmes d'information contenant les données à
caractère personnel~; de garantir que lors de l'utilisation d'un système de
traitement automatisé de données, les personnes autorisées ne puissent
accéder qu'aux données relevant de leur autorisation~; [\ldots] de
garantir que puisse être vérifiée et constatée a posteriori l'identité des
personnes ayant accès au système d'information contenant des données à
caractère personnel, la nature des données qui ont été introduites,
modifiées, altérées, copiées, effacées ou lues dans le système, ainsi que
la période de manipulation de ces données~; [\ldots] de sauvegarder les
données par la constitution de copies de sécurité protégées~; [\ldots] »}

\smallskip
\textbf{Portée.} L'article énumère dix obligations. La plus exigeante pour
l'audit est la traçabilité~: la loi vise les modifications, et aussi les
données « lues ». Un système qui ne journalise pas les
consultations ne permet pas de satisfaire l'article.
\end{boxjuris}

L'alinéa~2 du même article ajoute une obligation de compte rendu~: le
responsable « adresse à l'Autorité de protection des données à caractère
personnel un rapport annuel sur l'état de mise en œuvre des mesures de
sécurité contenues dans le référentiel des mesures techniques et
organisationnelles ». Tant que le référentiel n'est pas publié, le contenu
du rapport est un \NEK{} de critère~; l'existence d'un état annuel des
mesures, elle, peut être vérifiée.

% ============================================================================
\section{Dix Obligations, Dix Contrôles}
\label{sec:c07-controles}
% ============================================================================

\begin{tableaucours}{L{5.0cm}L{3.0cm}L{4.0cm}}{%
    \tableauHeader{Obligation (article~27)} &
    \tableauHeader{Mesure ISO/IEC 27002} &
    \tableauHeader{Test d'audit}
}
Disponibilité et résilience & 5.30, 8.14 & Plan de reprise exercé, durée
d'indisponibilité mesurée \\
Accès limité à l'autorisation & 5.15, 5.18, 8.3 & Comparer droits et
fonctions sur échantillon \\
Identité des destinataires vérifiée & 5.14 & Procédure de transmission,
échantillon d'envois \\
Traçabilité a posteriori & 8.15 & Journal des consultations, test de
reconstitution \\
Accès physique aux installations & 7.1, 7.2 & Visite, registre des accès \\
Supports protégés & 7.10 & Inventaire, chiffrement des supports amovibles \\
Introduction et modification non autorisées & 8.2, 8.5 & Comptes
privilégiés, validation des saisies \\
Transmission protégée & 5.14, 8.24 & Chiffrement des flux, configuration \\
Sauvegardes protégées & 8.13 & Test de restauration, sauvegarde hors
ligne \\
Autres mesures réglementaires & À préciser & \NEK{} de critère \\
\end{tableaucours}
\vspace{-0.8em}
\begin{center}{\small\itshape Tableau~12.1. L'article~27 et les mesures de l'ISO/IEC~27002:2022 (correspondances indicatives)}\end{center}

Les numéros de mesures renvoient à l'ISO/IEC~27002:2022~\cite{iso27002_2022}
et sont donnés à titre indicatif~; la correspondance n'est pas fixée par la
loi. L'ISO/IEC~27701~\cite{iso27701} complète ces mesures par des exigences
propres aux données personnelles (Chapitre~\ref{chap:normes}).

\subsection{Les droits d'accès}

L'audit des droits d'accès compare, pour un échantillon de comptes, les
droits accordés aux besoins de la fonction. Trois anomalies reviennent
partout~: les comptes de personnes parties (stagiaires, intérimaires,
salariés licenciés), les comptes partagés par une équipe, qui rendent la
traçabilité impossible, et les droits trop larges accordés « pour
dépanner » et jamais retirés.

\subsection{Les journaux de consultation}

Exploiter un journal de consultation, c'est répondre à une question précise~:
« Qui a consulté ce dossier, et chaque consultation avait-elle un lien avec
la finalité~? » On rapproche chaque ligne du journal d'une pièce qui la
justifie~: le planning du service, l'affectation du patient, l'acte facturé.
Une consultation sans pièce justificative n'est pas encore une faute~; c'est
une consultation à expliquer.

\begin{boxinfo}
\textbf{Un journal se lit avec ses limites.} Un compte partagé ne désigne pas
une personne~; une consultation par un compte ne prouve pas qui était au
clavier~; un journal qui ne garde que trente jours ne permet pas de
remonter plus loin. L'auditeur écrit ce que le journal établit (« le compte
X a ouvert le dossier à telle heure ») et ce qu'il n'établit pas (« la
personne derrière le compte »).
\end{boxinfo}

% ============================================================================
\section{Le Dossier de la Patiente}
\label{sec:c07-mbolo}
% ============================================================================

\begin{boxscenario}{Mission MBOLO. Les quatorze consultations, 11 et 12 mai 2026}
\begin{tableaucours}{L{3.0cm}cL{3.2cm}L{3.6cm}}{%
    \tableauHeader{Compte} &
    \tableauHeader{Consult.} &
    \tableauHeader{Fonction} &
    \tableauHeader{Pièce justificative}
}
Chirurgien & 4 & Bloc opératoire & Programme opératoire \\
Infirmière du service & 4 & Chirurgie & Planning, soins notés \\
Secrétaire maternité & 2 & Maternité & Aucune \\
Agent de facturation & 1 & Facturation & Facture émise le 12~mai \\
Compte « nuit\_brancard » & 2 & Partagé, 5 brancardiers & Aucune \\
Stagiaire & 1 & Stage terminé en mars & Aucune \\
\end{tableaucours}

\medskip
\textbf{Constatations.} Quatorze consultations. Neuf sont justifiées par une
pièce (chirurgien, infirmière, facturation). Cinq ne le sont pas~: deux par
la secrétaire de la maternité, deux par un compte partagé entre cinq
personnes, une par le compte d'un stagiaire parti depuis deux mois.

\textbf{Analyse} (confiance élevée sur les comptes, faible sur les
personnes). Cinq consultations sont sans lien établi avec la prise en
charge. Le journal ne permet pas de savoir lequel des cinq brancardiers a
utilisé le compte partagé, ni qui a utilisé le compte du stagiaire. La
proposition « la fuite vient d'une consultation du dossier » reste \NEK{}~:
le diagnostic a pu circuler oralement.
\end{boxscenario}

Le Dr~ABENA a désactivé le compte du stagiaire et le compte partagé le jour
même, et ouvert une procédure disciplinaire pour la secrétaire, qui a
reconnu la consultation. Elle a aussi posé à KORO AUDIT, mandaté en
octobre, une question plus large~: combien d'autres dossiers ont été
consultés de la même façon sans que personne ne s'en plaigne~?

\begin{boxlecture}{Le domaine D3 à MBOLO, octobre 2026}
\begin{tableaucours}{L{7.8cm}cL{3.2cm}}{%
    \tableauHeader{Exigence} &
    \tableauHeader{Cote} &
    \tableauHeader{Preuve}
}
Les consultations sont journalisées & \OUI{} & Journal MEDIKA \\
Les journaux sont conservés au moins un an & \NON{} & 90 jours \\
Chaque compte désigne une seule personne & \NON{} & 3 comptes partagés \\
Les comptes des personnes parties sont désactivés & \NON{} & 7 comptes
actifs sur 31 départs \\
Les droits correspondent aux fonctions & \NON{} & Facturation voit les
diagnostics \\
Les journaux sont revus & \NON{} & Aucune revue avant mai \\
Les sauvegardes sont restaurables & \OUI{} & Test de restauration du
22~octobre \\
Les flux de téléconsultation sont chiffrés de bout en bout & \NON{} &
Option non activée \\
Le prestataire applique les mesures convenues & \NEK{} & MEDISOFT non
audité \\
Le rapport annuel de l'article~27(2) est conforme & \NEK{} & Référentiel non
publié \\
\end{tableaucours}

\medskip
\textbf{Décompte.} Deux \OUI{}, six \NON{}, deux \NEK{} sur dix.
\hfill\lecture{0{,}20}{0{,}20}{1{,}37}

\medskip
\textbf{Lecture.} Domaine opaque au seuil ($U = 0{,}20$), mais six écarts
établis qui touchent tous la traçabilité et les droits d'accès. La
journalisation existe~; c'est son usage qui manque.
\end{boxlecture}

% ============================================================================
\section{La Sécurité à NYANGA, après l'Incident}
\label{sec:c07-nyanga}
% ============================================================================

À NYANGA CRÉDIT, la sécurité s'est renforcée depuis mars~2025~: la
passerelle a été remplacée, l'inventaire des équipements est tenu, les
journaux sont centralisés. L'audit le constate et cote \OUI{} les
exigences correspondantes. Il relève aussi ce que l'incident n'a pas
changé~: les agents des agences voient les comptes de tous les clients du
réseau, et non de leur seule agence~; les relevés de compte partent encore
par messagerie personnelle quand un client le demande. La sécurité
technique a progressé~; la limitation des accès à la finalité, non.

% ============================================================================
\section{Ce que la Sécurité demande à chaque Métier}
\label{sec:c07-metiers}
% ============================================================================

\begin{boxmetier}{RSSI}
Supprimer les comptes partagés, désactiver les comptes à la date de départ,
conserver les journaux de consultation assez longtemps pour enquêter, et
les revoir régulièrement sur les dossiers sensibles.
\end{boxmetier}

\begin{boxmetier}{Auditeur de la protection des données}
Tester les droits d'accès contre les fonctions et rapprocher les journaux de
pièces justificatives. Écrire ce que le journal établit sur un compte, et ce
qu'il n'établit pas sur une personne.
\end{boxmetier}

\begin{boxmetier}{Responsable de traitement}
Tirer les conséquences disciplinaires d'une consultation injustifiée, et
faire savoir au personnel que les consultations sont tracées et revues~:
la dissuasion fait partie de la mesure.
\end{boxmetier}

\begin{boxresume}
\begin{itemize}
    \item Les obligations de sécurité figurent aux articles~7, 15, 22, 24,
          25 et~27~; l'article~27 en donne dix, dont la traçabilité des
          données lues.
    \item Le responsable adresse un rapport annuel sur les mesures à
          l'Autorité (article~27, alinéa~2).
    \item Les menaces internes se contrôlent par les droits d'accès et
          l'exploitation des journaux de consultation.
    \item Un journal établit ce qu'un compte a fait, pas ce qu'une personne a
          fait~; les comptes partagés détruisent la traçabilité.
    \item À MBOLO, cinq consultations d'un dossier sensible sont sans lien
          établi avec la prise en charge~; six écarts sont établis dans le
          domaine de la sécurité.
\end{itemize}
\end{boxresume}

\begin{boxinfo}
\textbf{Points de contrôle.} Vous avez compris ce chapitre si vous savez~:
\begin{itemize}
    \item traduire une obligation de l'article~27 en test d'audit~;
    \item rapprocher un journal de consultation de pièces justificatives~;
    \item dire pourquoi un compte partagé est un écart à l'article~27.
\end{itemize}
\end{boxinfo}

\begin{boxexo}[title={\ding{45}\quad Exercice~12.1. Quel contrôle \niveaubase}]
Proposez un test d'audit pour chaque obligation~: (a) les sauvegardes sont
protégées~; (b) les supports amovibles ne sont pas lisibles par un tiers~;
(c) l'identité des destinataires est vérifiée avant transmission~;
(d) les données sont disponibles en cas de panne.
\end{boxexo}

\begin{boxexo}[title={\ding{45}\quad Exercice~12.2. Lire un journal \niveaumoyen}]
Le compte de l'agent de facturation a consulté le dossier une fois, le
12~mai à 16~h~04~; la facture a été émise à 16~h~11. Le compte a aussi accès
à l'onglet « diagnostic ». Quelles cotes donnez-vous aux propositions
« cette consultation était justifiée » et « les droits de ce compte sont
limités à sa fonction »~?
\end{boxexo}

\begin{boxexo}[title={\ding{45}\quad Exercice~12.3. Combien d'autres dossiers \niveaumoyen}]
Proposez une méthode pour répondre à la question du Dr~ABENA~: combien
d'autres dossiers ont été consultés sans lien avec la prise en charge~?
Précisez les données à extraire, les règles de rapprochement, l'échantillon,
et la manière de protéger les données des patients pendant l'analyse.
\end{boxexo}

\begin{boxexo}[title={\ding{45}\quad Exercice~12.4. Sécurisé mais non conforme \niveauexpert}]
Le DSI de NYANGA affirme~: « Depuis l'incident, notre sécurité est au
niveau des meilleures banques. » Montrez, à partir des constats de la
section~\ref{sec:c07-nyanga}, en quoi cette affirmation peut être vraie et
l'établissement non conforme à l'article~27. Rédigez la recommandation.
\end{boxexo}

\bigskip

\begin{boxinfo}
\textbf{Pour aller plus loin.}
\begin{itemize}
    \item L'ISO/IEC~27002:2022~\cite{iso27002_2022}, mesures de contrôle
          d'accès et de journalisation.
    \item L'ISO/IEC~27701~\cite{iso27701}, pour les mesures propres aux
          données personnelles.
\end{itemize}
\end{boxinfo}

\begin{boxinfo}
\textbf{Transition.} La Partie~II a montré comment auditer. Il faut
maintenant lire de près le cadre qui fournit les critères~: la loi
n°~2024/017 elle-même, son Autorité, ses sanctions, et le cadre africain
dans lequel elle s'inscrit. C'est l'objet du Chapitre~\ref{chap:loi}.
\end{boxinfo}


    \part{Pratique : de l'Écart à la Conformité}
    \label{part:pratique}

        % ============================================================================
% Fichier  : partie5/P05_C13_sous_traitants.tex
% Titre    : Sous-traitants et Responsables Conjoints
% Auteur   : MINKA MI NGUIDJOÏ Thierry Emmanuel
% Version  : 1.0 -- novembre 2026
% Corrigés : corriges/Corr_C13.tex
% ============================================================================
% Sources : loi n° 2024/017, art. 5, 16, 20, 23(3), 24, 26, 30, 31 (texte
% officiel) ; ISO/IEC 27701 et 27018 (paraphrasées).
% Faits nouveaux (bible) : contrat SOREC (2019), contrat de l'éditeur du
% logiciel bancaire, avenant MEDISOFT du 20/05/2026, convention avec SAWA
% ASSURANCES pour l'assurance des emprunteurs ; 17 au 21 août 2026.
% ============================================================================

\chapter{Sous-traitants et Responsables Conjoints}
\label{chap:soustraitants}
\chaptermark{Sous-traitants et responsables conjoints}

\epigraph{%
    « Confier ses données à un prestataire,\\
    ce n'est pas lui confier sa responsabilité. »%
}{Principe du cours}

\niveauChapitre{Intermédiaire}{%
    Chapitres~\ref{chap:donnees} et~\ref{chap:cartographie}
}

\begin{boxobjectifs}
\begin{itemize}
    \item Appliquer les obligations des articles~16, 26 et~30 à un contrat
          de sous-traitance.
    \item Contrôler les clauses d'un contrat contre une liste de vérification
          déclarée.
    \item Vérifier que le sous-traitant applique ce que le contrat prévoit.
    \item Reconnaître une situation de responsables conjoints et vérifier la
          convention de l'article~31.
    \item Tirer les conséquences de l'article~20~: mêmes obligations pour le
          responsable et le sous-traitant.
\end{itemize}
\end{boxobjectifs}

\bigskip

% ============================================================================
\section*{Avant de commencer : un contrat de deux pages}
% ============================================================================

Douala, lundi 17~août~2026. Mme~BELLO rencontre le gérant de SOREC, le
prestataire de recouvrement de NYANGA CRÉDIT. Chaque mois, SOREC reçoit par
messagerie un classeur des débiteurs en retard~: noms, téléphones,
adresses, employeurs, montants. Ses agents appellent, se déplacent,
parfois chez l'employeur.

Le contrat, signé en 2019, tient sur deux pages~: objet, rémunération au
pourcentage des sommes recouvrées, durée. Il ne dit rien des données~:
ni leur usage, ni leur sécurité, ni leur sort à la fin du dossier. Le gérant
reconnaît sans difficulté que les classeurs des cinq dernières années sont
conservés « au cas où », et que ses agents ont accès à tous les dossiers de
tous les clients de SOREC.

\begin{boxalert}
\textbf{Violation typique.} Tenir le prestataire pour seul responsable de ce
qu'il fait des données, ou tenir le contrat pour une preuve que le
prestataire protège les données.

\smallskip
\textbf{Règle~:} le responsable « veille au respect » des mesures par son
sous-traitant (article~16) et doit s'assurer qu'il offre des garanties
suffisantes (article~30). Le contrat fixe les obligations~; l'audit vérifie
qu'elles sont tenues.
\end{boxalert}

% ============================================================================
\section{Les Obligations de la Loi}
\label{sec:c10-obligations}
% ============================================================================

\begin{boxjuris}
\textbf{\loicam{n°~2024/017}, article~30}~\cite{loi2024017}~:
\textit{« (1) Lorsque le responsable du traitement recourt à un
sous-traitant, il s'assure que celui-ci dispose des garanties suffisantes
quant à la mise en œuvre des mesures techniques et organisationnelles
appropriées de manière à ce que le traitement respecte les exigences de la
présente loi. (2) Le traitement par un sous-traitant est régi par un
contrat qui lie le sous-traitant au responsable du traitement. (3) Le
contrat visé à l'alinéa 2 ci-dessus définit, notamment, le type de données
à caractère personnel et les catégories de personnes concernées, les
obligations et les droits du responsable du traitement, ainsi que du
sous-traitant. »}

\smallskip
\textbf{Portée.} Le contrat est obligatoire~; son contenu minimal est
court. Le mot « notamment » laisse au responsable le soin de le compléter
pour satisfaire l'alinéa~1.
\end{boxjuris}

Trois autres articles complètent le régime. L'article~16 oblige le
sous-traitant à « apporter des garanties suffisantes au respect des
mesures de sécurité » et le responsable à veiller à leur application~;
l'article~26 répète l'obligation de garanties pour la sécurité technique et
organisationnelle~; l'article~20 assujettit le responsable et le
sous-traitant « aux mêmes obligations ». L'article~23, alinéa~3, les rend
enfin « solidairement responsables de toute publication de données à
caractère personnel effectuée en l'absence du consentement de la personne
concernée ».

Le sous-traitant doit aussi tenir son propre registre (article~29) et
informer l'Autorité et les personnes d'une violation (article~22). Un
prestataire qui ignore ces obligations n'est pas seulement un risque pour
son client~: il est lui-même en infraction.

% ============================================================================
\section{Le Contrat de Sous-traitance}
\label{sec:c10-contrat}
% ============================================================================

La loi ne fixe que trois éléments du contrat. Pour satisfaire l'obligation
de garanties suffisantes, le cours retient une liste de vérification de
douze clauses, déclarée comme critère d'audit et inspirée des normes
ISO/IEC~27701 et~27018.

\begin{tableaucours}{cL{5.2cm}L{5.8cm}}{%
    \tableauHeader{N°} &
    \tableauHeader{Clause} &
    \tableauHeader{Fondement}
}
1 & Types de données et catégories de personnes & Article~30(3) \\
2 & Obligations et droits des deux parties & Article~30(3) \\
3 & Objet, finalité et durée du traitement & Article~11 \\
4 & Traitement sur instructions documentées seulement & Article~24(2) \\
5 & Confidentialité des personnes autorisées & Articles~24 et~25(3) \\
6 & Mesures de sécurité, renvoi à l'article~27 & Articles~16, 26, 27 \\
7 & Sous-traitance ultérieure soumise à accord & Article~30(1) \\
8 & Assistance pour les droits des personnes & Articles~37 à~47 \\
9 & Information sans délai en cas de violation & Article~22 \\
10 & Restitution ou destruction en fin de contrat & Article~13 \\
11 & Droit d'audit du responsable & Article~16(2) \\
12 & Lieu du traitement, transferts éventuels & Article~32 \\
\end{tableaucours}
\vspace{-0.8em}
\begin{center}{\small\itshape Tableau~13.1. Les douze clauses du cours (critère déclaré)}\end{center}

Les clauses~1 et~2 sont exigées par la lettre de la loi~: leur absence est
un \NON{} sans discussion. Les dix autres traduisent l'obligation de
garanties~: leur absence est un \NON{} au regard du critère déclaré, et le
rapport le précise. Le laboratoire APDC-LAB-06 exerce cette lecture sur
cinq contrats.

\begin{boxinfo}
\textbf{Le droit d'audit.} La clause~11 est la condition de toutes les
autres. Sans droit d'audit, le responsable ne peut pas vérifier les
garanties qu'il doit s'assurer d'obtenir~; il ne peut que croire le
prestataire. Un prestataire qui refuse tout audit oblige l'auditeur à
coter \NEK{} chaque exigence qui dépend de lui, et le responsable à en
tirer les conséquences.
\end{boxinfo}

% ============================================================================
\section{Vérifier l'Exécution}
\label{sec:c10-execution}
% ============================================================================

Un contrat parfait ne prouve rien de ce que fait le prestataire. L'audit des
sous-traitants procède en trois cercles, du moins coûteux au plus probant~:

\begin{enumerate}
    \item \textbf{Les documents du prestataire}~: registre, politique de
          sécurité, certificats, rapports d'audit indépendants. Un
          certificat ISO/IEC~27001 ou 27018 couvre un périmètre~; on vérifie
          qu'il inclut le service fourni.
    \item \textbf{Le questionnaire et l'entretien}, recoupés par des pièces~:
          liste des personnes ayant accès, journal des accès aux données du
          client, preuve des destructions.
    \item \textbf{L'audit sur place}, exercé au titre de la clause~11, pour
          les prestataires qui portent les traitements les plus risqués.
\end{enumerate}

Le choix du cercle suit le risque. Un fournisseur de papeterie qui ne
reçoit qu'une adresse de livraison ne justifie qu'un contrôle documentaire~;
un hébergeur de dossiers médicaux justifie un audit sur place.

% ============================================================================
\section{Les Responsables Conjoints}
\label{sec:c10-conjoints}
% ============================================================================

\begin{boxjuris}
\textbf{\loicam{n°~2024/017}, article~31}~\cite{loi2024017}~:
\textit{« (1) En cas de traitement conjoint, les responsables déterminent
les finalités, ainsi que les moyens du traitement et en sont solidairement
responsables. (2) Les responsables conjoints du traitement définissent de
manière transparente, par une convention rendue publique, leurs
obligations respectives aux fins d'assurer le respect des exigences de la
présente loi, notamment en ce qui concerne leurs rôles respectifs, leurs
relations vis-à-vis de la personne concernée, l'exercice des droits de la
personne concernée, ainsi que leurs obligations respectives quant à la
communication des informations visées à l'article 21 ci-dessus.
(3) Nonobstant les termes de la convention visée à l'alinéa 2 ci-dessus,
la personne concernée exerce les droits que lui confère la présente loi à
l'égard des responsables du traitement. »}

\smallskip
\textbf{Portée.} La convention doit être \textit{rendue publique}~: c'est
une exigence plus forte que celle du droit européen, qui n'impose que de
mettre « les grandes lignes » à disposition. La personne peut s'adresser à
l'un ou l'autre responsable, quelle que soit la répartition convenue.
\end{boxjuris}

Il y a responsabilité conjointe quand deux organisations décident ensemble
de la finalité et des moyens, et pas quand l'une exécute les instructions de
l'autre. NYANGA CRÉDIT et SAWA ASSURANCES\footnote{SAWA ASSURANCES est
fictive.} proposent ensemble une assurance des emprunteurs~: NYANGA
collecte le questionnaire de santé à l'agence, SAWA ASSURANCES décide de
l'acceptation du risque, et les deux ont conçu ensemble le formulaire et
le circuit. Ce sont des responsables conjoints, et le questionnaire de
santé est une donnée sensible.

% ============================================================================
\section{Les Contrats de NYANGA et de MBOLO}
\label{sec:c10-cas}
% ============================================================================

\begin{boxscenario}{Missions NYANGA et MBOLO. Cinq contrats, 17 au 21 août 2026}
\begin{tableaucours}{L{3.4cm}cccL{4.2cm}}{%
    \tableauHeader{Contrat} &
    \tableauHeader{Cl. 1--2} &
    \tableauHeader{Cl. 3--12} &
    \tableauHeader{Audit} &
    \tableauHeader{Observation}
}
SOREC (recouvrement) & 0/2 & 0/10 & Non & Classeurs de cinq ans
conservés \\
Éditeur du logiciel bancaire & 2/2 & 6/10 & Non & Hébergement à Abidjan,
pas de clause de transfert \\
Opérateur de SMS & 1/2 & 3/10 & Non & Contrat type de l'opérateur \\
MEDISOFT (avenant 2026) & 2/2 & 10/10 & Oui & Avenant du 20~mai, audit
prévu en octobre \\
SAWA ASSURANCES (conjoint) & sans objet & sans objet & sans objet &
Convention non rédigée \\
\end{tableaucours}

\medskip
\textbf{Constatations.} Sur quatre contrats de sous-traitance, un seul
contient les douze clauses (MEDISOFT, depuis l'avenant de mai). SOREC
n'a aucune clause. L'éditeur du logiciel bancaire n'accepte pas d'audit et
ne mentionne pas le pays d'hébergement. Aucune convention de
responsabilité conjointe n'existe avec SAWA ASSURANCES.

\textbf{Analyse} (confiance élevée sur les contrats, faible sur
l'exécution). Les écarts contractuels sont établis. Pour l'exécution, seul
SOREC a été vu~: la conservation de cinq ans de classeurs est un écart à
l'article~13, dont NYANGA répond aussi (article~20).
\end{boxscenario}

\begin{boxlecture}{Le domaine D4, sous-traitants, à NYANGA}
\begin{tableaucours}{L{7.8cm}cL{3.2cm}}{%
    \tableauHeader{Exigence} &
    \tableauHeader{Cote} &
    \tableauHeader{Preuve}
}
Chaque sous-traitant est lié par un contrat & \OUI{} & Trois contrats \\
Les contrats définissent données et personnes (art.~30(3)) & \NON{} &
SOREC, opérateur \\
Les contrats contiennent les douze clauses & \NON{} & Un sur trois au
mieux \\
Le droit d'audit est prévu & \NON{} & Aucun contrat \\
SOREC détruit les données des dossiers clos & \NON{} & Classeurs
conservés \\
L'éditeur applique les mesures de sécurité convenues & \NEK{} & Pas
d'audit, pas de rapport \\
L'opérateur de SMS limite l'usage des numéros & \NEK{} & Aucune pièce \\
La convention avec SAWA ASSURANCES est publique & \NON{} & Inexistante \\
\end{tableaucours}

\medskip
\textbf{Décompte.} Un \OUI{}, cinq \NON{}, deux \NEK{} sur huit.
\hfill\lecture{0{,}125}{0{,}25}{1{,}30}

\medskip
\textbf{Lecture.} Domaine opaque ($U = 0{,}25$)~: on ne sait pas ce que font
deux des trois prestataires. Le manque de droit d'audit explique les deux
\NEK{}~: corriger les contrats est la condition pour lever l'incertitude.
\end{boxlecture}

% ============================================================================
\section{Ce que la Sous-traitance demande à chaque Métier}
\label{sec:c10-metiers}
% ============================================================================

\begin{boxmetier}{Juriste conformité}
Rédiger un avenant type aux douze clauses, le négocier en commençant par les
prestataires les plus exposés, et tenir à jour la liste des sous-traitants
avec leur contrat.
\end{boxmetier}

\begin{boxmetier}{Responsable de traitement}
Refuser ou quitter un prestataire qui refuse le droit d'audit pour un
traitement sensible. La responsabilité solidaire et l'article~20 font que
ses manquements sont aussi les vôtres.
\end{boxmetier}

\begin{boxmetier}{Auditeur de la protection des données}
Distinguer dans le rapport les écarts contractuels (établis sur pièces) et
les écarts d'exécution (établis seulement là où l'on a pu vérifier), et
coter \NEK{} l'exécution des prestataires non vus.
\end{boxmetier}

\begin{boxresume}
\begin{itemize}
    \item Tout recours à un sous-traitant exige un contrat qui définit au
          moins les données, les personnes, les obligations et les droits
          des parties (article~30).
    \item Le responsable s'assure des garanties du sous-traitant et veille à
          leur application (articles~16 et~30)~; les deux ont les mêmes
          obligations (article~20).
    \item Le cours déclare un critère de douze clauses, dont le droit
          d'audit, condition de toutes les autres.
    \item Les responsables conjoints fixent leurs obligations dans une
          convention rendue publique (article~31).
    \item À NYANGA, aucun contrat ne prévoit d'audit, et SOREC conserve cinq
          ans de données de débiteurs.
\end{itemize}
\end{boxresume}

\begin{boxinfo}
\textbf{Points de contrôle.} Vous avez compris ce chapitre si vous savez~:
\begin{itemize}
    \item contrôler un contrat contre les douze clauses~;
    \item distinguer sous-traitant et responsable conjoint~;
    \item choisir le cercle de vérification adapté au risque d'un
          prestataire.
\end{itemize}
\end{boxinfo}

\begin{boxexo}[title={\ding{45}\quad Exercice~13.1. Sous-traitant ou conjoint \niveaubase}]
Qualifiez la relation~: (a) un cabinet comptable qui établit la paie de
MBOLO~; (b) un laboratoire d'analyses qui reçoit les prélèvements et
envoie les résultats au médecin~; (c) l'hébergeur de NOUN PAY~; (d) NYANGA
et SAWA ASSURANCES pour l'assurance des emprunteurs.
\end{boxexo}

\begin{boxexo}[title={\ding{45}\quad Exercice~13.2. Lire un contrat \niveaumoyen}]
Le contrat de l'opérateur de SMS prévoit que « l'opérateur pourra utiliser
les données d'acheminement pour améliorer ses services et proposer des
offres à ses abonnés ». Qualifiez l'opérateur pour cette utilisation, et
dites quelles clauses vous ajoutez.
\end{boxexo}

\begin{boxexo}[title={\ding{45}\quad Exercice~13.3. La convention publique \niveaumoyen}]
Rédigez les rubriques de la convention entre NYANGA et SAWA ASSURANCES
exigée par l'article~31. Qui recueille le consentement au questionnaire de
santé, et sur quel fondement~?
\end{boxexo}

\begin{boxexo}[title={\ding{45}\quad Exercice~13.4. Le refus d'audit \niveauexpert}]
L'éditeur du logiciel bancaire refuse tout audit, invoquant la sécurité de
ses autres clients. Proposez trois solutions qui permettraient à NYANGA de
s'assurer des garanties exigées par l'article~30, et dites comment vous
coterez la proposition « l'éditeur applique les mesures convenues » dans
chaque cas.
\end{boxexo}

\bigskip

\begin{boxinfo}
\textbf{Pour aller plus loin.}
\begin{itemize}
    \item L'ISO/IEC~27701~\cite{iso27701}, pour les obligations des
          sous-traitants.
    \item L'ISO/IEC~27018~\cite{iso27018}, pour l'hébergement en nuage.
\end{itemize}
\end{boxinfo}

\begin{boxinfo}
\textbf{Transition.} L'éditeur du logiciel bancaire héberge les comptes à
Abidjan, et l'hébergeur de NOUN PAY en Europe. Les données quittent le
Cameroun~: c'est un transfert, soumis à autorisation. C'est l'objet du
Chapitre~\ref{chap:transferts}.
\end{boxinfo}

        % ============================================================================
% Fichier  : partie5/P05_C14_transferts.tex
% Titre    : Les Transferts Internationaux
% Auteur   : MINKA MI NGUIDJOÏ Thierry Emmanuel
% Version  : 1.0 -- novembre 2026
% Corrigés : corriges/Corr_C14.tex
% ============================================================================
% Sources : loi n° 2024/017, art. 5 (transfert international, clauses
% contractuelles types), 21, 32, 53, 60, 69 (texte officiel) ; ISO/IEC 27018 ;
% recommandations 01/2020 du CEPD (comparaison) ; Minka Mi Nguidjoï,
% Massoussi Songo et Bialo (2026, manuscrit) pour la tension T3.
% Faits nouveaux (bible) : cartographie des flux de NOUN PAY, 21 au
% 25 septembre 2026.
% ============================================================================

\chapter{Les Transferts Internationaux}
\label{chap:transferts}
\chaptermark{Les transferts internationaux}

\epigraph{%
    « Une donnée hébergée à l'étranger a quitté le Cameroun,\\
    même si on la consulte depuis Bafoussam. »%
}{Principe du cours}

\niveauChapitre{Intermédiaire}{%
    Chapitres~\ref{chap:cartographie}, \ref{chap:loi} et~\ref{chap:soustraitants}
}

\begin{boxobjectifs}
\begin{itemize}
    \item Reconnaître un transfert international au sens de l'article~5,
          y compris l'hébergement et l'accès à distance.
    \item Appliquer les quatre vérifications de l'article~32.
    \item Cartographier les flux sortants d'une organisation à partir de
          sources indépendantes.
    \item Coter les transferts quand l'Autorité n'a publié ni liste de pays
          ni clauses types.
    \item Situer la tension entre transferts et coopération pénale.
\end{itemize}
\end{boxobjectifs}

\bigskip

% ============================================================================
\section*{Avant de commencer : « nos serveurs, c'est le nuage »}
% ============================================================================

Bafoussam, lundi 21~septembre~2026. Mme~KENFACK, responsable produit de NOUN
PAY, présente l'architecture de l'application à KORO AUDIT~: « Tout est dans
le nuage. » Où, exactement~? Elle cherche. L'hébergeur principal est
établi dans l'Union européenne~; la vérification d'identité par selfie est
confiée à un prestataire dont elle ne connaît pas le pays~; l'outil de
mesure d'audience et la messagerie d'entreprise sont américains~; le
support client de nuit est assuré par un centre d'appels au Maroc, qui
consulte les comptes à distance.

Aucun de ces flux n'a été autorisé. Mme~KENFACK pensait qu'un transfert
supposait d'« envoyer des fichiers à l'étranger »~; l'hébergement et la
consultation à distance ne lui semblaient pas en être.

\begin{boxalert}
\textbf{Violation typique.} Réduire le transfert à l'envoi de fichiers, et
oublier l'hébergement étranger, les services en ligne et la consultation
à distance depuis l'étranger.

\smallskip
\textbf{Règle~:} tout acheminement de données au-delà des frontières est un
transfert, quel qu'en soit le moyen. Il exige une autorisation préalable de
l'Autorité, après les vérifications de l'article~32.
\end{boxalert}

% ============================================================================
\section{Qu'est-ce qu'un Transfert}
\label{sec:c11-definition}
% ============================================================================

\begin{boxdef}
\termecle{Transfert international de données}~: selon l'article~5,
« acheminement des données à caractère personnel au-delà des frontières
d'un pays vers un autre, que ce soit vers un État tiers ou une organisation
internationale ».
\end{boxdef}

La définition vise l'\textit{acheminement}, pas le moyen. Sont des
transferts~: l'envoi d'un fichier à un partenaire étranger~; l'hébergement
d'une base chez un prestataire étranger~; l'utilisation d'un service en
ligne dont les serveurs sont à l'étranger (messagerie, formulaires, mesure
d'audience)~; la consultation à distance de données stockées au Cameroun
par une personne située à l'étranger, puisque les données traversent la
frontière pour s'afficher sur son écran. L'auditeur qui hésite sur un cas
le cote comme transfert et le signale~: la qualification restrictive
exposerait l'organisation aux articles~60 et~69.

% ============================================================================
\section{L'Article 32}
\label{sec:c11-article32}
% ============================================================================

\begin{boxjuris}
\textbf{\loicam{n°~2024/017}, article~32}~\cite{loi2024017}~:
\textit{« (1) Le transfert de données à caractère personnel vers un Etat
étranger, ou à une organisation internationale, est subordonné à
l'obtention d'une autorisation préalable délivrée par l'Autorité de
protection des données à caractère personnel dans les conditions
permettant de garantir l'exercice des droits de la personne concernée.
(2) A l'occasion de la délivrance de ladite autorisation, l'Autorité de
protection des données à caractère personnel doit préalablement
s'assurer~: que le pays de destination des données à caractère personnel
offre un niveau de protection suffisant~; de l'entrée en vigueur préalable
d'un instrument juridique signé avec le pays de destination des données à
caractère personnel transférées, en liaison avec les administrations et
organismes compétents~; de l'assujettissement de l'entité sollicitant
l'importation des données à caractère personnel aux règles contraignantes
de sécurité en matière de protection desdites données~; de la souscription
préalable, par les entités importatrices et exportatrices concernées, de
clauses contractuelles types en matière de transfert à l'international des
données à caractère personnel, édictées et publiées par l'Autorité de
protection des données à caractère personnel. (3) Les modalités de
délivrance de l'autorisation prévue à l'alinéa 1 ci-dessus sont fixées par
voie réglementaire. »}

\smallskip
\textbf{Portée.} Les quatre vérifications sont énumérées sans « ou »~: la
lecture littérale les rend cumulatives. Le deuxième point, un instrument
juridique « signé avec le pays de destination », est particulièrement
exigeant~; il suppose un accord d'État à État.
\end{boxjuris}

L'article~53 charge l'Autorité d'établir la liste des pays « reconnus comme
offrant un niveau de protection [\ldots] équivalent » et l'article~5
définit les \termecle{clauses contractuelles types} comme celles que
l'Autorité élabore et publie pour encadrer les transferts. Ni la liste ni
les clauses n'étaient publiées en novembre~2026.

\subsection{Coter sans liste ni clauses}

L'auditeur se trouve devant quatre conditions dont trois dépendent
d'instruments publics inexistants. Il décompose~:

\begin{itemize}
    \item « Le transfert est autorisé »~: \NON{} s'il n'existe pas
          d'autorisation, comme pour tout traitement.
    \item « Le pays de destination offre un niveau suffisant »~: \NEK{} de
          critère tant que la liste n'existe pas~; l'auditeur peut décrire,
          à titre d'information, le cadre du pays (une loi de protection,
          une autorité), sans en tirer de cote.
    \item « L'importateur est soumis à des règles contraignantes de
          sécurité »~: cotable sur pièces (contrat, certificat ISO/IEC~27018,
          rapport d'audit).
    \item « Les clauses contractuelles types sont souscrites »~: \NEK{} de
          critère tant qu'elles ne sont pas publiées.
    \item « L'organisation a documenté chaque transfert et préparé sa
          demande »~: cotable sur pièces.
\end{itemize}

\begin{boxinfo}
\textbf{Le modèle européen, en comparaison.} Le droit européen admet des
transferts fondés sur une décision d'adéquation, des clauses
contractuelles types adoptées par la Commission, des règles d'entreprise
contraignantes, ou des dérogations~; il demande d'évaluer, au cas par cas,
la législation du pays de destination~\cite{edpb2020transferts}. La loi
camerounaise ajoute à cet arsenal l'autorisation de l'Autorité et
l'instrument signé d'État à État. Les outils européens peuvent servir de
bonne pratique documentaire~; ils ne valent pas autorisation.
\end{boxinfo}

% ============================================================================
\section{Cartographier les Flux Sortants}
\label{sec:c11-flux}
% ============================================================================

On établit la liste des transferts à partir de trois sources
indépendantes~:

\begin{enumerate}
    \item \textbf{Les contrats et factures}~: chaque service payé à un
          prestataire étranger est un transfert probable~; le contrat dit
          parfois où sont les serveurs.
    \item \textbf{Les configurations et le code}~: adresses des services
          appelés par l'application, bibliothèques intégrées (mesure
          d'audience, vérification d'identité, notifications).
    \item \textbf{Les traces réseau}~: noms de domaine résolus et volumes
          échangés vers l'extérieur, rapprochés d'une base de localisation
          des adresses.
\end{enumerate}

La localisation d'une adresse IP est approximative~: un réseau de diffusion
de contenu répond depuis le point le plus proche, et un prestataire peut
stocker ailleurs qu'il ne répond. L'auditeur ne déduit pas le lieu de
stockage d'une adresse~; il le cherche dans le contrat ou la documentation
du prestataire, et cote \NEK{} quand il ne le trouve pas. Le laboratoire
APDC-LAB-07 exerce ce rapprochement.

% ============================================================================
\section{Les Flux de NOUN PAY}
\label{sec:c11-nounpay}
% ============================================================================

\begin{boxscenario}{Mission NOUN PAY. Six flux sortants, 21 au 25 septembre 2026}
\begin{tableaucours}{L{3.2cm}L{3.0cm}L{2.6cm}L{3.2cm}}{%
    \tableauHeader{Flux} &
    \tableauHeader{Données} &
    \tableauHeader{Pays} &
    \tableauHeader{Source de la preuve}
}
Hébergement principal & Toutes & Union européenne & Contrat, console \\
Vérification d'identité & Selfies, pièces d'identité & \NEK{} & Code de
l'application~; contrat muet \\
Mesure d'audience & Identifiants, parcours & États-Unis & Code, trafic \\
Messagerie d'entreprise & Échanges avec les clients & États-Unis &
Facture, documentation \\
Support de nuit & Consultation des comptes & Maroc & Contrat, journaux
d'accès \\
SMS & Numéros, messages & Cameroun & Contrat (pas un transfert) \\
\end{tableaucours}

\medskip
\textbf{Constatations.} Cinq transferts sur six flux sortants. Aucun n'est
autorisé, aucun n'est mentionné dans l'information des clients
(article~21). Le pays de traitement des selfies, données biométriques,
n'est établi par aucune pièce.

\textbf{Analyse} (confiance élevée). Le modèle d'entreprise de NOUN PAY
repose sur des services étrangers~: le rapatriement complet n'est pas
réaliste à court terme. La priorité est de documenter chaque transfert, de
préparer les demandes d'autorisation, d'informer les clients, et de lever
l'inconnu sur la biométrie.
\end{boxscenario}

\begin{boxlecture}{Les transferts de NOUN PAY}
\begin{tableaucours}{L{7.8cm}cL{3.2cm}}{%
    \tableauHeader{Exigence} &
    \tableauHeader{Cote} &
    \tableauHeader{Preuve}
}
Les transferts sont recensés & \OUI{} & Cartographie de l'audit \\
Les transferts sont autorisés & \NON{} & Aucune autorisation \\
Les clients sont informés des transferts & \NON{} & Notice de
l'application \\
Le pays de chaque transfert est connu & \NEK{} & Selfies, pays inconnu \\
Les importateurs sont soumis à des règles de sécurité & \OUI{} & Certificat
ISO/IEC 27018 de l'hébergeur \\
Le niveau de protection des pays est suffisant & \NEK{} & Liste non
publiée \\
Les clauses contractuelles types sont souscrites & \NEK{} & Clauses non
publiées \\
Les demandes d'autorisation sont préparées & \NON{} & Aucun dossier \\
\end{tableaucours}

\medskip
\textbf{Décompte.} Deux \OUI{}, trois \NON{}, trois \NEK{} sur huit.
\hfill\lecture{0{,}25}{0{,}375}{1{,}56}

\medskip
\textbf{Lecture.} $U = 0{,}375$~: le domaine est fortement opaque. Mais deux
des trois \NEK{} sont de critère et ne se lèveront qu'avec des textes~; un
seul est de preuve (le pays des selfies), et c'est lui qu'il faut lever en
priorité, parce qu'il porte sur des données biométriques.
\end{boxlecture}

\textbf{Décrire le pays de destination.} Pour le Maroc et pour la Côte
d'Ivoire, où se trouve l'éditeur de NYANGA, l'auditeur dispose des fiches
pays du Chapitre~\ref{chap:international}
(section~\ref{sec:c03-pays})~: une loi de protection des données, une
autorité, des engagements internationaux. Ces éléments entrent dans le
rapport à titre d'information. Ils ne lèvent pas le \NEK{} de critère sur
le niveau de protection, que seule la liste de l'Autorité camerounaise
pourra lever.

\subsection{Les données qui reviennent de l'Union européenne}

Un transfert a deux sens. Les 6\,840 clients de NOUN DIASPORA résident en
France et en Belgique. Leurs données sont collectées dans l'Union, puis
consultées depuis Douala par les équipes de NOUN PAY. Pour le droit
européen, cette consultation est un transfert vers un pays tiers, soumis
au chapitre~V du règlement. Le Cameroun ne bénéficie pas d'une décision
d'adéquation de la Commission. Le transfert doit donc s'appuyer sur une
autre garantie, en pratique les clauses contractuelles types de
2021~\cite{ue2021clauses}, complétées par une évaluation de la législation
du pays de destination~\cite{edpb2020transferts}.

L'auditeur cote alors deux exigences distinctes pour le même flux. Au
titre de la loi camerounaise, le traitement à Douala de données venues de
l'étranger reste un traitement soumis à l'article~19. Au titre du
règlement, l'exigence est « une garantie du chapitre~V couvre le
transfert vers le Cameroun », cotable sur pièces. À NOUN PAY, aucune
clause n'a été signée~: la cote est \NON{}. Cette double lecture est
celle du Chapitre~\ref{chap:rgpd} (section~\ref{sec:c04-double}).

% ============================================================================
\section{Transferts et Coopération Pénale}
\label{sec:c11-penal}
% ============================================================================

Les transferts ne concernent pas que les entreprises. Quand un officier de
police judiciaire camerounais demande à un fournisseur étranger de
conserver des données, ou quand une autorité étrangère demande au Cameroun
des données sur un suspect, des données personnelles franchissent une
frontière. La Convention de Budapest, à laquelle le Cameroun est Partie,
prévoit en son article~29 une procédure de conservation rapide entre
États. La loi n°~2024/017 soumet les transferts à autorisation préalable
(article~32), sous réserve des textes spécifiques à la justice
(article~72).

L'auteur de ce cours et ses coauteurs analysent cette rencontre comme une
« tension transfrontalière » dans un travail en cours~\cite{minka2026chainepenale}~:
elle n'est pas, selon eux, une incompatibilité, mais elle impose de
distinguer conservation, préservation, production, transmission et
transfert, qui peuvent relever de régimes différents. Le
Chapitre~\ref{chap:recherche} présente cette piste de recherche.

% ============================================================================
\section{Ce que les Transferts demandent à chaque Métier}
\label{sec:c11-metiers}
% ============================================================================

\begin{boxmetier}{Délégué à la protection des données}
Tenir une liste des transferts, avec pour chacun le pays, les données, le
prestataire, le contrat et l'état de la demande d'autorisation. Exiger de
chaque nouveau service en ligne la localisation de ses serveurs avant de
l'adopter.
\end{boxmetier}

\begin{boxmetier}{RSSI}
Fournir les traces des flux sortants et l'inventaire des bibliothèques
intégrées aux applications~; bloquer les services non déclarés.
\end{boxmetier}

\begin{boxmetier}{Juriste conformité}
Préparer les dossiers d'autorisation~: description du transfert, garanties
contractuelles, informations sur le pays, et, en attendant les clauses
types de l'Autorité, des clauses de protection équivalentes.
\end{boxmetier}

\begin{boxresume}
\begin{itemize}
    \item Tout acheminement de données au-delà des frontières est un
          transfert, y compris l'hébergement, les services en ligne et la
          consultation à distance.
    \item L'article~32 subordonne le transfert à une autorisation préalable,
          après quatre vérifications dont un instrument signé avec le pays
          de destination.
    \item Sans liste de pays ni clauses types publiées, deux de ces
          vérifications sont des \NEK{} de critère~; les autres se cotent sur
          pièces.
    \item Le lieu de stockage se prouve par le contrat ou la documentation,
          pas par la localisation d'une adresse IP.
    \item Le transfert sans autorisation est puni administrativement
          (article~60) et pénalement (article~69).
\end{itemize}
\end{boxresume}

\begin{boxinfo}
\textbf{Points de contrôle.} Vous avez compris ce chapitre si vous savez~:
\begin{itemize}
    \item dire si un service en ligne constitue un transfert~;
    \item énoncer les quatre vérifications de l'article~32~;
    \item distinguer un \NEK{} de critère d'un \NEK{} de preuve dans un audit
          des transferts.
\end{itemize}
\end{boxinfo}

\begin{boxexo}[title={\ding{45}\quad Exercice~14.1. Transfert ou non \niveaubase}]
(a) Un médecin de MBOLO consulte un dossier depuis un congrès à Paris~;
(b) NYANGA stocke ses sauvegardes chez un prestataire à Douala~; (c) NOUN
PAY envoie ses SMS par un opérateur camerounais~; (d) le logiciel bancaire
de NYANGA est hébergé à Abidjan~; (e) un formulaire de satisfaction en ligne
dont les serveurs sont aux États-Unis.
\end{boxexo}

\begin{boxexo}[title={\ding{45}\quad Exercice~14.2. Le dossier d'autorisation \niveaumoyen}]
Rédigez le dossier de demande d'autorisation du transfert d'hébergement de
NOUN PAY, rubrique par rubrique, en vous appuyant sur l'article~32.
Indiquez les pièces que vous joindrez.
\end{boxexo}

\begin{boxexo}[title={\ding{45}\quad Exercice~14.3. Les selfies \niveaumoyen}]
Proposez une démarche pour établir le pays de traitement des selfies.
Quelles pièces demandez-vous au prestataire~? Que recommandez-vous si
NOUN PAY ne parvient pas à l'établir~?
\end{boxexo}

\begin{boxexo}[title={\ding{45}\quad Exercice~14.4. Abidjan \niveauexpert}]
Le logiciel bancaire de NYANGA est hébergé en Côte d'Ivoire, qui dispose de
sa propre loi de protection des données. (a) Ce fait suffit-il à satisfaire
la première vérification de l'article~32~? (b) Quelles autres
vérifications restent à satisfaire~? (c) Comment cotez-vous la proposition
« le transfert vers Abidjan est conforme » au 30~septembre~2026~?
\end{boxexo}

\bigskip

\begin{boxinfo}
\textbf{Pour aller plus loin.}
\begin{itemize}
    \item Les recommandations 01/2020 du Comité européen de la protection des
          données~\cite{edpb2020transferts}, pour la méthode d'évaluation
          d'un pays de destination.
    \item L'ISO/IEC~27018~\cite{iso27018}, pour les obligations des
          hébergeurs.
\end{itemize}
\end{boxinfo}

\begin{boxinfo}
\textbf{Transition.} Une donnée qui sort sans autorisation est un écart~;
une donnée qui sort par effraction est une violation. Comment la détecter,
la qualifier et en informer qui de droit~? C'est l'objet du
Chapitre~\ref{chap:violations}.
\end{boxinfo}

        % ============================================================================
% Fichier  : partie5/P05_C15_violations.tex
% Titre    : Les Violations de Données
% Auteur   : MINKA MI NGUIDJOÏ Thierry Emmanuel
% Version  : 1.0 -- novembre 2026
% Corrigés : corriges/Corr_C15.tex
% ============================================================================
% Sources : loi n° 2024/017, art. 5 (violation), 22, 63, 73 (texte
% officiel) ; ISO/IEC 27035 ; NIST SP 800-61 r2 ; Minka Mi Nguidjoï,
% Massoussi Songo et Bialo (2026, manuscrit) pour la tension T2 et les
% articles 102 et 154 du Code de procédure pénale cités par eux.
% Faits nouveaux (bible) : lettre de l'OPJ du 20/03/2025 ; aucune
% information des clients ; MBOLO, consultations de mai 2026.
% ============================================================================

\chapter{Les Violations de Données}
\label{chap:violations}
\chaptermark{Les violations de données}

\epigraph{%
    « Une violation tenue secrète protège l'organisation un temps,\\
    et laisse la personne sans défense. »%
}{Principe du cours}

\niveauChapitre{Intermédiaire}{%
    Chapitres~\ref{chap:securite} et~\ref{chap:loi}~; réponse à incident
    (UE MSN-2S3-SEC07)
}

\begin{boxobjectifs}
\begin{itemize}
    \item Reconnaître une violation de données au sens de l'article~5,
          y compris l'accès interne non autorisé.
    \item Appliquer l'obligation d'information « sans délai » de
          l'article~22.
    \item Qualifier une violation du point de vue des personnes.
    \item Auditer le dispositif de gestion des violations et le registre
          qui les documente.
    \item Analyser la rencontre entre l'obligation d'informer et le secret
          d'une enquête pénale.
\end{itemize}
\end{boxobjectifs}

\bigskip

% ============================================================================
\section*{Avant de commencer : une lettre de l'OPJ}
% ============================================================================

Bafoussam, vendredi 21~août~2026. Dans le dossier de l'incident de
mars~2025, Mme~BELLO trouve une lettre de l'adjudant-chef NANA, officier
de police judiciaire, datée du 20~mars~2025~: « Dans l'intérêt de
l'enquête, il vous est demandé de ne pas informer vos clients de la
compromission avant instruction contraire. » Aucune instruction contraire
n'a été donnée. Le jugement a été rendu en février~2026. Le 21~août, les
18\,412 personnes dont les dossiers ont été copiés n'ont toujours reçu
aucune information.

Le directeur général s'en tient à la lettre~: « La police nous a demandé de
nous taire. » L'auditrice note que la lettre visait l'enquête, que
l'enquête est close, et que personne n'a vérifié depuis si la demande
tenait toujours.

\begin{boxalert}
\textbf{Violation typique.} Traiter une violation comme un problème
technique à corriger en silence, ou s'abriter indéfiniment derrière une
demande de discrétion devenue sans objet.

\smallskip
\textbf{Règle~:} une violation s'informe « sans délai » à l'Autorité et aux
personnes (article~22). Une raison de différer, si elle existe, se
documente, se limite dans le temps et se réexamine~; quand elle cesse,
l'information est due.
\end{boxalert}

% ============================================================================
\section{Qu'est-ce qu'une Violation}
\label{sec:c12-definition}
% ============================================================================

\begin{boxdef}
\termecle{Violation de données à caractère personnel}~: selon l'article~5,
« toute atteinte à la sécurité des données à caractère personnel
entraînant, de manière accidentelle ou illicite, la destruction, la perte,
l'altération, la divulgation non autorisée desdites données transmises,
conservées ou traitées d'une autre manière, ou l'accès non autorisé à de
telles données ».
\end{boxdef}

La définition est large. Une violation n'est pas forcément une attaque~:
une erreur d'envoi, un ordinateur portable perdu, une sauvegarde effacée
par erreur en sont. Elle n'est pas forcément externe~: l'accès non autorisé
d'un agent à un dossier, comme à MBOLO, en est une. Elle n'est pas
forcément une fuite~: la destruction ou l'indisponibilité de données
indispensables en est une aussi.

\begin{tableaucours}{L{3.0cm}L{4.6cm}L{4.4cm}}{%
    \tableauHeader{Atteinte} &
    \tableauHeader{Exemple} &
    \tableauHeader{Risque pour la personne}
}
Confidentialité & Fichier clients copié, dossier consulté sans raison,
courriel au mauvais destinataire & Fraude, chantage, discrimination,
atteinte à la réputation \\
Intégrité & Résultat d'analyse attribué au mauvais patient & Soin inadapté,
décision erronée \\
Disponibilité & Dossiers chiffrés par un rançongiciel, sauvegardes
détruites & Soin retardé, droits impossibles à exercer \\
\end{tableaucours}
\vspace{-0.8em}
\begin{center}{\small\itshape Tableau~15.1. Les trois atteintes et leurs effets}\end{center}

% ============================================================================
\section{L'Obligation d'Informer}
\label{sec:c12-article22}
% ============================================================================

\begin{boxjuris}
\textbf{\loicam{n°~2024/017}, article~22}~\cite{loi2024017}~:
\textit{« (1) Dès connaissance de toute violation des données à caractère
personnel, le responsable du traitement ou le sous-traitant est tenu
d'informer sans délai, l'Autorité de protection des données à caractère
personnel et la personne concernée. (2) En tout état de cause, et sans
préjudice des dispositions de l'alinéa 1 ci-dessus, le responsable du
traitement ou le sous-traitant est tenu de mettre en œuvre les mesures de
protection techniques et organisationnelles appropriées, telles que le
chiffrement, l'instauration des autorisations restrictives d'accès aux
données. »}

\smallskip
\textbf{Portée.} L'obligation est plus large que dans le droit européen sur
deux points~: elle vise \textit{toute} violation, sans seuil de risque, et
elle impose d'informer la personne dans tous les cas. Elle est moins
précise sur un troisième~: « sans délai » n'est pas chiffré, et le contenu
de l'information n'est pas fixé.
\end{boxjuris}

En l'absence de contenu fixé, le cours retient une information minimale,
déclarée comme critère~: la nature de la violation, les catégories et le
nombre approximatif de personnes et de données touchées, les conséquences
probables pour les personnes, les mesures prises et celles que la personne
peut prendre elle-même (changer un mot de passe, surveiller ses comptes),
et un contact. Pour « sans délai », l'audit mesure le temps écoulé entre la
connaissance de la violation et chaque information, et le rapproche des
raisons documentées d'un éventuel retard.

\begin{boxinfo}
\textbf{Le sous-traitant aussi.} L'article~22 vise « le responsable du
traitement ou le sous-traitant ». Le contrat doit organiser qui informe
qui, et dans quel ordre~: un prestataire qui découvre une violation doit
au moins prévenir immédiatement son client (clause~9 du
Chapitre~\ref{chap:soustraitants}).
\end{boxinfo}

% ============================================================================
\section{Gérer une Violation}
\label{sec:c12-gerer}
% ============================================================================

\begin{boxprocedure}
\textbf{Auditer le dispositif de gestion des violations.}
\begin{enumerate}
    \item Vérifier qu'une procédure écrite existe, connue des équipes et des
          sous-traitants, avec un point de contact unique.
    \item Vérifier la détection~: comment l'organisation apprend-elle une
          violation (alertes, revues de journaux, signalements)~?
    \item Obtenir le registre des violations~; à défaut, le reconstituer à
          partir des incidents de sécurité, des plaintes et des sinistres.
    \item Pour chaque violation, relever les dates de connaissance, de
          confinement, d'information de l'Autorité et des personnes.
    \item Vérifier la qualification~: données, personnes, conséquences
          pour elles.
    \item Vérifier le contenu des informations envoyées.
    \item Vérifier les enseignements tirés et les mesures correctives.
\end{enumerate}
\end{boxprocedure}

La loi n'impose pas explicitement un registre des violations. Il découle
cependant de deux obligations~: prouver qu'on a informé « sans délai »
suppose de dater la connaissance~; et la traçabilité de l'article~27
suppose de pouvoir reconstituer ce qui s'est passé. Un registre où chaque
violation, même non notifiée, est décrite avec sa date, sa qualification
et sa décision, est la seule manière de le démontrer.

% ============================================================================
\section{Informer pendant une Enquête Pénale}
\label{sec:c12-penal}
% ============================================================================

L'affaire de NYANGA pose une question que la loi ne règle pas. L'article~22
impose d'informer sans délai les personnes~; la procédure pénale protège
le secret de l'enquête. Dans leur travail en cours, l'auteur de ce cours et
ses coauteurs relèvent que le Code de procédure pénale impose le secret de
la procédure durant l'enquête de police judiciaire (article~102) et le
secret de l'information judiciaire (article~154)~\cite{minka2026chainepenale}.
Une information des personnes peut révéler l'existence de l'enquête,
l'identité d'un suspect ou une opération en cours.

Les auteurs qualifient cette rencontre de « tension procédurale », et non
d'incompatibilité~: les deux obligations peuvent être conciliées si l'on
organise le moment et les modalités de l'information. Ils proposent qu'un
mécanisme permette au magistrat compétent d'apprécier ce moment lorsque
l'information risque de compromettre l'enquête. En attendant un tel
mécanisme, l'auditeur vérifie quatre choses~:

\begin{enumerate}
    \item la demande de différer existe par écrit, émane d'une autorité
          compétente et vise l'enquête~;
    \item l'Autorité de protection a été informée, ou l'organisation a
          documenté pourquoi elle ne pouvait pas l'être~;
    \item l'organisation a réexaminé la demande à intervalles réguliers~;
    \item l'information a été donnée dès que la raison de différer a cessé.
\end{enumerate}

S'y ajoute une question de temps, que les mêmes auteurs nomment « tension
temporelle »~: la violation de NYANGA date du 14~mars~2025, après la
promulgation de la loi mais pendant le délai de mise en conformité de
l'article~73. L'obligation d'informer s'appliquait-elle déjà~? L'auditeur
ne tranche pas cette question de droit~; il constate en revanche qu'au
24~juin~2026, à l'échéance du délai, la violation n'était pas réparée du
point de vue des personnes, et qu'aucune raison documentée ne justifiait
plus le silence. Le Chapitre~\ref{chap:recherche} présente ces deux
tensions comme des pistes de recherche.

% ============================================================================
\section{Les Violations de NYANGA et de MBOLO}
\label{sec:c12-cas}
% ============================================================================

\begin{boxscenario}{Missions NYANGA et MBOLO. Deux violations, deux traitements}
\textbf{NYANGA.} Connaissance~: 14~mars~2025, 9~heures. Qualification
technique faite par l'enquête (3,2~Go, 18\,412 fichiers). Aucune
qualification du point de vue des personnes. Demande de l'OPJ du
20~mars~2025, jamais réexaminée. Jugement le 15~février~2026. Au
31~août~2026, aucune information des personnes ni de l'Autorité~; aucun
registre des violations.

\textbf{MBOLO.} Connaissance~: 13~mai~2026, par la plainte de la patiente.
Investigation le jour même (journal des consultations), comptes désactivés,
procédure disciplinaire. Patiente informée par écrit le 15~mai, avec la
liste des consultations non justifiées et les mesures prises. L'Autorité
n'a pas été informée~; la directrice a noté au registre, créé le 14~mai,
qu'aucun interlocuteur n'était identifié.

\textbf{Analyse} (confiance élevée). MBOLO a traité sa violation comme la
loi le demande pour ce qui dépendait de lui, et l'a documentée. NYANGA a
traité la sienne comme un dossier pénal, sans jamais la regarder comme une
violation de données envers ses clients.
\end{boxscenario}

\begin{boxlecture}{Le domaine D5, violations, à NYANGA}
\begin{tableaucours}{L{7.8cm}cL{3.2cm}}{%
    \tableauHeader{Exigence} &
    \tableauHeader{Cote} &
    \tableauHeader{Preuve}
}
Une procédure de gestion des violations existe & \NON{} & Aucune \\
Un registre des violations est tenu & \NON{} & Aucun \\
La violation de 2025 est qualifiée pour les personnes & \NON{} & Rapport
technique seul \\
La demande de différer est écrite et motivée & \OUI{} & Lettre du
20~mars~2025 \\
La demande a été réexaminée & \NON{} & Aucune démarche \\
Les personnes ont été informées quand la raison a cessé & \NON{} & Rien au
31~août~2026 \\
L'Autorité a été informée & \NEK{} & Autorité non organisée, aucune
démarche documentée \\
L'obligation d'informer s'appliquait en mars~2025 & \NEK{} & Question de
droit non tranchée \\
\end{tableaucours}

\medskip
\textbf{Décompte.} Un \OUI{}, cinq \NON{}, deux \NEK{} sur huit.
\hfill\lecture{0{,}125}{0{,}25}{1{,}30}

\medskip
\textbf{Lecture.} Les deux \NEK{} ne changent pas la recommandation~: quelle
que soit la réponse à la question de droit, la raison de se taire a cessé
en février~2026, et l'information des 18\,412 personnes est due. C'est la
première recommandation du rapport.
\end{boxlecture}

% ============================================================================
\section{Ce que les Violations demandent à chaque Métier}
\label{sec:c12-metiers}
% ============================================================================

\begin{boxmetier}{Responsable de traitement}
Désigner un point de contact pour les violations, décider de l'information
et la signer. Réexaminer par écrit toute demande de discrétion et lever le
silence dès qu'elle cesse.
\end{boxmetier}

\begin{boxmetier}{RSSI}
Alimenter le registre des violations depuis les incidents de sécurité, avec
les éléments techniques, et apporter la qualification du point de vue des
données touchées.
\end{boxmetier}

\begin{boxmetier}{Juriste conformité}
Articuler l'obligation d'informer avec les demandes des autorités
judiciaires~: obtenir des demandes écrites, limitées dans le temps, et
documenter chaque décision.
\end{boxmetier}

\begin{boxresume}
\begin{itemize}
    \item Une violation est toute atteinte à la confidentialité, à
          l'intégrité ou à la disponibilité de données personnelles,
          accidentelle ou illicite, interne ou externe.
    \item L'article~22 impose d'informer sans délai l'Autorité et la
          personne, pour toute violation.
    \item Un registre des violations est la seule manière de prouver le
          respect du « sans délai ».
    \item L'information peut rencontrer le secret de l'enquête pénale~: une
          demande de différer se documente, se limite et se réexamine.
    \item À NYANGA, 18\,412 personnes n'avaient toujours pas été informées
          en août~2026~; MBOLO a informé sa patiente en deux jours.
\end{itemize}
\end{boxresume}

\begin{boxinfo}
\textbf{Points de contrôle.} Vous avez compris ce chapitre si vous savez~:
\begin{itemize}
    \item reconnaître une violation interne~;
    \item rédiger l'information d'une personne après une violation~;
    \item auditer une demande de différer l'information.
\end{itemize}
\end{boxinfo}

\begin{boxexo}[title={\ding{45}\quad Exercice~15.1. Violation ou non \niveaubase}]
(a) un relevé de compte envoyé au mauvais client~; (b) une tentative
d'intrusion bloquée par le pare-feu~; (c) un serveur indisponible deux
jours, sans perte de données~; (d) un portable chiffré perdu dans un
taxi~; (e) un portable non chiffré volé.
\end{boxexo}

\begin{boxexo}[title={\ding{45}\quad Exercice~15.2. Informer M. TAGNE \niveaumoyen}]
Rédigez l'information que NYANGA doit adresser aux 18\,412 personnes,
selon le contenu minimal du cours. Que dites-vous du retard~?
\end{boxexo}

\begin{boxexo}[title={\ding{45}\quad Exercice~15.3. La chronologie \niveaumoyen}]
À partir des dates de la section~\ref{sec:c12-cas}, calculez pour NYANGA le
délai entre la connaissance de la violation et le jugement, puis entre le
jugement et le 31~août~2026. Rédigez la phrase du rapport sur le respect du
« sans délai ».
\end{boxexo}

\begin{boxexo}[title={\ding{45}\quad Exercice~15.4. La demande de l'OPJ \niveauexpert}]
Un OPJ demande aujourd'hui à MBOLO de ne pas informer des patientes dont les
dossiers ont été consultés par un agent sous enquête. Rédigez la réponse de
MBOLO à l'OPJ et la note au registre des violations, en organisant la
conciliation entre l'article~22 et le secret de l'enquête.
\end{boxexo}

\bigskip

\begin{boxinfo}
\textbf{Pour aller plus loin.}
\begin{itemize}
    \item L'ISO/IEC~27035~\cite{iso27035} et le NIST SP~800-61~\cite{nist80061r2},
          pour la gestion des incidents.
    \item Le cours de forensique de l'auteur (UE MSN-2S3-SEC07), pour
          l'enquête sur l'incident de NYANGA.
\end{itemize}
\end{boxinfo}

\begin{boxinfo}
\textbf{Transition.} Les constats s'accumulent~: registre, fondements,
droits, sécurité, sous-traitants, transferts, violations. Il reste à les
transformer en rapport et en programme de mise en conformité. C'est
l'objet du Chapitre~\ref{chap:programme}.
\end{boxinfo}

        % ============================================================================
% Fichier  : partie5/P05_C16_programme_rapport.tex
% Titre    : Programme de Conformité et Rapport d'Audit
% Auteur   : MINKA MI NGUIDJOÏ Thierry Emmanuel
% Version  : 3.0 -- novembre 2026
% Corrigés : corriges/Corr_C16.tex
% ============================================================================
% Sources : loi n° 2024/017, art. 13, 27(2), 28, 73 (texte officiel) ;
% ISO 19011:2018 (rapport, paraphrasé) ; NRAC C07 (observation en sept
% éléments) ; auditologie.
% Faits nouveaux (bible) : rapport de KORO AUDIT remis le 25/09/2026 ;
% réponse de la direction du 18/09/2026 ; synthèse des six domaines.
% ============================================================================

\chapter{Programme de Conformité et Rapport d'Audit}
\label{chap:programme}
\chaptermark{Programme et rapport}

\epigraph{%
    « Un rapport qui énumère des écarts sans les classer\\
    laisse la direction choisir les plus faciles. »%
}{Principe du cours}

\niveauChapitre{Intermédiaire}{%
    Chapitres~\ref{chap:audit} à~\ref{chap:violations}
}

\begin{boxobjectifs}
\begin{itemize}
    \item Rédiger une observation d'audit en sept éléments.
    \item Synthétiser les domaines d'un audit avec $\mu$, $U$ et $H$ et
          conclure selon la règle de décision.
    \item Prioriser les recommandations selon la gravité pour les
          personnes.
    \item Construire un programme de mise en conformité daté, avec des
          indicateurs.
    \item Auditer la conservation et la purge des données (articles~13
          et~28).
\end{itemize}
\end{boxobjectifs}

\bigskip

% ============================================================================
\section*{Avant de commencer : quarante-six exigences, une page}
% ============================================================================

Yaoundé, lundi 14~septembre~2026. Les papiers de travail de la mission
NYANGA remplissent deux classeurs~: quarante-six exigences cotées dans six
domaines, des dizaines de pièces. Le conseil d'administration, lui, attend
une page. M.~TCHOUA a prévenu~: « Dites-nous ce qui est grave, ce qui coûte
cher et ce que nous devons faire d'abord. »

Mme~BELLO sait qu'une page ne peut pas contenir quarante-six constats. Elle
sait aussi qu'une page qui ne dit pas ce qui reste inconnu tromperait le
conseil. Elle choisit de construire le rapport en trois couches~: une
synthèse d'une page, un constat par domaine avec ses cotes, et les
observations détaillées en annexe.

\begin{boxalert}
\textbf{Violation typique.} Remettre un rapport qui liste les écarts par
ordre d'apparition, sans les relier à des articles, sans dire ce qui reste
inconnu, et sans ordre de priorité.

\smallskip
\textbf{Règle~:} chaque constat se rattache à un article et à une preuve~;
la synthèse dit, par domaine, ce qui est établi et ce qui ne l'est pas~;
les recommandations sont classées par gravité pour les personnes, puis par
exposition de l'organisation.
\end{boxalert}

% ============================================================================
\section{L'Observation en Sept Éléments}
\label{sec:c13-observation}
% ============================================================================

\begin{tableaucours}{cL{2.8cm}L{8.4cm}}{%
    \tableauHeader{N°} &
    \tableauHeader{Élément} &
    \tableauHeader{Contenu}
}
1 & Norme & L'article de la loi, ou l'exigence de la norme, qui fonde le
critère \\
2 & Constatation & Le fait établi, chiffré, avec sa preuve et sa cote \\
3 & Risques & Le risque pour les personnes, puis pour l'organisation \\
4 & Contradictoire & La réponse de l'audité, datée \\
5 & Analyse & La cause de l'écart, avec un niveau de confiance \\
6 & Conclusion & La position de l'auditeur \\
7 & Recommandation & Ce qu'il faut faire, en traitant la cause \\
\end{tableaucours}
\vspace{-0.8em}
\begin{center}{\small\itshape Tableau~16.1. L'observation en sept éléments}\end{center}

\begin{boxscenario}{Mission NYANGA. Observation D2-06~: les oppositions non exécutées}
\textbf{Norme.} Article~40 (droit d'opposition à la prospection)~;
article~64 (délit de traitement maintenu malgré l'opposition).

\textbf{Constatation.} Cinq oppositions à la prospection reçues entre le
23~juin et le 31~août~2026, toutes acceptées par lettre. Trois des cinq
numéros figurent dans la liste d'appels de prospection du 28~août (pièce
D2-06-03). Cote \NON{}.

\textbf{Risques.} Pour les personnes~: sollicitations non voulues, perte de
confiance dans l'exercice de leurs droits. Pour l'établissement~: délit de
l'article~64.

\textbf{Contradictoire} (18~septembre). « Les agents n'ont pas reçu la mise
à jour de la liste. Nous corrigeons. »

\textbf{Analyse} (confiance élevée). La cause n'est pas l'agent~: la liste
d'appels est extraite chaque mois de la base des clients, et aucune
procédure ne reporte les oppositions dans la base.

\textbf{Conclusion.} L'exigence n'est pas satisfaite~; l'écart est
systémique.

\textbf{Recommandation.} Ajouter dans la base un indicateur d'opposition
qui exclut automatiquement le client des extractions de prospection~;
vérifier chaque mois, par sondage, l'absence des opposants dans les listes.
\end{boxscenario}

% ============================================================================
\section{La Synthèse par Domaine}
\label{sec:c13-synthese}
% ============================================================================

Cinq domaines ont été cotés dans les chapitres précédents~: D1
(Chapitre~\ref{chap:audit}), D2 (Chapitre~\ref{chap:droits}), D3
(Chapitre~\ref{chap:securite}), D4 (Chapitre~\ref{chap:soustraitants}) et D5
(Chapitre~\ref{chap:violations}). Il reste le sixième, la gouvernance, que
l'équipe cote en dernier parce qu'il dépend de tout ce qui précède.

\begin{boxlecture}{Le domaine D6, gouvernance, à NYANGA}
\begin{tableaucours}{L{7.8cm}cL{3.2cm}}{%
    \tableauHeader{Exigence} &
    \tableauHeader{Cote} &
    \tableauHeader{Preuve}
}
Le risque de chaque traitement est évalué (art.~25) & \NON{} & Aucune
évaluation \\
Une évaluation d'impact précède les traitements à risque élevé (art.~33) &
\NON{} & Douze traitements, aucune \\
Un dossier de conformité peut être présenté au contrôle (art.~35) & \NON{} &
Pièces dispersées, aucun dossier \\
Le rapport annuel sur les mesures est adressé à l'Autorité (art.~27(2)) &
\NEK{} & Autorité non organisée \\
\end{tableaucours}

\medskip
\textbf{Décompte.} Aucun \OUI{}, trois \NON{}, un \NEK{} sur quatre.
\hfill\lecture{0}{0{,}25}{0{,}81}

\medskip
\textbf{Lecture.} C'est le domaine le plus court et le plus net. Les
critères de l'évaluation d'impact relèvent d'un texte à venir, mais
l'absence de toute évaluation se constate sans lui~: la cote est \NON{}
et non \NEK{}. Le seul \NEK{} est de critère. Ce domaine explique les
cinq autres~: personne n'est chargé du sujet, donc rien n'est évalué,
tenu ni rendu.
\end{boxlecture}

Les quarante-six exigences sont réunies dans
l'Annexe~\ref{ann:exigences}, avec leur cote, leur preuve et la nature de
chaque \NEK{}.

\begin{boxlecture}{NYANGA CRÉDIT, synthèse des six domaines, 25 septembre 2026}
\begin{tableaucours}{cL{3.8cm}cccc}{%
    \tableauHeader{Code} &
    \tableauHeader{Domaine} &
    \tableauHeader{O / N / K} &
    \tableauHeader{$\mu$} &
    \tableauHeader{$U$} &
    \tableauHeader{$H$}
}
D1 & Recensement et fondements & 1 / 5 / 2 & 0,13 & 0,25 & 1,30 \\
D2 & Droits des personnes & 1 / 7 / 2 & 0,10 & 0,20 & 1,16 \\
D3 & Sécurité et confidentialité & 4 / 3 / 1 & 0,50 & 0,13 & 1,41 \\
D4 & Sous-traitants et transferts & 1 / 5 / 2 & 0,13 & 0,25 & 1,30 \\
D5 & Violations & 1 / 5 / 2 & 0,13 & 0,25 & 1,30 \\
D6 & Gouvernance & 0 / 3 / 1 & 0 & 0,25 & 0,81 \\
Total & & 8 / 28 / 10 & 0,17 & 0,22 & 1,35 \\
\end{tableaucours}

\medskip
\textbf{Lecture.} Cinq domaines sur six sont opaques au sens de la règle
($U \geq 0{,}20$)~; aucun ne peut être déclaré conforme. Le seul domaine
clair, D3, ne l'est pas non plus~: $\mu = 0{,}50$ atteint le seuil, mais
trois \NON{} portent sur la limitation des accès, et la règle exige
l'absence d'écart grave. Dans tous les domaines, les \NON{} dominent~:
l'opacité n'est pas ce qui empêche de conclure. L'établissement n'est pas
conforme, et l'audit sait pourquoi sur vingt-huit points précis.
\end{boxlecture}

La synthèse distingue trois messages qu'un conseil doit entendre
séparément~: ce qui est établi comme non conforme (vingt-huit écarts), ce
qui reste inconnu et pourquoi (dix \NEK{}, dont six de critère, liés aux
textes d'application et à l'Autorité), et ce qui fonctionne (huit
exigences, dont la moitié en sécurité). Fondre ces trois messages en une
note globale effacerait l'information utile.

\textbf{Une limite du plan, dite dans le rapport.} Le domaine D4
s'intitule « sous-traitants et transferts ». Ses huit exigences portent
sur les contrats. Le transfert vers Abidjan, où l'éditeur héberge le
logiciel bancaire, n'y apparaît que par l'absence de clause de transfert.
Le plan d'audit du 2~juillet ne l'avait pas isolé. Le rapport le signale
et donne la cote hors décompte~: « le transfert est autorisé » vaut
\NON{}, faute d'autorisation~; « le pays offre un niveau de protection
suffisant » vaut \NEK{} de critère, faute de liste
(Chapitre~\ref{chap:transferts}). La recommandation n°~9 en découle. Un
auditeur qui découvre en cours de mission une exigence oubliée ne modifie
pas son décompte après coup. Il l'écrit, la cote à part et l'inscrit au
plan de la mission suivante.

% ============================================================================
\section{Prioriser les Recommandations}
\label{sec:c13-priorites}
% ============================================================================

Les recommandations se classent d'abord par \textbf{gravité pour les
personnes}, ensuite par \textbf{exposition de l'organisation} (sanctions
encourues, contentieux), enfin par \textbf{effort}. L'ordre inverse, du
plus facile au plus difficile, est tentant et dangereux~: il traite d'abord
ce qui protège le moins.

\begin{tableaucours}{cL{6.0cm}L{2.8cm}L{2.4cm}}{%
    \tableauHeader{Rang} &
    \tableauHeader{Recommandation} &
    \tableauHeader{Articles} &
    \tableauHeader{Échéance}
}
1 & Informer les 18\,412 personnes de la violation de 2025 & 22 & 30 jours \\
2 & Supprimer le classeur des mauvais payeurs & 9, 36, 48, 66 & Immédiat \\
3 & Exécuter les oppositions et répondre aux demandes en attente & 39, 40,
56, 64 & 30 jours \\
4 & Désactiver le kit de mesure d'audience sans consentement & loi
2010/012, art.~47 & Immédiat \\
5 & Réécrire formulaires et notices, recueillir les consentements & 9, 21 &
90 jours \\
6 & Avenants aux contrats des sous-traitants, droit d'audit & 16, 30 & 90
jours \\
7 & Fixer et appliquer les durées de conservation, purger & 13, 28 & 6
mois \\
8 & Limiter les accès des agences à leurs clients & 27 & 6 mois \\
9 & Documenter les transferts, préparer les autorisations & 19, 32 & 6
mois \\
10 & Désigner un délégué à la protection des données (bonne pratique) &
aucun & 90 jours \\
\end{tableaucours}
\vspace{-0.8em}
\begin{center}{\small\itshape Tableau~16.2. Les dix premières recommandations à NYANGA CRÉDIT}\end{center}

La dixième recommandation n'a pas de fondement légal~; le rapport le dit.
Elle répond à une cause commune à presque tous les écarts~: personne,
dans l'établissement, n'est chargé de la protection des données.

% ============================================================================
\section{L'Exposition Juridique}
\label{sec:c16-exposition}
% ============================================================================

Le conseil demandera ce que l'établissement risque. Dès le 10~août,
M.~NGANDO a rapproché les constats de terrain des articles qui les
sanctionnent dans les deux lois (Chapitres~\ref{chap:loi}
et~\ref{chap:loi2010}).

\begin{boxscenario}{Mission NYANGA. La note juridique du 10 août 2026}
\textbf{Constatations} (rapprochement des constats d'audit et des
articles). Traitements sans autorisation~: 47 (article~55). Demandes
d'accès sans réponse~: 12 (article~56). Oppositions non exécutées~: 3
(article~64). Interconnexion sans obligation légale~: 1, le classeur des
mauvais payeurs (article~58). Données sensibles sans consentement exprès~:
au moins 4 traitements (article~66). Kit de mesure d'audience sans
consentement~: 1 (loi n°~2010/012, article~47).

\textbf{Analyse} (confiance moyenne). L'exposition théorique est élevée,
mais sa réalisation dépend de l'organisation de l'Autorité, de l'action
des personnes et du parquet. M.~NGANDO déconseille de chiffrer une
« amende probable »~: aucune décision n'existe encore pour en fonder une
estimation. Il classe les écarts par gravité pour les personnes, pas par
montant.
\end{boxscenario}

Cette note accompagne le rapport sans s'y substituer. Elle classe, elle ne
chiffre pas.

% ============================================================================
\section{La Conservation et la Purge}
\label{sec:c13-conservation}
% ============================================================================

L'article~13 limite la conservation à la durée nécessaire, et l'article~28
renvoie la fixation du délai maximum au référentiel de l'Autorité, « compte
tenu de la finalité du traitement ou de la nature des données collectées ».
En attendant le référentiel, l'organisation fixe elle-même ses durées dans
une politique de conservation, traitement par traitement, en distinguant
trois périodes~: la \textbf{base active} (le temps de la finalité), les
\textbf{archives intermédiaires} (le temps d'une obligation légale ou d'un
contentieux possible, avec accès restreint), et la \textbf{purge}
(effacement, ou anonymisation prouvée pour les statistiques).

L'audit de la conservation compare, pour chaque traitement, la date des
données les plus anciennes à la durée fixée. Le laboratoire APDC-LAB-09
exerce ce calcul sur un extrait de la base des comptes de NYANGA.

\begin{boxscenario}{Mission NYANGA. La politique de conservation proposée}
\begin{tableaucours}{L{3.6cm}L{2.8cm}L{3.2cm}L{2.6cm}}{%
    \tableauHeader{Traitement} &
    \tableauHeader{Base active} &
    \tableauHeader{Archive} &
    \tableauHeader{Justification}
}
Comptes clients & Durée du compte & 10 ans après la clôture & Obligations
comptables (dix ans en droit OHADA) \\
Prospects & 1 an sans contact & Aucune & Finalité seule \\
Candidats non retenus & 6 mois & Aucune & Finalité seule \\
Images des agences & 30 jours & Durée d'une procédure si incident &
Sécurité \\
Appels du service client & 6 mois & Durée d'un litige & Preuve (loi
2010/012, art.~45) \\
\end{tableaucours}

\medskip
\textbf{Constatation} (31~août). 22\,950 comptes clos avant 2016 dépassent
déjà la durée d'archive proposée~; la base des prospects contient des
entrées de 2017.
\end{boxscenario}

% ============================================================================
\section{Le Programme de Mise en Conformité}
\label{sec:c13-programme}
% ============================================================================

Un rapport d'audit se clôt par des recommandations~; un programme les
transforme en projet. Il fixe pour chaque recommandation un responsable,
une échéance, une preuve d'achèvement et un indicateur, puis il prévoit un
audit de suivi.

\begin{tableaucours}{L{5.6cm}L{6.4cm}}{%
    \tableauHeader{Indicateur} &
    \tableauHeader{Définition}
}
Couverture du registre & Traitements découverts et inscrits, sur
traitements existants \\
Traitements documentés pour autorisation & Dossiers complets, sur
traitements du registre \\
Délai médian de réponse aux demandes & Médiane du délai entre réception et
réponse complète \\
Exécution des oppositions & Opposants absents des listes de prospection,
par sondage mensuel \\
Contrats conformes & Contrats aux douze clauses, sur contrats de
sous-traitance \\
Violations informées dans les délais & Violations informées dans le délai
cible, sur violations du registre \\
Données au-delà de la durée & Enregistrements dépassant la durée fixée \\
\end{tableaucours}
\vspace{-0.8em}
\begin{center}{\small\itshape Tableau~16.3. Indicateurs d'un programme de conformité}\end{center}

L'article~27, alinéa~2, impose en outre un rapport annuel à l'Autorité sur
l'état des mesures de sécurité. Le programme peut en préparer la matière
dès maintenant~: les indicateurs de sécurité du programme en sont le
cœur, et le format sera fixé quand le référentiel sera publié.

% ============================================================================
\section{Ce que le Rapport demande à chaque Métier}
\label{sec:c13-metiers}
% ============================================================================

\begin{boxmetier}{Auditeur de la protection des données}
Rédiger chaque observation en sept éléments, synthétiser par domaine avec
les trois mesures, séparer établi, inconnu et satisfaisant, et classer les
recommandations par gravité pour les personnes.
\end{boxmetier}

\begin{boxmetier}{Responsable de traitement}
Répondre au contradictoire par écrit, accepter ou motiver le refus de
chaque recommandation, et nommer un responsable par action.
\end{boxmetier}

\begin{boxmetier}{Délégué à la protection des données}
Piloter le programme, mesurer les indicateurs chaque mois, et préparer
l'audit de suivi.
\end{boxmetier}

\begin{boxresume}
\begin{itemize}
    \item L'observation d'audit compte sept éléments, dont le
          contradictoire et une recommandation qui traite la cause.
    \item La synthèse par domaine sépare ce qui est établi non conforme, ce
          qui est inconnu et pourquoi, et ce qui fonctionne.
    \item Les recommandations se classent par gravité pour les personnes,
          puis par exposition, puis par effort.
    \item La conservation se gère en trois périodes~: base active, archive,
          purge~; l'audit compare l'âge des données aux durées fixées.
    \item À NYANGA, vingt-huit écarts établis, dix \NEK{} dont six de
          critère, et l'information des 18\,412 personnes en première
          recommandation.
\end{itemize}
\end{boxresume}

\begin{boxinfo}
\textbf{Points de contrôle.} Vous avez compris ce chapitre si vous savez~:
\begin{itemize}
    \item rédiger une observation en sept éléments~;
    \item conclure sur un domaine à partir de ses cotes~;
    \item ordonner des recommandations et justifier l'ordre.
\end{itemize}
\end{boxinfo}

\begin{boxexo}[title={\ding{45}\quad Exercice~16.1. Ordonner \niveaubase}]
Classez par gravité pour les personnes~: (a) le registre incomplet~; (b) les
selfies traités dans un pays inconnu~; (c) un contrat de papeterie sans
clause de protection des données~; (d) le pointage biométrique des
salariés~; (e) l'absence de délégué.
\end{boxexo}

\begin{boxexo}[title={\ding{45}\quad Exercice~16.2. Une observation \niveaumoyen}]
Rédigez l'observation D1-04 « le classeur des mauvais payeurs » en sept
éléments, à partir des Chapitres~\ref{chap:cartographie} et~\ref{chap:risque}.
Inventez de façon cohérente la réponse de la direction.
\end{boxexo}

\begin{boxexo}[title={\ding{45}\quad Exercice~16.3. La page du conseil \niveaumoyen}]
Rédigez la synthèse d'une page destinée au conseil d'administration de
NYANGA, sans jargon, en trois paragraphes~: ce qui est établi, ce qui est
inconnu, ce qu'il faut faire d'abord.
\end{boxexo}

\begin{boxexo}[title={\ding{45}\quad Exercice~16.4. L'audit de suivi \niveauexpert}]
Six mois après le rapport, NYANGA annonce que « toutes les recommandations
sont mises en œuvre ». Construisez le plan de l'audit de suivi~: exigences
retestées, preuves attendues, échantillons, et la manière de présenter
l'évolution de $\mu$, $U$ et $H$ par domaine.
\end{boxexo}

\bigskip

\begin{boxinfo}
\textbf{Pour aller plus loin.}
\begin{itemize}
    \item L'ISO~19011~\cite{iso19011_2018}, pour le rapport d'audit et le
          suivi.
    \item T.~E.~Minka Mi Nguidjoï~\cite{minka2026nek}, pour la lecture des
          trois mesures.
\end{itemize}
\end{boxinfo}

\begin{boxinfo}
\textbf{Transition.} Un traitement de NOUN PAY n'a pas encore été audité~:
la notation automatique qui décide des crédits. Profilage, décisions
automatisées, biométrie~: ce sont les défis de la Partie~V, et l'objet du
Chapitre~\ref{chap:profilage}.
\end{boxinfo}


    \part{Défis Avancés}
    \label{part:defis}

        % ============================================================================
% Fichier  : partie6/P06_C17_profilage_ia.tex
% Titre    : Profilage, Décisions Automatisées et Intelligence Artificielle
% Auteur   : MINKA MI NGUIDJOÏ Thierry Emmanuel
% Version  : 1.0 -- novembre 2026
% Corrigés : corriges/Corr_C17.tex
% ============================================================================
% Sources : loi n° 2024/017, art. 5 (profilage), 9(3), 17, 18, 21, 44, 48,
% 65 (texte officiel) ; ISO/IEC 27701 ; GCI C19 (gouvernance de l'IA).
% Faits nouveaux (bible) : NOUN SCORE ; refus de crédit de M. ATANGANA du
% 12/08/2026 ; selfies de vérification d'identité ; pointage biométrique de
% MBOLO ; audit de NOUN PAY, octobre 2026.
% ============================================================================

\chapter{Profilage, Décisions Automatisées et Intelligence Artificielle}
\label{chap:profilage}
\chaptermark{Profilage et IA}

\epigraph{%
    « Un modèle qui décide d'un crédit\\
    doit pouvoir être expliqué à celui qu'il refuse. »%
}{Principe du cours}

\niveauChapitre{Expert}{%
    Chapitres~\ref{chap:principes}, \ref{chap:risque}, \ref{chap:droits}
    et~\ref{chap:transferts}
}

\begin{boxobjectifs}
\begin{itemize}
    \item Reconnaître un profilage et une décision fondée exclusivement sur
          un traitement automatisé.
    \item Appliquer l'article~44 et vérifier les garanties dues à la
          personne~: intervention humaine, point de vue, contestation.
    \item Lire l'article~65 et dire ce que l'auditeur peut en conclure.
    \item Auditer un modèle de notation~: données d'entrée, finalité,
          information, non-discrimination.
    \item Traiter les données biométriques au regard des articles~5 et~48.
\end{itemize}
\end{boxobjectifs}

\bigskip

% ============================================================================
\section*{Avant de commencer : « l'ordinateur a dit non »}
% ============================================================================

Bafoussam, mercredi 12~août~2026. M.~ATANGANA, couturier, demande à NOUN
PAY un microcrédit de 150\,000~francs CFA pour acheter une machine. La
réponse arrive en quatre secondes~: refusé. Il appelle le service client,
qui lui explique que « l'ordinateur a dit non » et qu'aucun agent ne peut
revenir sur la décision. Il écrit pour demander pourquoi.

Personne, à NOUN PAY, ne sait le lui dire. Le modèle NOUN SCORE a été
entraîné en 2025 par un prestataire sur les historiques de 210\,000
utilisateurs. Il utilise une soixantaine de variables~: régularité des
recharges de crédit téléphonique, montants reçus, heure habituelle des
transactions, nombre de contacts dans le répertoire du téléphone, et
localisation approximative. L'équipe produit connaît les variables~; elle
ne sait pas lesquelles ont pesé pour M.~ATANGANA.

\begin{boxalert}
\textbf{Violation typique.} Laisser un modèle décider seul d'une question
importante pour une personne, sans pouvoir expliquer la décision ni la
faire réexaminer par un humain.

\smallskip
\textbf{Règle~:} la personne a le droit de s'opposer à une décision fondée
exclusivement sur un traitement automatisé, et, en tout cas, d'obtenir une
intervention humaine, d'exprimer son point de vue et de contester
(article~44). Elle doit avoir été informée de la logique sous-jacente
(article~21).
\end{boxalert}

% ============================================================================
\section{Profilage et Décision Automatisée}
\label{sec:c14-definitions}
% ============================================================================

\begin{boxdef}
\termecle{Profilage}~: selon l'article~5, « traitement automatisé de données
à caractère personnel consistant à les utiliser pour évaluer certains
aspects personnels relatifs à une personne physique, notamment sa santé,
ses préférences, sa localisation et sa situation économique ».
\end{boxdef}

Le profilage \textit{évalue}~; la décision automatisée \textit{décide}. Un
score de solvabilité est un profilage~; le refus de crédit prononcé par le
système sans intervention humaine est une décision automatisée. Les deux se
recouvrent souvent, mais pas toujours~: un score présenté à un agent qui
décide réellement est un profilage sans décision automatisée.

\begin{boxjuris}
\textbf{\loicam{n°~2024/017}, article~44}~\cite{loi2024017}~:
\textit{« (1) La personne concernée a le droit de s'opposer à toute décision
basée exclusivement sur le traitement automatisé de ses données à
caractère personnel, y compris le profilage. (2) Les dispositions de
l'alinéa 1 ci-dessus ne s'appliquent pas si~: la personne concernée est
informée de l'utilisation du système automatisé de prise de décision et a
donné son consentement préalable, explicite et éclairé~; le traitement est
autorisé par la loi, sous réserve qu'elle prévoit des mesures appropriées
pour la sauvegarde des droits et libertés et des intérêts légitimes de la
personne concernée. (3) La personne concernée a également le droit~:
d'obtenir une intervention humaine de la part du responsable du
traitement~; d'exprimer son point de vue~; de contester la décision fondée
sur le traitement automatisé. »}

\smallskip
\textbf{Portée.} L'alinéa~3 ne connaît pas d'exception~: même quand la
personne a consenti, elle conserve le droit à une intervention humaine et
à la contestation. Un système « non révisable » méconnaît l'article~44.
\end{boxjuris}

% ============================================================================
\section{L'Article 65}
\label{sec:c14-article65}
% ============================================================================

\begin{boxjuris}
\textbf{\loicam{n°~2024/017}, article~65}~\cite{loi2024017}~:
\textit{« Est puni d'un emprisonnement de trois (03) à dix (10) ans et
d'une amende de un million (1~000~000) à vingt millions (20~000~000) de
francs CFA, ou de l'une de ces deux peines seulement, le responsable de
traitement ou le sous-traitant qui procède ou fait procéder à un
traitement de données à caractère personnel à des fins de profilage. »}
\end{boxjuris}

Pris isolément, le texte punit tout profilage. Lu avec l'article~44, qui
admet les décisions automatisées, y compris le profilage, sous conditions,
il pourrait viser seulement le profilage réalisé hors de ces conditions.
Les deux lectures sont défendables~; seul un juge, ou un texte
d'application, les départagera.

L'auditeur ne tranche pas. Il fait trois choses. Il signale la difficulté
au commanditaire, en citant les deux articles. Il cote \NEK{} de critère la
proposition « le profilage pratiqué est licite au regard de l'article~65 ».
Et il vérifie, sur pièces, toutes les conditions qui rendraient le
profilage défendable dans la lecture la plus favorable~: autorisation
demandée, consentement préalable, explicite et éclairé, information sur la
logique, intervention humaine effective, absence de discrimination. Si l'une
manque, la recommandation ne dépend plus de la lecture de l'article~65.

% ============================================================================
\section{Auditer un Modèle de Notation}
\label{sec:c14-modele}
% ============================================================================

L'audit d'un modèle ne demande pas d'être statisticien~; il demande de
poser, sur chaque élément du modèle, les questions de la loi.

\begin{tableaucours}{L{3.2cm}L{4.8cm}L{4.0cm}}{%
    \tableauHeader{Élément} &
    \tableauHeader{Question d'audit} &
    \tableauHeader{Articles}
}
Données d'entraînement & Collectées pour cette finalité, ou détournées~?
Consentement~? & 9, 11, 67 \\
Variables & Chacune a-t-elle un lien avec la solvabilité~? Une variable
révèle-t-elle une donnée sensible ou une origine~? & 11, 17, 48 \\
Décision & Exclusivement automatisée~? Intervention humaine possible et
effective~? & 44 \\
Information & La personne est-elle informée de la logique et des
conséquences~? & 21 \\
Non-discrimination & Les taux de refus diffèrent-ils selon la région, la
langue, le sexe~? & 17 \\
Prestataire & Le modèle est-il hébergé ou réentraîné à l'étranger~? & 30,
32 \\
\end{tableaucours}
\vspace{-0.8em}
\begin{center}{\small\itshape Tableau~17.1. Les questions de la loi sur un modèle de notation}\end{center}

\subsection{Les variables qui révèlent}

Une variable anodine peut révéler une donnée sensible. La localisation
habituelle peut révéler l'origine régionale ou le lieu de culte~; l'heure
des transactions, la religion pendant un jeûne~; le répertoire des
contacts, des données sur des tiers qui n'ont jamais consenti. L'article~17
impose de traiter les données « sans tenir compte de la classe sociale,
de l'origine ethnique ou régionale, de l'appartenance syndicale, de
l'opinion politique ou de la conviction religieuse de la personne
concernée »~: une variable qui sert
d'intermédiaire à l'une d'elles fait entrer ce critère dans la décision.

\subsection{Mesurer une discrimination}

L'auditeur compare les taux de refus entre groupes de personnes. Si les
demandeurs d'une région sont refusés deux fois plus que ceux d'une autre à
revenus comparables, le modèle traite probablement une origine régionale
par intermédiaire. Cette comparaison exige de connaître le groupe des
personnes~; si l'organisation ne le connaît pas, et elle ne doit pas
collecter une donnée sensible pour le savoir, l'auditeur utilise un
indicateur indirect déclaré (la région de l'agence d'inscription) et cote
avec prudence. Le laboratoire APDC-LAB-10 conduit ce calcul.

% ============================================================================
\section{Les Données Biométriques}
\label{sec:c14-biometrie}
% ============================================================================

La biométrie figure dans la liste des données sensibles de l'article~5, et
l'article~48, alinéa~1, interdit le traitement des données relatives
« à la génétique et à la biométrie en matière de santé ». Le
Chapitre~\ref{chap:donnees} a relevé l'ambiguïté de cette fin de phrase.
Deux usages du cours la mettent à l'épreuve~:

\begin{itemize}
    \item le \textbf{selfie de vérification d'identité} de NOUN PAY, comparé
          automatiquement à la photographie de la pièce d'identité~;
    \item le \textbf{pointage par empreinte digitale} du personnel de MBOLO,
          en place depuis 2023, qui n'a rien de médical mais se déroule dans
          un établissement de santé.
\end{itemize}

Dans les deux cas, l'auditeur établit les faits (quelles données, où elles
sont stockées, combien de temps, qui y accède), vérifie le consentement
exprès que l'article~66 exige pour conserver des données sensibles, et cote
\NEK{} de critère la licéité au regard de l'article~48 tant qu'aucune
autorisation ni aucun texte ne l'éclaire. Il recommande dans tous les cas
une solution de rechange non biométrique et le stockage du gabarit sur un
support détenu par la personne plutôt que dans une base centrale.

% ============================================================================
\section{L'Intelligence Artificielle : Trois Textes à Connaître}
\label{sec:c17-ia}
% ============================================================================

La loi n°~2024/017 ne contient pas le mot « intelligence artificielle ».
Elle saisit ces systèmes par leurs données et par leurs décisions~: les
articles~44 et~65 pour la décision automatisée, les principes pour les
données d'entraînement. Trois textes extérieurs complètent ce regard.
Aucun n'est un critère obligatoire au Cameroun. L'auditeur doit pourtant
les connaître, parce que les organisations qu'il audite achètent des
modèles conçus sous leur empire.

\begin{tableaucours}{L{3.0cm}L{2.4cm}L{6.6cm}}{%
    \tableauHeader{Texte} &
    \tableauHeader{Nature} &
    \tableauHeader{Ce que l'auditeur en retient}
}
NIST AI RMF~1.0 (2023)~\cite{nist2023airmf} & Cadre volontaire & Quatre
fonctions~: gouverner, cartographier, mesurer, gérer. Une liste de
qualités attendues d'un système, dont l'équité, l'explicabilité et le
respect de la vie privée \\
Règlement (UE) 2024/1689~\cite{ue2024ia} & Texte obligatoire dans l'Union &
Une approche par niveau de risque~: pratiques interdites, systèmes à haut
risque soumis à des obligations, obligations de transparence. La notation
de solvabilité des personnes physiques figure parmi les usages à haut
risque \\
Stratégie continentale de l'Union africaine (2024)~\cite{ua2024ia} &
Orientation politique & Un appel aux États à se doter de cadres nationaux
et à protéger les données dans le développement de l'intelligence
artificielle \\
\end{tableaucours}
\vspace{-0.8em}
\begin{center}{\small\itshape Tableau~17.2. Trois textes sur l'intelligence artificielle}\end{center}

\subsection{Ce que ces textes changent pour l'audit}

Le règlement européen s'applique à qui met un système sur le marché de
l'Union ou l'y utilise. NOUN SCORE note des clients camerounais et n'en
relève pas. La question se poserait si le modèle servait à NOUN
DIASPORA~: l'équipe juridique devrait alors l'instruire, comme elle le
fait pour le RGPD (Chapitre~\ref{chap:rgpd}).

Le cadre du NIST fournit un vocabulaire et une liste de vérification. Il
aide à ne rien oublier~: qui répond du modèle, sur quelles données il a
appris, comment sa performance et ses écarts entre groupes sont mesurés,
qui décide de le retirer. L'auditeur peut l'emprunter comme critère
volontaire, à condition de le déclarer (Chapitre~\ref{chap:normes}).

La stratégie africaine ne crée aucune obligation pour une entreprise.
Elle annonce la direction que prendront les textes nationaux.

\begin{boxalert}
\textbf{Violation typique.} Écrire dans un rapport qu'un modèle « n'est
pas conforme au règlement européen sur l'IA » alors que ce règlement ne
s'applique pas à l'organisation.

\smallskip
\textbf{Règle~:} un texte étranger se cite pour ce qu'il est. S'il
s'applique, c'est un critère obligatoire et le rapport dit pourquoi.
S'il ne s'applique pas, c'est une référence empruntée, déclarée au plan,
qui ne fonde aucune non-conformité légale.
\end{boxalert}

\subsection{Les questions propres aux systèmes qui apprennent}

Un modèle qui apprend pose quatre questions que la notation classique ne
posait pas.

\begin{enumerate}
    \item \textbf{Les données d'entraînement.} D'où viennent-elles, et les
          personnes ont-elles été informées de cet usage (articles~9, 11
          et~21)~?
    \item \textbf{La dérive.} Le modèle est-il réentraîné, et qui vérifie
          que ses résultats ne changent pas en silence~?
    \item \textbf{L'explication.} L'organisation sait-elle dire à une
          personne pourquoi elle a été refusée, dans des mots qu'elle
          comprend (article~44)~?
    \item \textbf{Le fournisseur.} Si le modèle vient d'un tiers, que dit
          le contrat sur la réutilisation des données des clients pour
          améliorer le produit (articles~30 et~32)~?
\end{enumerate}

Quand l'organisation ne peut pas répondre, la cote est \NEK{} de preuve.
Elle se lève par des pièces~: la fiche du modèle, le journal des
réentraînements, le contrat.

% ============================================================================
\section{L'Audit de NOUN SCORE}
\label{sec:c14-nounpay}
% ============================================================================

\begin{boxscenario}{Mission NOUN PAY. La notation automatique, octobre 2026}
\textbf{Constatations.}
\begin{itemize}
    \item 100~\% des demandes de microcrédit sont décidées par NOUN SCORE
          sans intervention humaine~; aucune procédure de réexamen.
    \item Les conditions d'utilisation mentionnent « l'utilisation
          d'algorithmes » sans décrire la logique ni les conséquences~; le
          consentement est une case unique couvrant tous les services.
    \item Le modèle utilise le nombre de contacts du répertoire, collecté
          par l'application sans consentement distinct.
    \item Sur les 41\,200 demandes de juillet à septembre~2026, le taux de
          refus est de 38~\% pour les inscrits des agences de l'Ouest et de
          61~\% pour ceux de l'Extrême-Nord, à tranches de revenus déclarés
          comparables.
    \item Aucune demande d'autorisation, aucune évaluation d'impact.
\end{itemize}

\textbf{Analyse} (confiance moyenne sur la discrimination, élevée sur le
reste). L'écart de taux de refus, à revenus comparables, suggère que des
variables de localisation servent d'intermédiaire à une origine
régionale~; l'indicateur indirect retenu (l'agence d'inscription) ne permet
pas de l'établir avec certitude. Les autres écarts sont établis.
\end{boxscenario}

\begin{boxlecture}{NOUN SCORE, octobre 2026}
\begin{tableaucours}{L{7.8cm}cL{3.2cm}}{%
    \tableauHeader{Exigence} &
    \tableauHeader{Cote} &
    \tableauHeader{Preuve}
}
Une intervention humaine est possible (art.~44(3)) & \NON{} & Procédure
absente, refus à M.~ATANGANA \\
La personne est informée de la logique (art.~21) & \NON{} & Conditions
d'utilisation \\
Le consentement à la décision automatisée est explicite (art.~44(2)) &
\NON{} & Case unique \\
Les variables sont liées à la finalité (art.~11) & \NON{} & Répertoire des
contacts \\
Le modèle ne discrimine pas selon l'origine (art.~17) & \NEK{} & Écart de
refus, indicateur indirect \\
Une évaluation d'impact a été conduite (art.~33) & \NON{} & Aucune \\
Le traitement est autorisé (art.~19) & \NON{} & Aucune autorisation \\
Le profilage est licite au regard de l'article~65 & \NEK{} & Lecture non
tranchée \\
\end{tableaucours}

\medskip
\textbf{Décompte.} Zéro \OUI{}, six \NON{}, deux \NEK{} sur huit.
\hfill\lecture{0}{0{,}25}{0{,}81}

\medskip
\textbf{Lecture.} Aucune exigence n'est satisfaite. Les deux \NEK{} portent
sur les questions les plus graves, discrimination et délit, mais la
recommandation n'en dépend pas~: les six \NON{} suffisent à demander la
suspension des refus automatiques jusqu'à la mise en place d'une
intervention humaine.
\end{boxlecture}

% ============================================================================
\section{Ce que le Profilage demande à chaque Métier}
\label{sec:c14-metiers}
% ============================================================================

\begin{boxmetier}{Responsable de traitement}
Ne jamais mettre en production un modèle qui décide sans prévoir le
réexamen humain, l'explication et la contestation. M.~NJOYA a suspendu les
refus automatiques le 20~octobre~: tout refus est désormais revu par un
agent.
\end{boxmetier}

\begin{boxmetier}{Auditeur de la protection des données}
Poser les questions de la loi variable par variable, mesurer les écarts de
traitement entre groupes avec un indicateur déclaré, et signaler la
difficulté de l'article~65 sans la trancher.
\end{boxmetier}

\begin{boxmetier}{Juriste conformité}
Rédiger l'information sur la logique de la décision dans des termes que la
personne comprend, et prévoir le circuit de contestation.
\end{boxmetier}

\begin{boxresume}
\begin{itemize}
    \item Le profilage évalue~; la décision automatisée décide~; les deux
          se recouvrent souvent sans se confondre.
    \item L'article~44 garantit, même après consentement, l'intervention
          humaine, l'expression du point de vue et la contestation.
    \item L'article~65 punit le traitement à des fins de profilage~; sa
          portée exacte n'est pas tranchée, et l'auditeur le dit.
    \item Une variable anodine peut servir d'intermédiaire à une donnée
          sensible~; l'audit compare les taux de décision entre groupes.
    \item À NOUN PAY, aucune exigence n'est satisfaite~; les refus
          automatiques ont été suspendus.
\end{itemize}
\end{boxresume}

\begin{boxinfo}
\textbf{Points de contrôle.} Vous avez compris ce chapitre si vous savez~:
\begin{itemize}
    \item distinguer profilage et décision automatisée~;
    \item énoncer les droits de l'article~44(3)~;
    \item repérer une variable qui sert d'intermédiaire à une origine.
\end{itemize}
\end{boxinfo}

\begin{boxexo}[title={\ding{45}\quad Exercice~17.1. Profilage, décision, ni l'un ni l'autre \niveaubase}]
(a) Un score de solvabilité montré à un conseiller qui décide~; (b) un
refus de crédit automatique~; (c) une liste de clients triés par ancienneté~;
(d) une publicité ciblée selon la localisation~; (e) le blocage automatique
d'une transaction suspecte.
\end{boxexo}

\begin{boxexo}[title={\ding{45}\quad Exercice~17.2. Répondre à M. ATANGANA \niveaumoyen}]
Rédigez la réponse de NOUN PAY à M.~ATANGANA~: information sur la logique,
réexamen humain, possibilité de contester. Que lui dites-vous des variables
utilisées~?
\end{boxexo}

\begin{boxexo}[title={\ding{45}\quad Exercice~17.3. Les taux de refus \niveaumoyen}]
Calculez le rapport des taux de refus entre l'Extrême-Nord et l'Ouest. Que
faudrait-il connaître de plus pour passer la proposition « le modèle
discrimine selon l'origine » de \NEK{} à \OUI{} ou à \NON{}~? Quelles
données ne devez-vous pas collecter pour le savoir~?
\end{boxexo}

\begin{boxexo}[title={\ding{45}\quad Exercice~17.4. Le pointage biométrique \niveauexpert}]
Le Dr~ABENA veut conserver le pointage par empreinte « parce que les
salariés trichaient avec les badges ». Analysez le traitement au regard des
articles~5, 9, 48 et~66. Proposez deux solutions, et dites comment vous
coterez la licéité de chacune.
\end{boxexo}

\bigskip

\begin{boxinfo}
\textbf{Pour aller plus loin.}
\begin{itemize}
    \item L'ISO/IEC~27701~\cite{iso27701}, pour l'encadrement des
          décisions automatisées.
    \item Le cours de gouvernance des SI de l'auteur (UE MSN-2S3-MSI03),
          chapitre sur la gouvernance des données et de l'IA.
\end{itemize}
\end{boxinfo}

\begin{boxinfo}
\textbf{Transition.} La santé, la banque, la justice et l'état civil sont
renvoyés par la loi à des textes particuliers. Comment auditer dans ces
secteurs, quand le régime général et le régime spécial se rencontrent~?
C'est l'objet du Chapitre~\ref{chap:secteurs}.
\end{boxinfo}

        % ============================================================================
% Fichier  : partie6/P06_C18_secteurs_sensibles.tex
% Titre    : Les Secteurs Sensibles
% Auteur   : MINKA MI NGUIDJOÏ Thierry Emmanuel
% Version  : 1.0 -- novembre 2026
% Corrigés : corriges/Corr_C18.tex
% ============================================================================
% Sources : loi n° 2024/017, art. 4, 5, 9, 18, 19(2), 40(3), 45, 48, 72
% (texte officiel) ; Minka Mi Nguidjoï, Massoussi Songo et Bialo (2026,
% manuscrit) pour le régime spécial de la justice.
% Faits nouveaux (bible) : vidéosurveillance de la salle d'attente de
% pédiatrie de MBOLO ; transmission des résultats aux mutuelles ; audit de
% suivi d'octobre 2026.
% ============================================================================

\chapter{Les Secteurs Sensibles}
\label{chap:secteurs}
\chaptermark{Les secteurs sensibles}

\epigraph{%
    « Quand la loi renvoie à un texte particulier,\\
    l'auditeur cherche le texte avant de conclure. »%
}{Principe du cours}

\niveauChapitre{Expert}{%
    Chapitres~\ref{chap:donnees}, \ref{chap:principes} et~\ref{chap:loi}
}

\begin{boxobjectifs}
\begin{itemize}
    \item Identifier les traitements que la loi renvoie à des textes
          particuliers (articles~4 et~72).
    \item Auditer quand régime général et régime spécial se rencontrent.
    \item Appliquer les règles propres aux données de santé, aux
          transactions bancaires et aux données des mineurs.
    \item Auditer une vidéosurveillance~: fondement, information, durée,
          lieux filmés.
    \item Situer l'expert, l'OPJ et le magistrat au regard de la loi.
\end{itemize}
\end{boxobjectifs}

\bigskip

% ============================================================================
\section*{Avant de commencer : une caméra dans la salle d'attente}
% ============================================================================

Bertoua, mardi 6~octobre~2026. Pour l'audit de suivi du CENTRE MÉDICAL
MBOLO, Mme~BELLO visite la pédiatrie. Une caméra filme la salle d'attente,
où des mères allaitent et où des enfants attendent leurs vaccins. Les images
sont affichées à l'accueil et enregistrées trente jours. Aucun panneau ne
le signale. « C'est pour la sécurité, explique le surveillant, il y a eu des
vols de téléphones. »

Au même moment, le service de facturation transmet chaque mois aux
mutuelles des patients la liste des actes, avec le diagnostic en clair,
« parce que les mutuelles le demandent ». Aucun texte n'a été montré à
l'auditrice qui oblige à transmettre le diagnostic.

\begin{boxalert}
\textbf{Violation typique.} Justifier un traitement de données de santé ou de
mineurs par la seule bonne intention (« la sécurité », « le remboursement »)
sans vérifier le fondement, la proportion et les règles propres au secteur.

\smallskip
\textbf{Règle~:} dans un secteur sensible, chaque traitement cumule les
exigences de la loi générale, les règles propres aux données sensibles et,
quand il existe, le texte particulier du secteur. Ce qui n'est pas prévu
par un texte reste soumis à la loi générale.
\end{boxalert}

% ============================================================================
\section{Les Renvois à des Textes Particuliers}
\label{sec:c15-renvois}
% ============================================================================

Deux articles de la loi renvoient à d'autres textes.

\begin{boxjuris}
\textbf{\loicam{n°~2024/017}, article~4}~\cite{loi2024017}~:
\textit{« Tout traitement des données à caractère personnel effectué par les
autorités compétentes en matière de sécurité et de défense est régi par
des textes particuliers. »}

\medskip
\textbf{\loicam{n°~2024/017}, article~72}~\cite{loi2024017}~:
\textit{« Le traitement des données à caractère personnel relatif notamment
à la sécurité, à la défense, à la santé, à la justice et à l'état civil est
régi par des textes spécifiques, en tant que de besoin. »}

\smallskip
\textbf{Portée.} L'article~4 vise des \textit{autorités}~; l'article~72 des
\textit{domaines}, avec « notamment » et « en tant que de besoin ». Ce
dernier n'exclut pas ces domaines de la loi générale~: il annonce des
textes qui la complèteront ou l'adapteront. Tant qu'un texte spécifique
n'est pas identifié, l'auditeur applique la loi générale et signale la
réserve.
\end{boxjuris}

\begin{boxprocedure}
\textbf{Auditer un traitement d'un secteur visé par l'article~72.}
\begin{enumerate}
    \item Rechercher le texte spécifique au secteur et au traitement~; noter
          la recherche et son résultat.
    \item S'il existe, l'appliquer en même temps que la loi générale, et
          relever les points où il déroge.
    \item S'il n'existe pas, appliquer la loi générale, et coter \NEK{} de
          critère les seules exigences qu'un texte spécifique pourrait
          modifier.
    \item Appliquer dans tous les cas les règles des données sensibles
          (articles~5, 48, 66) et des mineurs (articles~9(3), 18, 19(2),
          40(3)).
    \item Vérifier les secrets professionnels du secteur, qui s'ajoutent à
          la loi.
\end{enumerate}
\end{boxprocedure}

% ============================================================================
\section{La Santé}
\label{sec:c15-sante}
% ============================================================================

Les données de santé sont sensibles (article~5). Le soin trouve un
fondement dans la troisième dérogation de l'article~9~: le traitement
« nécessaire~[\ldots] à la préservation de la santé de la personne
concernée ou d'une autre personne physique ». Ce fondement couvre le soin,
pas tout ce qui l'entoure~: la facturation, la recherche, la formation, la
communication avec les mutuelles ou les employeurs exigent chacun leur
propre fondement.

Les professionnels de santé sont en outre tenus au secret médical, protégé
par le droit pénal et par la déontologie médicale. Le secret ne se partage
qu'entre ceux qui participent à la prise en charge, et dans la mesure
nécessaire~: c'est la règle que l'auditeur applique aux journaux de
consultation (Chapitre~\ref{chap:securite}), et c'est elle qui interdit de
transmettre un diagnostic à une mutuelle sans texte ou sans consentement
du patient.

% ============================================================================
\section{La Banque et la Finance}
\label{sec:c15-banque}
% ============================================================================

Les transactions bancaires sont des données sensibles (article~5), et
l'article~48, alinéa~2, interdit leur traitement « sans autorisation
préalable des administrations et structures compétentes, dans les
conditions et suivant les modalités prévues par les lois en vigueur ». La
loi ne désigne pas ces administrations~; il peut s'agir des autorités de
régulation bancaire ou monétaire. Pour l'auditeur, une microfinance comme
NYANGA ou une fintech comme NOUN PAY doit pouvoir montrer, outre la demande
d'autorisation à l'Autorité de protection, les agréments et autorisations
de son régulateur sectoriel, et le texte qui encadre l'usage des données
dans son activité. Faute de désignation claire, l'exigence « autorisation
des administrations compétentes au titre de l'article~48(2) » est un
\NEK{} de critère.

Le secteur financier connaît aussi des obligations de conservation et de
transmission (lutte contre le blanchiment, déclarations aux autorités) qui
fondent des traitements au titre de l'obligation légale (article~9(2)).
L'auditeur demande le texte précis~; une « obligation réglementaire »
invoquée sans référence reste une déclaration.

% ============================================================================
\section{Les Mineurs}
\label{sec:c15-mineurs}
% ============================================================================

La loi protège les mineurs à quatre endroits~: le consentement d'un mineur
de dix-huit ans n'est valable qu'appuyé par celui de ses parents ou de son
représentant légal (article~9(3))~; les traitements de leurs données en vue
d'une offre de services doivent être « adéquat[s], pertinent[s] et
limité[s] » (article~18)~; l'interconnexion de fichiers de données
sensibles relatives aux mineurs exige une autorisation propre
(article~19(2))~; les parents peuvent s'opposer au traitement des données
de leur enfant collectées sans leur consentement (article~40(3)).

En pédiatrie, l'enfant est soigné sur le fondement de la préservation de sa
santé, et les parents consentent aux actes~; le consentement aux autres
traitements (photographies pour un dossier de formation, participation à
une étude) doit être recueilli auprès des parents, et l'enfant informé à sa
mesure.

% ============================================================================
\section{La Vidéosurveillance}
\label{sec:c15-video}
% ============================================================================

La vidéosurveillance d'un lieu ouvert au public filme des personnes
identifiables~: c'est un traitement de données personnelles. La loi n'en
fait pas un régime particulier~; elle s'applique donc entièrement, et la
difficulté vient du fondement (Chapitre~\ref{chap:principes})~: ni
l'intérêt légitime ni le contrat n'existent dans la loi, et le
consentement de chaque personne filmée est irréaliste. L'auditeur cote
\NEK{} le fondement tant qu'aucune autorisation ne l'a précisé, et vérifie
tout ce qui ne dépend pas du fondement~: la finalité, les lieux filmés,
l'information par panneau, la durée de conservation, les personnes qui
voient les images.

\begin{tableaucours}{L{5.2cm}L{6.8cm}}{%
    \tableauHeader{Question} &
    \tableauHeader{Ce que l'auditeur vérifie}
}
Finalité & Une seule, écrite~: sécurité des biens et des personnes~; pas de
contrôle du travail \\
Lieux filmés & Pas de lieu de soin, de toilettes, de vestiaire, d'espace
d'allaitement \\
Information & Panneau visible à l'entrée de la zone filmée \\
Durée & Courte, fixée, appliquée (quelques semaines au plus, sauf incident) \\
Accès aux images & Personnes désignées, consultation journalisée \\
Affichage & Écrans hors de la vue du public \\
\end{tableaucours}
\vspace{-0.8em}
\begin{center}{\small\itshape Tableau~18.1. Les questions d'audit d'une vidéosurveillance}\end{center}

% ============================================================================
\section{Les Télécommunications}
\label{sec:c18-telecom}
% ============================================================================

L'article~72 ne cite pas les télécommunications. Le secteur est pourtant
celui où les textes particuliers sont les plus anciens et les plus
nombreux (Chapitre~\ref{chap:loi2010}). Un opérateur détient trois
familles de données que peu d'organisations réunissent~: l'identité de
chaque abonné, ses données de trafic et sa localisation.

\begin{tableaucours}{L{3.2cm}L{4.4cm}L{4.4cm}}{%
    \tableauHeader{Donnée} &
    \tableauHeader{Texte particulier} &
    \tableauHeader{Question d'audit}
}
Identité de l'abonné & Décret n°~2015/3759 du 3~septembre~2015 sur
l'identification des abonnés et des terminaux & Quelles pièces sont
copiées, où sont-elles gardées, qui y accède dans les points de vente~? \\
Contenu des communications & Loi n°~2010/012, articles~41 à~48 & Le secret
est-il garanti, y compris envers les salariés de l'opérateur~? \\
Données de connexion et de trafic & Loi n°~2010/012, article~29 & Quelle
durée est appliquée, et qui peut interroger ces données~? \\
Droits de l'utilisateur & Directive n°~07/08-UEAC-133-CM-18 de la
CEMAC~\cite{cemac2008} & L'abonné est-il informé et peut-il s'opposer à la
prospection~? \\
\end{tableaucours}
\vspace{-0.8em}
\begin{center}{\small\itshape Tableau~18.2. Les données d'un opérateur et leurs textes}\end{center}

\subsection{L'identification des abonnés}

L'obligation d'identifier chaque abonné fonde le traitement au titre de
l'obligation légale (article~9(2) de la loi n°~2024/017). Elle ne règle
pas tout. L'auditeur vérifie quatre points que le texte particulier
laisse à la loi générale.

\begin{enumerate}
    \item \textbf{La minimisation.} L'opérateur collecte-t-il seulement ce
          que le texte demande, ou ajoute-t-il des données pour ses
          besoins commerciaux~?
    \item \textbf{Les points de vente.} Les revendeurs photographient les
          pièces d'identité avec leur propre téléphone. Ces copies
          sont-elles effacées après transmission~?
    \item \textbf{La réutilisation.} Les données d'identification
          servent-elles à la prospection ou à la notation, ce qui serait
          un détournement de finalité (article~11)~?
    \item \textbf{La fin du contrat.} Que deviennent les données d'un
          abonné qui résilie~?
\end{enumerate}

\subsection{L'opérateur comme sous-traitant}

Un opérateur est aussi le prestataire de ses clients. Celui qui envoie les
SMS de NYANGA CRÉDIT reçoit des numéros et des messages de relance. Pour
ces envois, il est sous-traitant de la microfinance. Pour l'acheminement
lui-même et pour les données de trafic qu'il conserve, il est responsable
de son propre traitement. Le contrat type qu'il impose ne distingue pas
les deux rôles. C'est l'une des raisons du \NEK{} relevé dans le domaine
D4 (Chapitre~\ref{chap:soustraitants}).

Deux autorités se partagent le secteur~: le régulateur des
télécommunications pour les obligations de l'opérateur, l'Autorité de
protection des données pour les droits des personnes. Tant que la seconde
n'est pas organisée, l'auditeur note auprès de quelle autorité chaque
exigence pourrait être vérifiée.

% ============================================================================
\section{La Justice et la Chaîne Pénale}
\label{sec:c15-justice}
% ============================================================================

La justice figure parmi les domaines de l'article~72. L'expert judiciaire,
l'officier de police judiciaire et le magistrat traitent tous des données
personnelles~: saisies, analyses, rapports, transmissions. Leur qualité au
regard de la loi (responsable, sous-traitant, autre acteur habilité) n'est
pas fixée par un texte.

Dans leur travail en cours, l'auteur de ce cours et ses
coauteurs~\cite{minka2026chainepenale} proposent de déterminer cette qualité
« opération par opération », et non selon la qualité institutionnelle~:
un expert qui collecte des données sur instruction d'une autorité
judiciaire peut être sous-traitant pour cette opération, et devenir
responsable pour la copie qu'il conserve pour ses propres besoins
techniques. Le Chapitre~\ref{chap:recherche} présente cette grille comme
une piste de recherche. Pour l'audit d'un cabinet d'expertise, elle
suggère une règle pratique~: tenir un registre de ses traitements, et
pour chaque mission, documenter qui a déterminé la finalité.

% ============================================================================
\section{L'Audit de Suivi de MBOLO}
\label{sec:c15-mbolo}
% ============================================================================

\begin{boxscenario}{Mission MBOLO. Les traitements sensibles, octobre 2026}
\textbf{Constatations.} La caméra de pédiatrie filme un espace où des mères
allaitent~; aucun panneau~; écran visible depuis le couloir. Le diagnostic
est transmis en clair aux mutuelles chaque mois, pour 1\,640 patients en
septembre~; aucun texte n'a été présenté, et les formulaires d'adhésion des
mutuelles ne mentionnent pas cette transmission. Les consentements
parentaux sont recueillis pour les actes, pas pour les photographies prises
pour le dossier de formation des internes. Le texte spécifique au secteur
de la santé annoncé par l'article~72 n'a pas été trouvé.
\end{boxscenario}

\begin{boxlecture}{Traitements sensibles de MBOLO, octobre 2026}
\begin{tableaucours}{L{7.8cm}cL{3.2cm}}{%
    \tableauHeader{Exigence} &
    \tableauHeader{Cote} &
    \tableauHeader{Preuve}
}
Le soin repose sur la dérogation de l'article~9(2) & \OUI{} & Finalité de
soin \\
La transmission du diagnostic aux mutuelles a un fondement & \NON{} &
Aucun texte, aucun consentement \\
Les lieux filmés excluent les espaces intimes & \NON{} & Salle
d'allaitement filmée \\
La vidéosurveillance est signalée & \NON{} & Aucun panneau \\
La vidéosurveillance a un fondement valable & \NEK{} & Fondement incertain \\
Les photographies d'enfants ont le consentement des parents & \NON{} &
Dossier de formation \\
Un texte spécifique à la santé modifie ces exigences & \NEK{} & Texte non
trouvé \\
\end{tableaucours}

\medskip
\textbf{Décompte.} Un \OUI{}, quatre \NON{}, deux \NEK{} sur sept.
\hfill\lecture{0{,}14}{0{,}29}{1{,}38}

\medskip
\textbf{Lecture.} Domaine opaque, mais les quatre \NON{} ne dépendent
d'aucun texte spécifique~: un texte de santé pourrait préciser la
transmission aux mutuelles, il ne rendrait pas licite une caméra dans un
espace d'allaitement.
\end{boxlecture}

% ============================================================================
\section{Ce que les Secteurs demandent à chaque Métier}
\label{sec:c15-metiers}
% ============================================================================

\begin{boxmetier}{Juriste conformité}
Tenir une veille des textes spécifiques annoncés par l'article~72, et une
liste, par traitement, des textes sectoriels qui l'encadrent.
\end{boxmetier}

\begin{boxmetier}{Auditeur de la protection des données}
Documenter la recherche du texte spécifique, appliquer la loi générale en
son absence, et ne coter \NEK{} de critère que les exigences qu'un tel
texte pourrait réellement modifier.
\end{boxmetier}

\begin{boxmetier}{Responsable de traitement}
Faire déplacer la caméra et retirer le diagnostic des listes transmises
aux mutuelles dès le rapport~: ces écarts ne dépendent d'aucun texte à
venir.
\end{boxmetier}

\begin{boxresume}
\begin{itemize}
    \item L'article~4 renvoie la sécurité et la défense à des textes
          particuliers~; l'article~72 annonce des textes spécifiques pour la
          santé, la justice, l'état civil et d'autres domaines.
    \item En l'absence de texte spécifique, la loi générale s'applique, et
          l'auditeur signale la réserve.
    \item La préservation de la santé fonde le soin, pas la facturation, la
          formation ou la transmission aux tiers.
    \item Les mineurs bénéficient de quatre protections propres~; la
          vidéosurveillance se vérifie lieu par lieu.
    \item À MBOLO, quatre écarts sont établis indépendamment de tout texte
          sectoriel.
\end{itemize}
\end{boxresume}

\begin{boxinfo}
\textbf{Points de contrôle.} Vous avez compris ce chapitre si vous savez~:
\begin{itemize}
    \item lire l'article~72 sans en faire une exclusion générale~;
    \item auditer une vidéosurveillance~;
    \item dire ce que couvre la dérogation de santé de l'article~9.
\end{itemize}
\end{boxinfo}

\begin{boxexo}[title={\ding{45}\quad Exercice~18.1. Quel fondement \niveaubase}]
Quel fondement pour~: (a) le dossier médical d'un patient~; (b) la
facturation d'un acte~; (c) une étude sur le paludisme à partir des
dossiers~; (d) la photographie d'une plaie pour la formation des
internes~?
\end{boxexo}

\begin{boxexo}[title={\ding{45}\quad Exercice~18.2. Les mutuelles \niveaumoyen}]
Proposez un circuit de transmission aux mutuelles conforme à la loi et au
secret médical~: données transmises, fondement, information des patients.
\end{boxexo}

\begin{boxexo}[title={\ding{45}\quad Exercice~18.3. La caméra \niveaumoyen}]
Rédigez la recommandation sur la vidéosurveillance de la pédiatrie, en
reprenant les questions du Tableau~18.1. Que reste-t-il \NEK{} après sa mise
en œuvre, et pourquoi~?
\end{boxexo}

\begin{boxexo}[title={\ding{45}\quad Exercice~18.4. Un cabinet d'expertise \niveauexpert}]
Un cabinet d'expertise judiciaire vous demande un audit. Pour trois de ses
opérations (recevoir un scellé, analyser une copie, conserver une copie
après la mission), proposez la qualification de l'expert au regard de la
loi, en suivant la grille opération par opération, et dites quelles
exigences vous audit\-erez pour chacune.
\end{boxexo}

\bigskip

\begin{boxinfo}
\textbf{Pour aller plus loin.}
\begin{itemize}
    \item La loi n°~2024/017~\cite{loi2024017}, articles~4, 9, 18, 48 et~72.
    \item T.~E.~Minka Mi Nguidjoï, R.~T.~Massoussi Songo et
          D.~Bialo~\cite{minka2026chainepenale}, pour la chaîne pénale.
\end{itemize}
\end{boxinfo}

\begin{boxinfo}
\textbf{Transition.} Tout au long du cours, des \NEK{} se sont accumulés~: de
preuve et de critère, de faits et de textes. Peut-on mesurer cette
incertitude, et que dit-elle de l'état du droit lui-même~? C'est la
question de recherche du Chapitre~\ref{chap:recherche}.
\end{boxinfo}


    \part{Recherche : Mesurer la Conformité, Articuler les Normes}
    \label{part:recherche}

        % ============================================================================
% Fichier  : partie7/P07_C19_mesurer_articuler.tex
% Titre    : Mesurer la Conformité, Articuler les Normes
% Auteur   : MINKA MI NGUIDJOÏ Thierry Emmanuel
% Version  : 1.0 -- novembre 2026
% Corrigés : corriges/Corr_C19.tex
% ============================================================================
% Sources : auditologie (hal-05701007, hal-05701270) ; Shannon (1948) ;
% Minka Mi Nguidjoï, Massoussi Songo et Bialo (2026, manuscrit en
% préparation, septembre 2026) : champ de la loi (art. 2, 4, 72), tensions
% T1 à T3, matrice en huit étapes, échelle compatibilité / tension /
% incompatibilité, grille de qualification des acteurs, quatre propositions.
% Les éléments du manuscrit sont présentés comme une piste de recherche en
% cours de l'auteur et de ses coauteurs.
% ============================================================================

\chapter{Mesurer la Conformité, Articuler les Normes}
\label{chap:recherche}
\chaptermark{Mesurer et articuler}

\epigraph{%
    « Ce que l'on ignore à cause des faits se cherche.\\
    Ce que l'on ignore à cause des textes se signale. »%
}{Principe du cours}

\niveauChapitre{Recherche}{%
    Tout le cours, en particulier les Chapitres~\ref{chap:audit},
    \ref{chap:loi} et~\ref{chap:programme}
}

\begin{boxobjectifs}
\begin{itemize}
    \item Mesurer l'état de conformité d'une organisation par $\mu$, $U$
          et $H$, et en suivre l'évolution.
    \item Distinguer et mesurer séparément l'incertitude de preuve et
          l'incertitude de critère.
    \item Présenter les pistes de recherche actuelles de l'auteur et de ses
          coauteurs sur l'articulation de la loi n°~2024/017 avec la chaîne
          pénale.
    \item Appliquer une matrice d'analyse des tensions normatives à un cas.
    \item Formuler une question de recherche sur la mesure de la
          conformité ou l'articulation des normes.
\end{itemize}
\end{boxobjectifs}

\bigskip

% ============================================================================
\section*{Avant de commencer : « dix inconnus, est-ce notre faute ? »}
% ============================================================================

Bafoussam, vendredi 25~septembre~2026. Le conseil d'administration de NYANGA
CRÉDIT a lu la synthèse~: vingt-huit écarts, huit exigences satisfaites,
dix inconnues. Un administrateur demande~: « Ces dix inconnues, est-ce notre
faute~? » Mme~BELLO répond que quatre le sont~: des pièces que
l'établissement n'a pas pu produire. Les six autres ne dépendent de
personne dans la salle~: elles attendent un décret, une liste de pays, un
référentiel, ou la décision d'un juge sur le sens d'un article.

La réponse ouvre deux questions de recherche. La première est de mesure~:
peut-on séparer, dans un rapport, ce que l'organisation ignore de ce que
le droit lui-même ne dit pas encore~? La seconde est d'articulation~: que
faire quand deux textes applicables à la même donnée se rencontrent~?

\begin{boxalert}
\textbf{Violation typique.} Imputer à l'organisation auditée toute
l'incertitude d'un audit, y compris celle qui vient de textes absents ou
contradictoires.

\smallskip
\textbf{Règle~:} l'incertitude se décompose selon sa source. Le rapport
distingue ce qui manque dans les preuves de ce qui manque dans les
critères, et il présente les secondes comme une information sur l'état du
droit, pas comme un défaut de l'organisation.
\end{boxalert}

% ============================================================================
\section{Mesurer l'État de Conformité}
\label{sec:c16-mesurer}
% ============================================================================

\begin{boxdef}
Pour un ensemble de $n$ exigences cotées~: $\mu = n_O/n$ (taux de
conformité), $U = n_K/n$ (taux d'incertitude), $H = -\sum_c p_c \log_2 p_c$
(entropie de la répartition des cotes, en bits, entre~0 et
$\log_2 3 \approx 1{,}585$)~\cite{minka2026nek, shannon1948}.
\end{boxdef}

\begin{boxlecture}{NYANGA CRÉDIT, deux instantanés}
\begin{tableaucours}{L{3.4cm}cccccc}{%
    \tableauHeader{Date} &
    \tableauHeader{$n$} &
    \tableauHeader{O / N / K} &
    \tableauHeader{$\mu$} &
    \tableauHeader{$U$} &
    \tableauHeader{$H$} &
    \tableauHeader{}
}
6~juillet~2026 & 8 & 2 / 2 / 4 & 0,25 & 0,50 & 1,50 & premier jour \\
25~septembre~2026 & 46 & 8 / 28 / 10 & 0,17 & 0,22 & 1,35 & rapport \\
\end{tableaucours}

\medskip
\textbf{Lecture.} Entre les deux dates, $U$ baisse de 0,50 à 0,22~:
l'audit a appris. $\mu$ baisse aussi, de 0,25 à 0,17, sans que
l'établissement ait régressé~: les \NEK{} du premier jour se sont résolus
surtout en \NON{}. Une baisse de $\mu$ au cours d'un audit est donc souvent
un signe de connaissance, pas de dégradation. Les deux instantanés ne
portent pas sur le même ensemble d'exigences~; la comparaison n'est
rigoureuse que sur les huit exigences communes.
\end{boxlecture}

La dernière remarque est une première question de recherche~: comment
comparer deux états de conformité mesurés sur des ensembles d'exigences
différents, ou découpés différemment~? Les mesures dépendent de la
granularité~; un domaine découpé en dix exigences pèse plus qu'un domaine
découpé en quatre.

% ============================================================================
\section{Deux Sortes d'Incertitude}
\label{sec:c16-deux-nek}
% ============================================================================

Le cours a distingué tout au long deux sources de \NEK{}~:
\begin{itemize}
    \item le \textbf{\NEK{} de preuve}~: le critère est clair, la pièce
          manque~; il se lève en cherchant (audit du prestataire, examen des
          sauvegardes)~;
    \item le \textbf{\NEK{} de critère}~: la pièce existerait, mais la règle
          qui permettrait de la juger n'est pas fixée~; il se lève avec un
          texte (décret, référentiel, liste de pays) ou une décision de
          justice.
\end{itemize}

Le cours propose, à titre de piste, de décomposer le taux d'incertitude en
deux termes~:
\[
U = U_p + U_c, \qquad U_p = \frac{n_{K,p}}{n}, \qquad U_c = \frac{n_{K,c}}{n},
\]
où $n_{K,p}$ et $n_{K,c}$ comptent les \NEK{} de preuve et de critère. À
NYANGA, sur 46 exigences, quatre \NEK{} sont de preuve et six de critère~:
$U_p \approx 0{,}09$ et $U_c \approx 0{,}13$.

$U_p$ mesure ce que l'organisation peut réduire. $U_c$ mesure ce que l'état
du droit laisse ouvert~; mesuré sur plusieurs organisations et plusieurs
années, il pourrait servir d'indicateur de la maturité du cadre normatif
lui-même, et décroître à mesure que les textes d'application paraissent.
Cette décomposition n'est pas encore validée~: elle suppose une règle sûre
pour classer chaque \NEK{}, et certains relèvent des deux sources à la
fois.

% ============================================================================
\section{Une Piste de Recherche Actuelle : la Loi et la Chaîne Pénale}
\label{sec:c16-piste}
% ============================================================================

\begin{boxinfo}
\textbf{Statut de cette section.} Elle présente un travail en cours de
l'auteur de ce cours et de ses coauteurs, M.~Robert Thierry MASSOUSSI SONGO
et M.~Dieudonné BIALO, intitulé « Protection des données personnelles et
chaîne pénale au Cameroun~: analyse des zones de tension normative autour
de la Loi n°2024/017 »~\cite{minka2026chainepenale}. Ce manuscrit n'est pas
encore publié~; ses analyses et ses propositions sont celles des auteurs,
et elles peuvent évoluer. Les cas qu'il étudie sont construits à des fins
d'analyse.
\end{boxinfo}

\subsection{La question}

Les auteurs partent d'un constat~: la loi n°~2024/017 institue un régime
général, mais ses articles~4 et~72 renvoient la sécurité, la défense et la
justice à des textes particuliers ou spécifiques. Leur problématique est
de savoir dans quelle mesure la loi constitue le régime de référence pour
les traitements de la chaîne pénale (police judiciaire, expertise,
magistrature), et quelles zones de tension subsistent avec les autres
instruments applicables~: le Code pénal, le Code de procédure pénale, la
loi n°~2010/012, la Convention de Budapest à laquelle le Cameroun est
Partie, et la Convention de Malabo, signée mais non ratifiée.

\subsection{Trois niveaux d'articulation}

Les auteurs distinguent trois situations entre deux normes applicables à
la même opération~:
\begin{itemize}
    \item la \textbf{compatibilité}~: elles s'appliquent simultanément sans
          difficulté~;
    \item la \textbf{tension}~: leur application simultanée crée une
          difficulté qui demande une interprétation ou des garanties~;
    \item l'\textbf{incompatibilité}~: les deux obligations ne peuvent pas
          juridiquement être satisfaites ensemble.
\end{itemize}

Cette distinction, selon eux, évite de transformer toute difficulté
d'articulation en conflit de normes.

\subsection{La matrice d'analyse}

\begin{tableaucours}{cL{3.2cm}L{8.0cm}}{%
    \tableauHeader{N°} &
    \tableauHeader{Étape} &
    \tableauHeader{Question}
}
1 & Applicabilité & Chaque norme est-elle applicable~? \\
2 & Qualification & Quelle est la qualification juridique de l'opération~? \\
3 & Concurrence & Les normes portent-elles sur la même opération~? \\
4 & Conciliabilité & Imposent-elles réellement des obligations
simultanées~? \\
5 & Interprétation & Peut-on les concilier par interprétation~? \\
6 & Qualification de la tension & Compatibilité, tension ou
incompatibilité~? \\
7 & Mécanisme & Quel mécanisme de résolution est envisageable~? \\
8 & Proposition & Quelle solution proposer~? \\
\end{tableaucours}
\vspace{-0.8em}
\begin{center}{\small\itshape Tableau~19.1. La matrice des auteurs (d'après~\cite{minka2026chainepenale})}\end{center}

\subsection{Trois zones de tension}

\begin{tableaucours}{cL{3.0cm}L{8.2cm}}{%
    \tableauHeader{Zone} &
    \tableauHeader{Nature} &
    \tableauHeader{Description}
}
T1 & Temporelle & Les données collectées avant la loi, puis conservées,
analysées ou produites en justice après son entrée en vigueur~: quelle loi
régit chaque opération~? \\
T2 & Procédurale & L'information « sans délai » de la personne
(article~22) et le secret de l'enquête et de l'information judiciaire
(articles~102 et~154 du Code de procédure pénale) \\
T3 & Transfrontalière & La conservation rapide de la Convention de Budapest
(article~29) et l'autorisation préalable des articles~19 et~32, sous
réserve de l'article~72 \\
\end{tableaucours}
\vspace{-0.8em}
\begin{center}{\small\itshape Tableau~19.2. Les trois zones de tension (d'après~\cite{minka2026chainepenale})}\end{center}

Pour la zone temporelle, les auteurs privilégient une solution où la
légalité de la collecte est régie par la loi en vigueur au moment de la
collecte, en distinguant ensuite conservation et traitement ultérieur.
Pour la zone transfrontalière, ils recommandent de distinguer
conservation, préservation, production, transmission et transfert, qui
peuvent relever de régimes différents. Leur conclusion est que ces zones
ne constituent pas, à ce stade, des incompatibilités irréductibles~: leur
résolution appelle, selon les cas, une clarification jurisprudentielle,
réglementaire ou législative.

\subsection{Qualifier les acteurs opération par opération}

\begin{tableaucours}{L{2.4cm}L{9.8cm}}{%
    \tableauHeader{Acteur} &
    \tableauHeader{Qualification possible selon les auteurs}
}
Magistrat & Responsable du traitement lorsqu'il détermine les finalités et
les principaux éléments du traitement, sous réserve des circonstances \\
OPJ & Responsable du traitement ou autre acteur habilité, selon le cadre
institutionnel, les finalités et son autonomie de décision \\
Expert & Sous-traitant possible lorsqu'il traite pour le compte d'un
responsable sans déterminer lui-même les finalités \\
\end{tableaucours}
\vspace{-0.8em}
\begin{center}{\small\itshape Tableau~19.3. La grille fonctionnelle (d'après~\cite{minka2026chainepenale})}\end{center}

\subsection{Quatre propositions}

Les auteurs formulent quatre propositions~: clarifier le champ
d'application de la loi aux traitements judiciaires~; distinguer les
opérations dans les échanges transfrontaliers~; prévoir un mécanisme qui
permette au magistrat compétent d'apprécier le moment et les modalités de
l'information quand elle risque de compromettre une enquête~; élaborer une
doctrine sur la temporalité des traitements. Ils annoncent des études
complémentaires sur l'effet d'une éventuelle ratification de la Convention
de Malabo, sur la jurisprudence à venir, et sur les textes d'application de
la loi lorsqu'ils seront publiés.

% ============================================================================
\section{Deux Échelles à Trois Niveaux}
\label{sec:c16-echelles}
% ============================================================================

Le cours rapproche, à titre de question ouverte, deux instruments de
l'auteur. Les cotes \OUI{}, \NON{}, \NEK{} portent sur des \textit{faits}~:
une exigence est-elle satisfaite, au vu des preuves~? L'échelle
compatibilité, tension, incompatibilité porte sur des \textit{normes}~:
deux textes peuvent-ils être satisfaits ensemble~?

Une hypothèse de recherche se dessine~: une exigence qui relève d'une zone
de tension non résolue produit, dans un audit, un \NEK{} de critère. La
compatibilité permettrait de coter normalement~; l'incompatibilité
obligerait à choisir un texte et à le dire~; la tension laisserait un
\NEK{} tant qu'aucune interprétation ne s'impose. Si cette hypothèse se
vérifiait, le taux $U_c$ d'un audit pourrait se rattacher aux zones de
tension recensées par l'analyse juridique. Rien de cela n'est encore
établi~; c'est un programme de travail.

% ============================================================================
\section{Questions Ouvertes}
\label{sec:c16-questions}
% ============================================================================

\begin{itemize}
    \item \textbf{La granularité.} Comment normaliser le découpage des
          exigences pour comparer des audits entre eux~?
    \item \textbf{La pondération.} Toutes les exigences ne pèsent pas
          autant pour les personnes. Faut-il pondérer $\mu$ et $U$ par la
          gravité, et comment éviter que la pondération ne devienne
          arbitraire~?
    \item \textbf{La classification des \NEK{}.} Quelle règle permet de
          classer sûrement un \NEK{} en preuve ou en critère~?
    \item \textbf{La maturité normative.} Le taux $U_c$ mesuré sur plusieurs
          organisations permet-il de suivre la mise en place des textes
          d'application~?
    \item \textbf{L'articulation des normes.} La matrice des auteurs
          peut-elle être outillée pour guider l'auditeur devant une exigence
          prise entre deux textes~?
\end{itemize}

\begin{boxinfo}
\textbf{Construire un sujet de mémoire.}
\begin{enumerate}
    \item Choisir un objet~: une organisation, un secteur, un type de
          traitement, une zone de tension.
    \item Fixer l'ensemble des exigences et la règle de cotation.
    \item Recueillir des cas réels anonymisés, avec l'accord des
          organisations, ou des cas construits et documentés.
    \item Mesurer $\mu$, $U_p$, $U_c$ et $H$, ou appliquer la matrice des
          tensions.
    \item Confronter les résultats à un jugement d'experts.
    \item Dire les limites~: taille de l'échantillon, subjectivité des
          cotes, état des textes à la date de l'étude.
\end{enumerate}
\end{boxinfo}

% ============================================================================
\section{Ce que la Mesure demande à chaque Métier}
\label{sec:c16-metiers}
% ============================================================================

\begin{boxmetier}{Auditeur de la protection des données}
Classer chaque \NEK{} en preuve ou en critère, publier $U_p$ et $U_c$ à
côté de $\mu$, et dire dans le rapport ce qui ne dépend pas de
l'organisation.
\end{boxmetier}

\begin{boxmetier}{Juriste conformité}
Tenir la liste des zones de tension qui touchent l'organisation, avec,
pour chacune, l'interprétation retenue et sa justification, en attendant
les textes ou la jurisprudence.
\end{boxmetier}

\begin{boxmetier}{Responsable de traitement}
Suivre d'un audit à l'autre l'évolution de $U_p$, qui mesure ce que
l'organisation peut encore apprendre sur elle-même.
\end{boxmetier}

\begin{boxresume}
\begin{itemize}
    \item $\mu$, $U$ et $H$ décrivent l'état de conformité~; une baisse de
          $\mu$ au cours d'un audit traduit souvent une connaissance
          nouvelle.
    \item Le cours propose de décomposer $U$ en incertitude de preuve et
          incertitude de critère~; à NYANGA, $U_p \approx 0{,}09$ et
          $U_c \approx 0{,}13$.
    \item L'auteur et ses coauteurs étudient, dans un travail en cours, les
          tensions temporelle, procédurale et transfrontalière entre la loi
          et la chaîne pénale, avec une matrice en huit étapes.
    \item Ils distinguent compatibilité, tension et incompatibilité, et
          qualifient les acteurs opération par opération.
    \item Le rapprochement de cette échelle et des cotes \OUI{}, \NON{},
          \NEK{} est une question de recherche ouverte.
\end{itemize}
\end{boxresume}

\begin{boxinfo}
\textbf{Points de contrôle.} Vous avez compris ce chapitre si vous savez~:
\begin{itemize}
    \item calculer $U_p$ et $U_c$ pour un ensemble de cotes~;
    \item appliquer la matrice des tensions à un cas simple~;
    \item formuler une question de recherche et ses limites.
\end{itemize}
\end{boxinfo}

\begin{boxexo}[title={\ding{45}\quad Exercice~19.1. Calculer \niveaubase}]
Pour un audit de 20 exigences~: 6 \OUI{}, 9 \NON{}, 3 \NEK{} de preuve et
2 \NEK{} de critère. Calculez $\mu$, $U$, $U_p$, $U_c$ et $H$.
\end{boxexo}

\begin{boxexo}[title={\ding{45}\quad Exercice~19.2. Classer les NEK \niveaumoyen}]
Classez en preuve ou en critère les \NEK{} suivants, rencontrés dans le
cours~: (a) l'exécution des mesures par l'éditeur du logiciel bancaire~;
(b) le niveau de protection de la Côte d'Ivoire~; (c) la licéité du
profilage au regard de l'article~65~; (d) le pays de traitement des
selfies~; (e) l'application de l'article~22 en mars~2025.
\end{boxexo}

\begin{boxexo}[title={\ding{45}\quad Exercice~19.3. La matrice \niveaumoyen}]
Appliquez les huit étapes de la matrice à la tension T2 dans l'affaire
NYANGA~: la demande de l'OPJ du 20~mars~2025 et l'article~22. À quelle
étape la situation change-t-elle en février~2026~?
\end{boxexo}

\begin{boxexo}[title={\ding{45}\quad Exercice~19.4. Un sujet de mémoire \niveauexpert}]
Choisissez l'une des questions ouvertes de la section~\ref{sec:c16-questions}
et rédigez une proposition de mémoire d'une page~: question, hypothèse,
données, méthode, critère de validation, limites.
\end{boxexo}

\bigskip

\begin{boxinfo}
\textbf{Pour aller plus loin.}
\begin{itemize}
    \item T.~E.~Minka Mi Nguidjoï, « Vers l'auditologie »~\cite{minka2026auditologie}.
    \item T.~E.~Minka Mi Nguidjoï, sur l'instrument trichotomique~\cite{minka2026nek}.
    \item T.~E.~Minka Mi Nguidjoï, R.~T.~Massoussi Songo et D.~Bialo,
          manuscrit en préparation~\cite{minka2026chainepenale}.
\end{itemize}
\end{boxinfo}

\begin{boxinfo}
\textbf{Fin du cours.} L'affaire a commencé par une lettre, le lendemain
d'une échéance~: « Sommes-nous en règle~? » Elle s'achève sur un rapport
qui répond non sur vingt-huit points, oui sur huit, et qui dit pourquoi il
ne peut pas répondre sur dix. Entre les deux, l'audit a recensé, apprécié,
vérifié, coté et recommandé. La conformité n'est pas un état que l'on
déclare~; c'est une preuve que l'on construit, et dont on dit les limites.
\end{boxinfo}


% --------------------------------------------------------------------------
% ANNEXES
% --------------------------------------------------------------------------
\appendix

    % ============================================================================
% Fichier  : annexes/Ann_A_glossaire.tex
% Version  : 1.0 -- novembre 2026
% ============================================================================

\chapter{Glossaire}
\label{ann:glossaire}

Ce glossaire reprend, sous une forme abrégée, les définitions données dans
les chapitres. Les définitions tirées de l'article~5 de la loi n°~2024/017
sont paraphrasées~; le texte exact figure dans le chapitre indiqué.

\begin{description}
    \item[Anonymisation] Modification ou suppression d'informations pour
          rendre les données non attribuables à des personnes~; elle doit
          être prouvée (Ch.~\ref{chap:donnees}).
    \item[Autorisation préalable] Acte de l'Autorité auquel la loi soumet
          tout traitement (art.~19) et tout transfert (art.~32)
          (Ch.~\ref{chap:loi}).
    \item[Autorité de protection] Organisme public indépendant chargé de
          veiller à l'application de la loi, créé et organisé par décret
          (art.~53) (Ch.~\ref{chap:loi}).
    \item[Clauses contractuelles types] Clauses élaborées et publiées par
          l'Autorité pour encadrer les transferts (Ch.~\ref{chap:transferts}).
    \item[Consentement] Manifestation de volonté expresse, non équivoque, non
          contrainte, libre et éclairée~; préalable, spécifique et univoque
          selon l'article~9 (Ch.~\ref{chap:principes}).
    \item[Constatation] Fait vérifiable, rattaché à une pièce
          (Ch.~\ref{chap:prologue}).
    \item[Critère de substitution] Critère déclaré dans le plan d'audit
          quand le texte applicable n'est pas publié (Ch.~\ref{chap:droits}).
    \item[Décision automatisée] Décision fondée exclusivement sur un
          traitement automatisé (art.~44) (Ch.~\ref{chap:profilage}).
    \item[Découverte des données] Recherche, dans les supports réels, des
          données personnelles que le registre ignore
          (Ch.~\ref{chap:cartographie}).
    \item[Délégué à la protection des données] Fonction de bonne pratique,
          non prévue par la loi, chargée du registre, du conseil et des
          demandes (Ch.~\ref{chap:prologue}).
    \item[Destinataire] Personne habilitée à recevoir une communication des
          données (Ch.~\ref{chap:donnees}).
    \item[Donnée à caractère personnel] Information permettant d'identifier
          une personne directement ou indirectement (Ch.~\ref{chap:donnees}).
    \item[Donnée sensible] Donnée relative notamment aux opinions, aux
          transactions bancaires, à l'origine ethnique, linguistique ou
          régionale, à la vie sexuelle, à la génétique, à la biométrie, à la
          santé, aux poursuites et sanctions (Ch.~\ref{chap:donnees}).
    \item[Entropie ($H$)] Dispersion des cotes d'un ensemble d'exigences, en
          bits (Ch.~\ref{chap:audit}).
    \item[Évaluation d'impact] Procédure qui analyse la probabilité et la
          gravité des risques pour les droits et libertés (art.~33)
          (Ch.~\ref{chap:risque}).
    \item[Fichier] Ensemble structuré de données personnelles accessibles
          selon des critères (Ch.~\ref{chap:donnees}).
    \item[Finalité] But d'un traitement~; elle le définit et le limite
          (art.~11) (Ch.~\ref{chap:principes}).
    \item[Interconnexion] Mise en relation de données traitées pour une
          finalité avec d'autres données (art.~36)
          (Ch.~\ref{chap:cartographie}).
    \item[\NEK{} de critère] Incertitude qui tient à une règle non fixée
          (texte non publié, lecture non tranchée) (Ch.~\ref{chap:audit}).
    \item[\NEK{} de preuve] Incertitude qui tient à une pièce manquante
          (Ch.~\ref{chap:audit}).
    \item[\OUI{}, \NON{}, \NEK{}] Cotes d'une exigence~: conformité établie,
          écart établi, connaissance insuffisante (Ch.~\ref{chap:prologue}).
    \item[Papier de travail] Document qui trace un constat~: exigence,
          critère, test, preuve, cote, auteur (Ch.~\ref{chap:audit}).
    \item[Personne concernée] Personne dont les données sont traitées
          (Ch.~\ref{chap:donnees}).
    \item[Portabilité] Droit de recevoir gratuitement ses données dans un
          format structuré et de les faire transmettre (art.~43)
          (Ch.~\ref{chap:droits}).
    \item[Profilage] Traitement automatisé qui évalue des aspects personnels
          (santé, préférences, localisation, situation économique)
          (Ch.~\ref{chap:profilage}).
    \item[Pseudonymisation] Remplacement des identifiants directs par des
          identifiants artificiels~; les données restent personnelles
          (Ch.~\ref{chap:donnees}).
    \item[Registre des activités de traitement] Registre exigé par
          l'article~29 (Ch.~\ref{chap:cartographie}).
    \item[Responsable du traitement] Personne qui détermine les finalités et
          les moyens d'un traitement (Ch.~\ref{chap:donnees}).
    \item[Responsables conjoints] Responsables qui déterminent ensemble
          finalités et moyens~; convention publique (art.~31)
          (Ch.~\ref{chap:soustraitants}).
    \item[Sous-traitant] Personne qui traite des données pour le compte et
          sur instruction du responsable (Ch.~\ref{chap:donnees}).
    \item[Taux de conformité ($\mu$)] Part des exigences cotées \OUI{}
          (Ch.~\ref{chap:audit}).
    \item[Taux d'incertitude ($U$)] Part des exigences cotées \NEK{}~;
          décomposable en $U_p$ et $U_c$ (Ch.~\ref{chap:recherche}).
    \item[Tension normative] Situation où l'application simultanée de deux
          normes demande une interprétation, entre la compatibilité et
          l'incompatibilité (Ch.~\ref{chap:recherche}).
    \item[Traitement] Toute opération appliquée à des données personnelles,
          de la collecte à la destruction (Ch.~\ref{chap:donnees}).
    \item[Transfert international] Acheminement de données au-delà des
          frontières (Ch.~\ref{chap:transferts}).
    \item[Violation de données] Atteinte à la sécurité entraînant
          destruction, perte, altération, divulgation ou accès non autorisé
          (Ch.~\ref{chap:violations}).
\end{description}

    % ============================================================================
% Fichier  : annexes/Ann_B_outils.tex
% Version  : 1.0 -- novembre 2026
% ============================================================================

\chapter{Outils}
\label{ann:outils}

Les outils servent l'audit~; ils ne le remplacent pas. Chacun produit des
pièces qu'il faut ensuite qualifier, coter et documenter. Les outils de
découverte et d'analyse s'utilisent avec l'accord écrit du commanditaire,
sur les seuls systèmes de l'organisation.

\begin{tableaucours}{L{3.6cm}L{4.4cm}L{4.0cm}}{%
    \tableauHeader{Besoin} &
    \tableauHeader{Outil ou méthode} &
    \tableauHeader{Chapitre}
}
Tenir le registre & Tableur structuré selon le Tableau~4.1, ou base
légère~; un modèle figure à l'annexe~\ref{ann:modeles} &
\ref{chap:cartographie} \\
Découvrir les données & Recherche de motifs par expressions régulières
(\texttt{grep}, \texttt{ripgrep}) sur les partages~; lecture des schémas
de bases & \ref{chap:cartographie} \\
Conduire une évaluation d'impact & Logiciel PIA de la CNIL (libre) ou
modèle tableur & \ref{chap:risque} \\
Analyser les flux d'une application & Mandataire d'interception sur un
appareil de test (\texttt{mitmproxy}), journaux DNS, \texttt{Wireshark} &
\ref{chap:transferts} \\
Repérer les traqueurs d'une application mobile & Analyse statique
(\texttt{MobSF}), bases publiques de traqueurs (Exodus Privacy) &
\ref{chap:cartographie} \\
Exploiter des journaux d'accès & Tableur, \texttt{Python} et
\texttt{pandas} pour le rapprochement avec les plannings &
\ref{chap:securite} \\
Mesurer $\mu$, $U$, $H$ & Tableur ou script~; formules du
Chapitre~\ref{chap:audit} & \ref{chap:audit} \\
Comparer des taux de décision & Tableaux croisés par groupe déclaré &
\ref{chap:profilage} \\
Protéger les données de l'audit & Conteneur chiffré, masquage des
identités dans les papiers & \ref{chap:normes} \\
\end{tableaucours}
\vspace{-0.8em}
\begin{center}{\small\itshape Tableau~B.1. Outils de l'audit de protection des données}\end{center}

\begin{boxinfo}
\textbf{Un outil a ses angles morts.} Une recherche de motifs ne voit pas
les images scannées~; un mandataire d'interception ne voit pas les flux
d'un serveur~; une analyse statique ne dit pas ce qu'une bibliothèque
envoie vraiment. L'auditeur écrit, pour chaque outil, ce qu'il a couvert
et ce qu'il n'a pas pu couvrir.
\end{boxinfo}

    % ============================================================================
% Fichier  : annexes/Ann_C_modeles.tex
% Version  : 1.0 -- novembre 2026
% ============================================================================

\chapter{Modèles de Documents}
\label{ann:modeles}

Ces modèles sont des points de départ~; chaque mission les adapte et le
dit.

\section*{C.1 Fiche de registre}

\begin{tableaucours}{L{4.6cm}L{7.6cm}}{%
    \tableauHeader{Rubrique} &
    \tableauHeader{Contenu}
}
Nom et finalité & Une finalité, en une phrase \\
Responsable, sous-traitants & Noms et coordonnées, contrats \\
Personnes concernées & Catégories \\
Données & Catégories, données sensibles signalées \\
Fondement & Consentement ou dérogation de l'article~9, pièce \\
Destinataires & Internes, externes \\
Transferts & Pays, garanties, état de l'autorisation \\
Durée de conservation & Base active, archive, purge \\
Mesures de sécurité & Renvoi aux mesures de l'article~27 \\
Autorisation & Numéro, ou état de la demande \\
Propriétaire, revue & Nom, date \\
\end{tableaucours}

\section*{C.2 Papier de travail}

\begin{tableaucours}{L{4.6cm}L{7.6cm}}{%
    \tableauHeader{Rubrique} &
    \tableauHeader{Contenu}
}
Référence & Domaine, numéro de l'exigence \\
Exigence & Question fermée, atomique \\
Critère & Article, norme, ou critère de substitution déclaré \\
Test & Ce qui a été fait, sur quel échantillon \\
Preuve & Référence des pièces, sans données personnelles en clair \\
Cote & \OUI{}, \NON{} ou \NEK{} (de preuve ou de critère) \\
Commentaire & Constatation, puis analyse avec niveau de confiance \\
Auteur, réviseur, dates & \\
\end{tableaucours}

\section*{C.3 Réponse à une demande d'accès}

Objet~: votre demande d'accès du [date]. Nous confirmons traiter des données
vous concernant. Vous trouverez ci-joint~: la liste de ces données et leur
origine~; les finalités des traitements~; les destinataires ou catégories
de destinataires~; la durée de conservation~; les transferts éventuels hors
du Cameroun. Vous pouvez demander la rectification, l'effacement ou la
limitation de ces données, et vous opposer à leur traitement, en écrivant à
[contact]. Vous pouvez également introduire une réclamation auprès de
l'Autorité de protection des données à caractère personnel.

\section*{C.4 Information des personnes après une violation}

Objet~: incident concernant vos données. Le [date], nous avons constaté
[nature de la violation]. Les données concernées sont [catégories]. Les
conséquences possibles pour vous sont [conséquences]. Nous avons pris les
mesures suivantes~: [mesures]. Nous vous recommandons de [mesures à la
portée de la personne]. Pour toute question~: [contact].

\section*{C.5 Plan d'audit}

Mandat (commanditaire, objectif, destinataires, date de remise)~;
critères (articles, textes publiés, norme, politiques internes, critères
de substitution déclarés)~; périmètre et exclusions~; domaines et
exigences~; tests et échantillons~; traitement des données de l'audit~;
règle de décision et seuils~; calendrier.

\section*{C.6 Clauses de sous-traitance}

Les douze clauses du Tableau~13.1 du Chapitre~\ref{chap:soustraitants}.

    % ============================================================================
% Fichier  : annexes/Ann_D_textes.tex
% Version  : 1.0 -- novembre 2026
% ============================================================================

% Évite un titre de section isolé en bas de page.
\providecommand{\gardetitre}{\par\vskip 0pt plus 9\baselineskip\penalty-300\vskip 0pt plus -9\baselineskip}

\chapter{Textes et Normes}
\label{ann:textes}

\gardetitre
\section*{Degré de vérification des textes cités}

\begin{tableaucours}{L{4.6cm}L{3.6cm}L{3.8cm}}{%
    \tableauHeader{Texte} &
    \tableauHeader{Source utilisée} &
    \tableauHeader{Citation dans le cours}
}
Loi n°~2024/017 du 23~décembre~2024 & Copie certifiée conforme du
Secrétariat général de la Présidence & Articles cités mot pour mot \\
Loi n°~2010/012 du 21~décembre~2010 & Texte officiel & Articles 41 à 47
cités mot pour mot \\
Code de procédure pénale, articles 102 et 154 & D'après
Minka Mi Nguidjoï et al. (2026) & Paraphrasés, non cités \\
Convention de Malabo (2014) & Texte de l'Union africaine~; statut d'après
Minka Mi Nguidjoï et al. (2026) & Paraphrasée \\
Convention de Budapest (2001) & Conseil de l'Europe~; statut d'après
Minka Mi Nguidjoï et al. (2026) & Paraphrasée \\
RGPD (UE) 2016/679 & Journal officiel de l'Union européenne & Paraphrasé,
pour comparaison \\
Normes ISO et ISO/IEC & Sous licence & Paraphrasées, jamais
reproduites \\
Méthode PIA de la CNIL & Publication libre & Paraphrasée \\
\end{tableaucours}

\gardetitre
\section*{Textes internationaux et étrangers cités}

\begin{tableaucours}{L{4.4cm}L{2.2cm}L{5.4cm}}{%
    \tableauHeader{Texte} &
    \tableauHeader{Nature} &
    \tableauHeader{Place dans le cours}
}
Lignes directrices de l'OCDE (1980, révisées en 2013)~\cite{ocde1980} &
Recommandation & Origine des principes, Chapitre~\ref{chap:international} \\
Principes directeurs des Nations Unies (1990)~\cite{onu1990} & Résolution &
Chapitre~\ref{chap:international} \\
Convention~108 (1981) et protocole de 2018~\cite{coe1981conv108,coe2018conv108plus} &
Traité du Conseil de l'Europe & Chapitre~\ref{chap:international}~; fiches
pays \\
Convention de Malabo (2014)~\cite{ua2014malabo} & Traité de l'Union
africaine & Chapitre~\ref{chap:international} \\
Acte additionnel de la CEDEAO (2010)~\cite{cedeao2010acte} & Acte
communautaire & Comparaison, Chapitre~\ref{chap:international} \\
Directive n°~07/08-UEAC-133-CM-18 de la CEMAC (2008)~\cite{cemac2008} &
Directive communautaire & Chapitres~\ref{chap:loi2010} et~\ref{chap:secteurs} \\
Convention de Budapest (2001)~\cite{budapest2001} & Traité du Conseil de
l'Europe & Chapitres~\ref{chap:international}, \ref{chap:transferts}
et~\ref{chap:recherche} \\
RGPD (UE) 2016/679~\cite{ue2016rgpd} & Règlement européen &
Chapitre~\ref{chap:rgpd}~; double cote aux Chapitres~\ref{chap:transferts}
et~\ref{chap:violations} \\
Clauses contractuelles types de 2021~\cite{ue2021clauses} & Décision de la
Commission & Chapitre~\ref{chap:transferts} \\
Règlement (UE) 2024/1689 sur l'intelligence artificielle~\cite{ue2024ia} &
Règlement européen & Chapitre~\ref{chap:profilage} \\
Loi ivoirienne n°~2013-450~\cite{ci2013loi} & Loi étrangère & Fiche pays,
Chapitre~\ref{chap:international} \\
Loi marocaine n°~09-08~\cite{ma2009loi} & Loi étrangère & Fiche pays,
Chapitre~\ref{chap:international} \\
NIST Privacy Framework~1.0~\cite{nist2020pf} & Cadre volontaire &
Chapitre~\ref{chap:normes} \\
NIST AI RMF~1.0~\cite{nist2023airmf} & Cadre volontaire &
Chapitre~\ref{chap:profilage} \\
Stratégie continentale de l'Union africaine sur l'intelligence
artificielle (2024)~\cite{ua2024ia} & Orientation politique &
Chapitre~\ref{chap:profilage} \\
\end{tableaucours}

Aucun de ces textes n'est un critère obligatoire pour une organisation qui
ne traite que des données de personnes situées au Cameroun, à l'exception
des traités auxquels le Cameroun est Partie et des actes de la CEMAC. Le
Chapitre~\ref{chap:loi2010} donne la méthode pour choisir ses critères.

\gardetitre
\section*{Textes d'application attendus}

À la date du cours, l'auteur n'a pas trouvé de texte publié pour les
dispositions suivantes, qui renvoient à la voie réglementaire ou à un
référentiel de l'Autorité~: organisation de l'Autorité (art.~53)~;
modalités d'autorisation (art.~19, 32)~; critères d'évaluation du risque et
de l'évaluation d'impact (art.~25, 33)~; durées de conservation (art.~13,
28)~; délais d'examen des demandes (art.~47)~; référentiel des mesures
(art.~24, 27, 57)~; mécanisme de certification (art.~34)~; contrôle
(art.~35)~; textes spécifiques des secteurs (art.~4, 72). Le lecteur
vérifiera l'état de ces textes au moment où il lit.

\gardetitre
\section*{Correspondance des articles et des chapitres}

\begin{tableaucours}{L{4.0cm}L{8.2cm}}{%
    \tableauHeader{Articles} &
    \tableauHeader{Chapitres}
}
1 à 5 & \ref{chap:donnees} \\
6 à 18 & \ref{chap:principes} \\
19 & \ref{chap:loi} \\
20, 30, 31 & \ref{chap:soustraitants} \\
21, 37 à 47 & \ref{chap:principes}, \ref{chap:droits} \\
22 & \ref{chap:violations} \\
24 à 28 & \ref{chap:securite}, \ref{chap:programme} \\
29, 36 & \ref{chap:cartographie} \\
32 & \ref{chap:transferts} \\
33 & \ref{chap:risque} \\
34 & \ref{chap:loi}, \ref{chap:normes} \\
44, 65 & \ref{chap:profilage} \\
48 à 52, 72 & \ref{chap:donnees}, \ref{chap:secteurs} \\
53 à 71 & \ref{chap:loi} \\
73 & \ref{chap:prologue}, \ref{chap:recherche} \\
\end{tableaucours}

    % ============================================================================
% Fichier  : annexes/Ann_E_bibliographie.tex
% Version  : 1.0 -- novembre 2026
% ============================================================================

\chapter{Bibliographie Commentée}
\label{ann:bibliographie}

\begin{description}
    \item[Les textes.] La loi n°~2024/017~\cite{loi2024017} est la source
          première du cours~; on la lit en entier, une fois, avant
          d'auditer. La loi n°~2010/012~\cite{loi2010012} reste utile pour
          la vie privée dans les communications.
    \item[Le cadre africain.] La Convention de Malabo~\cite{ua2014malabo}
          et l'acte additionnel de la CEDEAO~\cite{cedeao2010acte} situent
          la loi camerounaise dans son environnement régional.
    \item[La comparaison européenne.] Le RGPD~\cite{ue2016rgpd} et les
          recommandations du Comité européen~\cite{edpb2020transferts}
          servent de référence, à condition de ne pas les confondre avec la
          loi camerounaise.
    \item[Les normes.] L'ISO/IEC~29100~\cite{iso29100} pour le vocabulaire,
          l'ISO/IEC~27701~\cite{iso27701} pour le système de management,
          l'ISO/IEC~29134~\cite{iso29134} pour l'évaluation d'impact,
          l'ISO~19011~\cite{iso19011_2018} pour la conduite de l'audit.
    \item[La méthode PIA de la CNIL.]~\cite{cnil2018pia} La plus accessible
          des méthodes d'évaluation d'impact~; ses échelles sont reprises au
          Chapitre~\ref{chap:risque}.
    \item[La vie privée.] D.~J.~Solove~\cite{solove2008} et
          H.~Nissenbaum~\cite{nissenbaum2010}, pour comprendre ce que la
          protection des données protège.
    \item[La recherche de l'auteur.] « Vers l'auditologie »~\cite{minka2026auditologie}
          et l'instrument trichotomique~\cite{minka2026nek}~; le manuscrit
          en préparation sur la chaîne pénale~\cite{minka2026chainepenale}.
\end{description}

    % ============================================================================
% Fichier  : annexes/Ann_F_contacts.tex
% Version  : 1.0 -- novembre 2026
% ============================================================================

\chapter{Contacts et Ressources}
\label{ann:contacts}

\section*{Au Cameroun}

\begin{tableaucours}{L{4.6cm}L{7.6cm}}{%
    \tableauHeader{Organisme} &
    \tableauHeader{Rôle pour le cours}
}
Autorité de protection des données à caractère personnel & Autorisations,
plaintes, référentiel, liste des pays (art.~53)~; organisation par décret
à vérifier \\
Agence nationale des technologies de l'information et de la
communication (ANTIC) & Sécurité des systèmes d'information~; agents
habilités au titre de la loi n°~2010/012 \\
Régulateurs sectoriels & Autorisations prévues par l'article~48(2) pour
les transactions bancaires \\
Juridictions & Référé et réparation (art.~62)~; poursuites pénales \\
\end{tableaucours}

\section*{Ressources internationales}

\begin{tableaucours}{L{4.6cm}L{7.6cm}}{%
    \tableauHeader{Ressource} &
    \tableauHeader{Adresse}
}
Méthode et logiciel PIA (CNIL) & \url{https://www.cnil.fr/} \\
Comité européen de la protection des données & \url{https://www.edpb.europa.eu/} \\
Union africaine, traités & \url{https://au.int/} \\
Base de traqueurs Exodus Privacy & \url{https://exodus-privacy.eu.org/} \\
\end{tableaucours}

\section*{Le cours}

Dépôt du cours~: \url{https://github.com/MaletYon/cours-apdc}. NDAP SUKULU en
ligne~: \url{https://maletyon.github.io/cours-apdc/ndap-sukulu/}. Compte de
démonstration~: matricule \texttt{ETU001}, mot de passe \texttt{eTU001}.
Contact~: \href{mailto:maletyon@proton.me}{maletyon@proton.me}.

    % ============================================================================
% Fichier  : annexes/Ann_G_exigences.tex
% Version  : 1.0 -- novembre 2026
% Objet    : les quarante-six exigences de la mission NYANGA, par domaine,
%            avec leur cote, leur preuve et la nature de chaque NEK.
%            Les lignes reprennent mot pour mot les encadrés des chapitres
%            8, 11, 12, 13, 15 et 16.
% ============================================================================

% Évite un titre de section isolé en bas de page.
\providecommand{\gardetitre}{\par\vskip 0pt plus 9\baselineskip\penalty-300\vskip 0pt plus -9\baselineskip}

\chapter{Les Quarante-six Exigences de la Mission NYANGA}
\label{ann:exigences}

Cette annexe réunit les exigences cotées par KORO AUDIT à NYANGA CRÉDIT du
24~juin au 25~septembre~2026. Elle permet de refaire chaque calcul du
livre~: 8 + 10 + 8 + 8 + 8 + 4 = 46 exigences, 8~\OUI{}, 28~\NON{},
10~\NEK{}. La dernière colonne donne la nature de chaque \NEK{}~: de
\textit{preuve} quand une pièce manque, de \textit{critère} quand c'est
le texte ou la référence qui manque.

\gardetitre
\section*{D1. Recensement et fondements (Chapitre~\ref{chap:audit})}

\begin{tableaucours}{cL{6.2cm}cL{3.6cm}}{%
    \tableauHeader{N°} &
    \tableauHeader{Exigence} &
    \tableauHeader{Cote} &
    \tableauHeader{Preuve, nature du NEK}
}
1 & Les familles de données sont recensées & \OUI{} & Inventaire du 7~juillet \\
2 & Un registre des activités de traitement est tenu & \NON{} & Aucun registre \\
3 & Chaque traitement a un fondement valable & \NON{} & Formulaires, cahiers \\
4 & L'autorisation de l'Autorité est demandée & \NON{} & Attestation \\
5 & Une durée de conservation est fixée & \NON{} & Aucune règle \\
6 & L'information de l'article~21 est fournie & \NON{} & Formulaire 2016 \\
7 & L'identification des clients repose sur une obligation légale & \NEK{} &
Texte non produit. \textit{Critère} \\
8 & La vidéosurveillance repose sur un fondement valable & \NEK{} &
Fondement incertain. \textit{Critère} \\
\end{tableaucours}

\gardetitre
\section*{D2. Droits des personnes (Chapitre~\ref{chap:droits})}

\begin{tableaucours}{cL{6.2cm}cL{3.6cm}}{%
    \tableauHeader{N°} &
    \tableauHeader{Exigence} &
    \tableauHeader{Cote} &
    \tableauHeader{Preuve, nature du NEK}
}
9 & Les personnes sont informées de leurs droits & \NON{} & Formulaire 2016 \\
10 & Une procédure de traitement des demandes existe & \NON{} & Aucune \\
11 & Un registre des demandes est tenu & \NON{} & Reconstitué par l'audit \\
12 & Les demandes d'accès reçoivent une réponse & \NON{} & 12 sur 31 \\
13 & Les réponses respectent le critère de 30 jours & \NON{} & 8 sur 12 \\
14 & Les oppositions sont exécutées dans les systèmes & \NON{} & 3 numéros
encore appelés \\
15 & La portabilité est possible & \NON{} & Refus du 14~août \\
16 & L'identité est vérifiée de façon proportionnée & \NEK{} & Pratiques non
documentées. \textit{Preuve} \\
17 & Les effacements acceptés sont exécutés & \NEK{} & Sauvegardes non
examinées. \textit{Preuve} \\
18 & La procédure de succession respecte l'article~45 & \OUI{} & Note de
2025, 3 dossiers \\
\end{tableaucours}

\gardetitre
\section*{D3. Sécurité et confidentialité (Chapitre~\ref{chap:securite})}

\begin{tableaucours}{cL{6.2cm}cL{3.6cm}}{%
    \tableauHeader{N°} &
    \tableauHeader{Exigence} &
    \tableauHeader{Cote} &
    \tableauHeader{Preuve, nature du NEK}
}
19 & L'équipement à l'origine de l'incident est remplacé et tenu à jour &
\OUI{} & Bon de livraison, état des correctifs \\
20 & Un inventaire des équipements est tenu & \OUI{} & Inventaire de
mars~2026 \\
21 & Les journaux sont centralisés et conservés & \OUI{} & Console, treize
mois en ligne \\
22 & Les sauvegardes sont restaurables & \OUI{} & Test du 3~septembre \\
23 & Chaque agent n'accède qu'aux clients de son agence & \NON{} & Test sur
six comptes \\
24 & Les relevés sont transmis par un canal maîtrisé & \NON{} & Messagerie
personnelle \\
25 & Les habilitations sont revues & \NON{} & Aucune revue depuis 2023 \\
26 & Les mesures suivent le référentiel de l'Autorité (art.~24) & \NEK{} &
Référentiel non publié. \textit{Critère} \\
\end{tableaucours}

\gardetitre
\section*{D4. Sous-traitants et transferts (Chapitre~\ref{chap:soustraitants})}

\begin{tableaucours}{cL{6.2cm}cL{3.6cm}}{%
    \tableauHeader{N°} &
    \tableauHeader{Exigence} &
    \tableauHeader{Cote} &
    \tableauHeader{Preuve, nature du NEK}
}
27 & Chaque sous-traitant est lié par un contrat & \OUI{} & Trois contrats \\
28 & Les contrats définissent données et personnes (art.~30(3)) & \NON{} &
SOREC, opérateur \\
29 & Les contrats contiennent les douze clauses & \NON{} & Un sur trois au
mieux \\
30 & Le droit d'audit est prévu & \NON{} & Aucun contrat \\
31 & SOREC détruit les données des dossiers clos & \NON{} & Classeurs
conservés \\
32 & L'éditeur applique les mesures de sécurité convenues & \NEK{} & Ni
audit ni rapport. \textit{Preuve} \\
33 & L'opérateur de SMS limite l'usage des numéros & \NEK{} & Aucune pièce.
\textit{Preuve} \\
34 & La convention avec SAWA ASSURANCES est publique & \NON{} & Inexistante \\
\end{tableaucours}

\gardetitre
\section*{D5. Violations (Chapitre~\ref{chap:violations})}

\begin{tableaucours}{cL{6.2cm}cL{3.6cm}}{%
    \tableauHeader{N°} &
    \tableauHeader{Exigence} &
    \tableauHeader{Cote} &
    \tableauHeader{Preuve, nature du NEK}
}
35 & Une procédure de gestion des violations existe & \NON{} & Aucune \\
36 & Un registre des violations est tenu & \NON{} & Aucun \\
37 & La violation de 2025 est qualifiée pour les personnes & \NON{} &
Rapport technique seul \\
38 & La demande de différer est écrite et motivée & \OUI{} & Lettre du
20~mars~2025 \\
39 & La demande a été réexaminée & \NON{} & Aucune démarche \\
40 & Les personnes ont été informées quand la raison a cessé & \NON{} &
Rien au 31~août~2026 \\
41 & L'Autorité a été informée & \NEK{} & Autorité non organisée.
\textit{Critère} \\
42 & L'obligation d'informer s'appliquait en mars~2025 & \NEK{} & Question
de droit non tranchée. \textit{Critère} \\
\end{tableaucours}

\gardetitre
\section*{D6. Gouvernance (Chapitre~\ref{chap:programme})}

\begin{tableaucours}{cL{6.2cm}cL{3.6cm}}{%
    \tableauHeader{N°} &
    \tableauHeader{Exigence} &
    \tableauHeader{Cote} &
    \tableauHeader{Preuve, nature du NEK}
}
43 & Le risque de chaque traitement est évalué (art.~25) & \NON{} & Aucune
évaluation \\
44 & Une évaluation d'impact précède les traitements à risque élevé
(art.~33) & \NON{} & Douze traitements, aucune \\
45 & Un dossier de conformité peut être présenté au contrôle (art.~35) &
\NON{} & Aucun dossier \\
46 & Le rapport annuel sur les mesures est adressé à l'Autorité
(art.~27(2)) & \NEK{} & Autorité non organisée. \textit{Critère} \\
\end{tableaucours}

\gardetitre
\section*{Récapitulation}

\begin{tableaucours}{cL{4.2cm}ccccc}{%
    \tableauHeader{Code} &
    \tableauHeader{Domaine} &
    \tableauHeader{\OUI{}} &
    \tableauHeader{\NON{}} &
    \tableauHeader{\NEK{}} &
    \tableauHeader{de preuve} &
    \tableauHeader{de critère}
}
D1 & Recensement et fondements & 1 & 5 & 2 & 0 & 2 \\
D2 & Droits des personnes & 1 & 7 & 2 & 2 & 0 \\
D3 & Sécurité et confidentialité & 4 & 3 & 1 & 0 & 1 \\
D4 & Sous-traitants et transferts & 1 & 5 & 2 & 2 & 0 \\
D5 & Violations & 1 & 5 & 2 & 0 & 2 \\
D6 & Gouvernance & 0 & 3 & 1 & 0 & 1 \\
Total & & 8 & 28 & 10 & 4 & 6 \\
\end{tableaucours}

Les quatre \NEK{} de preuve se lèvent par le travail de l'auditeur ou de
l'établissement~: une procédure écrite, l'examen des sauvegardes, un audit
des deux prestataires. Les six \NEK{} de critère ne dépendent ni de l'un ni
de l'autre. Ils attendent un texte d'application, l'organisation de
l'Autorité ou une décision de justice
(Chapitre~\ref{chap:recherche}). Le transfert vers Abidjan n'entre pas dans
ce décompte~: le Chapitre~\ref{chap:programme} dit pourquoi et le cote à
part.


% --------------------------------------------------------------------------
% ÉLÉMENTS FINAUX
% --------------------------------------------------------------------------
\backmatter

    \nocite{*}
    \printbibliography[heading=bibintoc, title={Références bibliographiques}]

    \printindex

% ==============================================================================
% 4E DE COUVERTURE
% ==============================================================================
\clearpage
\thispagestyle{empty}
\pagecolor{CouleurPrimaire}

\begin{tikzpicture}[remember picture, overlay]
    \fill[CouleurAccent]
        (current page.south west) rectangle
        ([yshift=1.8cm] current page.south east);
    \fill[CouleurAccent]
        ([yshift=-1.4cm] current page.north west) rectangle
        ([yshift=-1.7cm] current page.north east);
\end{tikzpicture}

\vspace*{1.8cm}

\begin{center}
    {\color{CouleurAccent}\rule{0.6\linewidth}{1pt}}\\[1.5ex]
    {\color{white}\Large\bfseries À propos de ce cours}\\[1.5ex]
    {\color{CouleurAccent}\rule{0.6\linewidth}{1pt}}
\end{center}

\vspace{1cm}

\begin{center}
\begin{minipage}{0.75\textwidth}
\color{white!85!gray}\setstretch{1.4}
Depuis le 23 juin 2026, les organisations camerounaises doivent appliquer la
loi n°~2024/017 relative à la protection des données à caractère personnel.
Dire qu'elles le font demande une preuve. Ce cours apprend à la chercher~:
recenser les traitements, apprécier le risque pour les personnes, vérifier
les droits, la sécurité, les sous-traitants et les transferts, puis conclure
en disant ce que l'on sait et ce que l'on ignore.

\medskip

Il progresse en sept parties, précédées d'un prologue~: les fondements, le
contexte international (RGPD, NIST, normes ISO), le contexte national (lois
de 2024 et de 2010), la méthode d'audit, la pratique, les défis avancés et
la recherche sur la mesure de la conformité.

\medskip

Il s'adresse aux étudiants du Master Professionnel en Sécurité des Systèmes
Numériques et aux professionnels qui exercent comme auditeurs de la
protection des données, délégués à la protection des données, juristes
conformité, RSSI ou responsables de traitement.
\end{minipage}
\end{center}

\vfill

\begin{center}
    {\color{CouleurAccent}\rule{0.7\linewidth}{1.2pt}}\\[1.2ex]
    {\color{white}\large\bfseries MINKA MI NGUIDJOÏ Thierry Emmanuel}\\[0.4ex]
    {\color{white!80!gray}\small
        CISA~·~CISM~·~CGEIT~·~CRISC~·~CDPSE~·~ISO 27001~·~ISO 22301~·~ISO 9001}\\[0.3ex]
    {\color{white!70!gray}\footnotesize\itshape
        Chercheur au laboratoire d'Ingénierie Mathématiques et des Systèmes d'Information}\\[0.3ex]
    {\color{white!70!gray}\footnotesize\itshape
        Chercheur~|~Cryptologue~|~Auditeur Senior~|~Expert Judiciaire~|~GRC \& Sécurité de l'Information}\\[1.2ex]
    {\color{CouleurAccent}\rule{0.7\linewidth}{1.2pt}}
\end{center}

\vspace{1cm}

\afterpage{\nopagecolor}
\end{document}
